\documentclass[11pt,a4paper]{article}
\usepackage[T1]{fontenc}
\usepackage[utf8]{inputenc}
\usepackage{lmodern,microtype}
\usepackage[a4paper,margin=25mm,headheight=14pt]{geometry}
\usepackage{amsmath,amssymb,amsthm,mathtools,mathrsfs}
\usepackage{booktabs,array,enumitem,xcolor,fancyhdr,longtable}
\usepackage[hyphens]{url}
\usepackage[colorlinks=true,linkcolor=blue!45!black,citecolor=blue!45!black,urlcolor=blue!45!black]{hyperref}
\usepackage{bookmark}
\hypersetup{pdftitle={Exact Recurrences, Algebraic Independence, and Certified Statistical Recovery},pdfauthor={KEIO}}
\newtheorem{theorem}{Theorem}[section]
\newtheorem{lemma}[theorem]{Lemma}
\newtheorem{proposition}[theorem]{Proposition}
\newtheorem{corollary}[theorem]{Corollary}
\newtheorem{conjecture}[theorem]{Conjecture}
\theoremstyle{definition}
\newtheorem{definition}[theorem]{Definition}
\newtheorem{example}[theorem]{Example}
\theoremstyle{remark}
\newtheorem{remark}[theorem]{Remark}
\newcommand{\C}{\mathbb C}
\newcommand{\Q}{\mathbb Q}
\newcommand{\Z}{\mathbb Z}
\newcommand{\R}{\mathbb R}
\newcommand{\N}{\mathbb N}
\newcommand{\T}{\mathbb T}
\newcommand{\Qbar}{\overline{\mathbb Q}}
\newcommand{\E}{\mathbb E}
\newcommand{\PP}{\mathbb P}
\newcommand{\fl}[1]{\lfloor #1\rfloor}
\newcommand{\cl}[1]{\lceil #1\rceil}
\newcommand{\ff}[2]{#1^{\underline{#2}}}
\DeclareMathOperator{\ord}{ord}
\DeclareMathOperator{\Res}{Res}
\DeclareMathOperator{\Sol}{Sol}
\DeclareMathOperator{\sgn}{sgn}
\DeclareMathOperator{\SL}{SL}
\DeclareMathOperator{\Gal}{Gal}
\DeclareMathOperator{\spn}{span}
\DeclareMathOperator{\lc}{lc}
\DeclareMathOperator{\tr}{tr}
\DeclareMathOperator{\gr}{gr}
\DeclareMathOperator{\Sp}{Sp}
\DeclareMathOperator{\rank}{rank}
\DeclareMathOperator{\supp}{supp}
\DeclareMathOperator{\diag}{diag}
\DeclareMathOperator{\Var}{Var}
\DeclareMathOperator{\Cov}{Cov}
\DeclareMathOperator{\CT}{CT}
\numberwithin{equation}{section}
\setlength{\emergencystretch}{3em}
\setlength{\parskip}{3pt}
\allowdisplaybreaks[2]
\setcounter{tocdepth}{1}
\pagestyle{fancy}\fancyhf{}
\fancyhead[L]{\small KEIO}
\fancyhead[R]{\small Exact recurrences and certified recovery}
\fancyfoot[C]{\thepage}
\begin{document}
\title{Exact Recurrences, Algebraic Independence,\\
and Certified Statistical Recovery\\[6pt]
\large Vertex balance, Bessel modules, and entropy transport}
\author{KEIO}
\date{20 October 2026}
\maketitle
\begin{abstract}
We study exact differential equations and polynomial recurrences for
ensemble traces, Bessel determinants and moments, together with arithmetic
classifications and certified mixture estimation. The analytic method
combines vertex visibility, a Wronskian balance identity and Laurent
operator division. For circular ensembles it gives minimal differential
order $n+1$ for every $\beta>0$ and nonnegative integer twist, and exact
moment-recurrence invariants. Jacobi traces, orthogonal-group moments,
even-parameter vertex integrals, decaying Bessel moments and planar-walk
moments are treated with explicit parameter restrictions.

At every positive even Dyson parameter, the normalized circular integral
generates the same irreducible Bessel symmetric power. An invertible
Laurent jet gauge identifies its function and arithmetic-value fields
across parameters. Its unipotent block gives minimum recurrence degree
$n-1$ for every integer twist, including exceptional parameters missed
by exponent divisibility. A distinct three-particle family at
$\beta=2q-1$ has degree one; minimum-order degrees are classified for
all three-particle parameters. Jacobi degree bounds are extended through
zero indicial parameters and beyond a positivity condition, with a full
one- and two-particle classification. Exponential-tail moment uniqueness
also proves optimality on every real-shifted parity lattice for decaying
Bessel products.

The arithmetic part determines algebraicity under sampling and
quasipolynomial weighting, multipoint jet relations and field laws.
The statistical part gives known-support binomial and Poisson recovery
rates, entropy--transport bounds and finite-grid likelihood certificates.
Weighted Franel recurrences, Ap\'ery limits and split-beta orthogonality
provide further applications. Construction, minimality, arithmetic
specialization and computational approximation are distinguished
throughout; noninteger-twist minimality and residual resonant degree
problems remain explicitly open.
\end{abstract}
\noindent\textbf{Status.} Research manuscript prepared with AI assistance
and internal mathematical checks. Statements are subject to their
explicit hypotheses. External peer review and a comprehensive
historical-priority assessment have not been completed.
\clearpage
\begingroup\small\setlength{\parskip}{0pt}\tableofcontents\endgroup
\clearpage

\section{Scope and organization}\label{fusion:introduction}
This volume has four parts. Part I develops scalar equations and minimal
recurrences, beginning with determinants and then extending the method
to positive-parameter ensembles, orthogonal groups and Bessel moments.
Part II develops differential Galois theory, local monodromy,
algebraicity and arithmetic specialization. Part III treats
entropy--transport bounds and certified mixture estimation. Part IV
treats weighted Franel sequences and arithmetic classifications.
The analytical and statistical strands use different invariants:
exponential factors and local monodromy in the former, posterior
relative entropy and likelihood scores in the latter.

\subsection*{General-parameter recurrences and their refinements}
The six ensemble and moment theorems are stated together in
Section~\ref{V:overview}. Their proofs establish visibility of every
relevant vertex, absence of nonzero apparent singularities, and the
transfer of indicial factors to recurrence coefficients. The following
refinements use this structure together with the arithmetic part.
\begin{enumerate}[leftmargin=*]
\item \textbf{Even-parameter modules and arithmetic fields.}
Theorems~\ref{evb:module}--\ref{evb:arithmetic} identify the circular
integral at every $\beta=2m$, $m\ge1$, and integer $q\ge0$ with the
same Bessel symmetric power. The complete initial jets are related by
invertible Laurent matrices, with an explicit determinant. Their
function fields and numerical jet fields at nonzero algebraic arguments
are independent of $m,q$ at fixed dimension.
\item \textbf{Exceptional twist degrees.}
Theorem~\ref{vtd:even-degree} proves minimum coefficient degree $n-1$
for every positive even $\beta$ and every integer twist, among eventual
recurrences of arbitrary order. Theorem~\ref{vtd:odd-degree-drop}
proves an infinite family with minimum degree one at $n=3$,
$\beta=2q-1$. Theorem~\ref{vtd:three-particle-classification}
classifies degrees at minimum order for every three-particle parameter.
The latter is not a claim about every higher-order recurrence at the
remaining non-even exceptional parameters.
\item \textbf{Jacobi resonance strata.}
Theorem~\ref{V:thm:jacobi-resonances} replaces the positivity condition
by an explicit finite resonance set, gives a multiplicity-sensitive
degree lower bound everywhere, and proves the exact first-resonance
drop. Corollary~\ref{V:cor:jacobi-small-rank} classifies the one- and
two-particle cases, including an explicit exceptional density equation.
\item \textbf{Shifted Bessel moment lattices.}
Theorem~\ref{V:new:mellin-parity} proves minimum width
$\lceil n/2\rceil$ and minimum degree $n+1$ for
$\int_0^\infty t^{2N+r}f(t)\,dt$, every real $r>-1$, and every nonzero
decaying linear combination of degree-$n$ Bessel products. Its
leading-factor conclusion holds for arbitrary eventual recurrences.
\item \textbf{A partial result for real twists.}
Proposition~\ref{V:prop:real-twist-compact} constructs a compact
distribution transform for every real $q\ge0$ and locates its possible
singularities. Nonvanishing at all candidate vertices, and hence
minimality for general noninteger $q$, is not concluded.
\end{enumerate}

\subsection*{Arithmetic and statistical results}
Theorem~\ref{fusion:observable-dichotomy} gives the algebraicity
dichotomy for holomorphic fixed-parity homogeneous Laurent Bessel
observables, including arithmetic sampling and quasipolynomial
weighting. Theorems~\ref{mpj:product}--\ref{mpj:alljets} determine
multipoint jet relations, and Theorem~\ref{mpj:fields} gives the
componentwise field laws. In statistics,
Theorems~\ref{fg:binomial} and~\ref{fg:poisson} and
Corollary~\ref{fg:rates} give deterministic finite grids that retain
the stated mixing-law recovery rates; Theorem~\ref{fg:gaussian}
gives the corresponding Gaussian likelihood guarantee.
Classical antecedents are cited where used. An application of a
classical theorem is not counted as an independent new method.

\subsection*{Coverage and parameter restrictions}
\begin{center}\small
\renewcommand{\arraystretch}{1.2}
\begin{tabular}{@{}>{\raggedright\arraybackslash}p{0.26\textwidth}>{\raggedright\arraybackslash}p{0.67\textwidth}@{}}
\toprule
Family & Conclusions and restrictions\\
\midrule
Circular ensembles & Shape for $\beta>0$, real $q>-1/2$; minimal
differential order for integer $q\ge0$. Untwisted recurrence width
and degree for every $\beta>0$ (Theorems~\ref{V:thm:A}, \ref{V:thm:B}).\\
Twisted circular moments & Width for every $\beta>0$ and integer
$q\ge0$; degree outside the stated finite exceptional set, and at
every positive even $\beta$ (Theorem~\ref{vtd:even-degree}).\\
Jacobi traces & Minimal transform order for every $a,b>-1$, $\beta>0$.
Degree and density refinements according to the resonance set
(Theorem~\ref{V:thm:jacobi-resonances}).\\
Orthogonal groups & Minimal moment recurrences for types $B,D$;
the type $B$ path interpretation restricts zero steps at the boundary
(Theorem~\ref{V:thm:D}).\\
Even-parameter vertex integrals & Exact sorted origin orders and
explicit solution basis (Theorem~\ref{V:thm:E}).\\
Bessel and planar-walk moments & Minimal widths and degrees, with
leading-factor divisibility (Theorem~\ref{V:thm:F}).\\
Weighted Franel sequences & Minimal widths, including arithmetic
subsequences; no general minimum-degree classification
(Theorem~\ref{E:ac:subsequence}).\\
\bottomrule
\end{tabular}
\end{center}

\subsection*{Notation}
A recurrence is eventual unless its index range is stated. Its width
is the difference between extreme nonzero shifts; its degree is the
maximum degree of its polynomial coefficients. The two minima are
distinct optimization questions. Symbols are local to their sections:
in statistics $N$ is sample size, $t$ replication, and $M$ a known
Poisson support endpoint. Algebraicity of a generating function means
algebraicity of its germ over $\mathbb C(z)$; it is distinct from
algebraicity of a numerical value.


\clearpage\part{Vertex balance, scalar equations, and minimal recurrences}

\section{The three combinatorial families and principal results}\label{B:families:main}


\subsection{Three sequences}
For $\ell\ge1$ and $n\ge0$ let
\begin{itemize}
\item $u_\ell(n)$ be the number of permutations of $[n]=\{1,\dots,n\}$ whose longest increasing subsequence has length at most $\ell$;
\item $y_k(n)$ be the number of standard Young tableaux with $n$ cells and at most $k$ rows; by the Robinson--Schensted correspondence, $y_k(n)$ is also the number of involutions of $[n]$ with no decreasing subsequence of length $k+1$;
\item $m_\ell(n)$ be the number of perfect matchings of $[2n]$ with no $(\ell+1)$-crossing (equivalently, with no $(\ell+1)$-nesting) \cite{CDDSY07}.
\end{itemize}
By the Robinson--Schensted correspondence, $u_\ell(n)=\sum_{\lambda}(f^\lambda)^2$ over partitions $\lambda\vdash n$ with at most $\ell$ rows, where $f^\lambda$ is the number of standard tableaux of shape $\lambda$. In random matrix language,
\begin{equation}\label{B:eq:moments}
u_\ell(n)=\E_{U(\ell)}|\tr U|^{2n},\qquad m_\ell(n)=\E_{\Sp(2\ell)}(\tr U)^{2n},
\end{equation}
with Haar measure on the unitary and the compact symplectic group. The first identity is due to Gessel and Rains \cite{Ges90,Rai98}. The second follows from the classical identification of $(\ell+1)$-noncrossing perfect matchings with oscillating tableaux of height at most $\ell$, which count the $\Sp(2\ell)$-invariants in $(\C^{2\ell})^{\otimes 2n}$ \cite{Ber86,Sun90,CDDSY07}.

Write $b_m(t)=I_m(2t)=\sum_{j\ge0}t^{2j+|m|}/(j!\,(j+|m|)!)$ for the modified Bessel functions. The basic generating functions are Bessel determinants:
\begin{align}
\sum_{n\ge0}u_\ell(n)\frac{t^{2n}}{(n!)^2}&=D_\ell(t):=\det\bigl[b_{i-j}(t)\bigr]_{i,j=1}^{\ell}=\E_{U(\ell)}e^{2t\,\mathrm{Re}\,\tr U},\label{B:eq:Gessel}\\
Y_{2r}(t)&=\det\bigl[b_{i-j}+b_{i+j-1}\bigr]_{i,j=1}^{r},\qquad Y_{2r+1}(t)=e^{t}\det\bigl[b_{i-j}-b_{i+j}\bigr]_{i,j=1}^{r},\label{B:eq:Gordon}\\
\sum_{n\ge0}m_\ell(n)\frac{t^{2n}}{(2n)!}&=\E_{\Sp(2\ell)}e^{t\tr U}=\det\bigl[b_{i-j}-b_{i+j}\bigr]_{i,j=1}^{\ell}=e^{-t}\,Y_{2\ell+1}(t),\label{B:eq:Sp}
\end{align}
where $Y_k(t)=\sum_n y_k(n)t^n/n!$. Here \eqref{B:eq:Gessel} is Gessel's theorem \cite{Ges90}, \eqref{B:eq:Gordon} is due to Gordon and Gessel \cite{Ges90}, and \eqref{B:eq:Sp} follows from Weyl's integration formula for $\Sp(2\ell)$ and Andreief's identity (Lemma~\ref{B:lem:jacobi-int}), together with \eqref{B:eq:Gordon}. In particular all three sequences are P-recursive.

\subsection{Conjectures and claims}
Bergeron, Favreau and Krob \cite{BFK95} computed differential equations and recurrences for small heights and conjectured their general shape. For tableaux, the conjecture states that $Y_k$ satisfies an equation of order $\fl{k/2}+1$ with coefficients of degree at most $\fl{k/2}$, leading coefficient $t^{\fl{k/2}}$, and the coefficient of the $m$-th derivative divisible by $t^{m-1}$. It was checked for $k\le11$ by Bergeron and Gascon \cite{BG00}, listed as open in Mishna's survey \cite[Conj.~9]{Mis19}, and proved in the manuscript \cite{Kei26}. Theorem~\ref{B:thm:jacobi} below recalls that result, and Appendix~\ref{B:app:jacobi} contains a complete proof.

For pairs of tableaux, Bergeron and Gascon \cite{BG00} wrote down the equations satisfied by
\[
G_\ell(x):=D_\ell(\sqrt x)=\sum_{n\ge0}u_\ell(n)\frac{x^n}{(n!)^2}
\]
for $\ell\le7$, of orders $3,\dots,8$, each with leading term $x^\ell\,\partial_x^{\ell+1}$. They state that these "settle the corresponding conjectures for the cases $k\le7$". We formulate this as follows.

\begin{conjecture}[Bergeron--Favreau--Krob, Bergeron--Gascon; pairs]\label{B:conj:pairs}
For every $\ell\ge1$ there are polynomials $q_0,\dots,q_{\ell+1}$ of degree at most $\ell$ with $q_{\ell+1}=x^\ell$ and $x^{m-1}\mid q_m$ for $m\ge1$, such that $\sum_{m=0}^{\ell+1}q_m(x)G_\ell^{(m)}(x)=0$.
\end{conjecture}

Forrester and Wei \cite{FW26} recently derived scalar differential equations for unitary matrix integrals $\langle(\det U)^q e^{x^{1/2}\tr(U+U^\dagger)}\rangle_{U(\ell)}$ from a first-order matrix system. For $q=0$ the integral is $G_\ell(x)$ \cite[(4.6)]{FW26}. Their Proposition~4.1 is stated for a differential-difference system that covers these integrals and their $\beta$-analogues. For $q=0$ it asserts an equation $(x^\ell\partial^{\ell+1}+q_\ell(x)x^{\ell-1}\partial^\ell+\dots+q_0(x))G_\ell=0$ in which "the degree in $x$ of the coefficient of each $d^j/dx^j$ is $[(\ell-1+j)/2]$". The elimination procedure behind it is described rather than carried out in full. Appendix~B of \cite{FW26}, due to Bornemann, argues that the order $\ell+1$ is optimal, using Siegel's theorem on the algebraic independence of $I_0$ and $I_0'$.

For the sequence itself, Forrester and Wei note that the conjectured recurrence form of \cite{BFK95}, with $\fl{\ell/2}+1$ terms beyond the first, is not sharp: for $\ell=4$ a recurrence with only three terms holds. They state the corrected form as Proposition~4.3, obtained "(by inspection)": the sequence satisfies a recurrence of order $\fl{(\ell-1)/2}+1$. From the cases $\ell=2,\dots,5$ they observe "an obvious pattern" that the degree of the polynomial coefficients equals $\ell-1$, which they call "difficult to anticipate". Here we prove the recurrence structure and the corresponding order and degree lower bounds, also treating tableaux and matchings.

\subsection{Main results}
Throughout, $\theta$ denotes the Euler operator $x\,d/dx$ or $t\,d/dt$, and for an operator written in the form $\sum_k x^kP_k(\theta)$ we call the polynomials $P_k$ its \emph{$\theta$-coefficients}. We write $K=\cl{\ell/2}$, and for $0\le b\le \ell$
\begin{equation}\label{B:eq:gamma}
\lambda_b=2(\ell-2b),\qquad \gamma_b=\tfrac12\bigl((\ell-b)^2+b^2\bigr).
\end{equation}

\begin{theorem}[The Bessel--Toeplitz operator]\label{B:thm:A}
Let $\ell\ge1$.
\begin{enumerate}
\item[(a)] $G_\ell$ satisfies no linear differential equation of order at most $\ell$ over $\C(x)$. Its minimal operator, normalized to have coprime polynomial coefficients, is
\[
\begin{gathered}
L_\ell=\sum_{m=0}^{\ell+1}q_m(x)\,\partial_x^m,\qquad q_{\ell+1}=x^\ell,\\
x^{m-1}\mid q_m\ (1\le m\le\ell+1),\qquad
\deg q_m=\Bigl\lfloor\frac{\ell-1+m}{2}\Bigr\rfloor .
\end{gathered}
\]
Every operator of order $\ell+1$ with polynomial coefficients annihilating $G_\ell$ is a polynomial multiple of $L_\ell$.
\item[(b)] $L_\ell$ has no singular point in $\C\setminus\{0\}$. One has $xL_\ell=\sum_{k=0}^{K}x^kP_k(\theta)$ with $\deg P_k=\ell+1-2k$, where
\[
P_0(\theta)=\prod_{b=0}^{\ell}\bigl(\theta+b(\ell-b)\bigr),\qquad \sum_{k=0}^{K}\lc(P_k)\,\mu^{\ell+1-2k}=\prod_{b=0}^{\ell}(\mu-\ell+2b),
\]
and, if $\ell$ is even, $P_K(\theta)=\lc(P_K)\,(\theta+\ell^2/8)$. In particular the local exponents of $L_\ell$ at $x=0$ are the $\ell+1$ integers $-b(\ell-b)$, $0\le b\le \ell$.
\item[(c)] At $x=\infty$ the operator $L_\ell$ has a basis of formal solutions $e^{\lambda_b\sqrt x}\,x^{-\gamma_b/2}\bigl(1+O(x^{-1/2})\bigr)$, $0\le b\le\ell$.
\end{enumerate}
\end{theorem}

The same statements hold in the variable $t=\sqrt x$ for $D_\ell(t)$; see Theorem~\ref{B:thm:Dt}. There the minimal operator has leading coefficient $t^\ell$, coefficients of degree at most $\ell$, and exponents $-2b(\ell-b)$ at $t=0$.

\begin{corollary}\label{B:cor:FW41}
The shape asserted in Proposition~4.1 of \cite{FW26} holds for $q=0$, with the exact degrees of the coefficients; the order $\ell+1$ is minimal, which is Bornemann's claim; and Conjecture~\ref{B:conj:pairs} holds for every $\ell\ge1$. The integrality statement is completed constructively in Section~\ref{B:sec:integral-continuant}.
\end{corollary}

\begin{theorem}[Permutations with bounded increasing subsequences]\label{B:thm:B}
Let $\ell\ge1$, $K=\cl{\ell/2}$, and let $P_0,\dots,P_K$ be as in Theorem~\ref{B:thm:A}. Put
\[
R_0(n)=\prod_{b=1}^{\ell-1}\bigl(n+b(\ell-b)\bigr),\qquad R_k(n)=P_k(n-k)\,\bigl[(n-1)(n-2)\cdots(n-k+1)\bigr]^2\quad(1\le k\le K).
\]
\begin{enumerate}
\item[(a)] $\sum_{k=0}^{K}R_k(n)\,u_\ell(n-k)=0$ for all $n\ge1$, with $u_\ell(m)=0$ for $m<0$. Each $R_k$ has degree exactly $\ell-1$, and
\[
R_K(n)=(-1)^K\prod_{j=0}^{K-1}(\ell-2j)^2\cdot\bigl[(n-1)\cdots(n-K+1)\bigr]^2\cdot\begin{cases}1,&\ell\text{ odd},\\ n-K+\ell^2/8,&\ell\text{ even}.\end{cases}
\]
\item[(b)] The characteristic polynomial is $\sum_{k=0}^{K}\lc(R_k)\lambda^{K-k}=\prod_{j=0}^{K-1}\bigl(\lambda-(\ell-2j)^2\bigr)$.
\item[(c)] Every linear recurrence with polynomial coefficients satisfied by $u_\ell(n)$ for all large $n$ has order at least $K$ and coefficients of degree at least $\ell-1$. Its leading coefficient, written as the coefficient of $u_\ell(n)$ in a relation among $u_\ell(n),u_\ell(n-1),\dots$, is divisible by $R_0(n)$. Every such recurrence of order $K$ is a polynomial multiple of the one in (a).
\end{enumerate}
\end{theorem}

For $\ell=4$ the recurrence in (a) is
\[
(n+3)^2(n+4)\,u_4(n)-2(10n^3+31n^2+11n-12)\,u_4(n-1)+64n(n-1)^2\,u_4(n-2)=0,
\]
which is \cite[(4.7)]{FW26}. Table~\ref{B:tab:perm} lists the cases $\ell\le7$.

\begin{corollary}\label{B:cor:FW43}
Proposition~4.3 of \cite{FW26} holds. The minimal order of a recurrence for $u_\ell(n)$ is exactly $\fl{(\ell-1)/2}+1=\cl{\ell/2}$, the minimal degree is exactly $\ell-1$, and the primitive minimal recurrence is unique up to a nonzero constant factor. In particular a recurrence of the form conjectured in \cite{BFK95}, with $\fl{\ell/2}+1$ terms beyond the first, exists; its order is minimal exactly when $\ell$ is odd, and for even $\ell$ the minimal order is one less.
\end{corollary}

\subsection*{Twisted integrals}
For $q\in\Z$ put
\[
D_{\ell,q}(t)=\det\bigl[b_{i-j+q}(t)\bigr]_{i,j=1}^{\ell}=\E_{U(\ell)}(\det U)^q\,e^{2t\,\mathrm{Re}\tr U}.
\]
Since $D_{\ell,-q}=D_{\ell,q}$, we assume $q\ge0$. By Lemma~\ref{B:lem:twisted-expansion},
\[
G_{\ell,q}(x):=x^{-q\ell/2}D_{\ell,q}(\sqrt x)=\sum_{n\ge0}u_{\ell,q}(n)\frac{x^n}{n!\,(n+q\ell)!},\qquad u_{\ell,q}(n)=\sum_{\mu}f^{\mu}f^{\mu+(q^\ell)},
\]
where $\mu$ runs over the partitions of $n$ with at most $\ell$ parts and $\mu+(q^\ell)$ adds $q$ to each of the $\ell$ parts. Thus $G_{\ell,0}=G_\ell$. Up to a constant factor, $G_{\ell,q}$ is the hypergeometric function ${}_0F_1(q+\ell;(x)^\ell)$ with Jack parameter $1$ \cite[(2.4), (2.5)]{FW26}. For $\beta=2$ and $n=\ell$, the function in \cite[Prop.~4.1]{FW26} is a constant multiple of ${}_0F_1(2\ell+a+b;(x)^\ell)$ \cite[(3.10) and Rem.~3.4]{FW26}; the choice $a+b=q-\ell$, that is $B_p=p(\ell+q-p)$ in the system \cite[(3.6)]{FW26}, gives $G_{\ell,q}$. The case $q=\ell$ arises in the random matrix theory of joint moments of characteristic polynomials and their derivatives \cite{KW23}, \cite[\S4.3]{FW26}.

\begin{theorem}[Twisted integrals]\label{B:thm:twisted}
Let $\ell\ge1$, $q\ge0$ and $K=\cl{\ell/2}$.
\begin{enumerate}
\item[(a)] $G_{\ell,q}$ satisfies no linear differential equation of order at most $\ell$ over $\C(x)$. Its minimal operator $L_{\ell,q}=\sum_mq_m(x)\partial_x^m$, normalized to have coprime polynomial coefficients, has the shape of Theorem~\ref{B:thm:A}(a): $q_{\ell+1}=x^\ell$, $x^{m-1}\mid q_m$ and $\deg q_m=\fl{(\ell-1+m)/2}$. It has no singular point in $\C\setminus\{0\}$, and $xL_{\ell,q}=\sum_{k=0}^{K}x^kP_k(\theta)$ with $\deg P_k=\ell+1-2k$, $\sum_k\lc(P_k)\mu^{\ell+1-2k}=\prod_b(\mu-\ell+2b)$,
\[
P_0(\theta)=\prod_{i=0}^{\ell}\bigl(\theta+i(\ell+q-i)\bigr),
\]
and, for even $\ell$, $P_K(\theta)=\lc(P_K)(\theta+\ell^2/8+q\ell/2)$. At infinity $L_{\ell,q}$ has the formal solutions $e^{\lambda_b\sqrt x}x^{-(\gamma_b+q\ell)/2}(1+O(x^{-1/2}))$.
\item[(b)] The minimal operator of $D_{\ell,q}(t)$ has order $\ell+1$ and no singular point in $\C\setminus\{0\}$. Its exponents at $t=0$ are $q(\ell-2i)-2i(\ell-i)$, $0\le i\le\ell$. Its leading coefficient is $t^\ell$ if one of these exponents is $0$, and $t^{\ell+1}$ otherwise.
\item[(c)] $u_{\ell,q}$ satisfies a recurrence of order $K$ whose coefficients have degree $\ell-1$, with leading coefficient $\prod_{i=1}^{\ell-1}(n+i(\ell+q-i))$, and it satisfies no recurrence of smaller order.
\end{enumerate}
\end{theorem}

\begin{corollary}\label{B:cor:FW41q}
For $\beta=2$ and every $q\in\Z_{\ge0}$, the shape asserted in Proposition~4.1 of \cite{FW26} holds for ${}_0F_1(q+\ell;(x)^\ell)$, with the exact degrees, and the order $\ell+1$ is minimal. In particular this holds for $q=\ell$, and Bornemann's optimality claim \cite[App.~B]{FW26} also holds for the Hankel determinant $w_\ell$ treated there. Section~\ref{B:sec:integral-continuant} also establishes integrality in this exact normalization.
\end{corollary}

For tableaux let $r=\fl{k/2}$, $\varepsilon=k-2r$, $e_b=b^2-2br-\varepsilon b$ and $\lambda^{(k)}_j=\varepsilon+2r-4j$ for $0\le j\le r$.

\begin{theorem}[Tableaux of bounded height]\label{B:thm:C}
Let $k\ge2$ and put $R=r+1-[4\mid k]$, where $[4\mid k]$ is $1$ if $4$ divides $k$ and $0$ otherwise. There are polynomials $c_0,\dots,c_R$ of degree at most $r$ with
\[
\sum_{i=0}^{R}c_i(N)\,y_k(N-i)=0\quad(N\ge1),\qquad c_0(N)=\prod_{b=1}^{r}\bigl(N+b(k-b)\bigr),
\]
and $c_R(N)=\prod_{j=0}^{r}(-\lambda^{(k)}_j)\cdot(N-1)(N-2)\cdots(N-r)$ if $4\nmid k$, while $c_R(N)=\prod_{j\ne r/2}(-\lambda^{(k)}_j)\cdot(N-1)\cdots(N-r+1)\,(N-r+r^2/2)$ if $4\mid k$. The characteristic polynomial is $\prod_{j:\lambda^{(k)}_j\ne0}(\lambda-\lambda^{(k)}_j)$. Every recurrence satisfied by $y_k$ for large $N$ has order at least $R$ and degree at least $r$, and its leading coefficient is divisible by $c_0$. Every recurrence of order $R$ is a polynomial multiple of this one.
\end{theorem}

\begin{theorem}[Noncrossing matchings]\label{B:thm:D}
Let $\ell\ge1$ and $K=\cl{\ell/2}$. There are polynomials $d_0,\dots,d_K$ of degree exactly $\ell$ with
\[
\sum_{i=0}^{K}d_i(N)\,m_\ell(N-i)=0\quad(N\ge1),\qquad d_0(N)=\prod_{b=1}^{\ell}\bigl(2N+b(2\ell+1-b)\bigr),
\]
and $d_K(N)=\prod_{j=0}^{\ell}(-\lambda_j)\cdot(2N-1)(2N-2)\cdots(2N-\ell)$ if $\ell$ is odd, while $d_K(N)=\kappa_\ell(2N-1)\cdots(2N-\ell+1)\bigl(2N-\ell+\tfrac{\ell(\ell+1)}2\bigr)$ with $\kappa_\ell=\prod_{j\ne\ell/2}(-\lambda_j)$ if $\ell$ is even. The characteristic polynomial is $\prod_{j<\ell/2}(\lambda-4(\ell-2j)^2)$. Every recurrence satisfied by $m_\ell$ for large $N$ has order at least $K$ and degree at least $\ell$, and its leading coefficient is divisible by $d_0$. Every recurrence of order $K$ is a polynomial multiple of this one.
\end{theorem}

In Theorems~\ref{B:thm:C} and~\ref{B:thm:D} the coefficients are normalized as in Section~\ref{B:sec:CD}; only their ratios matter. Finally, the recurrences translate into differential equations for the underlying distributions.

\begin{theorem}[Distributions of traces]\label{B:thm:E}
Let $\ell\ge2$ and $K=\cl{\ell/2}$. The distribution of $|\tr U|^2$, $U$ Haar distributed in $U(\ell)$, has a density $\rho_\ell$ on $[0,\ell^2]$ that is annihilated, in the sense of distributions on $\R$, by
\[
\mathcal L_\ell=\sum_{j=0}^{K}w^j\,R_{K-j}(K-j-1-\theta),\qquad \theta=w\frac{d}{dw}.
\]
The operator $\mathcal L_\ell$ has order $\ell-1$ and is Fuchsian on $\mathbb P^1$. One has $\mathcal L_\ell=w^{K-1}\mathcal L^\circ_\ell$, where $\mathcal L^\circ_\ell$ has polynomial coefficients and leading coefficient $c\,w^{\fl{\ell/2}}\prod_{j=0}^{K-1}(w-(\ell-2j)^2)$ with $c\ne0$. No nonzero operator of order less than $\ell-1$ annihilates $\rho_\ell$ in the sense of distributions on $\R$, and among such annihilators of order $\ell-1$ the operator $\mathcal L_\ell$ is unique up to a factor in $\C(w)$. Its exponents at $w=0$ are $0,0,1,1,\dots,K-2,K-2$, together with $\ell^2/8-1$ when $\ell$ is even. Its exponents at $w=\infty$, in the coordinate $1/w$, are $1-b(\ell-b)$, $1\le b\le\ell-1$. Analogous statements hold for $\tr U$ on $\Sp(2\ell)$, with order $\ell$, and for tableaux, with order $\fl{k/2}$; see Section~\ref{B:sec:densities}.
\end{theorem}

\subsection{Method}
The proof of Theorem~\ref{B:thm:A} has four steps.

\emph{Step 1: the solution space.} We write $D_\ell(t)$ as a symmetric multilinear Toeplitz form in solutions of Bessel's equation. Since the differential Galois group of Bessel's equation is $\SL_2$, the solution space of the minimal operator is spanned by $\ell+1$ \emph{vertex integrals} $\Psi_b$. In $\Psi_b$, $b$ of the $\ell$ eigenvalue variables run over the negative real axis instead of the unit circle.

\emph{Step 2: infinity.} Laplace's method at the critical point with $b$ coordinates equal to $-1$ and $\ell-b$ equal to $1$ shows that $\Psi_b\sim C_be^{\lambda_bt}t^{-\gamma_b}$, with pairwise distinct rates $\lambda_b$.

\emph{Step 3: the origin.} Here the vertex integrals are genuinely hard to evaluate: pairs of them share the same leading exponent, and the leading terms cancel in complicated ways. We only prove the crude bound $\Psi_b=O(t^{-2\fl{b/2}(\ell-\fl{b/2})}|\log t|^b)$. It comes from expanding a mixed Toeplitz determinant and from an elementary inequality for the total displacement of a permutation.

\emph{Step 4: the squeeze.} The crude bound gives a lower bound for the sum of the local exponents at $0$. Abel's identity for the Wronskian, combined with the asymptotics at infinity, gives the corresponding upper bound, with a deficit equal to the total contribution of the apparent singularities. The two bounds coincide because of the identity
\[
\sum_{b=0}^{\ell}\gamma_b-2\sum_{b=0}^{\ell}b(\ell-b)=\binom{\ell+1}{2}.
\]
This forces both the absence of apparent singularities and the exact exponents. The shape of the operator then follows as in \cite{Kei26}. For the twisted determinants the crude bound involves the order statistics of $2i(\ell-i)-q(\ell-2i)$, whose sum does not depend on $q$, so the same squeeze applies.

The passage to recurrences uses two simple principles, Lemmas~\ref{B:lem:types} and~\ref{B:lem:filtration}. First, every recurrence of order $M$ yields an annihilating operator whose \emph{span} in the variable is $M$, and such an operator can carry at most $M$ (for $t$: $2M$) solutions with distinct nonzero exponential types. Second, the lowest $\theta$-coefficient of that operator is essentially the leading coefficient of the recurrence. It must vanish at every negative exponent of the minimal operator at the origin, with multiplicity. The exponents $-b(\ell-b)$ thus become the factor $R_0$ of every recurrence for $u_\ell$, and the $\ell-1$ negative exponents force the degree bound $\ell-1$.

The same scheme applies verbatim to the Hankel (orthogonal and symplectic) Bessel determinants \eqref{B:eq:Gordon}--\eqref{B:eq:Sp}, whose minimal operators were determined in \cite{Kei26}. Table~\ref{B:tab:summary} summarizes the data that drive all the results.

\begin{table}[ht]
\centering
\small
\renewcommand{\arraystretch}{1.3}
\begin{tabular}{@{}lccc@{}}
\toprule
 & $D_\ell(t)$ & $Y_k(t)$ & $e^{-t}Y_{2\ell+1}(t)$\\
 & (unitary) & (tableaux) & (symplectic)\\
\midrule
order of minimal operator & $\ell+1$ & $r+1$ & $\ell+1$\\
rates at $\infty$ & $2(\ell-2b)$ & $\varepsilon+2r-4j$ & $2(\ell-2j)$\\
exponents $\gamma$ at $\infty$ ($e^{\lambda t}t^{-\gamma}$) & $\gamma_b$, see \eqref{B:eq:gamma} & $\gamma^{(k)}_j$, see \eqref{B:eq:jacobi-data} & $\gamma^{(2\ell+1)}_j$\\
exponents at $0$ & $-2b(\ell-b)$ & $-b(k-b)$ & $-b(2\ell+1-b)$\\
minimal recurrence: order & $\cl{\ell/2}$ & $r+1-[4\mid k]$ & $\cl{\ell/2}$\\
minimal recurrence: degree & $\ell-1$ & $r$ & $\ell$\\
\bottomrule
\end{tabular}
\medskip
\caption{Local data of the three Bessel determinants, with $k=2r+\varepsilon$. The indices run over $0\le b\le\ell$ for $D_\ell$, over $0\le b,j\le r$ for $Y_k$, and over $0\le b,j\le\ell$ in the symplectic case. For $Y_k$ and the symplectic case the exponents at $0$ are distinct; for $D_\ell$ the values for $b$ and $\ell-b$ coincide.}\label{B:tab:summary}
\end{table}

\subsection{Scope and prior work}
Table~\ref{B:tab:closed} lists the statements from the literature that are settled here.

\begin{table}[ht]
\centering
\small
\renewcommand{\arraystretch}{1.25}
\begin{tabular}{@{}>{\raggedright\arraybackslash}p{0.36\textwidth}>{\raggedright\arraybackslash}p{0.14\textwidth}>{\raggedright\arraybackslash}p{0.25\textwidth}>{\raggedright\arraybackslash}p{0.15\textwidth}@{}}
\toprule
Statement & Source & Previous status & Here\\
\midrule
Shape of the equation for $G_\ell$ (pairs conjecture) & \cite{BFK95,BG00} & equations computed for $\ell\le7$ & Cor.~\ref{B:cor:FW41}\\
Shape and degrees $\fl{(\ell-1+j)/2}$, $q=0$ & \cite[Prop.~4.1]{FW26} & proof sketched & Thm.~\ref{B:thm:A}\\
The same for $\beta=2$ and all $q\in\Z_{\ge0}$ & \cite[Prop.~4.1]{FW26} & proof sketched; $q=\ell$: $\ell\le5$ listed & Thm.~\ref{B:thm:twisted}\\
Order $\ell+1$ is optimal ($q=0$ and $q=\ell$) & \cite[App.~B]{FW26} & argued & Cor.~\ref{B:cor:FW41}, \ref{B:cor:FW41q}\\
Recurrence of order $\fl{(\ell-1)/2}+1$ for $u_\ell$ & \cite[Prop.~4.3]{FW26} & ``by inspection'' & Thm.~\ref{B:thm:B}\\
Degree $\ell-1$ of that recurrence & \cite{FW26} & observed for $\ell\le5$ & Thm.~\ref{B:thm:B}\\
Recurrence form with $\fl{\ell/2}+1$ terms & \cite{BFK95} & conjectured & Cor.~\ref{B:cor:FW43}\\
Minimal recurrences for $y_k$ and $m_\ell$ & & not determined & Thms.~\ref{B:thm:C}, \ref{B:thm:D}\\
\bottomrule
\end{tabular}
\medskip
\caption{Statements established here and their relation to the cited literature. The table records mathematical coverage, not a comprehensive priority claim.}\label{B:tab:closed}
\end{table}

The existence of differential equations of order $\ell+1$ for $G_\ell$ and of the recurrences of order $\cl{\ell/2}$ was known in the sense of \cite{BG00,FW26}. The results proved here include:
\begin{itemize}
\item a complete proof of the shape of the minimal operator, including the exact coefficient degrees, the exponents, and the absence of apparent singularities;
\item minimality of the order through an explicit differential-Galois argument, complementing the Bessel-module construction of \cite[Appendix~B]{FW26};
\item the recurrence of order $\cl{\ell/2}$, with explicit extreme coefficients, minimality in order and degree, and uniqueness;
\item the same shape and minimality for the twisted integrals with $(\det U)^q$, $q\in\Z$. For $q=\ell$, Keating and Wei derived an equation of order $\ell+1$ \cite{KW23}, and \cite[Prop.~4.6]{FW26} lists the equations for $\ell\le5$.
\end{itemize}
For tableaux, the minimal operators are taken from \cite{Kei26}. The equations of the right shape also follow from the Forrester--Kumar system for Jacobi traces \cite{FK22}, as explained in \cite{Kei26}. The recurrence results for tableaux and matchings include minimality and explicit extreme coefficients; recurrences for matchings with small crossing number were computed in \cite{CQRZ08}. For the $\ell$-step planar random walk, whose generating function is $I_0(2t)^\ell$, the recurrence of order $\cl{\ell/2}$ and degree $\ell-1$ was given in general by Verrill \cite{Ver04}; see \cite{BNSW11}. Remark~\ref{B:rem:walks} proves order minimality.
Theorem~\ref{rev:verrill-min-degree} proves degree minimality by combining
a local splitting lemma with the iterated Borel bound; the numerical
exponents at zero alone do not supply this conclusion.

For every $\beta>0$ and real $q\ge0$, Corollary~\ref{B:cor:ic-beta} constructs the Jack-family scalar equation. Theorem~\ref{V:thm:A} extends the coefficient-degree statement to $q>-1/2$ and proves minimality for nonnegative integer twists. Theorem~\ref{V:thm:C} treats general Jacobi parameters. The even-parameter Galois and arithmetic refinements are in Theorems~\ref{evb:module}--\ref{evb:arithmetic}; noninteger-twist minimality remains open here.

\subsection*{Organization}
Section~\ref{B:sec:ops} collects facts on linear differential operators, including the filtration and squeeze lemmas. Section~\ref{B:sec:dict} records the dictionaries between recurrences and operators. Section~\ref{B:sec:unitary} proves Theorem~\ref{B:thm:A}, Section~\ref{B:sec:twisted} extends it to the twisted determinants (Theorem~\ref{B:thm:twisted}), and Section~\ref{B:sec:perm} proves Theorem~\ref{B:thm:B}. Section~\ref{B:sec:CD} recalls the Hankel case and proves Theorems~\ref{B:thm:C} and~\ref{B:thm:D}. Section~\ref{B:sec:densities} treats the distributions (Theorem~\ref{B:thm:E}) and explains the transport principle behind them. Appendix~\ref{B:app:comp} describes the computer checks.

\section{Linear differential operators}\label{B:sec:ops}

Let $\partial=d/dt$ and $\theta=t\partial$, and write $\ff{x}{m}=x(x-1)\cdots(x-m+1)$ for the falling factorial, so that $t^m\partial^m=\ff{\theta}{m}$. An operator $L=\sum_{m=0}^np_m(t)\partial^m$ with $p_m\in\C[t]$ and $p_n\ne0$ is \emph{primitive} if the $p_m$ have no common zero. For a primitive operator, a point $\tau\in\C$ is singular if $p_n(\tau)=0$; it is \emph{apparent} if all local solutions are holomorphic at $\tau$.

\subsection{\texorpdfstring{$\theta$}{theta}-forms}
After left multiplication by a power of $t$, every operator with polynomial coefficients becomes a finite sum $\sum_k t^kQ_k(\theta)$ with $Q_k\in\C[\theta]$. Since $t^k\theta^i=\sum_{m\le i}S(i,m)\,t^{k+m}\partial^m$, where $S(i,m)$ are Stirling numbers of the second kind, the $\partial$-order of such a sum is $\max_k\deg Q_k$.

\begin{definition}
A \emph{$\theta$-form} is an operator $A=\sum_{k=k_0}^{k_1}t^kQ_k(\theta)$ with $Q_k\in\C[\theta]$, $k_0\le k_1$ integers and $Q_{k_0}Q_{k_1}\ne0$. Its \emph{span} is $k_1-k_0$. Its \emph{lowest coefficient} is $Q_{k_0}$ and its \emph{top coefficient} is $Q_{k_1}$.
\end{definition}

A $\theta$-form is Fuchsian at $0$ if and only if $\deg Q_{k_0}=\max_k\deg Q_k$. In that case $Q_{k_0}$ is, up to a constant factor, the indicial polynomial at $0$; this is Fuchs' criterion \cite{Inc56}. Similarly, $A$ is Fuchsian at $\infty$ if and only if $\deg Q_{k_1}=\max_k\deg Q_k$.

\subsection{Expansions at the origin and the filtration lemma}
Let $\mathcal N$ be the space of functions $f$ on a slit punctured disc $\{0<|t|<\delta\}\setminus(-\delta,0]$ that can be written as
\begin{equation}\label{B:eq:nilsson}
f=\sum_{\nu\ge\nu_0,\ \nu\in\Z}t^\nu P_\nu(\log t),
\end{equation}
with polynomials $P_\nu$ of bounded degree and the series converging absolutely on the slit disc. Equivalently, $f$ is a finite sum of products of a power of $\log t$ and a convergent Laurent series with finitely many negative powers. For $0\ne f\in\mathcal N$ let $\ord f=\min\{\nu:P_\nu\ne0\}$; the term $t^{\ord f}P_{\ord f}(\log t)$ is the \emph{leading part}. If $|f(t)|\le Ct^{-a}|\log t|^c$ for small $t>0$, then $\ord f\ge -a$, because the leading part dominates all other terms as $t\to0^+$.

\begin{lemma}[filtration]\label{B:lem:filtration}
Let $A=\sum_{k=k_0}^{k_1}t^kQ_k(\theta)$ be a $\theta$-form, and let $V\subset\mathcal N$ be a finite-dimensional space of solutions of $Af=0$. For $\nu\in\Z$ put $V_{\ge\nu}=\{f\in V:f=0\text{ or }\ord f\ge\nu\}$. Then
\[
\dim V_{\ge\nu}/V_{\ge\nu+1}\le \operatorname{mult}_\nu(Q_{k_0}),
\]
the multiplicity of $\nu$ as a root of $Q_{k_0}$, which is $0$ if $Q_{k_0}(\nu)\ne0$. Consequently $\dim V_{\ge\nu}$ is at most the number of roots of $Q_{k_0}$ in $\Z_{\ge\nu}$, counted with multiplicity.
\end{lemma}

\begin{proof}
Write $\Lambda=\log t$ and $D=d/d\Lambda$ acting on polynomials in $\Lambda$. Then $\theta(t^\nu P(\Lambda))=t^\nu(\nu+D)P(\Lambda)$. Let $f\in V_{\ge\nu}$ with expansion \eqref{B:eq:nilsson}. In $Af$, the terms $t^kQ_k(\theta)f$ with $k>k_0$ involve only powers $t^{\mu}$ with $\mu\ge\nu+k>\nu+k_0$. Hence the coefficient of $t^{\nu+k_0}$ in $Af=0$ gives $Q_{k_0}(\nu+D)P_\nu=0$. Let $m=\operatorname{mult}_\nu(Q_{k_0})$. Taylor expansion gives $Q_{k_0}(\nu+D)=D^m\,(c+DH(D))$ with $c\ne0$ and a polynomial $H$. The factor $c+DH(D)$ is invertible on polynomials, because $D$ is nilpotent there. Therefore $D^mP_\nu=0$, that is, $\deg P_\nu<m$. The linear map $V_{\ge\nu}\to\C[\Lambda]_{<m}$, $f\mapsto P_\nu$, has kernel $V_{\ge\nu+1}$, which proves the first claim. For the second, note that $V_{\ge\mu}=0$ for large $\mu$, because a finite-dimensional space of nonzero elements of $\mathcal N$ has bounded orders. Summing the first claim over $\mu\ge\nu$ gives the bound.
\end{proof}

\begin{lemma}[squeeze]\label{B:lem:squeeze}
Let $L$ be an operator of order $n$ that is Fuchsian at $0$, with indicial polynomial $I$, and suppose that its local solution space $V$ at $0$ is contained in $\mathcal N$. Then the roots of $I$ are integers, $\nu_1\ge\dots\ge\nu_n$. If $f_1,\dots,f_n$ is any basis of $V$ with $\ord f_1\ge\dots\ge\ord f_n$, then $\ord f_i\le\nu_i$ for all $i$. In particular
\[
\sum_{i=1}^n\ord f_i\le\sum_{i=1}^n\nu_i,
\]
with equality if and only if $\ord f_i=\nu_i$ for all $i$. In the case of equality, $\dim V_{\ge\nu}=\#\{i:\nu_i\ge\nu\}$ for every $\nu\in\Z$.
\end{lemma}

\begin{proof}
Apply Lemma~\ref{B:lem:filtration} to the $\theta$-form of $t^{n-\ord_0p_n}L$, whose lowest coefficient is $I$. The span of $\{f_i:\ord f_i\ge\nu\}$ lies in $V_{\ge\nu}$, so
\[
\#\{i:\ord f_i\ge\nu\}\le\dim V_{\ge\nu}\le\#\{\text{roots of $I$ in $\Z_{\ge\nu}$}\}.
\]
Taking $\nu\le\min_i\ord f_i$ shows that $I$, which has degree $n$, has $n$ integer roots. For a given $i$, choosing $\nu=\ord f_i$ gives $i\le\#\{j:\nu_j\ge\ord f_i\}$, hence $\nu_i\ge\ord f_i$. The remaining assertions follow, since equality of the sums forces $\ord f_i=\nu_i$, and then the two outer quantities in the display coincide.
\end{proof}

\subsection{Abel's identity and the top degree}
The next two lemmas are standard; we recall short proofs from \cite{Kei26}.

\begin{lemma}[Abel]\label{B:lem:abel}
Let $L=\sum_{m=0}^np_m\partial^m$ be primitive, and let $W$ be the Wronskian of a basis of solutions on a simply connected domain. Then $W'/W=-p_{n-1}/p_n$, and:
\begin{enumerate}
\item[(a)] if $\tau$ is a regular singular point with exponents $\rho_1,\dots,\rho_n$, the residue of $W'/W$ at $\tau$ is $\sum_i\rho_i-\binom n2$;
\item[(b)] if $p_n(\tau)=0$ and all solutions are holomorphic at $\tau$, the residue of $W'/W$ at $\tau$ is a positive integer.
\end{enumerate}
\end{lemma}

\begin{proof}
The identity is Abel's. For (a), write $(t-\tau)^{n-A}L=\sum_m s_m(t)(t-\tau)^m\partial^m$ with $A=\ord_\tau p_n$ and $s_n(\tau)\ne0$. The indicial polynomial is $\sum_ms_m(\tau)\ff{\rho}{m}=s_n(\tau)\bigl(\rho^n-(\binom n2-s_{n-1}(\tau)/s_n(\tau))\rho^{n-1}+\cdots\bigr)$, and the residue of $-p_{n-1}/p_n$ is $-s_{n-1}(\tau)/s_n(\tau)$. For (b), $W$ is holomorphic at $\tau$, so the residue is $\ord_\tau W\ge0$. If $W(\tau)\ne0$, Cramer's rule applied to the linear system that the coefficients $p_m/p_n$ satisfy shows that all $p_m/p_n$ are holomorphic at $\tau$. Then $(t-\tau)$ divides every $p_m$, contradicting primitivity.
\end{proof}

\begin{lemma}[top degree]\label{B:lem:top}
Let $L=\sum_{m=0}^np_m\partial^m$, $d=\max_m\deg p_m$, and $\chi(\lambda)=\sum_m[t^d]p_m\,\lambda^m$. Suppose there are solutions $y_1,\dots,y_n$, pairwise distinct $\lambda_1,\dots,\lambda_n$, constants $C_j\ne0$ and $\gamma_j$ with
\[
y_j^{(m)}(t)=C_je^{\lambda_jt}t^{-\gamma_j}\bigl(\lambda_j^m+O(t^{-1/2})\bigr)\qquad(t\to+\infty,\ 0\le m\le n).
\]
Then $\chi(\lambda)=[t^d]p_n\cdot\prod_j(\lambda-\lambda_j)$; in particular $\deg p_n=d$.
\end{lemma}

\begin{proof}
We have $0=C_j^{-1}e^{-\lambda_jt}t^{\gamma_j}Ly_j=t^d\chi(\lambda_j)+O(t^{d-1/2})$. So $\chi$ is a nonzero polynomial of degree at most $n$ with the $n$ distinct roots $\lambda_j$. Hence $\deg\chi=n$ and $[t^d]p_n\ne0$.
\end{proof}

\subsection{Exponential types bound the span}
The following lemma converts exponential types at infinity into lower bounds for the order of recurrences.

\begin{lemma}[types]\label{B:lem:types}
Let $A=\sum_{k=k_0}^{k_1}t^kQ_k(\theta)$ be a $\theta$-form. Suppose $y_1,\dots,y_N$ are solutions of $Ay=0$ on an interval $(T,\infty)$ such that, for some $C_j\ne0$ and $\gamma_j\in\R$,
\[
y_j^{(m)}(t)=C_je^{\lambda_jt}t^{-\gamma_j}\bigl(\lambda_j^m+O(t^{-1/2})\bigr)\qquad(t\to+\infty,\ 0\le m\le\max_k\deg Q_k),
\]
where $\lambda_1,\dots,\lambda_N$ are pairwise distinct and \emph{nonzero}. Then $N\le k_1-k_0$.
\end{lemma}

\begin{proof}
Write $d_k=\deg Q_k$ for $Q_k\ne0$. From $\theta^i=\sum_{m\le i}S(i,m)t^m\partial^m$ and $S(i,i)=1$ we get $\theta^iy_j=C_je^{\lambda_jt}t^{-\gamma_j}\bigl(\lambda_j^it^i+O(t^{i-1/2})\bigr)$. Hence
\[
t^kQ_k(\theta)y_j=C_je^{\lambda_jt}t^{-\gamma_j}\bigl(\lc(Q_k)\lambda_j^{d_k}t^{k+d_k}+O(t^{k+d_k-1/2})\bigr).
\]
Let $w=\max_k(k+d_k)$ and $S=\{k:Q_k\ne0,\ k+d_k=w\}$. Then $Ay_j=0$ gives $h(\lambda_j)=0$ for $h(\lambda)=\sum_{k\in S}\lc(Q_k)\lambda^{w-k}$. Now $h(\lambda)=\lambda^{w-\max S}g(\lambda)$, where $g(\lambda)=\sum_{k\in S}\lc(Q_k)\lambda^{\max S-k}$ has degree $\max S-\min S$ and $g(0)\ne0$. The $\lambda_j$ are $N$ distinct nonzero roots of $g$, so $N\le\max S-\min S\le k_1-k_0$.
\end{proof}

\subsection{Galois orbits}
Let $E$ be a Picard--Vessiot extension of $\C(t)$ with differential Galois group $G$.

\begin{lemma}\label{B:lem:galois-orbit}
For a D-finite element $f\in E$ over $\C(t)$, the $\C$-span $U$ of the orbit $\{\sigma(f):\sigma\in G\}$ is finite-dimensional, and it is the solution space in $E$ of the minimal monic operator $L_f\in\C(t)[\partial]$ annihilating $f$. In particular $\ord L_f=\dim U$.
\end{lemma}

\begin{proof}
Since $L_f$ has coefficients in $\C(t)$, it annihilates every $\sigma(f)$, so $U\subseteq\Sol_E(L_f)$ and $\dim U\le\ord L_f$. Conversely, $U$ is finite-dimensional and $G$-stable. The monic operator $L_U$ with solution space $U$, built from Wronskians of a basis, is therefore fixed by $G$, so its coefficients lie in $\C(t)$ by the Galois correspondence \cite{vdPS03}. Since $L_Uf=0$, minimality gives $\ord L_f\le\dim U$.
\end{proof}

\section{Recurrences and operators}\label{B:sec:dict}

We record the three dictionaries used later. Let $a:\Z_{\ge0}\to\C$ be a sequence and suppose
\begin{equation}\label{B:eq:rec}
\sum_{s=0}^{M}c_s(n)\,a(n+s)=0\qquad(n\ge n_0),
\end{equation}
with polynomials $c_s$ of degree at most $d$ and $c_0c_M\ne0$. We call $M$ the \emph{order}, $d$ the \emph{degree} and $c_M$ the \emph{leading coefficient} of the recurrence. In the backward normalization $\sum_{i=0}^{M}c_{M-i}(N-M)\,a(N-i)=0$, the leading coefficient is the coefficient of $a(N)$.

\begin{lemma}[exponential generating function]\label{B:lem:dict-egf}
Let $Y(t)=\sum_na(n)t^n/n!$ and assume \eqref{B:eq:rec}. Then $Y$ is annihilated by the $\theta$-form
\[
A=\sum_{s=0}^{M}t^{M-s}\,\ff{(\theta-s)}{n_0}\,c_s(\theta-s)\,\ff{\theta}{s},
\]
which has span $M$, $\theta$-degree at most $d+M+n_0$, lowest coefficient $\ff{\theta}{M+n_0}\,c_M(\theta-M)$ and top coefficient $\ff{\theta}{n_0}\,c_0(\theta)$.
\end{lemma}

\begin{proof}
Since $\partial^sY=\sum_na(n+s)t^n/n!$ and $\theta$ acts on $t^n/n!$ by multiplication with $n$, the operator $B=\sum_sc_s(\theta)\partial^s$ satisfies $BY=\sum_n\bigl(\sum_sc_s(n)a(n+s)\bigr)t^n/n!$. This is a polynomial of degree less than $n_0$, so it is killed by $\ff{\theta}{n_0}$. Using $\partial^s=t^{-s}\ff{\theta}{s}$ and $t^M\ff{\theta}{n_0}=\ff{(\theta-M)}{n_0}t^M$, we obtain $t^M\ff{\theta}{n_0}B=A$. Finally, $\ff{(\theta-M)}{n_0}\ff{\theta}{M}=\ff{\theta}{M+n_0}$.
\end{proof}

\begin{lemma}[double exponential generating function]\label{B:lem:dict-degf}
Let $G(x)=\sum_na(n)x^n/(n!)^2$, $\theta=x\,d/dx$, and assume \eqref{B:eq:rec}. Then $G$ is annihilated by
\[
O=\sum_{k=0}^{M}x^k\,\ff{(\theta+k)}{M+n_0}\,c_{M-k}(\theta+k-M)\,\bigl(\ff{\theta}{M-k}\bigr)^2,
\]
a $\theta$-form of span $M$ with lowest coefficient $\ff{\theta}{M+n_0}\,c_M(\theta-M)\,\bigl(\ff{\theta}{M}\bigr)^2$. Conversely, if $\sum_{k=0}^{M}x^kP_k(\theta)$ annihilates $G$, then
\begin{equation}\label{B:eq:degf-rec}
\sum_{k=0}^{M}P_k(N-k)\,\bigl(\ff{N}{k}\bigr)^2\,a(N-k)=0\qquad\text{for all }N\ge0,
\end{equation}
where terms with $N-k<0$ are omitted.
\end{lemma}

\begin{proof}
For the converse, the coefficient of $x^N$ in $x^kP_k(\theta)G$ is $P_k(N-k)a(N-k)/((N-k)!)^2$, and $(N!/(N-k)!)^2=(\ff{N}{k})^2$. For the direct statement, put $P_k(\theta)=c_{M-k}(\theta+k-M)(\ff{\theta}{M-k})^2$ and $O_0=\sum_kx^kP_k(\theta)$. For $N\ge M$, $\ff{(N-k)}{M-k}=(N-k)!/(N-M)!$, so the coefficient of $x^N$ in $O_0G$ equals
\[
\frac{1}{((N-M)!)^2}\sum_{k=0}^{M}c_{M-k}(N-M)\,a(N-k),
\]
which vanishes for $N-M\ge n_0$ by \eqref{B:eq:rec}. For $N<M$, every term contains a factor $\ff{(N-k)}{M-k}$ with $N-k\ge0>N-M$, which vanishes. Hence $O_0G$ is a polynomial of degree less than $M+n_0$ and is killed by $\ff{\theta}{M+n_0}$. Finally, $\ff{\theta}{M+n_0}x^k=x^k\ff{(\theta+k)}{M+n_0}$.
\end{proof}

\begin{lemma}[moments]\label{B:lem:dict-mom}
Let $\rho$ be a compactly supported distribution on $\R$ with moments $\mu_n=\langle\rho,w^n\rangle$. For polynomials $P_k$ the following are equivalent:
\begin{enumerate}
\item[(i)] $\sum_{k=0}^{M}w^kP_k(\theta)\rho=0$ in the sense of distributions, where $\theta=w\,d/dw$;
\item[(ii)] $\sum_{k=0}^{M}P_k(-n-k-1)\,\mu_{n+k}=0$ for all $n\ge0$.
\end{enumerate}
\end{lemma}

\begin{proof}
The formal adjoint of $\theta$ is $-\theta-1$. Hence $\langle w^kP_k(\theta)\rho,w^n\rangle=\langle\rho,P_k(-\theta-1)w^{n+k}\rangle=P_k(-n-k-1)\mu_{n+k}$. A compactly supported distribution whose moments all vanish is zero, since its Fourier transform is an entire function whose Taylor coefficients vanish.
\end{proof}

\section{The unitary Bessel--Toeplitz operator}\label{B:sec:unitary}

Fix $\ell\ge1$ and recall $K=\cl{\ell/2}$, $\lambda_b=2(\ell-2b)$ and $\gamma_b=\frac12((\ell-b)^2+b^2)$ from \eqref{B:eq:gamma}.

\subsection{Two contour sequences}
For $t>0$ and $m\in\Z$ put
\[
\begin{aligned}
i_m(t)&=\frac1{2\pi}\int_{-\pi}^{\pi}e^{2t\cos\varphi}e^{im\varphi}\,d\varphi=I_m(2t),\\
k_m(t)&=\int_0^\infty e^{-t(s+s^{-1})}(-s)^m\,\frac{ds}{s}=2(-1)^mK_m(2t).
\end{aligned}
\]
see \cite[10.32.3, 10.32.10]{DLMF}. Both are integrals $\int_\gamma e^{t(z+z^{-1})}z^m\,\frac{dz}{z}$, up to the factor $(2\pi i)^{-1}$ in the first case: $\gamma$ is the unit circle for $i_m$, and the negative real axis traversed from $0$ to $-\infty$ for $k_m$.

\begin{lemma}\label{B:lem:seq}
For $h\in\{i,k\}$ and $m\in\Z$ one has $h_{-m}=h_m$, $h_m'=h_{m-1}+h_{m+1}$ and $t(h_{m-1}-h_{m+1})=m\,h_m$. Consequently there are rational functions $A_m,B_m\in\Q(t)$, the same for $h=i$ and $h=k$, such that $h_m=A_mh_0+B_mh_0'$. Moreover $h_0$ satisfies Bessel's equation
\begin{equation}\label{B:eq:T}
Th:=t\,h''+h'-4t\,h=0,
\end{equation}
and $i_0,k_0$ form a basis of its solutions on $(0,\infty)$.
\end{lemma}

\begin{proof}
Differentiation under the integral sign multiplies the integrand by $z+z^{-1}$, which gives the second relation. Next, $\int_\gamma d\bigl(e^{t(z+z^{-1})}z^m\bigr)=0$: the circle is closed, and on the negative axis $e^{t(z+z^{-1})}$ vanishes to infinite order at both ends. This gives $t(h_{m+1}-h_{m-1})+mh_m=0$. The substitution $z\mapsto z^{-1}$ maps each contour to itself, reversing the orientation and changing the sign of $dz/z$, so $h_{-m}=h_m$. The two relations give $h_1=h_{-1}=h_0'/2$, and they express every $h_m$ through $h_0$ and its derivatives. They also give $h_1'=h_0+h_2$ and $h_0-h_2=h_1/t$, hence $h_0''=4h_0-h_0'/t$, which is \eqref{B:eq:T}; then all $h_m$ lie in $\Q(t)h_0+\Q(t)h_0'$. Finally, $I_0(2t)$ grows and $K_0(2t)$ decays as $t\to\infty$, so $i_0,k_0$ are independent.
\end{proof}

Equation \eqref{B:eq:T} is the modified Bessel equation of order $0$ in the variable $2t$. By Kolchin's theorem \cite{Kol68} (see also \cite{vdPS03}), its differential Galois group over $\C(t)$ acts on the solution space $\Sol(T)$ as the full group $\SL(\Sol(T))\cong\SL_2(\C)$. In particular it acts transitively on $\Sol(T)\setminus\{0\}$.

\subsection{The Toeplitz form}
For $g\in\Sol(T)$ let $g_\bullet=(g_m)_{m\in\Z}$, $g_m=A_mg+B_mg'$; thus $(i_0)_\bullet=i$ and $(k_0)_\bullet=k$. For sequences $h^{(1)},\dots,h^{(\ell)}$ of this form put
\[
F\bigl(h^{(1)},\dots,h^{(\ell)}\bigr)=\frac1{\ell!}\sum_{\sigma,\tau\in S_\ell}\sgn(\sigma)\sgn(\tau)\prod_{j=1}^{\ell}h^{(j)}_{\sigma(j)-\tau(j)} .
\]
This is $\C$-multilinear and symmetric, and $F(h,\dots,h)=\det[h_{p-p'}]_{p,p'=1}^{\ell}$. Indeed, substituting $j=\tau^{-1}(p')$ turns the double sum into $\ell!$ copies of the Leibniz expansion. For $g\in\Sol(T)$ put
\[
\Phi(g)=F(g_\bullet,\dots,g_\bullet)=\det\bigl[A_{p-p'}g+B_{p-p'}g'\bigr]_{p,p'=1}^{\ell};
\]
it is a homogeneous polynomial of degree $\ell$ in $g,g'$ with coefficients in $\Q(t)$. By \eqref{B:eq:Gessel}, $\Phi(i_0)=D_\ell$.

For $0\le b\le\ell$ define the \emph{vertex integrals}
\[
\Psi_b=F(\underbrace{k,\dots,k}_{b},\underbrace{i,\dots,i}_{\ell-b}).
\]

\begin{lemma}\label{B:lem:Psi-int}
Let $z_j=-s_j$ for $j\le b$ and $z_j=e^{i\varphi_j}$ for $j>b$, $f(z)=\sum_{j=1}^{\ell}(z_j+z_j^{-1})$ and $\Delta(z)=\prod_{p<p'}(z_{p'}-z_p)$. Then
\[
\Psi_b(t)=\frac1{\ell!}\int_{(0,\infty)^b\times(-\pi,\pi)^{\ell-b}}e^{tf(z)}\,\Delta(z)\Delta(z^{-1})\prod_{j\le b}\frac{ds_j}{s_j}\prod_{j>b}\frac{d\varphi_j}{2\pi},
\]
and $\Psi_b$ is real on $(0,\infty)$. Moreover, $\Phi(\alpha i_0+\beta k_0)=\sum_{b=0}^{\ell}\binom{\ell}{b}\alpha^{\ell-b}\beta^b\,\Psi_b$ for $\alpha,\beta\in\C$.
\end{lemma}

\begin{proof}
Expand $\Delta(z)\Delta(z^{-1})=\det[z_j^{p-1}]\det[z_j^{1-p}]=\sum_{\sigma,\tau}\sgn(\sigma\tau)\prod_jz_j^{\sigma(j)-\tau(j)}$ and apply Fubini to the integral representations of $i_m$ and $k_m$. Since the $i_m$ and $k_m$ are real and $F$ has integer coefficients, $\Psi_b$ is real. The last identity follows from multilinearity and symmetry of $F$, because $(\alpha i_0+\beta k_0)_\bullet=\alpha i+\beta k$.
\end{proof}

\begin{proposition}\label{B:prop:solspace}
Let $L^t_\ell$ be the minimal monic operator over $\C(t)$ annihilating $D_\ell(t)$. Then its solution space is
\[
\Sol(L^t_\ell)=V:=\spn_\C\{\Psi_0,\dots,\Psi_\ell\}.
\]
\end{proposition}

\begin{proof}
Let $E$ be the Picard--Vessiot extension of $\C(t)$ for \eqref{B:eq:T}, realized by functions on $(0,\infty)$, and $G$ its differential Galois group. Since $\Phi$ is a differential polynomial with coefficients in $\C(t)$, $D_\ell=\Phi(i_0)\in E$ and $\sigma(D_\ell)=\Phi(\sigma i_0)$ for $\sigma\in G$. By Kolchin's theorem $\{\sigma i_0\}=\Sol(T)\setminus\{0\}$. Hence the span $U$ of the orbit of $D_\ell$ is the span of $\{\Phi(g):g\in\Sol(T)\}$. By the polarization identity
\[
F(h^{(1)},\dots,h^{(\ell)})=\frac1{\ell!}\sum_{\emptyset\ne S\subseteq[\ell]}(-1)^{\ell-|S|}F\bigl(h_S,\dots,h_S\bigr),\qquad h_S=\sum_{j\in S}h^{(j)},
\]
$U$ contains every $F(g^{(1)}_\bullet,\dots,g^{(\ell)}_\bullet)$ with $g^{(j)}\in\Sol(T)$, in particular every $\Psi_b$. Conversely, expanding $\Phi(\alpha i_0+\beta k_0)$ by Lemma~\ref{B:lem:Psi-int} shows $U\subseteq V$. Now apply Lemma~\ref{B:lem:galois-orbit}.
\end{proof}

\subsection{Asymptotics at infinity}

\begin{lemma}[vertex asymptotics]\label{B:lem:vertex}
For $0\le b\le\ell$ and $m\ge0$, as $t\to+\infty$,
\[
\Psi_b^{(m)}(t)=C_b\,e^{\lambda_bt}t^{-\gamma_b}\bigl(\lambda_b^m+O(t^{-1})\bigr),\qquad
C_b=\frac{(-1)^{b(b-1)/2}\,4^{b(\ell-b)}}{\ell!\,(2\pi)^{\ell-b}}\;\mathcal M_{\ell-b}\,\mathcal M_b\ne0,
\]
where $\mathcal M_n=\int_{\R^n}e^{-|x|^2}\Delta(x)^2\,dx>0$ is Mehta's integral. More precisely, $\Psi_b$ has an asymptotic expansion $e^{\lambda_bt}t^{-\gamma_b}\sum_{i\ge0}c_{b,i}t^{-i}$ with $c_{b,0}=C_b$, and the expansions of the derivatives are the termwise derivatives.
\end{lemma}

\begin{proof}
On the domain of Lemma~\ref{B:lem:Psi-int} we have $z+z^{-1}=2\cos\varphi\le2$ on the circle and $z+z^{-1}=-(s+s^{-1})\le-2$ on the negative axis. So $f\le\lambda_b$, with equality only at the vertex $v^*$ where $s_j=1$ for $j\le b$ and $\varphi_j=0$ for $j>b$. Introduce the global coordinates
\[
x_j=2\sin(\varphi_j/2)\in(-2,2)\quad(j>b),\qquad y_j=s_j^{1/2}-s_j^{-1/2}\in\R\quad(j\le b).
\]
In these coordinates the phase is \emph{exactly} quadratic: $2\cos\varphi_j=2-x_j^2$ and $-(s_j+s_j^{-1})=-2-y_j^2$, so $f=\lambda_b-|x|^2-|y|^2$. Also $d\varphi_j=(1-x_j^2/4)^{-1/2}dx_j$ and $ds_j/s_j=2s_j^{1/2}(1+s_j)^{-1}dy_j$, with Jacobian factors equal to $1$ at $v^*$.

Next consider the factors of $\Delta(z)\Delta(z^{-1})=\prod_{p<p'}(2-z_p/z_{p'}-z_{p'}/z_p)$. A circle--circle factor equals $4\sin^2((\varphi_p-\varphi_{p'})/2)$. It is an analytic function of $(x_p,x_{p'})$ near $0$ that vanishes to second order on $x_p=x_{p'}$, so it equals $(x_p-x_{p'})^2$ times an analytic function with value $1$ at $0$. Similarly, a factor with both variables on the axis equals $-(s_p-s_{p'})^2/(s_ps_{p'})=-(y_p-y_{p'})^2\cdot(\text{analytic, }1\text{ at }0)$. A mixed factor equals $2+e^{i\varphi}/s+se^{-i\varphi}$, which is $4$ at $v^*$. Hence, near $v^*$,
\[
e^{tf}\Delta(z)\Delta(z^{-1})\,d(\cdots)=e^{\lambda_bt}e^{-t(|x|^2+|y|^2)}\,\Delta(x)^2\Delta(y)^2\,H(x,y)\,dx\,dy,
\]
with $H$ analytic and $H(0)=(-1)^{b(b-1)/2}4^{b(\ell-b)}(2\pi)^{b-\ell}/\ell!$.

Fix $\delta>0$ so small that $H$ is analytic on the ball $B_\delta=\{|w|\le\delta\}$, $w=(x,y)$. Outside $B_\delta$ we have $f\le\lambda_b-\delta^2$. We estimate this part in the original variables $(s,\varphi)$, where the integrand of Lemma~\ref{B:lem:Psi-int} is $e^{tf}$ times a function bounded by a polynomial in the $s_j^{\pm1}$. Writing $e^{tf}=e^{(t-1)f}e^{f}$ and using $\int e^{f}\prod_j(s_j+s_j^{-1})^N\prod_j\frac{ds_j}{s_j}<\infty$, we see that the part outside $B_\delta$ contributes $O(e^{(\lambda_b-\delta^2)t})$ for $t\ge1$. On $B_\delta$ write $H$ as its Taylor polynomial of degree $2J$ plus a remainder bounded by $C|w|^{2J+1}$; the remainder contributes $O(t^{-\gamma_b-J-1/2})$. A monomial $w^\alpha$ contributes
\[
\int_{\R^\ell}e^{-t|w|^2}\Delta(x)^2\Delta(y)^2w^\alpha\,dw+O(e^{-\delta^2t/2})=t^{-\gamma_b-|\alpha|/2}\int_{\R^\ell}e^{-|w|^2}\Delta(x)^2\Delta(y)^2w^\alpha\,dw+O(e^{-\delta^2t/2}).
\]
Here we used that $\Delta(x)^2\Delta(y)^2$ is homogeneous of degree $(\ell-b)(\ell-b-1)+b(b-1)=2\gamma_b-\ell$. The integral vanishes for $|\alpha|$ odd, by the symmetry $w\mapsto-w$. This gives the expansion in integral powers of $t^{-1}$, with leading coefficient $H(0)\mathcal M_{\ell-b}\mathcal M_b=C_b$.

For $\Psi_b^{(m)}$ the amplitude is multiplied by $f^m=(\lambda_b-|w|^2)^m$. Monomial by monomial, the resulting expansion is the termwise derivative, since $\int|w|^2e^{-t|w|^2}(\cdots)=-\frac{d}{dt}\int e^{-t|w|^2}(\cdots)$. Its leading term is $\lambda_b^mC_b$.
\end{proof}

\begin{corollary}\label{B:cor:dim}
The functions $\Psi_0,\dots,\Psi_\ell$ are linearly independent, so $\ord L^t_\ell=\dim V=\ell+1$.
\end{corollary}

\begin{proof}
The rates $\lambda_b$ are pairwise distinct and $C_b\ne0$.
\end{proof}

\subsection{The origin}\label{B:ssec:origin}
Every element of $V$ is a polynomial in $i_0,i_0',k_0,k_0'$ with coefficients in $\Q(t)$. Moreover $i_0$ is entire and $k_0=-2(\log t)\,i_0+(\text{entire})$ \cite[10.31.2]{DLMF}. Hence $V\subset\mathcal N$, every element of $V$ has moderate growth at $0$, and $0$ is a regular singular point of $L^t_\ell$ (Fuchs).

\begin{lemma}[displacement]\label{B:lem:disp}
\begin{enumerate}
\item[(a)] Let $S=\{s_1<\dots<s_j\}\subseteq[\ell]$ and let $\sigma$ be a permutation of $S$. Then $\sum_{s\in S}|\sigma(s)-s|\le2p(\ell-p)$ with $p=\fl{j/2}$.
\item[(b)] Let $\pi\in S_\ell$ and $S\subseteq[\ell]$ with $|S|=j$. Then
\[
\sum_{s\in S}|\pi(s)-s|-\sum_{s\notin S}|\pi(s)-s|\le 2p(\ell-p),\qquad p=\fl{j/2}.
\]
\end{enumerate}
\end{lemma}

\begin{proof}
(a) We have
\[
\sum_s|\sigma(s)-s|=\sum_s\max(s,\sigma(s))-\sum_s\min(s,\sigma(s)).
\]
The $2j$ numbers $\max(s,\sigma(s))$ and $\min(s,\sigma(s))$ form the multiset in which every element of $S$ occurs twice. So the difference is at most the sum of the $j$ largest elements of that multiset minus the sum of the $j$ smallest, which is $2\sum_{i=1}^{p}(s_{j+1-i}-s_i)$. Since $s_{j+1-i}\le\ell+1-i$ and $s_i\ge i$, this is at most $2\sum_{i=1}^p(\ell+1-2i)=2p(\ell-p)$.

(b) For $s\in S$, follow the cycle of $\pi$ through $s$ until it returns to $S$: $s\to\pi(s)\to\dots\to\pi^r(s)=:\sigma(s)\in S$, with $\pi^i(s)\notin S$ for $0<i<r$. This defines a permutation $\sigma$ of $S$. By the triangle inequality,
\[
|\pi(s)-s|\le|\sigma(s)-s|+\sum_{0<i<r}|\pi^{i+1}(s)-\pi^i(s)|.
\]
Every $s'\notin S$ occurs at most once as an intermediate point $\pi^i(s)$. Summing over $s\in S$ therefore gives $\sum_{s\in S}|\pi(s)-s|\le\sum_{s\in S}|\sigma(s)-s|+\sum_{s'\notin S}|\pi(s')-s'|$. Now apply (a).
\end{proof}

\begin{lemma}[origin bound]\label{B:lem:origin}
There is a constant $C$ such that for $0<t\le1$ and $0\le b\le\ell$
\[
|\Psi_b(t)|\le C\,t^{-2\fl{b/2}(\ell-\fl{b/2})}\,(1+|\log t|)^{b}.
\]
Consequently $\ord\Psi_b\ge-2\fl{b/2}(\ell-\fl{b/2})$.
\end{lemma}

\begin{proof}
On $(0,1]$ we have $|i_m(t)|\le e\,t^{|m|}$, since $|I_m(2t)|\le t^{|m|}e^{t^2}/|m|!$. We also have $|k_m(t)|\le C\,t^{-|m|}(1+|\log t|)$. Indeed, $k_{-m}=k_m$, and for $m>0$ splitting the integral at $s=1$ and using $e^{-t(s+s^{-1})}\le e^{-ts}$ for $s\ge1$ gives $|k_m(t)|\le 1/m+\Gamma(m)t^{-m}$. For $m=0$ the symmetry $s\mapsto s^{-1}$ gives $k_0(t)=2\int_1^\infty e^{-t(s+s^{-1})}ds/s\le2\int_1^\infty e^{-ts}ds/s$, which is $O(1+|\log t|)$. Let $T_I=[i_{p-p'}]$ and $T_K=[k_{p-p'}]$. By Lemma~\ref{B:lem:Psi-int}, $\binom{\ell}{b}\Psi_b$ is the coefficient of $\alpha^{\ell-b}\beta^b$ in $\det(\alpha T_I+\beta T_K)$. Expanding this determinant column by column,
\[
\binom{\ell}{b}\Psi_b=\sum_{|S|=b}\ \sum_{\pi\in S_\ell}\sgn(\pi)\prod_{s\in S}k_{\pi(s)-s}\prod_{s\notin S}i_{\pi(s)-s}.
\]
Each term is bounded by $C\,t^{-E(\pi,S)}(1+|\log t|)^b$, where $E(\pi,S)=\sum_{s\in S}|\pi(s)-s|-\sum_{s\notin S}|\pi(s)-s|$. By Lemma~\ref{B:lem:disp}(b), $E(\pi,S)\le2\fl{b/2}(\ell-\fl{b/2})$.
\end{proof}

\subsection{The squeeze and the shape of the operator}

\begin{theorem}\label{B:thm:Dt}
Let $\ell\ge1$ and let $L^t_\ell=\sum_{m=0}^{\ell+1}p_m(t)\partial^m$ be the minimal operator of $D_\ell$, normalized to be primitive. Then:
\begin{enumerate}
\item[(i)] $\ord L^t_\ell=\ell+1$; after scaling, $p_{\ell+1}=t^\ell$, $\deg p_m\le\ell$, and $t^{m-1}\mid p_m$ for $1\le m\le\ell+1$. Moreover $p_m(-t)=(-1)^{m+1}p_m(t)$.
\item[(ii)] $L^t_\ell$ has no singular point in $\C\setminus\{0\}$.
\item[(iii)] The exponents at $t=0$ are the integers $-2b(\ell-b)$, $0\le b\le\ell$, and $\ord\Psi_b=-2\fl{b/2}(\ell-\fl{b/2})$ for each $b$.
\item[(iv)] At $t=\infty$, $L^t_\ell$ has the basis of formal solutions $e^{\lambda_bt}t^{-\gamma_b}\sum_{i\ge0}c_{b,i}t^{-i}$, $0\le b\le\ell$. Moreover $\sum_m[t^\ell]p_m\,\lambda^m=\prod_{b=0}^{\ell}(\lambda-\lambda_b)$.
\item[(v)] $tL^t_\ell=\sum_{k=0}^{K}t^{2k}Q_{2k}(\theta)$ with $\deg Q_{2k}=\ell+1-2k$ and $\lc(Q_{2k})=[\lambda^{\ell+1-2k}]\chi(\lambda)$, where $\chi(\lambda)=\prod_b(\lambda-\lambda_b)$. The subleading coefficients are given by
\begin{equation}\label{B:eq:eta}
\sum_{k=0}^{K}[\theta^{\ell-2k}]Q_{2k}\cdot\lambda^{\ell-2k}=\sum_{b=0}^{\ell}\gamma_b\prod_{b'\ne b}(\lambda-\lambda_{b'})-\frac{\lambda}{2}\chi''(\lambda).
\end{equation}
\end{enumerate}
\end{theorem}

\begin{proof}
By Proposition~\ref{B:prop:solspace} and Corollary~\ref{B:cor:dim}, $n:=\ord L^t_\ell=\ell+1$ and $\Sol(L^t_\ell)=V$ has basis $\Psi_0,\dots,\Psi_\ell$.

\emph{Singular points.} Every element of $V$ extends analytically along all paths in $\C\setminus\{0\}$. Hence every zero $\tau\ne0$ of $p_n$ is an apparent singularity. By Lemma~\ref{B:lem:abel}(b) the residue $\kappa_\tau$ of $W'/W$ at $\tau$ is a positive integer, where $W$ is the Wronskian of $\Psi_0,\dots,\Psi_\ell$.

\emph{Abel at infinity.} The function $R=W'/W=-p_{n-1}/p_n$ is rational. Its poles are at $0$ and at the apparent points, and all of them are simple. At $0$ this is Fuchs' criterion; at an apparent point it holds because $W$ is holomorphic there. Hence
\[
R=Q(t)+\frac{\kappa_0}{t}+\sum_{\tau}\frac{\kappa_\tau}{t-\tau},\qquad W=c\,e^{\int Q}\,t^{\kappa_0}\prod_\tau(t-\tau)^{\kappa_\tau},
\]
with $Q\in\C[t]$ and $\kappa_0=\sum\nu_i-\binom n2$ by Lemma~\ref{B:lem:abel}(a), where $\nu_1,\dots,\nu_n$ are the exponents at $0$. On the other hand, Lemma~\ref{B:lem:vertex} gives
\[
W(t)=\Bigl(\prod_bC_b\Bigr)e^{(\sum_b\lambda_b)t}\,t^{-\sum_b\gamma_b}\bigl(\textstyle\prod_{b<b'}(\lambda_{b'}-\lambda_b)+O(t^{-1})\bigr),
\]
and $\sum_b\lambda_b=0$. Comparing, $Q=0$ and $\kappa_0+\sum_\tau\kappa_\tau=-\sum_b\gamma_b$, that is,
\begin{equation}\label{B:eq:abel-sum}
\sum_{\tau\ne0}\kappa_\tau=\binom{n}{2}-\sum_b\gamma_b-\sum_i\nu_i .
\end{equation}

\emph{The squeeze.} By Lemma~\ref{B:lem:origin} and Lemma~\ref{B:lem:squeeze},
\[
\sum_i\nu_i\ \ge\ \sum_b\ord\Psi_b\ \ge\ -2\sum_{b=0}^{\ell}\fl{b/2}\bigl(\ell-\fl{b/2}\bigr)=-2\sum_{b=0}^{\ell}b(\ell-b)=-\frac{(\ell-1)\ell(\ell+1)}{3}.
\]
The middle equality holds because the multisets $\{\fl{b/2}(\ell-\fl{b/2})\}_{b=0}^{\ell}$ and $\{b(\ell-b)\}_{b=0}^{\ell}$ coincide: both consist of the values $p(\ell-p)$ for $0\le p<\ell/2$, each twice, together with $\ell^2/4$ once if $\ell$ is even. Since $\sum_b\gamma_b=\sum_bb^2=\ell(\ell+1)(2\ell+1)/6$, the right-hand side of \eqref{B:eq:abel-sum} is at most
\[
\frac{\ell(\ell+1)}2-\frac{\ell(\ell+1)(2\ell+1)}6+\frac{(\ell-1)\ell(\ell+1)}3=\frac{\ell(\ell+1)}6\bigl(3-(2\ell+1)+2(\ell-1)\bigr)=0 .
\]
As each $\kappa_\tau\ge1$, there are no apparent singularities, which proves (ii). All the inequalities above are therefore equalities. In particular $\ord\Psi_b=-2\fl{b/2}(\ell-\fl{b/2})$ for every $b$, because each term was bounded separately. By the equality case of Lemma~\ref{B:lem:squeeze}, the exponents are the numbers $\ord\Psi_b$, that is, the multiset $\{-2b(\ell-b)\}$. This proves (iii).

\emph{Shape.} By (ii), $p_n=ct^A$ with $c\ne0$. By Lemma~\ref{B:lem:vertex} and Lemma~\ref{B:lem:top}, $\deg p_m\le A$ for all $m$ and $\sum_m[t^A]p_m\lambda^m=c\prod_b(\lambda-\lambda_b)$. Fuchs' criterion at $0$ gives $\ord_0p_m\ge A-n+m$, and the indicial polynomial is $I(\rho)=\sum_m[t^{A-n+m}]p_m\,\ff{\rho}{m}$. If $A\le n-2$, the terms with $m<n-A$ vanish, so $\ff{\rho}{n-A}$ divides $I$ and $1$ would be an exponent; this contradicts (iii), as all exponents are $\le0$. Hence $A\ge n-1$. Since $0$ is an exponent, $[t^{A-n}]p_0=I(0)=0$. Hence $t^{A-n+1}$ divides every $p_m$, and primitivity gives $A\le n-1$. So $A=\ell$, which proves the first part of (i) and the second part of (iv).

\emph{Formal solutions.} By Lemma~\ref{B:lem:vertex}, applying $L^t_\ell$ termwise to the asymptotic expansion of $\Psi_b$ gives the asymptotic expansion of $L^t_\ell\Psi_b=0$. So each expansion is a formal solution, and they are independent because the $\lambda_b$ are distinct. This proves (iv).

\emph{Parity.} Since $I_m(-2t)=(-1)^mI_m(2t)$, conjugation by $\mathrm{diag}((-1)^p)$ shows $D_\ell(-t)=D_\ell(t)$. The operator obtained from $L^t_\ell$ by $t\mapsto-t$, namely $\sum_mp_m(-t)(-1)^m\partial^m$, is primitive of order $n$ and annihilates $D_\ell$. By uniqueness of the minimal operator it is a constant multiple of $L^t_\ell$, and comparing leading terms the constant is $-1$. This gives $p_m(-t)=(-1)^{m+1}p_m(t)$.

\emph{The $\theta$-coefficients.} Write $tL^t_\ell=\sum_mr_m(t)\ff{\theta}{m}$ with $r_m=p_m/t^{m-1}$ for $m\ge1$ and $r_0=tp_0$, so that $\deg r_m\le\ell+1-m$. Thus $tL^t_\ell=\sum_kt^kQ_k(\theta)$ with $\deg Q_k\le\ell+1-k$. By the parity just proved, $tL^t_\ell$ is invariant under $t\mapsto-t$, and $\theta$ is invariant too; hence $Q_k=0$ for odd $k$. Comparing the coefficients of $t^{\ell+1}\partial^m$ gives $[t^\ell]p_m=[\theta^m]Q_{\ell+1-m}$. Hence $\chi(\lambda)=\sum_k[\theta^{\ell+1-k}]Q_k\,\lambda^{\ell+1-k}$. The polynomial $\chi(\lambda)=\lambda^{[\ell\text{ even}]}\prod_{b<\ell/2}(\lambda^2-\lambda_b^2)$ has nonzero coefficients at all powers $\lambda^{\ell+1-2k}$, $0\le k\le K$: up to sign they are elementary symmetric functions of the positive numbers $\lambda_b^2$. Hence $\deg Q_{2k}=\ell+1-2k$, with $\lc(Q_{2k})=[\lambda^{\ell+1-2k}]\chi$.

For \eqref{B:eq:eta}, substitute a formal solution $y=e^{\lambda t}t^{-\gamma}g$, with $\lambda=\lambda_b$, $\gamma=\gamma_b$ and $g=1+O(t^{-1})$, into $tL^t_\ell$. Since $\theta(e^{\lambda t}t^{-\gamma}g)=e^{\lambda t}t^{-\gamma}(\lambda t-\gamma+\theta)g$, we must expand $Q_{2k}(\lambda t+\theta-\gamma)$ applied to $g$. Write $d_k=\ell+1-2k$. The terms of order $t^{d_k}$ and $t^{d_k-1}$ come from the leading monomial and from the subleading one. In the leading monomial $(\lambda t+\theta-\gamma)^{d_k}$, the contribution with one factor $\theta-\gamma$ in position $p$ acting on $t^{d_k-p}g$ gives $\lambda^{d_k-1}t^{d_k-1}(d_k-p-\gamma)$. Summing over $p$ gives $\lambda^{d_k-1}t^{d_k-1}\bigl(\binom{d_k}2-d_k\gamma\bigr)$. Collecting the coefficient of $t^{\ell}$ and using $\chi(\lambda_b)=0$, we get, for every $b$,
\[
\sum_k\lambda_b^{d_k-1}[\theta^{d_k-1}]Q_{2k}+\frac{\lambda_b}{2}\chi''(\lambda_b)-\gamma_b\chi'(\lambda_b)=0 .
\]
The polynomial obtained from the left side of \eqref{B:eq:eta} by adding $\frac\lambda2\chi''(\lambda)$ has degree at most $\ell$. At the $\ell+1$ points $\lambda_b$ it takes the values $\gamma_b\chi'(\lambda_b)=\gamma_b\prod_{b'\ne b}(\lambda_b-\lambda_{b'})$. Lagrange interpolation gives \eqref{B:eq:eta}.
\end{proof}

\begin{remark}
The squeeze shows a little more than (iii): the crude bound of Lemma~\ref{B:lem:origin} is sharp in the power of $t$ for every single vertex integral (the logarithmic factor need not be). Pairs of vertex integrals share the same order at $0$, namely $\Psi_{2p}$ and $\Psi_{2p+1}$. Computing their leading terms directly would require controlling the cancellations among the $\ell!^2$ terms of the Toeplitz expansion. The squeeze avoids this entirely. High-precision evaluation confirms the orders for $\ell\le6$ (Appendix~\ref{B:app:comp}).
\end{remark}

\subsection{Proof of Theorem~\ref{B:thm:A}}
Put $x=t^2$, so that $\theta_t=2\theta_x$ and $G_\ell(x)=D_\ell(t)$. By Theorem~\ref{B:thm:Dt}(v),
\[
\Theta:=tL^t_\ell=\sum_{k=0}^{K}x^k\,Q_{2k}(2\theta_x)
\]
is an operator in $x$ annihilating $G_\ell$. If $G_\ell$ satisfied an equation of order at most $\ell$ in $x$, then $D_\ell$ would satisfy one in $t$, contradicting Theorem~\ref{B:thm:Dt}(i). So the minimal operator of $G_\ell$ has order $\ell+1$, and $\Theta=xL_\ell$ up to normalization. We normalize so that $P_k(\theta):=2^{-\ell-1}Q_{2k}(2\theta)$ gives a monic $P_0$. Then $\deg P_k=\ell+1-2k$ and $\lc(P_k)=2^{-2k}\lc(Q_{2k})$, so
\[
\sum_k\lc(P_k)\mu^{\ell+1-2k}=2^{-\ell-1}\chi(2\mu)=\prod_b(\mu-\ell+2b).
\]
The polynomial $P_0$ is the indicial polynomial at $x=0$. By Theorem~\ref{B:thm:Dt}(iii) its roots are half the $t$-exponents, so $P_0(\theta)=\prod_b(\theta+b(\ell-b))$.

If $\ell$ is even, the rate $\lambda_{\ell/2}$ vanishes, and $\Psi_{\ell/2}(\sqrt x)=C\,x^{-\gamma_{\ell/2}/2}(1+O(x^{-1/2}))$ is a solution. Applying $\Theta$, the top term $x^KP_K(\theta)x^{-\gamma_{\ell/2}/2}=P_K(-\gamma_{\ell/2}/2)\,x^{K-\gamma_{\ell/2}/2}$ must vanish, because all other terms are of lower order. So $P_K(\theta)=\lc(P_K)(\theta+\ell^2/8)$, as $\gamma_{\ell/2}=\ell^2/4$. This proves (b), except for the absence of singular points in $\C^*$, which follows from Theorem~\ref{B:thm:Dt}(ii) because $\tau\mapsto\tau^2$ maps $\C^*$ onto $\C^*$. Part (c) is Theorem~\ref{B:thm:Dt}(iv) rewritten in $x$.

It remains to prove (a). Write $L_\ell=\sum_mq_m\partial_x^m$. Since $x^k\theta^i=\sum_{m\le i}S(i,m)x^{k+m}\partial_x^m$, the polynomial $xq_m$ is a combination of monomials $x^{k+m}$ coming from $x^kP_k(\theta)$ with $\deg P_k\ge m$. Hence $x^{m-1}\mid q_m$, $q_{\ell+1}=x^\ell$, and $\deg q_m\le k^*+m-1=\fl{(\ell-1+m)/2}$, where $k^*=\fl{(\ell+1-m)/2}$ is the largest $k$ with $\ell+1-2k\ge m$. The coefficient of $x^{k^*+m}$ in $xq_m$ comes only from $P_{k^*}$. It equals $\lc(P_{k^*})\ne0$ if $\deg P_{k^*}=m$, and $\binom{m+1}{2}\lc(P_{k^*})+[\theta^m]P_{k^*}$ if $\deg P_{k^*}=m+1$.

In the second case put $d=m+1=\ell+1-2k^*$. By \eqref{B:eq:eta} and $\lc(Q_{2k})=[\lambda^{\ell+1-2k}]\chi$, this coefficient equals $2^{d-1-\ell-1}$ times the coefficient of $\lambda^{d-1}$ in $\frac\lambda2\chi''(\lambda)+\sum_b\gamma_b\prod_{b'\ne b}(\lambda-\lambda_{b'})$. Write $\varepsilon_\ell=[\ell\text{ even}]$, $c_j=\lambda_j^2>0$ for $j<K$, and $\chi=\lambda^{\varepsilon_\ell}\prod_{j<K}(\lambda^2-c_j)$. Pairing $b$ with $\ell-b$ (and keeping $b=\ell/2$ alone if $\ell$ is even), a direct computation shows that this coefficient is
\[
(-1)^{k^*}\Bigl(\tbinom{d}{2}\sigma_{k^*}(c)+2\sum_{j<K}\gamma_j\,\sigma_{k^*}(c_{\ne j})+\varepsilon_\ell\,\gamma_{\ell/2}\,\sigma_{k^*}(c)\Bigr),
\]
where $\sigma_i$ denotes the elementary symmetric functions and $c_{\ne j}$ the list with $c_j$ removed. All terms are nonnegative. If $k^*<K$, the middle term is positive; if $k^*=K$, then $d=1$, so $\varepsilon_\ell=1$ and the last term is positive. Hence $\deg q_m=\fl{(\ell-1+m)/2}$ exactly.

Finally, $q_1(0)=P_0(1)=\prod_b(1+b(\ell-b))\ne0$ and $q_{\ell+1}=x^\ell$ has no zero in $\C^*$. Hence $L_\ell$ is primitive, and every operator of order $\ell+1$ with polynomial coefficients annihilating $G_\ell$ is $r\,L_\ell$ with $r\in\C(x)$, necessarily a polynomial. \qed

\begin{proof}[Proof of Corollary~\ref{B:cor:FW41}]
For $q=0$ the function in \cite[Prop.~4.1]{FW26} is $G_\ell$ \cite[(4.6)]{FW26}, and Theorem~\ref{B:thm:A}(a) gives exactly the asserted shape. The optimality of the order is the first statement of Theorem~\ref{B:thm:A}(a). Conjecture~\ref{B:conj:pairs} follows from (a), since $\fl{(\ell-1+m)/2}\le\ell$.
\end{proof}

\begin{remark}
Appendix~B of \cite{FW26}, due to Bornemann, discusses optimality using algebraic independence of the Bessel functions. Here minimality follows from Proposition~\ref{B:prop:solspace} and Corollary~\ref{B:cor:dim}; it also follows from irreducibility of the $\ell$th symmetric power of the standard representation of $\SL_2$.
\end{remark}

\section{Twisted Toeplitz determinants}\label{B:sec:twisted}

Fix $\ell\ge1$ and $q\ge0$. The arguments of Section~\ref{B:sec:unitary} extend to $D_{\ell,q}$. The main new ingredient is an inequality for displacements with a drift, Lemma~\ref{B:lem:disp-q}. In addition, the leading coefficient in the variable $t$ can now be $t^{\ell+1}$, and the gauge factor $x^{q\ell/2}$ produces an extra term in the exact-degree argument.

\subsection{The functions}

\begin{lemma}\label{B:lem:twisted-expansion}
For $q\in\Z$ one has $D_{\ell,-q}=D_{\ell,q}$ and $D_{\ell,q}(t)=\E_{U(\ell)}(\det U)^q\,e^{2t\,\mathrm{Re}\tr U}$. For $q\ge0$,
\[
D_{\ell,q}(t)=\sum_{n\ge0}u_{\ell,q}(n)\,\frac{t^{2n+q\ell}}{n!\,(n+q\ell)!},\qquad u_{\ell,q}(n)=\sum_{\mu}f^{\mu}f^{\mu+(q^\ell)},
\]
where $\mu$ runs over the partitions of $n$ with at most $\ell$ parts. In particular $G_{\ell,q}$ is a power series with $G_{\ell,q}(0)=f^{(q^\ell)}/(q\ell)!>0$, and $D_{\ell,q}(-t)=(-1)^{q\ell}D_{\ell,q}(t)$.
\end{lemma}

\begin{proof}
Transposition and $b_{-m}=b_m$ give $D_{\ell,-q}=D_{\ell,q}$. By Heine's identity, that is, Andreief's identity on the unit circle (see \cite{BR01}), $\E_{U(\ell)}\prod_j\varphi(z_j)=\det[\hat\varphi_{i-j}]_{i,j=1}^{\ell}$, where $z_1,\dots,z_\ell$ are the eigenvalues of $U$ and $\hat\varphi_m$ is the $m$-th Fourier coefficient of $\varphi$ on the unit circle. For $\varphi(z)=z^{-q}e^{t(z+z^{-1})}$ we have $\hat\varphi_m=b_{m+q}$. Hence $\E(\det U)^{-q}e^{2t\,\mathrm{Re}\tr U}=D_{\ell,q}$, and replacing $q$ by $-q$ gives the stated formula.

For the expansion, write $e^{2t\,\mathrm{Re}\tr U}=\sum_{a,c\ge0}\frac{t^{a+c}}{a!\,c!}(\tr U)^a\,\overline{\tr U}^{\,c}$ and use $(\tr U)^a=\sum_{\lambda\vdash a}f^\lambda s_\lambda(U)$, where $s_\lambda(U)=0$ if $\lambda$ has more than $\ell$ parts. Since $(\det U)^qs_\mu(U)=s_{\mu+(q^\ell)}(U)$ and $\overline{\det U}=(\det U)^{-1}$, Schur orthogonality gives $\E\,(\det U)^{-q}s_\lambda(U)\overline{s_\mu(U)}=\delta_{\lambda,\mu+(q^\ell)}$. Collecting terms in $\E(\det U)^{-q}e^{2t\,\mathrm{Re}\tr U}$ gives the expansion. The parity follows from $b_m(-t)=(-1)^mb_m(t)$ after conjugation by $\mathrm{diag}((-1)^i)$.
\end{proof}

\subsection{Vertex integrals with a twist}
For sequences $h^{(1)},\dots,h^{(\ell)}$ as in Section~\ref{B:sec:unitary} put
\[
F_q\bigl(h^{(1)},\dots,h^{(\ell)}\bigr)=\frac1{\ell!}\sum_{\sigma,\tau\in S_\ell}\sgn(\sigma)\sgn(\tau)\prod_{j=1}^{\ell}h^{(j)}_{\sigma(j)-\tau(j)+q},
\]
$\Phi_q(g)=F_q(g_\bullet,\dots,g_\bullet)$ for $g\in\Sol(T)$, and $\Psi^{(q)}_b=F_q(k,\dots,k,i,\dots,i)$ with $b$ copies of $k$. As before, $F_q$ is multilinear and symmetric, $F_q(h,\dots,h)=\det[h_{p-p'+q}]_{p,p'=1}^{\ell}$, and $\Phi_q(i_0)=D_{\ell,q}$. Let $L^t_{\ell,q}$ be the minimal operator of $D_{\ell,q}$ over $\C(t)$, normalized to be primitive.

\begin{proposition}\label{B:prop:twisted-basic}
\begin{enumerate}
\item[(a)] $\Psi^{(q)}_b$ is given by the integral of Lemma~\ref{B:lem:Psi-int} with the additional factor $\prod_jz_j^q$ in the integrand. It is real on $(0,\infty)$, and $\Phi_q(\alpha i_0+\beta k_0)=\sum_b\binom{\ell}{b}\alpha^{\ell-b}\beta^b\,\Psi^{(q)}_b$.
\item[(b)] $\Sol(L^t_{\ell,q})=V_q:=\spn_\C\{\Psi^{(q)}_0,\dots,\Psi^{(q)}_\ell\}$.
\item[(c)] Lemma~\ref{B:lem:vertex} holds for $\Psi^{(q)}_b$, with the same $\lambda_b$ and $\gamma_b$ and with $C_b$ replaced by $(-1)^{qb}C_b$. In particular $\ord L^t_{\ell,q}=\dim V_q=\ell+1$.
\item[(d)] $V_q\subset\mathcal N$, and $0$ is a regular singular point of $L^t_{\ell,q}$.
\end{enumerate}
\end{proposition}

\begin{proof}
Parts (a), (b) and (d) are proved like Lemma~\ref{B:lem:Psi-int}, Proposition~\ref{B:prop:solspace} and the first paragraph of Section~\ref{B:ssec:origin}; $\Psi^{(q)}_b$ is real because the $i_m$ and $k_m$ are real. For (c), the additional factor $\prod_{j\le b}(-s_j)^q\prod_{j>b}e^{iq\varphi_j}$ is real-analytic near the vertex $v^*$, takes the value $(-1)^{qb}$ there, and is bounded by a polynomial in the $s_j^{\pm1}$. After absorbing it into $H$, the proof of Lemma~\ref{B:lem:vertex} applies without change; the vanishing of the odd Gaussian moments used there does not depend on $H$. The rates $\lambda_b$ are pairwise distinct, which gives the last assertion as in Corollary~\ref{B:cor:dim}.
\end{proof}

\subsection{Displacement with drift}

\begin{lemma}\label{B:lem:disp-q}
For $0\le i\le\ell$ put $f(i)=2i(\ell-i)-q(\ell-2i)$, and let $f_{(0)}\le f_{(1)}\le\dots\le f_{(\ell)}$ be these $\ell+1$ numbers in increasing order. Let $\pi\in S_\ell$ and $S\subseteq[\ell]$ with $|S|=b$, put $d_s=\pi(s)-s+q$, and
\[
E_q(\pi,S)=\sum_{s\in S}|d_s|-\sum_{s\notin S}|d_s| .
\]
Then $E_q(\pi,S)\le f_{(b)}$. Moreover $\sum_{b=0}^{\ell}f_{(b)}=\frac{(\ell-1)\ell(\ell+1)}3$.
\end{lemma}

For $q=0$ we have $f_{(b)}=2\fl{b/2}(\ell-\fl{b/2})$, so the lemma contains Lemma~\ref{B:lem:disp}(b).

\begin{proof}
Let $S_+=\{s\in S:d_s>0\}$, $S_-=S\setminus S_+$ and $j=|S_+|$, and put $A=\sum_{s\in S_+}d_s$ and $B=-\sum_{s\in S_-}d_s\ge0$. Then $\sum_{s\in S}|d_s|=A+B$ and $\sum_{s\in S}d_s=A-B$. Since $\sum_{s=1}^{\ell}d_s=q\ell$, we have $\sum_{s\notin S}|d_s|\ge|q\ell-A+B|$, and hence
\[
E_q(\pi,S)\le A+B-|q\ell-A+B|=\min(2A-q\ell,\;2B+q\ell).
\]
For a set $T\subseteq[\ell]$ with $m$ elements, both $\sum_{s\in T}(\pi(s)-s)$ and $\sum_{s\in T}(s-\pi(s))$ are at most $\sum_{i=1}^m\bigl((\ell+1-i)-i\bigr)=m(\ell-m)$. Applied to $T=S_+$ and $T=S_-$, this gives $A\le j(\ell-j)+jq$ and $B\le(b-j)(\ell-b+j)-(b-j)q$. Substituting, one finds
\[
2A-q\ell\le f(j),\qquad 2B+q\ell\le f(j'),\qquad j'=j+\ell-b .
\]
So $E_q(\pi,S)\le c:=\min(f(j),f(j'))$, where $0\le j\le j'\le\ell$. Since $f$ is a concave quadratic function, $f(i)\ge c$ for each of the $\ell-b+1$ integers $i\in[j,j']$. Hence at most $b$ of the numbers $f(0),\dots,f(\ell)$ are smaller than $c$, which means $f_{(b)}\ge c$. Finally, $\sum_bf_{(b)}=\sum_if(i)=2\sum_ii(\ell-i)$ because $\sum_{i=0}^{\ell}(\ell-2i)=0$.
\end{proof}

\begin{lemma}\label{B:lem:origin-q}
There is a constant $C$ such that $|\Psi^{(q)}_b(t)|\le C\,t^{-f_{(b)}}(1+|\log t|)^b$ for $0<t\le1$ and $0\le b\le\ell$. Consequently $\ord\Psi^{(q)}_b\ge-f_{(b)}$.
\end{lemma}

\begin{proof}
By Proposition~\ref{B:prop:twisted-basic}(a), $\binom{\ell}b\Psi^{(q)}_b$ is the coefficient of $\alpha^{\ell-b}\beta^b$ in $\det(\alpha T^{(q)}_I+\beta T^{(q)}_K)$, where $T^{(q)}_I=[i_{p-p'+q}]$ and $T^{(q)}_K=[k_{p-p'+q}]$. Expanding as in the proof of Lemma~\ref{B:lem:origin}, every term is bounded by $C\,t^{-E_q(\pi,S)}(1+|\log t|)^b$, and Lemma~\ref{B:lem:disp-q} applies.
\end{proof}

\subsection{The shape}

\begin{theorem}\label{B:thm:Dtq}
Let $n=\ell+1$, write $L^t_{\ell,q}=\sum_{m=0}^{n}p_m(t)\partial^m$, and put $e_i=q(\ell-2i)-2i(\ell-i)$ for $0\le i\le\ell$.
\begin{enumerate}
\item[(i)] $L^t_{\ell,q}$ has order $\ell+1$ and no singular point in $\C\setminus\{0\}$.
\item[(ii)] The exponents at $t=0$ are $e_0,\dots,e_\ell$, and $\ord\Psi^{(q)}_b=-f_{(b)}$ for every $b$.
\item[(iii)] After scaling, $p_n=t^A$, where $A=\ell$ if $0\in\{e_0,\dots,e_\ell\}$ and $A=\ell+1$ otherwise. Moreover $\deg p_m\le A$ and $\ord_0p_m\ge m+A-n$ for all $m$, and $\sum_m[t^A]p_m\lambda^m=\chi(\lambda)=\prod_{b=0}^{\ell}(\lambda-\lambda_b)$.
\item[(iv)] At $t=\infty$, $L^t_{\ell,q}$ has a basis of formal solutions $e^{\lambda_bt}t^{-\gamma_b}\sum_{i\ge0}c^{(q)}_{b,i}t^{-i}$, $0\le b\le\ell$, with $c^{(q)}_{b,0}=(-1)^{qb}C_b$.
\item[(v)] $t^{n-A}L^t_{\ell,q}=\sum_{k=0}^{K}t^{2k}Q_{2k}(\theta)$ with $Q_0(\theta)=\prod_{i=0}^{\ell}(\theta-e_i)$, $\deg Q_{2k}=\ell+1-2k$ and $\lc(Q_{2k})=[\lambda^{\ell+1-2k}]\chi(\lambda)$, and \eqref{B:eq:eta} holds for these $Q_{2k}$.
\end{enumerate}
\end{theorem}

For $q=0$ this is Theorem~\ref{B:thm:Dt}. For $q\ge1$ the exponent $0$ occurs only when $q(\ell-2i)=2i(\ell-i)$ for some $1\le i<\ell/2$. For $\ell\le4$ the only such pairs $(\ell,q)$ are $(3,4)$ and $(4,3)$; further examples are $(5,12)$, $(6,8)$ and $(8,6)$.

\begin{proof}
By Proposition~\ref{B:prop:twisted-basic}, $\Sol(L^t_{\ell,q})=V_q$ has the basis $\Psi^{(q)}_0,\dots,\Psi^{(q)}_\ell$, with the same rates $\lambda_b$ and exponents $\gamma_b$ at infinity as for $q=0$. The parts \emph{Singular points}, \emph{Abel at infinity} and \emph{The squeeze} of the proof of Theorem~\ref{B:thm:Dt} therefore apply, with Lemma~\ref{B:lem:origin-q} in place of Lemma~\ref{B:lem:origin}. They give
\[
\sum_{\tau\ne0}\kappa_\tau=\binom n2-\sum_b\gamma_b-\sum_i\nu_i\le\binom n2-\sum_b\gamma_b+\sum_bf_{(b)}=0,
\]
by Lemma~\ref{B:lem:disp-q} and the computation in that proof. Hence there are no apparent singularities, $\ord\Psi^{(q)}_b=-f_{(b)}$ for all $b$, and, by the equality case of Lemma~\ref{B:lem:squeeze}, the exponents at $0$ form the multiset $\{-f_{(b)}\}=\{-f(i)\}=\{e_i\}$. This proves (i) and (ii).

For (iii), $p_n=ct^A$ by (i). Lemma~\ref{B:lem:top}, applied with the asymptotics of Proposition~\ref{B:prop:twisted-basic}(c), gives $\deg p_m\le A$ and the formula for $\chi$. By Fuchs' criterion $\ord_0p_m\ge A-n+m$, and the indicial polynomial is $I(\rho)=\sum_m[t^{A-n+m}]p_m\,\ff{\rho}{m}$. All $e_i$ are congruent to $q\ell$ modulo $2$. Suppose first that $0$ is not an exponent. If $A<n$, the term $m=0$ of $I$ vanishes and $I(0)=0$, which is impossible; so $A\ge n$. Then $t^{A-n}$ divides every $p_m$, and primitivity gives $A=n$. Suppose now that $0$ is an exponent. Then $q\ell$ is even, so $1$ is not an exponent. If $A\le n-2$, the terms $m\le1$ of $I$ vanish and $\ff{\rho}{2}$ divides $I$, which is impossible; so $A\ge n-1$. Since $I(0)=0$, the power $t^{A-n+1}$ divides every $p_m$, and primitivity gives $A\le n-1$.

Part (iv) follows from Proposition~\ref{B:prop:twisted-basic}(c) as in the proof of Theorem~\ref{B:thm:Dt}.

For (v), put $\Theta_q=t^{n-A}L^t_{\ell,q}$. By (iii), $\Theta_q=\sum_mr_m(t)\ff{\theta}{m}$ with $r_m=t^{n-A-m}p_m\in\C[t]$ of degree at most $n-m$. So $\Theta_q=\sum_kt^kQ_k(\theta)$ with $\deg Q_k\le n-k$, and $Q_0(\theta)=\sum_mr_m(0)\ff{\theta}{m}=I(\theta)$. By Lemma~\ref{B:lem:twisted-expansion}, $D_{\ell,q}(-t)=\pm D_{\ell,q}(t)$. As in the proof of Theorem~\ref{B:thm:Dt}, the operator obtained from $L^t_{\ell,q}$ by the substitution $t\mapsto-t$ is a constant multiple of $L^t_{\ell,q}$. Comparing the terms with $\partial^n$, the constant is $(-1)^{A+n}$, so $\Theta_q$ is invariant under $t\mapsto-t$. Hence $Q_k=0$ for odd $k$. The remaining assertions follow as in the proof of Theorem~\ref{B:thm:Dt}(v). Comparing coefficients of $t^n\partial^m$ gives $[t^A]p_m=[\theta^m]Q_{n-m}$, which determines the degrees and leading coefficients from $\chi$. The computation with the formal solutions of (iv) gives \eqref{B:eq:eta}.
\end{proof}

\begin{remark}
The bound of Lemma~\ref{B:lem:disp-q} is attained for every $b$. Indeed, $E_q(\pi,S)$ is an integer, so if every term in the proof of Lemma~\ref{B:lem:origin-q} had $E_q(\pi,S)<f_{(b)}$, we would get $\ord\Psi^{(q)}_b>-f_{(b)}$, contradicting Theorem~\ref{B:thm:Dtq}(ii). We also confirmed this by enumeration for $\ell\le6$ and $q\le8$.
\end{remark}

\subsection{Proof of Theorem~\ref{B:thm:twisted}}
Part (b) is Theorem~\ref{B:thm:Dtq}(i)--(iii).

\emph{(a)} Put $x=t^2$. By definition $G_{\ell,q}(t^2)=t^{-q\ell}D_{\ell,q}(t)$, and by Lemma~\ref{B:lem:twisted-expansion} this is a power series in $t^2$. Conjugation by $t^{q\ell}$ replaces $\theta_t$ by $\theta_t+q\ell$, and $\theta_t=2\theta_x$. By Theorem~\ref{B:thm:Dtq}(v) the $\theta$-form
\[
\Theta'=\sum_{k=0}^{K}x^k\,Q_{2k}(2\theta_x+q\ell)
\]
annihilates $G_{\ell,q}$. As in the proof of Theorem~\ref{B:thm:A}, $G_{\ell,q}$ satisfies no equation of order at most $\ell$, because $D_{\ell,q}(t)=t^{q\ell}G_{\ell,q}(t^2)$ would then satisfy one. Put $P_k(\theta)=2^{-\ell-1}Q_{2k}(2\theta+q\ell)$. Then $\deg P_k=\ell+1-2k$, the leading coefficients of the $P_k$ are those of Theorem~\ref{B:thm:A}(b), and
\[
P_0(\theta)=\prod_{i=0}^{\ell}\Bigl(\theta+\frac{q\ell-e_i}2\Bigr)=\prod_{i=0}^{\ell}\bigl(\theta+i(\ell+q-i)\bigr).
\]
For even $\ell$, the solution $x^{-q\ell/2}\Psi^{(q)}_{\ell/2}(\sqrt x)=C\,x^{-(\gamma_{\ell/2}+q\ell)/2}\bigl(1+O(x^{-1/2})\bigr)$ gives $P_K(\theta)=\lc(P_K)(\theta+\ell^2/8+q\ell/2)$, exactly as in the proof of Theorem~\ref{B:thm:A}. The formal solutions at infinity are those of Theorem~\ref{B:thm:Dtq}(iv), multiplied by $x^{-q\ell/2}$.

All roots of $P_0$ are $\le0$, and $P_0(0)=0$. Hence the last three paragraphs of the proof of Theorem~\ref{B:thm:A} apply to $L_{\ell,q}:=x^{-1}\sum_kx^kP_k(\theta)$. They give $q_{\ell+1}=x^\ell$, so that $L_{\ell,q}$ has no singular point in $\C\setminus\{0\}$, as well as $x^{m-1}\mid q_m$, $\deg q_m\le\fl{(\ell-1+m)/2}$, and primitivity, because $q_1(0)=P_0(1)=\prod_i(1+i(\ell+q-i))\ne0$. For the exact degrees, the only change concerns the case $\deg P_{k^*}=m+1=:d$. Write $Q_{2k^*}(u)=a_du^d+a_{d-1}u^{d-1}+\cdots$. Then $[\theta^{d-1}]P_{k^*}=2^{d-\ell-2}(a_{d-1}+d\,q\ell\,a_d)$, so the coefficient $\binom{d}2\lc(P_{k^*})+[\theta^{d-1}]P_{k^*}$ is $2^{d-\ell-2}$ times the coefficient of $\lambda^{d-1}$ in
\[
\frac\lambda2\chi''(\lambda)+\sum_b\gamma_b\prod_{b'\ne b}(\lambda-\lambda_{b'})+q\ell\,\chi'(\lambda).
\]
Compared with the case $q=0$, the additional term is $q\ell\,d\,a_d=(-1)^{k^*}d\,q\ell\,\sigma_{k^*}(c)$, in the notation of that proof. It has the same sign $(-1)^{k^*}$ as the other terms there, so the coefficient does not vanish, and $\deg q_m=\fl{(\ell-1+m)/2}$ in all cases.

\emph{(c)} Write $g(n)=u_{\ell,q}(n)/(n!\,(n+q\ell)!)$ for the Taylor coefficients of $G_{\ell,q}$. Comparing coefficients of $x^N$ in $xL_{\ell,q}G_{\ell,q}=0$ gives $\sum_kP_k(N-k)g(N-k)=0$. Multiply by $N!\,(N+q\ell)!$ and, for $N\ge1$, divide by $N(N+q\ell)$, which are the factors of $P_0(N)$ with $i=0$ and $i=\ell$. This gives
\[
\sum_{k=0}^{K}R^{(q)}_k(N)\,u_{\ell,q}(N-k)=0\quad(N\ge1),\qquad R^{(q)}_0(N)=\prod_{i=1}^{\ell-1}\bigl(N+i(\ell+q-i)\bigr),
\]
with $R^{(q)}_k(N)=P_k(N-k)\,\ff{(N-1)}{k-1}\,\ff{(N+q\ell-1)}{k-1}$ for $k\ge1$; all $R^{(q)}_k$ have degree $\ell-1$. Conversely, suppose that $u_{\ell,q}$ satisfies a recurrence of order $M$. Since $(n+q\ell)!/n!$ is a polynomial in $n$, the sequence $a(n)=(n!)^2g(n)=n!\,u_{\ell,q}(n)/(n+q\ell)!$ satisfies a recurrence of order $M$ as well. By Lemma~\ref{B:lem:dict-degf}, $G_{\ell,q}=\sum_na(n)x^n/(n!)^2$ is annihilated by a $\theta$-form of span $M$. In the variable $t$, after conjugation by $t^{q\ell}$, this is a $\theta$-form of span $2M$ annihilating $D_{\ell,q}$. As in the proof of Theorem~\ref{B:thm:B}(c), Lemma~\ref{B:lem:types} and Proposition~\ref{B:prop:twisted-basic}(c) give $2K\le2M$. \qed

\begin{proof}[Proof of Corollary~\ref{B:cor:FW41q}]
By \cite[(2.4), (2.5)]{FW26}, the function ${}_0F_1(q+\ell;(x)^\ell)$ with Jack parameter $\alpha=1$ is a constant multiple of $x^{-q\ell/2}D_{\ell,q}(\sqrt x)=G_{\ell,q}(x)$, and by \cite[(3.10)]{FW26} it is, for $\beta=2$, $n=\ell$ and $a+b=q-\ell$, a constant multiple of the function in \cite[Prop.~4.1]{FW26}. Theorem~\ref{B:thm:twisted}(a) gives the shape asserted there, with exact degrees, and shows that the order $\ell+1$ is minimal. For the Hankel determinant $w_\ell(z)=c_\ell(-1)^{\ell(\ell-1)/2}z^{-\ell^2}\det[b_{j+k+1}(z)]_{j,k=0}^{\ell-1}$ of \cite[App.~B]{FW26}, reversing the order of the columns gives $w_\ell(z)=c_\ell z^{-\ell^2}D_{\ell,\ell}(z)=c_\ell\,G_{\ell,\ell}(z^2)$. If $w_\ell$ satisfied an equation of order at most $\ell$ in $z$, so would $D_{\ell,\ell}$, contradicting Theorem~\ref{B:thm:Dtq}(i).
\end{proof}

\section{Permutations with bounded increasing subsequences}\label{B:sec:perm}

\begin{proof}[Proof of Theorem~\ref{B:thm:B}]
\emph{(a)} By Theorem~\ref{B:thm:A}(b), the operator $xL_\ell=\sum_{k=0}^{K}x^kP_k(\theta)$ annihilates the double exponential generating function $G_\ell$ of $u_\ell$. The converse part of Lemma~\ref{B:lem:dict-degf} gives
\[
\sum_{k=0}^{K}P_k(N-k)\bigl(\ff{N}{k}\bigr)^2u_\ell(N-k)=0\qquad(N\ge0).
\]
For $N\ge1$ divide by $N^2$. Since $P_0(\theta)=\theta^2\prod_{b=1}^{\ell-1}(\theta+b(\ell-b))$, the term $k=0$ becomes $R_0(N)u_\ell(N)$. For $k\ge1$ we have $(\ff{N}{k})^2/N^2=(\ff{(N-1)}{k-1})^2$. This is the recurrence in (a). The degree of $R_k$ is $\deg P_k+2k-2=\ell-1$, by Theorem~\ref{B:thm:A}(b). The leading coefficient of $P_K$ is the coefficient of $\mu^{\ell+1-2K}=\mu^{[\ell\text{ even}]}$ in
\[
\prod_{b=0}^{\ell}(\mu-\ell+2b)=\mu^{[\ell\text{ even}]}\prod_{j=0}^{K-1}\bigl(\mu^2-(\ell-2j)^2\bigr),
\]
namely $(-1)^K\prod_{j<K}(\ell-2j)^2$. Together with $P_K(\theta)=\lc(P_K)(\theta+\ell^2/8)$ for even $\ell$, this gives the formula for $R_K$.

\emph{(b)} We have $\lc(R_k)=\lc(P_k)$. The identity just displayed, divided by $\mu^{[\ell\text{ even}]}$, gives
\[
\sum_k\lc(P_k)\,\mu^{2(K-k)}=\prod_{j<K}\bigl(\mu^2-(\ell-2j)^2\bigr).
\]
Now put $\lambda=\mu^2$.

\emph{(c), order.} Suppose $u_\ell$ satisfies a recurrence of order $M$ for $n\ge n_0$. By Lemma~\ref{B:lem:dict-degf}, $G_\ell$ is annihilated by a $\theta$-form $O=\sum_{k=0}^{M}x^kO_k(\theta_x)$ of span $M$. Substituting $x=t^2$, $\theta_x=\theta_t/2$, we obtain the $\theta$-form $O_t=\sum_kt^{2k}O_k(\theta_t/2)$ of span $2M$, which annihilates $D_\ell(t)$. By minimality of $L^t_\ell$ we have $O_t=QL^t_\ell$ for some operator $Q$ over $\C(t)$, so all $\Psi_b$ are solutions of $O_t$. By Lemma~\ref{B:lem:vertex} the $\Psi_b$ with $\lambda_b\ne0$ satisfy the hypotheses of Lemma~\ref{B:lem:types}; there are $2K$ of them. Hence $2K\le2M$.

\emph{(c), degree and leading coefficient.} Write the recurrence as in \eqref{B:eq:rec}, with leading coefficient $c_M$. By Lemma~\ref{B:lem:dict-degf}, the lowest coefficient of $O$ is
\[
O_0(\theta)=\ff{\theta}{M+n_0}\,c_M(\theta-M)\,\bigl(\ff{\theta}{M}\bigr)^2 .
\]
Let $V_x$ be the local solution space of $L_\ell$ at $x=0$. It is spanned by the functions $\Psi_b(\sqrt x)$, and $\ord_x\Psi_b(\sqrt x)=\frac12\ord_t\Psi_b$. Here $\Psi_b(\sqrt x)$ lies in the space $\mathcal N$ of Section~\ref{B:sec:ops}, taken in the variable $x$. Indeed, $tL^t_\ell=\sum_kt^{2k}Q_{2k}(\theta)$ contains only even powers of $t$, so the part of an element of $V$ formed by the odd powers of $t$ is again a solution. Its leading exponent would be an odd root of $Q_0$, but all roots of $Q_0$ are even, so this part vanishes. By Theorem~\ref{B:thm:Dt}(iii) and the equality case of Lemma~\ref{B:lem:squeeze} (applied in the variable $x$, where the exponents are the roots of $P_0$), we get $\dim(V_x)_{\ge\nu}/(V_x)_{\ge\nu+1}=\operatorname{mult}_\nu(P_0)$ for all $\nu\in\Z$. Since $O=Q'L_\ell$, the space $V_x$ consists of solutions of $O$. Lemma~\ref{B:lem:filtration} then gives $\operatorname{mult}_\nu(P_0)\le\operatorname{mult}_\nu(O_0)$. For $\nu<0$ the factors $\ff{\theta}{M+n_0}$ and $(\ff{\theta}{M})^2$ do not vanish, so
\[
\operatorname{mult}_\nu\bigl(c_M(\theta-M)\bigr)\ge\operatorname{mult}_\nu(P_0)\qquad(\nu<0).
\]
The negative roots of $P_0$, with multiplicity, are $-b(\ell-b)$ for $1\le b\le\ell-1$. Therefore $\prod_{b=1}^{\ell-1}(\theta+b(\ell-b))$ divides $c_M(\theta-M)$. Equivalently, $R_0(n+M)$ divides $c_M(n)$, that is, $R_0(N)$ divides the coefficient of $u_\ell(N)$ in the backward normalization. In particular $\deg c_M\ge\ell-1$.

\emph{(c), uniqueness.} Let $\mathcal R=\sum_{s=0}^Kc_s(n)S^s$ be the recurrence operator of (a), written in forward form with shift $S$, and let $\mathcal R'=\sum_{s=0}^Kc'_s(n)S^s$ be another recurrence of order $K$, valid for large $n$. The operator $c_K\mathcal R'-c'_K\mathcal R$ has order less than $K$ and annihilates $u_\ell$ for large $n$. By the order bound it is zero, so $c_Kc'_s=c'_Kc_s$ for all $s$. The coefficients $c_s$ have no common factor: dividing by one would produce a recurrence of degree $\ell-2$, contradicting the degree bound. Hence $c'_K/c_K$ is a polynomial, and $\mathcal R'=(c'_K/c_K)\mathcal R$.
\end{proof}

\begin{proof}[Proof of Corollary~\ref{B:cor:FW43}]
Theorem~\ref{B:thm:B} gives a recurrence of order $K=\cl{\ell/2}=\fl{(\ell-1)/2}+1$ and degree $\ell-1$, and shows that neither number can be lowered. The minimal recurrence is unique up to a constant by (c). For odd $\ell$, $K=\fl{\ell/2}+1$. For even $\ell$, $K=\ell/2<\fl{\ell/2}+1$, so a recurrence of order $\fl{\ell/2}+1$, as in \cite[(4.8)]{FW26}, exists but is not minimal.
\end{proof}

\begin{table}[ht]
\centering
\renewcommand{\arraystretch}{1.25}
\small
\begin{tabular}{@{}cl@{}}
\toprule
$\ell$ & $R_0(n),\ R_1(n),\dots,R_K(n)$ in $\sum_kR_k(n)u_\ell(n-k)=0$\\
\midrule
$2$ & $n+1,\quad -2(2n-1)$\\
$3$ & $(n+2)^2,\quad -(10n^2+2n-3),\quad 9(n-1)^2$\\
$4$ & $(n+3)^2(n+4),\quad -2(10n^3+31n^2+11n-12),\quad 64n(n-1)^2$\\
$5$ & $(n+4)^2(n+6)^2,\quad -(35n^4+322n^3+843n^2+400n-375),$\\
 & $(n-1)^2(259n^2+622n+45),\quad -225(n-2)^2(n-1)^2$\\
$6$ & $(n+5)^2(n+8)^2(n+9),\quad -4(14n^5+273n^4+1840n^3+4757n^2+2758n-2352),$\\
 & $4(n-1)^2(196n^3+1608n^2+3167n+444),\quad -1152(n-2)^2(n-1)^2(2n+3)$\\
$7$ & $(n+6)^2(n+10)^2(n+12)^2,$\\
 & $-(84n^6+2940n^5+38958n^4+238440n^3+634756n^2+429858n-343035),$\\
 & $(n-1)^2(1974n^4+36336n^3+213240n^2+407840n+82425),$\\
 & $-2(n-2)^2(n-1)^2(6458n^2+42646n+49875),\quad 11025(n-3)^2(n-2)^2(n-1)^2$\\
\bottomrule
\end{tabular}
\medskip
\caption{The minimal recurrences of Theorem~\ref{B:thm:B} for $2\le\ell\le7$, valid for $n\ge1$ with $u_\ell(m)=0$ for $m<0$. They were checked exactly against the values $u_\ell(n)=\sum_\lambda(f^\lambda)^2$ computed from the hook length formula (Appendix~\ref{B:app:comp}). For $\ell=2$ one recovers the Catalan recurrence; the case $\ell=4$ is \cite[(4.7)]{FW26}.}\label{B:tab:perm}
\end{table}

\begin{remark}
The leading coefficients in Table~\ref{B:tab:perm} illustrate Theorem~\ref{B:thm:B}(c): the factor $R_0(n)=\prod_{b=1}^{\ell-1}(n+b(\ell-b))$ is a square times $(n+\ell^2/4)$ when $\ell$ is even, and a square when $\ell$ is odd. Its roots record the exponents at $0$ of the Bessel--Toeplitz operator. The trailing coefficient $R_K$ records the exponent $\ell^2/8$ of the non-exponential solution at infinity when $\ell$ is even. The characteristic roots $(\ell-2j)^2$ are the squared critical values of $|\tr U|$ on $U(\ell)$, attained at matrices with $j$ eigenvalues $-1$ and $\ell-j$ eigenvalues $+1$, up to a common phase.
\end{remark}

\begin{remark}\label{B:rem:asymp}
We note without proof the asymptotic structure that one expects from the formal solutions of $L_\ell$ at infinity. Each characteristic root $(\ell-2b)^2$ with $b<\ell/2$ carries a formal solution $(\ell-2b)^{2n}n^{(1-(\ell-b)^2-b^2)/2}(1+O(n^{-1}))$. This is what the double Borel transform of the formal solution $e^{2(\ell-2b)\sqrt x}x^{-\gamma_b/2}$ of $L_\ell$ predicts. For $b=0$ it is consistent with Regev's asymptotics $u_\ell(n)\sim c_\ell\,\ell^{2n}n^{-(\ell^2-1)/2}$ \cite{Reg81}. Nothing in this paper depends on it.
\end{remark}

\section{The Hankel case: tableaux and matchings}\label{B:sec:CD}

\subsection{Jacobi integrals}
Let $\Delta(u)=\prod_{p<q}(u_q-u_p)$.

\begin{lemma}\label{B:lem:jacobi-int}
Let $k=2r+\varepsilon$ with $r\ge1$, $\varepsilon\in\{0,1\}$, and $w(u)=(2-u)^{\varepsilon-\frac12}(2+u)^{\frac12}$ on $(-2,2)$. Then
\[
Y_k(t)=\frac{e^{\varepsilon t}}{r!\,(2\pi)^r}\int_{[-2,2]^r}e^{t(u_1+\cdots+u_r)}\,\Delta(u)^2\prod_{p=1}^rw(u_p)\,du .
\]
Moreover, for $\ell\ge1$, $\E_{\Sp(2\ell)}e^{t\tr U}=\det[b_{i-j}-b_{i+j}]_{i,j=1}^{\ell}$, which by \eqref{B:eq:Gordon} equals $e^{-t}Y_{2\ell+1}(t)$.
\end{lemma}

\begin{proof}
The first formula is proved in \cite[Part~I]{Kei26}. It follows from \eqref{B:eq:Gordon} and the integral $b_m=\frac1\pi\int_0^\pi e^{2t\cos\varphi}\cos(m\varphi)\,d\varphi$. One then uses the product formulas for $\cos((i-j)\varphi)\pm\cos((i+j-1+\varepsilon)\varphi)$, Andreief's identity and the substitution $u=2\cos\varphi$. For the second, Weyl's integration formula expresses the eigenvalue angles $\varphi_1,\dots,\varphi_\ell\in[0,\pi]$ of $U\in\Sp(2\ell)$ through the density proportional to $\prod_{p<q}(\cos\varphi_p-\cos\varphi_q)^2\prod_p\sin^2\varphi_p$, and $\tr U=\sum_p2\cos\varphi_p$. The identities $\det[\sin(i\varphi_p)]=c\prod_p\sin\varphi_p\prod_{p<q}(\cos\varphi_q-\cos\varphi_p)$ and $b_{i-j}-b_{i+j}=\frac2\pi\int_0^\pi e^{2t\cos\varphi}\sin(i\varphi)\sin(j\varphi)\,d\varphi$, together with Andreief's identity, give the determinant \cite{BR01,Joh97}. The normalization follows from the value $1$ at $t=0$.
\end{proof}

In particular $\E_{\Sp(2\ell)}(\tr U)^{N}$ is the $N$-th moment of the law of $\sum_p2\cos\varphi_p$, and these moments vanish for odd $N$ since $-I\in\Sp(2\ell)$.

\subsection{The minimal operator of \texorpdfstring{$Y_k$}{Yk}}
We recall the theorem of \cite{Kei26}, in the form needed here. Let $k=2r+\varepsilon\ge2$ and put
\begin{equation}\label{B:eq:jacobi-data}
e_b=b^2-2br-\varepsilon b=-b(k-b),\qquad \lambda^{(k)}_j=\varepsilon+2r-4j,\qquad \gamma^{(k)}_j=(r-j)\bigl(r-j+\varepsilon-\tfrac12\bigr)+j\bigl(j+\tfrac12\bigr),
\end{equation}
for $0\le b,j\le r$.

\begin{theorem}[\cite{Kei26}]\label{B:thm:jacobi}
Let $k\ge2$. The minimal operator $L_k=\sum_{m=0}^{r+1}p_m(t)\partial^m$ of $Y_k$, normalized to be primitive, has order $r+1$, with $p_{r+1}=t^r$, $\deg p_m\le r$ and $t^{m-1}\mid p_m$. It has no singular point in $\C\setminus\{0\}$. Its solution space has a basis $y_0,\dots,y_r$, with $y_0$ a nonzero multiple of $Y_k$, such that:
\begin{enumerate}
\item[(i)] $y_b$ has an expansion in $\mathcal N$ with leading part $C'_bt^{e_b}$, $C'_b\ne0$; thus the exponents at $0$ are the distinct integers $e_0=0>e_1>\dots>e_r$;
\item[(ii)] $y_j^{(m)}(t)=C_je^{\lambda^{(k)}_jt}t^{-\gamma^{(k)}_j}\bigl((\lambda^{(k)}_j)^m+O(t^{-1})\bigr)$ as $t\to+\infty$ for all $m\ge0$, with $C_j\ne0$; and
\[
\sum_m[t^r]p_m\,\lambda^m=\chi_k(\lambda):=\prod_{j=0}^{r}\bigl(\lambda-\lambda^{(k)}_j\bigr).
\]
\end{enumerate}
\end{theorem}

\begin{proof}[Outline]
A complete proof, following \cite{Kei26}, is given in Appendix~\ref{B:app:jacobi}. It follows the scheme of Section~\ref{B:sec:unitary}. The one-variable functions are the Laplace transforms of $w$ over $[-2,2]$ and of $(2-u)^{\varepsilon-\frac12}|2+u|^{\frac12}$ over $(-\infty,-2]$. They solve a second-order equation gauge equivalent to \eqref{B:eq:T}, whose Galois group is $\SL_2$. The Toeplitz form is replaced by the Hankel form $F(g_1,\dots,g_r)=\sum_{\sigma,\tau}\sgn(\sigma\tau)\prod_pg_p^{(\sigma(p)+\tau(p)-2)}$, and the vertex integrals are taken over $(-\infty,-2]^b\times[-2,2]^{r-b}$. Laplace's method at the vertices gives (ii). The substitution $u_p=-s_p/t$ on the rays gives (i) directly, with $C'_b\ne0$; here no squeeze is needed, because the exponents are distinct. Abel's identity together with
\[
\sum_{j=0}^r\gamma^{(k)}_j+\sum_{b=0}^re_b=\binom{r+1}{2}
\]
excludes apparent singularities, and the shape follows as in the proof of Theorem~\ref{B:thm:Dt}.
\end{proof}

\subsection{Proof of Theorem~\ref{B:thm:C}}
Write $p_{m,i}=[t^i]p_m$. Comparing the coefficients of $t^M/M!$ in $L_kY_k=0$, and using $[t^M/M!]\,t^i\partial^mY_k=\ff{M}{i}\,y_k(M+m-i)$, we get
\begin{equation}\label{B:eq:egf-rec}
\sum_{s}C_s(M)\,y_k(M+s)=0\quad(M\ge0),\qquad C_s(M)=\sum_{m-i=s}p_{m,i}\,\ff{M}{i}.
\end{equation}
Since $m-1\le i\le r$, the shifts $s=m-i$ lie in $[-r,1]$, and $\deg C_s\le r$.

The top shift $s=1$ collects the terms with $i=m-1$. The indicial polynomial at $0$ is $I(\rho)=\sum_mp_{m,m-1}\ff{\rho}{m}=\prod_{b=0}^r(\rho-e_b)$; it is monic because $p_{r+1}=t^r$, and its roots are the $e_b$ by Theorem~\ref{B:thm:jacobi}(i). Since $\ff{(M+1)}{m}=(M+1)\ff{M}{m-1}$,
\[
C_1(M)=\frac{I(M+1)}{M+1}=\prod_{b=1}^{r}\bigl(M+1+b(k-b)\bigr).
\]
The bottom shift $s=-r$ has $C_{-r}(M)=p_{0,r}\ff{M}{r}=\chi_k(0)\ff{M}{r}$, and $\chi_k(0)=\prod_j(-\lambda^{(k)}_j)$ vanishes exactly when some $\lambda^{(k)}_j=0$, that is, when $\varepsilon=0$ and $r$ is even, i.e.\ $4\mid k$.

In that case the bottom shift is $s=1-r$, with
\[
C_{1-r}(M)=p_{0,r-1}\ff{M}{r-1}+p_{1,r}\ff{M}{r}=\ff{M}{r-1}\bigl(p_{0,r-1}+p_{1,r}(M-r+1)\bigr).
\]
To identify the linear factor, write $tL_k=\sum_lt^lQ_l(\theta)$ as in the proof of Theorem~\ref{B:thm:Dt}. Then $Q_{r+1}=p_{0,r}=0$ and $Q_r(\theta)=p_{0,r-1}+p_{1,r}\theta$. The solution $y_{r/2}$ has rate $0$ and behaves like $C\,t^{-\gamma}(1+O(t^{-1}))$ with $\gamma=\gamma^{(k)}_{r/2}=r^2/2$. In $tL_ky_{r/2}$ the top power $t^{r-\gamma}$ has coefficient $Q_r(-\gamma)$, so $Q_r(-r^2/2)=0$. Moreover $p_{1,r}=[\lambda]\chi_k=\prod_{j\ne r/2}(-\lambda^{(k)}_j)\ne0$. Setting $N=M+1$ and $c_i(N)=C_{1-i}(N-1)$ gives the recurrence of Theorem~\ref{B:thm:C}, with $c_0$ and $c_R$ as stated.

The leading coefficient of $C_s$ in $M$ is $p_{s+r,r}=[\lambda^{s+r}]\chi_k$. Hence $\sum_s\lc(C_s)\lambda^{s+r}=\chi_k(\lambda)$, and dividing by $\lambda$ when $4\mid k$ gives the characteristic polynomial.

\emph{Order.} Let a recurrence of order $M'$ hold for large $N$. Lemma~\ref{B:lem:dict-egf} gives a $\theta$-form of span $M'$ annihilating $Y_k$. All $y_j$ are solutions of it, by minimality of $L_k$. The number of $j$ with $\lambda^{(k)}_j\ne0$ is $r+1-[4\mid k]=R$, so Lemma~\ref{B:lem:types} gives $M'\ge R$.

\emph{Degree and leading coefficient.} By Lemma~\ref{B:lem:dict-egf}, the lowest coefficient of that $\theta$-form is
\[
\ff{\theta}{M'+n_0}\,c_{M'}(\theta-M'),
\]
where $c_{M'}$ is the leading coefficient in forward normalization. By Lemma~\ref{B:lem:filtration} applied to $y_b$, which has leading exponent $e_b$, every $e_b$ is a root of the lowest coefficient. For $b\ge1$, $e_b\le1-2r-\varepsilon<0$, so $e_b$ is a root of $c_{M'}(\theta-M')$. The $e_b$ are distinct, so $\prod_{b=1}^r(\theta-e_b)$ divides $c_{M'}(\theta-M')$. In the backward normalization this says that $c_0(N)=\prod_b(N-e_b)$ divides the coefficient of $y_k(N)$. Uniqueness follows as in the proof of Theorem~\ref{B:thm:B}(c). \qed

\begin{example}
For $k=8$ one has $r=4$, $\lambda^{(8)}=(8,4,0,-4,-8)$ and $e=(0,-7,-12,-15,-16)$. The minimal recurrence has order $4$ and degree $4$, with
\[
c_0(N)=(N+7)(N+12)(N+15)(N+16),\qquad c_4(N)=1024\,(N-1)(N-2)(N-3)(N+4),
\]
and characteristic polynomial $(\lambda-8)(\lambda-4)(\lambda+4)(\lambda+8)$. The factor $N+4=N-r+r^2/2$ comes from the rate-zero solution, whose exponent at infinity is $\gamma^{(8)}_2=8$.
\end{example}

\subsection{Proof of Theorem~\ref{B:thm:D}}
Let $Z_\ell=\E_{\Sp(2\ell)}e^{t\tr U}=e^{-t}Y_{2\ell+1}$, and let $\mu_N=\E(\tr U)^N$, so that $\mu_{2n}=m_\ell(n)$ by \eqref{B:eq:moments} and $\mu_{2n+1}=0$. If $L_{2\ell+1}(\partial)$ annihilates $Y_{2\ell+1}$, then $L'=L_{2\ell+1}(\partial+1)=\sum_mp_m(t)(\partial+1)^m$ annihilates $Z_\ell$.

The operator $L'$ has the same shape as $L_{2\ell+1}$. It also has the same indicial polynomial at $0$, because $[t^{m-1}]$ of the coefficient of $\partial^m$ is unchanged: $p_{m'}$ is divisible by $t^{m'-1}$. Its solutions are $e^{-t}y_j$, so the rates are $\lambda^{(2\ell+1)}_j-1=2(\ell-2j)=\lambda_j$, with the same $\gamma^{(2\ell+1)}_j$. The exponents at $0$ are $e_b=-b(2\ell+1-b)$, and $\sum_m[t^\ell]p'_m\lambda^m=\chi'(\lambda):=\prod_{j=0}^{\ell}(\lambda-\lambda_j)$.

As in \eqref{B:eq:egf-rec}, $\sum_{s=s_{\min}}^{1}C'_s(M)\mu_{M+s}=0$ for $M\ge0$. Here $C'_1(M)=\prod_{b=1}^{\ell}(M+1+b(2\ell+1-b))$, and $s_{\min}=-\ell$ if $\ell$ is odd, while $s_{\min}=1-\ell$ if $\ell$ is even (then $\chi'(0)=0$). In both cases $s_{\min}$ is odd.

Take $M=2N-1$ odd. Then the terms with even $s$ involve odd moments and vanish, and we obtain
\[
\sum_{i=0}^{K}d_i(N)\,m_\ell(N-i)=0,\qquad d_i(N)=C'_{1-2i}(2N-1),\qquad K=\tfrac{1-s_{\min}}{2}=\cl{\ell/2}.
\]
The extreme coefficients are computed exactly as in the proof of Theorem~\ref{B:thm:C}. The rate-zero solution for even $\ell$ has exponent $\gamma^{(2\ell+1)}_{\ell/2}=\ell(\ell+1)/2$ at infinity, which produces the factor $2N-\ell+\ell(\ell+1)/2$.

Since $\chi'(-\lambda)=(-1)^{\ell+1}\chi'(\lambda)$, the nonzero coefficients of $\chi'$ sit at the powers $\lambda^{s+\ell}$ with $s$ odd. Hence $\lc(C'_{1-2i})=[\lambda^{1-2i+\ell}]\chi'\ne0$; up to sign these are elementary symmetric functions of the numbers $4(\ell-2j)^2$. So each $d_i$ has degree exactly $\ell$, and the characteristic polynomial is $\prod_{j<\ell/2}(\lambda-4(\ell-2j)^2)$.

For the lower bounds, a recurrence of order $M'$ and degree $d$ for $m_\ell$, valid for large $N$, yields a recurrence of order $2M'$ and degree $d$ for $\mu$: replace $N$ by $M''/2$, and note that for odd $M''$ all terms vanish. Lemma~\ref{B:lem:dict-egf} and Lemma~\ref{B:lem:types}, applied to the $\ell+1-[\ell\text{ even}]$ solutions $e^{-t}y_j$ with $\lambda_j\ne0$, give $2M'\ge\ell+1-[\ell\text{ even}]$, i.e.\ $M'\ge\cl{\ell/2}$. Lemma~\ref{B:lem:filtration} applied to the $\ell$ distinct negative exponents $e_b$ shows that $\prod_b(M''-e_b)$ divides the leading coefficient of the $\mu$-recurrence. In terms of $N=M''/2$, this is $\prod_b(2N-e_b)=d_0(N)$. Uniqueness is proved as before. \qed

\begin{example}
For $\ell=4$ (perfect matchings without $5$-crossings) the minimal recurrence is
\[
\prod_{b=1}^{4}\bigl(2N+b(9-b)\bigr)\,m_4(N)+d_1(N)\,m_4(N-1)+4096\,(N-1)(N+3)(2N-1)(2N-3)\,m_4(N-2)=0,
\]
with $d_1$ of degree $4$, $d_0(N)=(2N+8)(2N+14)(2N+18)(2N+20)$, and characteristic polynomial $(\lambda-64)(\lambda-16)$.
\end{example}

\section{Distributions of traces and the transport principle}\label{B:sec:densities}

\subsection{The transport principle}
The recurrences of Sections~\ref{B:sec:perm} and~\ref{B:sec:CD} are statements about moments, and by Lemma~\ref{B:lem:dict-mom} they are equivalent to differential equations for the underlying measures. The following classical observation explains why the exponential rates of the Laplace side become the singular points of the measure side.

\begin{proposition}\label{B:prop:transport}
Let $\rho$ be a finite measure on $\R$ with compact support and $Z(t)=\int e^{tw}\,d\rho(w)$. Let $\phi$ be the ring isomorphism of Weyl algebras with $\phi(t)=-\partial_w$ and $\phi(\partial_t)=w$. Then for every $A\in\C[t]\langle\partial_t\rangle$,
\[
AZ=0\iff\phi(A)\rho=0\quad\text{(in the sense of distributions)}.
\]
If $A=\sum_ia_i(t)\partial_t^i$ with $\max_i\deg a_i=d$, then $\phi(A)$ has order at most $d$ in $\partial_w$, and its coefficient of $\partial_w^d$ is $(-1)^d\sum_i[t^d]a_i\,w^i$.
\end{proposition}

\begin{proof}
For monomials, $t^a\partial_t^iZ=\int w^i\,\partial_w^a(e^{tw})\,d\rho=(-1)^a\int e^{tw}\,d\bigl(\partial_w^a(w^i\rho)\bigr)=\int e^{tw}\,d\bigl(\phi(t^a\partial_t^i)\rho\bigr)$. The Fourier--Laplace transform is injective on compactly supported distributions. The last assertion follows from $\phi(t^a\partial^i)=(-1)^a\partial_w^aw^i$.
\end{proof}

Thus, when the minimal operator of $Z$ has leading coefficient a power of $t$ and coefficients of degree at most $d$, the measure is annihilated by an operator of order $d$. The leading coefficient of that operator is the polynomial $\sum_i[t^d]a_i\,w^i$, whose roots are the exponential rates $\lambda$ of $Z$ at infinity. For example, $D_\ell$ is the Laplace transform of the law of $2\,\mathrm{Re}\tr U$. By Theorem~\ref{B:thm:Dt} and Proposition~\ref{B:prop:transport}, this law is annihilated by an operator of order $\ell$ with leading coefficient $\pm\prod_{b=0}^{\ell}(w-2(\ell-2b))$, whose zeros are exactly the critical values of $2\,\mathrm{Re}\tr U$. For the law of $|\tr U|^2$ we argue directly from Theorem~\ref{B:thm:B}.

\subsection{Proof of Theorem~\ref{B:thm:E}}
Write the recurrence of Theorem~\ref{B:thm:B} in forward form: for $n\ge0$, $\sum_{j=0}^{K}c_j(n)u_\ell(n+j)=0$ with $c_j(n)=R_{K-j}(n+K)$. By \eqref{B:eq:moments}, $u_\ell(n)$ is the $n$-th moment of the law of $|\tr U|^2$, a probability measure on $[0,\ell^2]$. For $\ell\ge2$ it is absolutely continuous, since $|\tr U|^2$ is a nonconstant real-analytic function of the eigenvalue angles; its density is $\rho_\ell$. Lemma~\ref{B:lem:dict-mom} with $P_j(s)=c_j(-s-j-1)=R_{K-j}(K-j-1-s)$ shows that $\mathcal L_\ell\rho_\ell=0$.

All $P_j$ have degree $\ell-1$, so $\mathcal L_\ell$ has order $\ell-1$. Since $\theta^{\ell-1}=w^{\ell-1}\partial_w^{\ell-1}+\cdots$, its leading coefficient is
\[
w^{\ell-1}\sum_j\lc(P_j)\,w^j=(-1)^{\ell-1}w^{\ell-1}\sum_j\lc(R_{K-j})\,w^j=(-1)^{\ell-1}w^{\ell-1}\prod_{i=0}^{K-1}\bigl(w-(\ell-2i)^2\bigr),
\]
by Theorem~\ref{B:thm:B}(b). The lowest and top $\theta$-coefficients $P_0$ and $P_K$ both have the maximal degree $\ell-1$, so $\mathcal L_\ell$ is Fuchsian at $0$ and $\infty$. The remaining singular points are simple zeros of the leading coefficient, hence regular singular. For $j<K$ the polynomial $R_{K-j}(n)$ contains the factor $[(n-1)\cdots(n-K+j+1)]^2$, so $P_j(\theta)=R_{K-j}(K-j-1-\theta)$ is divisible by $\ff{\theta}{K-j-1}=w^{K-j-1}\partial_w^{K-j-1}$. Hence every term $w^jP_j(\theta)$ is left divisible by $w^{K-1}$, which gives the factorization $\mathcal L_\ell=w^{K-1}\mathcal L^\circ_\ell$ and the leading coefficient of $\mathcal L^\circ_\ell$.

Let $A$ be a nonzero operator of order $d$ with polynomial coefficients such that $A\rho_\ell=0$ in the sense of distributions on $\R$. Then $w^dA$ is a $\theta$-form of $\theta$-degree $d$ annihilating $\rho_\ell$, and Lemma~\ref{B:lem:dict-mom} yields a recurrence of degree at most $d$ for $u_\ell$. Hence $d\ge\ell-1$ by Theorem~\ref{B:thm:B}(c). If $A$ and $B$ are two such annihilators of order $\ell-1$, then $\lc(B)A-\lc(A)B$ annihilates $\rho_\ell$ and has smaller order, so it vanishes.

The exponents at $0$ are the roots of $P_0(\theta)=R_K(K-1-\theta)$. By the formula for $R_K$, these are $0,1,\dots,K-2$, each twice, and $\ell^2/8-1$ when $\ell$ is even. The roots of the top coefficient $P_K(\theta)=R_0(-1-\theta)=\prod_b(b(\ell-b)-1-\theta)$ are $b(\ell-b)-1$. In the local coordinate $1/w$ at infinity, this means that the exponents are $1-b(\ell-b)$, $1\le b\le\ell-1$.

For the symplectic case, the recurrence $\sum_{s=s_{\min}}^{1}C'_s(M)\mu_{M+s}=0$ of Section~\ref{B:sec:CD} has order $1-s_{\min}=\ell+1-[\ell\text{ even}]$ and degree $\ell$ for the full moment sequence of $\tr U$. By Lemma~\ref{B:lem:dict-mom}, the law of $\tr U$ on $\Sp(2\ell)$ is annihilated, in the sense of distributions, by $A_\ell=\sum_kw^kP_k(\theta)$ with $P_k(\theta)=C'_{k+s_{\min}}(-\theta-k-1-s_{\min})$. As above, $A_\ell$ has order $\ell$, it is Fuchsian, and its leading coefficient is $\pm w^{-s_{\min}}\chi'(w)$. For $s\le0$ every term of $C'_s(M)=\sum_{m-i=s}p'_{m,i}\ff{M}{i}$ contains the factor $\ff{M}{-s}$, so $P_k(\theta)$ is divisible by $\ff{\theta}{-s_{\min}-k}$ for $k\le-s_{\min}$. Hence $A_\ell=w^{-s_{\min}}A^\circ_\ell$, where $A^\circ_\ell$ has polynomial coefficients and leading coefficient $\pm\chi'(w)=\pm\prod_{j=0}^{\ell}\bigl(w-2(\ell-2j)\bigr)$. Thus the finite singular points of $A^\circ_\ell$ lie among the critical values $2(\ell-2j)$ of $\tr U$; in particular $w=0$ is an ordinary point when $\ell$ is odd. No operator of smaller order annihilates the law, because every recurrence for $(\mu_N)$ has degree at least $\ell$, by the proof of Theorem~\ref{B:thm:D}. For tableaux, the same argument applied to \eqref{B:eq:egf-rec} shows the following. The law of $\varepsilon+\sum_pu_p$ under the probability measure proportional to $\Delta(u)^2\prod_pw(u_p)\,du$ on $[-2,2]^r$ is annihilated by a Fuchsian operator of order $r$ with polynomial coefficients and leading coefficient $\pm\chi_k(w)=\pm\prod_{j=0}^{r}(w-\lambda^{(k)}_j)$, and by no operator of smaller order. In particular $w=0$ is an ordinary point unless $4\mid k$. \qed

\begin{remark}
For $\ell=2$ the theorem says that $\rho_2$ satisfies a first-order equation with exponents $-\frac12$ at $0$ and $\frac12$ at $4$. Indeed $\rho_2(w)=\frac1{2\pi}\sqrt{(4-w)/w}$ on $[0,4]$, whose moments are the Catalan numbers $u_2(n)$. For $\ell=3$, $\rho_3$ satisfies a second-order Fuchsian equation with singular points $0,1,9,\infty$ and exponents $0,0$ at the origin.
\end{remark}

\section{A denominator-free construction and arithmetic refinements}
\label{B:sec:integral-continuant}

This section completes the coefficient-integrality question left aside in
Corollaries~\ref{B:cor:FW41} and~\ref{B:cor:FW41q}.  Integrality itself was already
asserted in \cite[Proposition~4.1]{FW26}; the point here is an explicit finite
formula, with no elimination denominators, and its consequences for the
operators and recurrences of this paper.  The construction also gives the
coefficient shape for the positive-parameter Jack family outside $\beta=2$.
Minimality for nonnegative integer twists in that larger family is proved by visibility in Theorem~\ref{V:thm:A}.

\subsection{An Ore continuant and a matching formula}

Fix an integer $\ell\ge1$, write $\theta=x\partial_x$, and take parameters
$B_0=0,B_1,\ldots,B_\ell$.  Put
\[
 w_j=(j+1)(\ell-j)\qquad(0\le j<\ell).
\]
In the Ore algebra defined by $\theta x=x(\theta+1)$ define
\begin{equation}\label{B:eq:ic-recursion}
 \mathcal C_{\ell+1}=1,\qquad \mathcal C_\ell=\theta+B_\ell,
 \qquad
 \mathcal C_j=(\theta+B_j)\mathcal C_{j+1}
                 -w_jx\mathcal C_{j+2}\quad(0\le j<\ell).
\end{equation}
A matching $M$ of the path with vertices $0,1,\ldots,\ell$ is a set of
pairwise disjoint edges $\{j,j+1\}$.  Let $V(M)$ be its incident vertices and,
for an unmatched vertex $i$, put
\[
 r_M(i)=\#\{\{j,j+1\}\in M:j>i\},\qquad
 w(M)=\prod_{\{j,j+1\}\in M}(j+1)(\ell-j).
\]
Empty products equal $1$.

\begin{theorem}[Integral matching formula]\label{B:thm:ic-matching}
With $K=\lceil\ell/2\rceil$, one has
\begin{equation}\label{B:eq:ic-matching}
 \mathcal C_0=\sum_{k=0}^{K}x^k\mathcal P_k(\theta),\qquad
 \mathcal P_k(z)=(-1)^k
 \sum_{\substack{M\text{ a path matching}\\|M|=k}}
 w(M)\prod_{i\notin V(M)}\bigl(z+B_i+r_M(i)\bigr).
\end{equation}
In particular:
\begin{enumerate}
\item $\mathcal P_k\in\mathbb Z[B_1,\ldots,B_\ell,z]$,
$\deg_z\mathcal P_k=\ell+1-2k$, and
\[
 \mathcal P_0(z)=\prod_{i=0}^{\ell}(z+B_i),\qquad
 \sum_{k=0}^K\operatorname{lc}(\mathcal P_k)\mu^{\ell+1-2k}
       =\prod_{i=0}^{\ell}(\mu-\ell+2i).
\]
\item The operator $\mathcal L=x^{-1}\mathcal C_0$ belongs to
$\mathbb Z[B_1,\ldots,B_\ell,x]\langle\partial_x\rangle$.
Its order is $\ell+1$, its leading coefficient is $x^\ell$, and, if
$\mathcal L=\sum_m a_m(x)\partial_x^m$, then
\[
 x^{m-1}\mid a_m\quad(m\ge1),\qquad
 \deg_x a_m\le\left\lfloor\frac{\ell-1+m}{2}\right\rfloor.
\]
\item If every $B_i\ge0$, all the degree bounds in (2) are equalities.
Moreover $(-1)^k\mathcal P_k(z)$ has nonnegative coefficients as a polynomial
in $z,B_1,\ldots,B_\ell$.
\end{enumerate}
\end{theorem}

\begin{proof}
Expand \eqref{B:eq:ic-recursion} recursively.  At the leftmost remaining
vertex $j$, either retain $\theta+B_j$, or select the edge $\{j,j+1\}$,
contributing $-w_jx$ and deleting both endpoints.  Thus the resulting terms
are indexed precisely by path matchings.  Moving every $x$ to the left
increases an unmatched factor $\theta+B_i$ by one for each selected edge
strictly to its right.  This proves \eqref{B:eq:ic-matching} and integrality.
Every summand for a $k$-matching is monic of degree $\ell+1-2k$, with sign
$(-1)^k$ and positive weight.  Consequently that degree is exact.

The polynomial $\sum_k\operatorname{lc}(\mathcal P_k)\mu^{\ell+1-2k}$
is the ordinary tridiagonal determinant with zero diagonal and products
$w_j$ of opposite off-diagonal entries.  It is the characteristic polynomial
of $v\partial_u+u\partial_v$ on homogeneous polynomials of degree $\ell$
in $u,v$: in the monomial basis the opposite entries have product
$(j+1)(\ell-j)$.  The basis
$(u+v)^{\ell-i}(u-v)^i$ diagonalizes this operator, with eigenvalues
$\ell-2i$.  This proves the claimed characteristic polynomial.

Since $B_0=0$, $\mathcal P_0(0)=0$.  The identity
$\theta^d=\sum_{m=0}^d S(d,m)x^m\partial_x^m$, with integral Stirling
numbers $S(d,m)$, now proves that $x^{-1}\mathcal C_0$ has polynomial
coefficients.  Indeed its only potentially negative power comes from
$x^{-1}\mathcal P_0(\theta)$, whose constant term is zero.  The bounds
in (2) follow directly from $\deg\mathcal P_k=\ell+1-2k$.

For exactness let $k_* =\lfloor(\ell+1-m)/2\rfloor$.  The coefficient
of the largest possible power $x^{k_*+m-1}$ in $a_m$ comes solely from
$\mathcal P_{k_*}$.  If its degree is $m$, that coefficient is its nonzero
leading coefficient.  If its degree is $m+1$ and $m\ge1$, its leading
term contributes $S(m+1,m)=\binom{m+1}{2}>0$ after multiplication by
$(-1)^{k_*}$, and all other contributions have the same nonnegative sign.
The remaining case is $m=0$, $\ell$ even.  A maximal matching leaving
vertex $0$ unmatched contributes a positive multiple of $z+K$ to
$(-1)^K\mathcal P_K(z)$, so its constant coefficient is positive.
This proves (3).
\end{proof}

\begin{proposition}[Scalarization of the published matrix system]
\label{B:prop:ic-system}
Suppose $I_0,\ldots,I_\ell$ satisfy
\begin{equation}\label{B:eq:ic-system}
 (\ell-p)xI_{p+1}=(\theta+B_p)I_p-pI_{p-1}
       \qquad(0\le p\le\ell),
\end{equation}
where the terms with nonexistent indices are omitted.  Then
$\mathcal L I_\ell=0$ for the operator of Theorem~\ref{B:thm:ic-matching}.
For $B_p=p(\ell+q-p)$, $q\in\mathbb Z_{\ge0}$, this operator equals
$L_{\ell,q}$ of Theorem~\ref{B:thm:twisted}.
\end{proposition}

\begin{proof}
Set $J_p=(\ell!/p!)I_p$.  Descending induction in \eqref{B:eq:ic-system}
gives
\[
 J_p=\mathcal C_{p+1}I_\ell\qquad(0\le p\le\ell).
\]
The equation with $p=0$ is $\theta J_0-\ell xJ_1=0$, hence
$\mathcal C_0I_\ell=0$.  System \eqref{B:eq:ic-system} is the
Forrester--Wei system \cite[(3.6)]{FW26}.  For the stated specialization,
its last component is a nonzero constant multiple of $G_{\ell,q}$;
for $q=0$ use the modified initial vector in \cite[Remark~3.4(1)]{FW26}.
Theorem~\ref{B:thm:twisted} gives minimal order $\ell+1$.
Both $\mathcal L$ and $L_{\ell,q}$ have that order and leading
coefficient $x^\ell$, so they coincide.
\end{proof}

\begin{corollary}[Polynomial dependence and recurrence integrality]
\label{B:cor:ic-integrality}
For every $\ell\ge1$ the coefficients of $L_{\ell,q}$ belong to
$\mathbb Z[q,x]$.  More precisely, $P_k(\theta)$ of
Theorem~\ref{B:thm:twisted} belongs to $\mathbb Z[q,\theta]$ and
$(-1)^kP_k$ has nonnegative coefficients in $(q,\theta)$.
The recurrence coefficients
\[
 R_0^{(q)}(N)=\prod_{i=1}^{\ell-1}(N+i(\ell+q-i)),\qquad
 R_k^{(q)}(N)=P_k(N-k)
  \ff{(N-1)}{k-1}\ff{(N+q\ell-1)}{k-1}\quad(k\ge1)
\]
belong to $\mathbb Z[q,N]$ in the normalization of
Theorem~\ref{B:thm:twisted}(c).  These statements hold without an unspecified
common denominator or an additional rescaling of the leading coefficient.
\end{corollary}

\begin{proof}
Substitute $B_i=i(\ell-i)+iq$ into
Theorem~\ref{B:thm:ic-matching} and use
Proposition~\ref{B:prop:ic-system}.  The displayed recurrence formula is
the coefficient extraction already established in the proof of
Theorem~\ref{B:thm:twisted}(c).
\end{proof}

\begin{corollary}[Positive-parameter Jack family]\label{B:cor:ic-beta}
Let $\beta>0$, $q\ge0$, and
\[
 F_{\ell,\beta,q}(x)
 ={}_0F_1^{(2/\beta)}\left(q+1+\frac{\beta}{2}(\ell-1);(x)^\ell\right).
\]
With $B_p=p(q+\frac{\beta}{2}(\ell-p))$, the operator $\mathcal L$
of Theorem~\ref{B:thm:ic-matching} annihilates $F_{\ell,\beta,q}$.
It has order $\ell+1$, leading coefficient $x^\ell$, and the exact
coefficient degrees in Theorem~\ref{B:thm:A}(a).  Its coefficients are
polynomials over $\mathbb Z$ in $q,\beta/2,x$; hence they are integral
whenever $q,\beta/2$ are nonnegative integers.
This corollary concerns existence and shape; Theorem~\ref{V:thm:A} supplies order minimality for nonnegative integer twists.
\end{corollary}

\begin{proof}
The identification of the last component of \eqref{B:eq:ic-system}
with the displayed Jack hypergeometric function is
\cite[Proposition~3.3]{FW26}, with the explicit substitution
$n=\ell$ and $a+b=q-1-\frac\beta2(\ell-1)$.
Its zero-$q$ normalization is as in Remark~3.4(1).
All $B_p$ are nonnegative.  Apply
Theorem~\ref{B:thm:ic-matching} and Proposition~\ref{B:prop:ic-system}.
\end{proof}

\subsection{Exactly which twists lower the leading power in the Bessel variable?}

\begin{theorem}[Pythagorean classification of the exceptional twists]
\label{B:thm:ic-pythagorean}
Let $\ell\ge1$ and $q\in\mathbb Z_{\ge0}$.  In the normalization of
Theorem~\ref{B:thm:Dtq}, the leading coefficient of the minimal operator of
$D_{\ell,q}(t)$ is $t^\ell$ precisely when
\[
 q=0\qquad\hbox{or}\qquad q>0\ \hbox{and}\ \ell^2+q^2\ \hbox{is a square}.
\]
In the latter case the unique index carrying exponent zero is
\[
 i_0=\frac{\ell+q-\sqrt{\ell^2+q^2}}{2},\qquad 0<i_0<\frac\ell2.
\]
Consequently all positive exceptional pairs are exactly
\[
 (\ell,q)=d(m^2-n^2,2mn)\quad\hbox{or}\quad
 (\ell,q)=d(2mn,m^2-n^2),
\]
where $d\ge1$, $m>n\ge1$, $\gcd(m,n)=1$, and $m,n$ have opposite parity.
For every other positive twist the leading coefficient is $t^{\ell+1}$.
\end{theorem}

\begin{proof}
Theorem~\ref{B:thm:Dtq} says that the lower power occurs exactly when
\[
 q(\ell-2i)-2i(\ell-i)=0
\]
for some integer $i\in[0,\ell]$.  For $q=0$, the indices $0,\ell$ do
this.  For $q>0$, the roots of this quadratic in $i$ are
\[
 i_\pm=\frac{\ell+q\pm\sqrt{\ell^2+q^2}}2.
\]
The plus root is larger than $\ell$, and the minus root belongs to
$(0,\ell/2)$.  Therefore an allowed integer index exists exactly when
the square root is an integer: necessity follows directly from
$\sqrt{\ell^2+q^2}=\ell+q-2i$; for sufficiency, if $s^2=\ell^2+q^2$
then $s\equiv\ell+q\pmod2$, so $i_-\in\mathbb Z$.  This also proves
uniqueness.  The displayed parametrization is the elementary classification
of integer right triangles.  For completeness, divide a triple by its
common divisor; one leg is even and the other odd.  Writing the even leg
as $2ab$ in the factorization $(s-u)(s+u)=v^2$, the two coprime factors
$(s+u)/2$ and $(s-u)/2$ are squares $m^2,n^2$.  Thus
$u=m^2-n^2$, $v=2mn$, with coprime opposite-parity $m,n$; restoring
the common divisor gives exactly the stated pairs.
\end{proof}

\begin{example}
For $\ell=3$, the only positive exceptional twist is $q=4$; for
$\ell=4$, it is $q=3$.  For $\ell=5$, it is $q=12$, while for
$\ell=8$ the positive exceptional twists are $q=6$ and $q=15$.
The classification concerns the primitive scalar operator in $t$;
it does not change the uniform leading coefficient $x^\ell$ of the
normalized operator $L_{\ell,q}$ in $x=t^2$.
\end{example}

\section{Vertex balance and general-parameter ensembles}\label{V:overview}

\subsection{The integrals}
Let $\beta>0$ and let $U$ be distributed according to the circular $\beta$-ensemble $\mathrm{C}\beta\mathrm{E}_n$, that is, its eigenvalues $e^{i\varphi_1},\dots,e^{i\varphi_n}$ have density proportional to $\prod_{j<k}|e^{i\varphi_j}-e^{i\varphi_k}|^\beta$ on $\T^n$. For $q\in\Z_{\ge0}$ Forrester and Wei \cite[(2.7)]{FW26} write
\begin{equation}\label{V:eq:FW27}
\E\bigl[(\det U)^q\,e^{t\tr(U+U^{-1})}\bigr]=c\,t^{qn}\,G^\beta_{n,q}(t^2),\qquad G^\beta_{n,q}(x)={}_0F_1^{(2/\beta)}\Bigl(q+1+\frac{(n-1)\beta}2;(x)^n\Bigr),
\end{equation}
where ${}_0F_1^{(\alpha)}(c;(x)^n)$ is the hypergeometric function of $n$ equal arguments with Jack parameter $\alpha$ \cite{Kan93}, and $c\ne0$. For $\beta=2$ this is Gessel's Bessel--Toeplitz determinant, whose minimal equations were determined in Sections~\ref{B:sec:unitary}--\ref{B:sec:twisted}; for $q=0$ it is the generating function $\sum_N\mu_Nx^N/(N!)^2$ of the moments $\mu_N=\E|\tr U|^{2N}$. As $\beta\to0$ the eigenvalues become independent and $G^0_{n,0}(t^2)=I_0(2t)^n$.

For the Jacobi ensemble with parameters $a,b>-1$, whose eigenvalues $s_1,\dots,s_n\in[0,1]$ have density proportional to $\prod_js_j^a(1-s_j)^b\prod_{j<k}|s_j-s_k|^\beta$, the Laplace transform of the trace is \cite{Kan93}, \cite[(2.12)]{FK22}
\begin{equation}\label{V:eq:FK212}
H(x)=\E\,e^{-x\sum_js_j}={}_1F_1^{(2/\beta)}\Bigl(\frac\beta2(n-1)+a+1;\ \beta(n-1)+a+b+2;\ (-x)^n\Bigr).
\end{equation}
For $\beta=1$ the trace is Pillai's $V$ statistic of multivariate analysis; for $\beta=2$ and $a,b\in\{\pm\frac12\}$ it is, up to an affine change, the trace of a Haar-distributed orthogonal or symplectic matrix (Section~\ref{V:sec:walks}).

\subsection{What was known}
Forrester and Kumar \cite[Cor.~2.5]{FK22} showed that $H$ is, up to normalization, a component of a solution of an $(n+1)$-dimensional first-order system, and Forrester and Wei \cite[Prop.~3.3]{FW26} obtained the analogous system for ${}_0F_1^{(2/\beta)}$ by a confluent limit. Eliminating, \cite[Prop.~4.1]{FW26} asserts, for all $\beta$ and all parameters, a scalar equation
\[
\Bigl(x^n\frac{d^{n+1}}{dx^{n+1}}+q_n(x)x^{n-1}\frac{d^n}{dx^n}+\dots+q_1(x)\frac{d}{dx}+q_0(x)\Bigr)I_n(x)=0
\]
in which ``the degree in $x$ of the coefficient of each $d^j/dx^j$ is $[(n-1+j)/2]$'', with integer coefficients when the parameters $B_j$ of the system are integers. The proof displays the first three elimination steps and then states that the procedure can be repeated. Bornemann \cite[App.~B]{FW26} argued that the order $n+1$ is optimal for $\beta=2$ and $q\in\{0,n\}$; no statement about the order is made for $\beta\ne2$. For the moments, \cite[Prop.~4.4]{FW26} shows that the elimination produces, for general parameters, a recurrence with $\fl{(n-1)/2}+2$ terms, and for $\beta=2$ they give the corresponding form ``by inspection'' \cite[Prop.~4.3]{FW26}; from the cases they computed they observe that the degree of the coefficients is $\ell-1$, which they call difficult to anticipate. Minimality is not addressed. Forrester and Kumar \cite[Prop.~1.2]{FK22} give a vector recurrence for the moments of the Jacobi trace, and in \cite[Remark~3.8]{FK22} they note that eliminating their system for the density of the trace gives a linear differential equation of order $N$, written out for $N=2$.

The determinant method in Sections~\ref{B:sec:unitary}--\ref{B:sec:twisted} extends to arbitrary positive Dyson parameters through visibility and an explicit Kac recursion. The results below also cover Jacobi trace laws, orthogonal-group moments, and decaying Bessel products.

\subsection{Main results}
Throughout, $\theta$ is the Euler operator of the variable at hand, $K=\cl{n/2}$, and for an operator $\sum_kx^kP_k(\theta)$ we call the polynomials $P_k$ its $\theta$-coefficients. Put, for $0\le p,b\le n$,
\begin{equation}\label{V:eq:data}
\varepsilon_p=p\Bigl(q+\frac\beta2(n-p)\Bigr),\qquad\lambda_b=2(n-2b),\qquad\gamma_b=\frac n2+\frac\beta2\Bigl(\binom{n-b}2+\binom b2\Bigr).
\end{equation}

\begin{theorem}[circular ensembles]\label{V:thm:A}
Let $\beta>0$, $n\ge1$ and $q>-\frac12$ real, and $G=G^\beta_{n,q}$.
\begin{enumerate}
\item[(a)] $G$ is annihilated by an operator $L=\sum_{m=0}^{n+1}q_m(x)\partial_x^m$ with
\[
q_{n+1}=x^n,\qquad x^{m-1}\mid q_m\ (1\le m\le n+1),\qquad\deg q_m=\Bigl\lfloor\frac{n-1+m}2\Bigr\rfloor\ (0\le m\le n+1).
\]
$L$ has no singular point in $\C\setminus\{0\}$, and it is primitive if $q\ge0$. The coefficients of the $q_m$ lie in $\Z[B_1,\dots,B_n]$, where $B_p=p(q+1+\frac\beta2(n-p))$; in particular they are integers when all $B_p$ are.
\item[(b)] $xL=\sum_{k=0}^{K}x^kP_k(\theta)$ with $\deg P_k=n+1-2k$,
\[
P_0(\theta)=\prod_{p=0}^{n}(\theta+\varepsilon_p),\qquad\sum_{k=0}^{K}\lc(P_k)\,\mu^{n+1-2k}=\prod_{b=0}^{n}(\mu-n+2b),
\]
and, for even $n$, $P_K(\theta)=\lc(P_K)\bigl(\theta+\frac12(\gamma_{n/2}+qn)\bigr)$. The exponents of $L$ at $x=0$ are $-\varepsilon_0,\dots,-\varepsilon_n$.
\item[(c)] At $x=\infty$, $L$ has the formal solutions $e^{\lambda_b\sqrt x}\,x^{-(\gamma_b+qn)/2}\bigl(1+O(x^{-1/2})\bigr)$, $0\le b\le n$.
\item[(d)] If $q\in\Z_{\ge0}$, then $G$ satisfies no linear differential equation of order at most $n$ over $\C(x)$; $L$ is its minimal operator, and every operator of order $n+1$ with polynomial coefficients annihilating $G$ is a polynomial multiple of $L$. The minimal operator of $\E[(\det U)^qe^{t\tr(U+U^{-1})}]$ in $t$ has order $n+1$, no singular point in $\C\setminus\{0\}$, and exponents $q(n-2p)-\beta p(n-p)$, $0\le p\le n$, at $t=0$.
\end{enumerate}
\end{theorem}

\begin{corollary}\label{V:cor:FW41}
For every $\beta>0$ and all parameters $a,b$ with $q':=\frac\beta2(n-1)+a+b+1>-\frac12$, the function $I_n$ of \cite[Prop.~3.3]{FW26}, which is a constant multiple of $G^\beta_{n,q'}$ \cite[(3.10)]{FW26}, satisfies the equation of \cite[Prop.~4.1]{FW26} with the coefficient of $d^j/dx^j$ of degree exactly $\fl{(n-1+j)/2}$, and with integer coefficients when the $B_j$ of \cite[(3.3)]{FW26} are integers; the operator obtained there by elimination is the operator $L$ of Theorem~\ref{V:thm:A} for $q=q'$. If $q'\in\Z_{\ge0}$, which is the case relevant to the circular ensembles \cite[(3.13), Remark~3.4]{FW26}, the order $n+1$ is minimal. Outside the range $q'>-\frac12$ the degree assertion of \cite[Prop.~4.1]{FW26} can fail (Remark~\ref{V:rem:FWfalse}).
\end{corollary}

\begin{theorem}[moments]\label{V:thm:B}
Let $\beta>0$, $n\ge2$, $q=0$, and let $P_0,\dots,P_K$ be as in Theorem~\ref{V:thm:A}. Put $R_0(N)=\prod_{i=1}^{n-1}\bigl(N+\frac\beta2i(n-i)\bigr)$ and $R_k(N)=P_k(N-k)\bigl[(N-1)(N-2)\cdots(N-k+1)\bigr]^2$ for $1\le k\le K$.
\begin{enumerate}
\item[(a)] $\sum_{k=0}^{K}R_k(N)\,\mu_{N-k}=0$ for all $N\ge1$, with $\mu_m=0$ for $m<0$. Each $R_k$ has degree exactly $n-1$.
\item[(b)] The characteristic polynomial $\sum_k\lc(R_k)\lambda^{K-k}$ is $\prod_{j<n/2}(\lambda-(n-2j)^2)$, independent of $\beta$, and $R_K(N)=(-1)^K\prod_{j<n/2}(n-2j)^2\cdot[(N-1)\cdots(N-K+1)]^2$ for odd $n$, while for even $n$ the last factor is multiplied by $N-K+\frac12\gamma_{n/2}$.
\item[(c)] Every recurrence satisfied by $(\mu_N)$ for all large $N$ has order at least $K$ and degree at least $n-1$, and its leading coefficient, written as the coefficient of $\mu_N$, is divisible by $R_0(N)$. Every recurrence of order $K$ is a polynomial multiple of the one in (a).
\item[(d)] Let $q\in\Z_{\ge1}$ and $a_q(N)=\E\bigl[(\det U)^q(\tr U)^N\,\overline{\tr U}^{\,N+qn}\bigr]$. Then $a_q$ satisfies a recurrence of order $K$ and degree $n-1$ with leading coefficient $\prod_{i=1}^{n-1}(N+\varepsilon_i)$, and no recurrence of smaller order. If no $\varepsilon_i$, $1\le i\le n-1$, lies in $\{1,2,\dots,qn\}$ (this excludes only finitely many $\beta$ for given $n,q$), then every recurrence has degree at least $n-1$ and leading coefficient divisible by $\prod_{i=1}^{n-1}(N+\varepsilon_i)$. If no $\varepsilon_i$ lies in $\{qn-K+1,\dots,qn\}$, then every recurrence of order $K$ has degree at least $n-1$, and the recurrence of order $K$ is unique up to a polynomial factor.
\end{enumerate}
\end{theorem}

For $\beta=2$, parts (a)--(c) recover Theorem~\ref{B:thm:B}. Part (d) gives the generic-parameter twisted degree bound. Theorem~\ref{vtd:even-degree} below also covers every positive even $\beta$, including all exceptional integer twists, while Theorem~\ref{vtd:odd-degree-drop} gives an infinite family where the degree drops. Corollary~\ref{V:cor:density} below gives the minimal equation, of order $n-1$, for the density of $|\tr U|^2$.

Let $\varepsilon^J_p=p\bigl(a+b+1+\frac\beta2(2n-p-1)\bigr)$ and $\gamma^J_k=(n-k)(a+1)+k(b+1)+\beta\bigl(\binom{n-k}2+\binom k2\bigr)$, and let $(\star)$ denote the condition $a+b+1+\frac\beta2(n-1)>0$.

\begin{theorem}[Jacobi ensembles]\label{V:thm:C}
Let $\beta>0$, $a,b>-1$, $n\ge1$, and let $m_N=\E(\sum_js_j)^N$ and $\rho^J$ be the moments and the density of the trace.
\begin{enumerate}
\item[(a)] $H$ satisfies no linear differential equation of order at most $n$. It is annihilated by $L^J=\sum_{m=0}^{n+1}q_m\partial_x^m$ with $q_{n+1}=x^n$, $x^{m-1}\mid q_m$ and $\deg q_m=n$ for $1\le m\le n+1$, and $\deg q_0=n-1$; $L^J$ has no singular point in $\C\setminus\{0\}$, and $xL^J=\sum_{k=0}^{n}x^kP^J_k(\theta)$ with $P^J_0=\prod_{p=0}^n(\theta+\varepsilon^J_p)$, $\sum_k\lc(P^J_k)\mu^{n+1-k}=\mu\prod_{k=1}^n(\mu+k)$ and $P^J_n=n!(\theta+\gamma^J_0)$. Its exponents at $0$ are $-\varepsilon^J_p$, and at infinity it has the formal solutions $e^{-kx}x^{-\gamma^J_k}(1+O(x^{-1}))$, $0\le k\le n$. Under $(\star)$, $L^J$ is primitive; it is then the minimal operator of $H$, and every operator of order $n+1$ with polynomial coefficients annihilating $H$ is a polynomial multiple of it. Without $(\star)$, $L^J$ may be divisible by $x$.
\item[(b)] With $c_0(N)=\prod_{p=1}^n(N+\varepsilon^J_p)$ and $c_k(N)=(-1)^kP^J_k(N-k)\,(N-1)\cdots(N-k+1)$,
\[
\sum_{k=0}^{n}c_k(N)\,m_{N-k}=0\qquad(N\ge1),
\]
all $c_k$ have degree $n$, and the characteristic polynomial is $\prod_{k=1}^n(\lambda-k)$. Every recurrence for $(m_N)$ has order at least $n$. Under $(\star)$, every recurrence has degree at least $n$, its leading coefficient is divisible by $c_0(N)$, and the recurrence of order $n$ is unique up to a polynomial factor.
\item[(c)] $\rho^J$ is annihilated by an explicit operator of order $n$ with leading coefficient $\pm w^n\prod_{k=1}^n(w-k)$. Under $(\star)$, no nonzero operator of order less than $n$ annihilates $\rho^J$.
\end{enumerate}
\end{theorem}

Under $(\star)$, part (c) shows that the equation of order $N$ obtained by elimination in \cite[Remark~3.8]{FK22} has minimal order (Remark~\ref{V:rem:FK38}).

For $\ell\ge1$ let $w_B(N)=\E_{\mathrm{SO}(2\ell+1)}(\tr U)^N$ and $w_D(N)=\E_{\mathrm{SO}(2\ell)}(\tr U)^{2N}$. By the Pieri rules, $w_B(N)$ is the number of sequences of dominant weights $0=\lambda^0,\lambda^1,\dots,\lambda^N=0$ of $\mathfrak{so}_{2\ell+1}$ with steps in $\{\pm e_1,\dots,\pm e_\ell\}$ and, when $\lambda_\ell>0$, also the step $0$; and $w_D(N)$ counts such sequences of length $2N$ for $\mathfrak{so}_{2\ell}$ with steps $\pm e_i$. The representation-theoretic connection with chamber paths is discussed by Grabiner and Magyar \cite{GM93}. Here the type $B$ count always uses the stated restriction on the zero step. Recurrences for individual small $\ell$ can be found by guessing, for instance the case $\mathrm{SO}(5)$ in \cite[A095922]{OEIS}.

\begin{theorem}[orthogonal groups and Weyl chamber walks]\label{V:thm:D}
\begin{enumerate}
\item[(a)] For $\ell\ge1$, $w_B$ satisfies a recurrence of order $\ell+1$ and degree $\ell$ whose leading coefficient is $\prod_{p=1}^{\ell}(N+p(2\ell-p))$, whose last coefficient is a constant times $(N-1)\cdots(N-\ell)$, and whose characteristic polynomial is $\prod_{j=0}^{\ell}(\lambda-2\ell-1+4j)$. Every recurrence for $w_B$ has order at least $\ell+1$ and degree at least $\ell$, with leading coefficient divisible by $\prod_p(N+p(2\ell-p))$, and every recurrence of order $\ell+1$ is a polynomial multiple of this one.
\item[(b)] For $\ell\ge2$, $w_D$ satisfies a recurrence of order $\cl{\ell/2}$ and degree $\ell$ whose leading coefficient is $\prod_{p=1}^{\ell}(2N+p(2\ell-1-p))$ and whose characteristic polynomial is $\prod_{j<\ell/2}(\lambda-4(\ell-2j)^2)$. Every recurrence for $w_D$ has order at least $\cl{\ell/2}$ and degree at least $\ell$, with leading coefficient divisible by $\prod_p(2N+p(2\ell-1-p))$, and every recurrence of order $\cl{\ell/2}$ is a polynomial multiple of this one.
\end{enumerate}
\end{theorem}

For $\ell=1$, (a) is the Riordan recurrence $(N+1)w(N)=(N-1)(2w(N-1)+3w(N-2))$; for $\ell=2$, (b) is the recurrence of the squared Catalan numbers. Table~\ref{V:tab:walks} lists the cases $\ell\le4$. The symplectic case is Theorem~\ref{B:thm:D}, and it is recovered by the same argument.

\begin{theorem}[even $\beta$: the vertex integrals]\label{V:thm:E}
Let $\beta=2m$ with $m\in\Z_{\ge1}$ and $q\in\Z_{\ge0}$. The minimal operator of $\E[(\det U)^qe^{t\tr(U+U^{-1})}]$ has a basis of solutions $\Psi^{(m,q)}_0,\dots,\Psi^{(m,q)}_n$ (the vertex integrals of Section~\ref{V:sec:even}), with $\Psi^{(m,q)}_b\sim C_be^{\lambda_bt}t^{-\gamma_b}$ as $t\to+\infty$, and
\[
\ord_{t=0}\Psi^{(m,q)}_b=-f_{m,(b)}\qquad(0\le b\le n),
\]
where $f_{m,(0)}\le\dots\le f_{m,(n)}$ are the numbers $2mi(n-i)-q(n-2i)$, $0\le i\le n$, in increasing order. This is proved without the first-order systems of \cite{FK22,FW26}; the argument also reproves, independently of those systems, the minimality and the statements in the variable $t$ in Theorem~\ref{V:thm:A}(d) for these $\beta$ and $q$.
\end{theorem}

For the Bessel moments $c_{a,b}(k)=\int_0^\infty t^kI_0(t)^aK_0(t)^b\,dt$ with integers $b>a\ge0$, put $n=a+b$. Bailey, Borwein, Borwein and Crandall conjectured \cite[Conj.~1]{BBBC07} (see also \cite{BBC06,BBBG08}), and Borwein and Salvy \cite{BS08} proved, that these moments satisfy a recursion with $\cl{n/2}+1$ terms in steps of two, obtained from the $n$-th symmetric power of Bessel's equation. The same symmetric power governs the even moments $W_n(2N)=\sum_{k_1+\dots+k_n=N}\binom{N}{k_1,\dots,k_n}^2$ of the planar random walk with $n$ unit steps, for which Verrill \cite{Ver04} found a recurrence of order $\cl{n/2}$ and degree $n-1$.

\begin{theorem}[Bessel moments]\label{V:thm:F}
Let $\sum_{k}t^{2k}P^{(n)}_k(\theta)$ be the symmetric power operator of Section~\ref{V:sec:bessel}, with $P^{(n)}_0=\theta^{n+1}/n!$.
\begin{enumerate}
\item[(a)] $\sum_{i=0}^{\cl{n/2}}P^{(n)}_i(-k-1-2i)\,c_{a,b}(k+2i)=0$ for all $k\ge0$; the coefficient of $c_{a,b}(k)$ is $(-k-1)^{n+1}/n!$, and the coefficient of $c_{a,b}(k+2i)$ has degree $n+1-2i$.
\item[(b)] If $\sum_{s=0}^{M}p_s(k)\,c_{a,b}(k+s)=0$ for all large $k$, with polynomials $p_s$ and $p_0p_M\ne0$, then $M\ge2\cl{n/2}$ and $(k+1)^{n+1}$ divides $p_0(k)$. Every such recursion with $M=2\cl{n/2}$ is a polynomial multiple of (a). In particular (a) has minimal order and minimal degree $n+1$.
\item[(c)] Every recurrence satisfied by $(W_n(2N))_N$ has order at least $\cl{n/2}$ and degree at least $n-1$, and its leading coefficient, written as the coefficient of $W_n(2N)$, is divisible by $N^{n-1}$. Verrill's recurrence is, up to a constant factor, the unique recurrence of order $\cl{n/2}$ and degree $n-1$.
\end{enumerate}
\end{theorem}

\subsection{The vertex balance method}\label{V:ssec:method}
All the integrals above, except the Bessel moments, are Laplace transforms $F(t)=\int e^{tw}\rho(w)\,dw$ of the law $\rho$ of a trace $w=f(z)$ over a compact domain: the torus $\T^n$, or the cube $[0,1]^n$. The trace has finitely many critical configurations on the domain and its faces, the \emph{vertices}; for $\T^n$ these are the points with all coordinates in $\{\pm1\}$, and for the cube its corners. The method consists of four moves.

\emph{1. Vertex data.} Each vertex $v$ produces an exponential solution $e^{f(v)t}t^{-\gamma_v}(1+\cdots)$ of the minimal operator at infinity, with $\gamma_v$ computed from the local homogeneity of the weight at $v$ (a Mehta--Selberg exponent). In the presence of a first-order system the same exponents appear as diagonal entries of the residue of a Kac matrix in its eigenbasis, which we compute with $\mathfrak{sl}_2$ (Proposition~\ref{V:prop:cbe-formal}).

\emph{2. Balance} (Theorem~\ref{V:thm:balance}). If a basis of solutions has these exponents at infinity and satisfies crude bounds $\ord_0\Psi_v\ge-\beta_v$ at the regular singular point, then the \emph{balance defect}
\[
\Delta=\binom N2-\sum_v\gamma_v+\sum_v\beta_v
\]
bounds the total weight of the apparent singularities. When $\Delta=0$ there are none, and every crude bound is attained. For the circular ensembles the vanishing of $\Delta$ is the identity $\sum_b\gamma_b-\sum_i(\beta i(n-i)-q(n-2i))=\binom{n+1}2$, which is the trace of the residue of the Kac system (Remark~\ref{V:rem:trace}).

\emph{3. Visibility} (Section~\ref{V:sec:vis}). Each vertex leaves a point at which $\rho$ is not smooth: localizing the Fourier transform at the critical value, all non-vertex points contribute $O(\tau^{-\infty})$ by integration by parts along translations of clusters of coinciding coordinates, which leave the singular weight invariant, and the vertex contributes $c\,\tau^{-\gamma_v}e^{i\tau f(v)}$ with $c\ne0$, by a product of two one-sided Erd\'elyi expansions. Dually, if $L$ annihilates $F$, then $\rho$ is analytic off the roots of the top-degree polynomial $\chi_L$, whose degree is at most $\ord L$ (Lemma~\ref{V:lem:dual}). Hence the minimal order is at least the number of vertex values, and the span of every annihilating $\theta$-form, which is the order of the corresponding recurrence, is at least the number of nonzero vertex values.

\emph{4. Transfer} (Lemma~\ref{V:lem:transfer}). Once the minimal operator is known to have no singular point in $\C\setminus\{0\}$, every annihilating $\theta$-form is a left multiple of it by a $\theta$-form. So its lowest $\theta$-coefficient, which is essentially the leading coefficient of the recurrence, is divisible by the indicial polynomial at $0$. Each negative exponent therefore contributes one to the degree of every recurrence.

For even $\beta$ the integrals are polynomial in Bessel functions, the solution space is spanned by explicit vertex integrals, and moves 1, 2 and 4 suffice. For general $\beta$ the first-order systems of \cite{FK22,FW26} give an explicit operator through a three-term \emph{Kac recursion} (Section~\ref{V:sec:sys}), which has no apparent singularities by construction, and move 3 proves that it is minimal. For Bessel moments, move 4 is applied to the Mellin transform.

\subsection*{Organization}
Section~\ref{V:sec:ops} collects the facts on operators, including the transfer lemma and the balance theorem. Section~\ref{V:sec:vis} proves visibility. Section~\ref{V:sec:sys} derives the Kac recursions and the formal solutions at infinity. Sections~\ref{V:sec:cbe}--\ref{V:sec:bessel} prove Theorems~\ref{V:thm:A}--\ref{V:thm:F}, and Section~\ref{V:sec:open} discusses what remains open. Appendix~\ref{V:app:comp} describes the computations.
\section{Operators, \texorpdfstring{$\theta$}{theta}-forms and the balance theorem}\label{V:sec:ops}

Throughout, $\theta$ denotes the Euler operator of the variable at hand ($\theta=x\,d/dx$ or $\theta=t\,d/dt$), $\ff{z}{m}=z(z-1)\cdots(z-m+1)$, and an operator $L=\sum_{m=0}^{r}p_m\partial^m$ with polynomial coefficients and $p_r\ne0$ is \emph{primitive} if the $p_m$ have no common zero. The \emph{minimal operator} of a function (or formal series) $y$ is the monic operator of least order in $\C(x)[\partial]$ annihilating $y$; after clearing denominators it is a primitive operator, unique up to a constant factor. Every operator over $\C(x)$ annihilating $y$ is a left multiple of it.

\subsection{\texorpdfstring{$\theta$}{theta}-forms}
A \emph{$\theta$-form} is an operator $A=\sum_{k=k_0}^{k_1}x^kQ_k(\theta)$ with $Q_k\in\C[\theta]$, $k_0\le k_1$ in $\Z$, and $Q_{k_0}Q_{k_1}\ne0$. Its \emph{span} is $k_1-k_0$, its \emph{lowest coefficient} is $Q_{k_0}$ and its \emph{top coefficient} is $Q_{k_1}$. Since $\theta=x\partial$, $\partial=x^{-1}\theta$ and $x^k\theta^i=\sum_{m\le i}S(i,m)x^{k+m}\partial^m$ (Stirling numbers of the second kind), the nonzero elements of $\C[x,x^{-1}][\partial]$ are exactly the $\theta$-forms, and the $\partial$-order of $A$ is $\max_k\deg Q_k$. The commutation rule
\begin{equation}\label{V:eq:mult}
x^jM(\theta)\cdot x^kQ(\theta)=x^{j+k}M(\theta+k)Q(\theta)
\end{equation}
shows that the product of two $\theta$-forms $\sum_{j=j_0}^{j_1}x^jM_j(\theta)$ and $\sum_{k=k_0}^{k_1}x^kQ_k(\theta)$ is a $\theta$-form with lowest coefficient $M_{j_0}(\theta+k_0)Q_{k_0}(\theta)$ and top coefficient $M_{j_1}(\theta+k_1)Q_{k_1}(\theta)$. In particular spans add.

\begin{lemma}[Laurent division]\label{V:lem:laurent}
Let $L\in\C[x][\partial]$ have order $r$ and leading coefficient $cx^A$ with $c\ne0$. If $O\in\C[x,x^{-1}][\partial]$ and $O=QL$ with $Q\in\C(x)[\partial]$, then $Q\in\C[x,x^{-1}][\partial]$.
\end{lemma}

\begin{proof}
The leading coefficient of $L$ is a unit of $\C[x,x^{-1}]$, so right division by $L$ can be carried out in $\C[x,x^{-1}][\partial]$: $O=Q'L+R$ with $Q',R\in\C[x,x^{-1}][\partial]$ and $\ord R<r$. By uniqueness of right division in $\C(x)[\partial]$, $Q'=Q$ and $R=0$.
\end{proof}

\begin{lemma}[transfer]\label{V:lem:transfer}
Let $y$ have minimal operator $L$, normalized to have polynomial coefficients, and suppose that its leading coefficient is a monomial $cx^A$, so that $L$ has no singular point in $\C\setminus\{0\}$. Write $x^eL=\sum_{k=k_0}^{k_1}x^kP_k(\theta)$ for some $e\in\Z$. If a $\theta$-form $O$ annihilates $y$, then $O=Q\cdot x^eL$ for a $\theta$-form $Q=\sum_{j=j_0}^{j_1}x^jM_j(\theta)$. Consequently:
\begin{enumerate}
\item[(a)] $\operatorname{span}O\ge k_1-k_0$;
\item[(b)] the lowest coefficient of $O$ is $M_{j_0}(\theta+k_0)P_{k_0}(\theta)$, and the top coefficient is $M_{j_1}(\theta+k_1)P_{k_1}(\theta)$;
\item[(c)] if $\operatorname{span}O=k_1-k_0$, then $O=x^{j_0}M(\theta)\,x^eL=\sum_kx^{j_0+k}M(\theta+k)P_k(\theta)$ for one polynomial $M$.
\end{enumerate}
\end{lemma}

\begin{proof}
By minimality, $O=\tilde QL$ with $\tilde Q\in\C(x)[\partial]$, and $\tilde Q\in\C[x,x^{-1}][\partial]$ by Lemma~\ref{V:lem:laurent}. Put $Q=\tilde Qx^{-e}$ and apply \eqref{V:eq:mult}.
\end{proof}

\subsection{Recurrences}
A sequence $a:\Z_{\ge0}\to\C$ satisfies a recurrence of \emph{order} $M$ and \emph{degree} $d$ if
\begin{equation}\label{V:eq:rec}
\sum_{s=0}^{M}c_s(N)\,a(N+s)=0\qquad(N\ge n_0)
\end{equation}
with polynomials $c_s$ of degree at most $d$ and $c_0c_M\ne0$. The \emph{leading coefficient} is $c_M$; in the backward form $\sum_{i=0}^{M}c_{M-i}(N-M)\,a(N-i)=0$ it is the coefficient of $a(N)$. Part (a) of the following lemma is Lemma~\ref{B:lem:dict-egf} ; part (b) sharpens Lemma~\ref{B:lem:dict-degf}, whose lowest coefficient contains $(\ff\theta M)^2$ instead of $\ff\theta M$, and allows a shift $r$.

\begin{lemma}[dictionaries]\label{V:lem:dict}
\begin{enumerate}
\item[(a)] Let $Y(t)=\sum_Na(N)t^N/N!$. If \eqref{V:eq:rec} holds, then $Y$ is annihilated by a $\theta$-form of span $M$ with lowest coefficient $\ff{\theta}{M+n_0}\,c_M(\theta-M)$. Conversely, if $\sum_{k=0}^{M}t^kP_k(\theta)$ annihilates $Y$, then $\sum_kP_k(N-k)\ff{N}{k}\,a(N-k)=0$ for all $N\ge0$, terms with $N-k<0$ being omitted.
\item[(b)] Let $r\in\Z_{\ge0}$ and $G(x)=\sum_Na(N)x^N/(N!\,(N+r)!)$. If \eqref{V:eq:rec} holds, then $G$ is annihilated by a $\theta$-form of span $M$ with lowest coefficient $\ff{\theta}{M+n_0}\,\ff{(\theta+r)}{M}\,c_M(\theta-M)$. Conversely, if $\sum_{k=0}^{M}x^kP_k(\theta)$ annihilates $G$, then $\sum_kP_k(N-k)\ff Nk\ff{(N+r)}k\,a(N-k)=0$ for all $N\ge0$.
\end{enumerate}
\end{lemma}

\begin{proof}
For (a) see Lemma~\ref{B:lem:dict-egf}. For (b) put $F(\theta)=\ff\theta{M+n_0}\ff{(\theta+r)}M$ and
\[
Q_i(\theta)=\frac{F(\theta+i)\,c_{M-i}(\theta+i-M)}{\ff{(\theta+i)}{i}\;\ff{(\theta+i+r)}{i}}\qquad(0\le i\le M).
\]
Since $i\le M$, $\ff{(\theta+i)}i$ divides $\ff{(\theta+i)}{M+n_0}$ and $\ff{(\theta+i+r)}i$ divides $\ff{(\theta+i+r)}M$, so the $Q_i$ are polynomials, with $Q_0=F(\theta)c_M(\theta-M)$ and $Q_M\ne0$. Because $\ff Ni\,(N-i)!=N!$ and $\ff{(N+r)}i\,(N+r-i)!=(N+r)!$, the coefficient of $x^N$ in $\sum_ix^iQ_i(\theta)G$ is
\[
\sum_{i=0}^{M}Q_i(N-i)\frac{a(N-i)}{(N-i)!\,(N-i+r)!}=\frac{F(N)}{N!\,(N+r)!}\sum_{s=0}^{M}c_s(N-M)\,a(N-M+s).
\]
For $N<M$, terms with negative sequence indices are omitted; each retained term vanishes because its factorial denominators are nonzero and $F(N)=0$. For $N\ge M$, the displayed expression vanishes for $N-M\ge n_0$ by \eqref{V:eq:rec}, and for $N<M+n_0$ because $F(N)=0$. The converse statements follow by comparing coefficients.
\end{proof}

The transfer lemma turns the lowest coefficient of the minimal operator into a lower bound for the degree of every recurrence. We record the form in which it will be used.

\begin{corollary}\label{V:cor:degree}
In the situation of Lemma~\ref{V:lem:dict}(a), or of Lemma~\ref{V:lem:dict}(b) with $r=0$, let the minimal operator $L$ of $Y$ \textup{(}resp.\ $G$\textup{)} be as in Lemma~\ref{V:lem:transfer}, with $k_1-k_0=K$. Then every recurrence \eqref{V:eq:rec} satisfied by $a$ has order $M\ge K$, and the polynomial $\prod_{\nu}(\theta-\nu)^{m_\nu}$, taken over the roots $\nu\notin\Z_{\ge0}$ of $P_{k_0}$ with their multiplicities $m_\nu$, divides $c_M(\theta-M)$. In particular $\deg c_M$ is at least the number of roots of $P_{k_0}$ outside $\Z_{\ge0}$, counted with multiplicity.
\end{corollary}

\begin{proof}
The operator $O$ of Lemma~\ref{V:lem:dict} has span $M$, so $M\ge K$ by Lemma~\ref{V:lem:transfer}(a). By Lemma~\ref{V:lem:transfer}(b), $P_{k_0}(\theta)$ divides the lowest coefficient of $O$. The factors $\ff{\theta}{M+n_0}$ and $\ff{\theta}{M}$ vanish only on $\Z_{\ge0}$.
\end{proof}

\subsection{The balance theorem}
Let $\mathcal N$ be the space of functions on a slit punctured disc $\{0<|t|<\delta\}\setminus(-\delta,0]$ of the form $\sum_{\nu\ge\nu_0}t^\nu P_\nu(\log t)$ with $\nu\in\Z$, polynomials $P_\nu$ of bounded degree, and absolute convergence; for $f\ne0$ in $\mathcal N$, $\ord f$ is the least $\nu$ with $P_\nu\ne0$. We use three facts from Section~\ref{B:sec:ops}, proved in Lemmas~\ref{B:lem:squeeze}, \ref{B:lem:abel}, and~\ref{B:lem:top}. (i) \emph{Squeeze}: if $L$ of order $N$ is Fuchsian at $0$ with local solution space in $\mathcal N$, its exponents $\nu_1\ge\dots\ge\nu_N$ are integers, and for every basis $f_1,\dots,f_N$ with $\ord f_1\ge\dots\ge\ord f_N$ one has $\ord f_i\le\nu_i$; hence $\sum_i\ord f_i\le\sum_i\nu_i$ with equality only if $\ord f_i=\nu_i$ for all $i$. (ii) \emph{Abel}: if $W$ is a Wronskian of $L=\sum p_m\partial^m$ (primitive), then $W'/W=-p_{N-1}/p_N$; at a regular singular point the residue is $\sum_i\nu_i-\binom N2$; at a zero $\tau$ of $p_N$ where all solutions are holomorphic, the residue is an integer $\kappa_\tau\ge1$. (iii) \emph{Top degree}: if $d=\max\deg p_m$ and $L$ has solutions $y_j$ with $y_j^{(m)}=C_je^{\lambda_jt}t^{-\gamma_j}(\lambda_j^m+O(t^{-1/2}))$, $C_j\ne0$, for $N$ distinct $\lambda_j$, then $[t^d]p_N\ne0$ and $\sum_m[t^d]p_m\lambda^m=[t^d]p_N\prod_j(\lambda-\lambda_j)$.

\begin{theorem}[balance]\label{V:thm:balance}
Let $L=\sum_{m=0}^{N}p_m(t)\partial^m$ be primitive, and assume:
\begin{enumerate}
\item[(B1)] every solution of $L$ continues analytically along every path in $\C\setminus\{0\}$;
\item[(B2)] $L$ is Fuchsian at $0$, with local solution space contained in $\mathcal N$;
\item[(B3)] $L$ has a basis of solutions $\Psi_1,\dots,\Psi_N$ on some interval $(T,\infty)$ with
\[
\Psi_j^{(m)}(t)=C_je^{\lambda_jt}t^{-\gamma_j}\bigl(\lambda_j^m+O(t^{-1/2})\bigr)\qquad(0\le m\le N),
\]
where $C_j\ne0$, $\gamma_j\in\R$, and the $\lambda_j$ are pairwise distinct;
\item[(B4)] $\ord\Psi_j\ge-\beta_j$ for integers $\beta_j$, where $\Psi_j$ is continued along $(0,\infty)$ to a slit disc at $0$.
\end{enumerate}
Define the \emph{balance defect}
\[
\Delta=\binom N2-\sum_{j=1}^{N}\gamma_j+\sum_{j=1}^{N}\beta_j .
\]
Then $\sum_{\tau}\kappa_\tau\le\Delta$, the sum running over the zeros $\tau\ne0$ of $p_N$. In particular $\Delta\ge0$, and $p_N$ has at most $\Delta$ distinct zeros in $\C\setminus\{0\}$. If $\Delta=0$, then $p_N=ct^A$, $\ord\Psi_j=-\beta_j$ for every $j$, and the exponents of $L$ at $0$ are $-\beta_1,\dots,-\beta_N$.
\end{theorem}

\begin{proof}
By (B1), every zero $\tau\ne0$ of $p_N$ is an apparent singularity, since a disc around $\tau$ avoiding $0$ is simply connected. Let $W$ be the Wronskian of $\Psi_1,\dots,\Psi_N$. The rational function $W'/W=-p_{N-1}/p_N$ has simple poles only, at $0$ by Fuchs' criterion and at the apparent points by Abel's lemma, so $W'/W=Q(t)+\kappa_0/t+\sum_\tau\kappa_\tau/(t-\tau)$ with $Q\in\C[t]$ and $\kappa_0=\sum_i\nu_i-\binom N2$, where $\nu_i$ are the exponents at $0$. Hence $W=c\,e^{\int Q}\,t^{\kappa_0}\prod_\tau(t-\tau)^{\kappa_\tau}$. By (B3), expanding the determinant,
\[
W(t)=\Bigl(\prod_jC_j\Bigr)e^{(\sum_j\lambda_j)t}\,t^{-\sum_j\gamma_j}\Bigl(\prod_{i<j}(\lambda_j-\lambda_i)+O(t^{-1/2})\Bigr),
\]
and the Vandermonde product is nonzero. Comparing the two expressions as $t\to+\infty$ gives $\int Q=(\sum_j\lambda_j)t+\mathrm{const}$ and $\kappa_0+\sum_\tau\kappa_\tau=-\sum_j\gamma_j$. Therefore
\[
\sum_\tau\kappa_\tau=\binom N2-\sum_j\gamma_j-\sum_i\nu_i\le\binom N2-\sum_j\gamma_j-\sum_j\ord\Psi_j\le\Delta,
\]
by the squeeze (i) and (B4). If $\Delta=0$, all inequalities are equalities: there are no apparent singularities, so $p_N$ has no zero in $\C\setminus\{0\}$; $\ord\Psi_j=-\beta_j$ for every $j$, because each term was bounded separately; and by the equality case of the squeeze the exponents are the numbers $\ord\Psi_j$.
\end{proof}

We call a family of solutions satisfying (B3)--(B4) with $\Delta=0$ \emph{balanced}. The balance theorem is the abstract form of the squeeze of Sections~\ref{B:sec:unitary}--\ref{B:sec:twisted}: crude bounds at the regular singular point and exact data at the irregular point determine the local exponents and exclude apparent singularities as soon as the defect vanishes. The following two lemmas convert this information into the shape of the operator.

\begin{lemma}[shape at the origin]\label{V:lem:shape-t}
Let $L=\sum_{m=0}^{N}p_m\partial^m$ be primitive with $p_N=ct^A$, Fuchsian at $0$. If $0$ is not an exponent at $0$, then $A=N$. If $0$ is an exponent and $1$ is not, then $A=N-1$. In both cases $\ord_0p_m\ge m+A-N$ for all $m$.
\end{lemma}

\begin{proof}
By Fuchs' criterion $\ord_0p_m\ge m+A-N$, and the indicial polynomial is
\[
I(\rho)=\sum_m[t^{m+A-N}]p_m\,\ff{\rho}{m}.
\] Primitivity gives some $p_m(0)\ne0$, hence $A\le N$. If $A\le N-1$, the term $m=0$ of $I$ vanishes, so $I(0)=0$; if $A\le N-2$, also the term $m=1$ vanishes and $I(1)=0$. This proves the first assertion, and that $A\ge N-1$ in the second case. If $0$ is an exponent and $A=N$, then $\ord_0p_m\ge m$ for all $m$ and $p_0(0)=I(0)=0$, so all $p_m$ vanish at $0$, contradicting primitivity.
\end{proof}

\begin{lemma}[the $\eta$-identity]\label{V:lem:eta}
Let $\Theta=\sum_{k=0}^{K}t^{2k}Q_{2k}(\theta)$ with $\deg Q_{2k}\le N-2k$, and suppose
\[
\chi(\lambda):=\sum_{k}[\theta^{N-2k}]Q_{2k}\,\lambda^{N-2k}=c\prod_{b=1}^{N}(\lambda-\lambda_b),\qquad c\ne0,
\]
with pairwise distinct $\lambda_b$. If for every $b$ the operator $\Theta$ has a formal solution $e^{\lambda_bt}t^{-g_b}(1+\sum_{i\ge1}c_{b,i}t^{-i})$, then
\begin{equation}\label{V:eq:eta-gen}
\sum_{k}[\theta^{N-2k-1}]Q_{2k}\,\lambda^{N-2k-1}=c\sum_{b=1}^{N}g_b\prod_{b'\ne b}(\lambda-\lambda_{b'})-\frac\lambda2\chi''(\lambda).
\end{equation}
\end{lemma}

\begin{proof}
Since $\theta(e^{\lambda t}t^{-g}h)=e^{\lambda t}t^{-g}(\lambda t-g+\theta)h$, we must expand $Q_{2k}(\lambda t+\theta-g)$, applied to $h=1+c_1t^{-1}+\cdots$. Write $d_k=N-2k$ and $Q_{2k}=a_k\theta^{d_k}+s_k\theta^{d_k-1}+\cdots$. In $(\lambda t+\theta-g)^{d_k}$ the terms with exactly one factor $\theta-g$, in position $p$ from the left, contribute $\lambda^{d_k-1}t^{d_k-1}(d_k-p-g)$; their sum is $\lambda^{d_k-1}t^{d_k-1}(\binom{d_k}2-d_kg)$. Collecting the coefficient of $t^{N-1}$ in $e^{-\lambda t}t^{g}\,\Theta(e^{\lambda t}t^{-g}h)$, and noting that the contribution of $c_1$ is $c_1\chi(\lambda)$, we obtain for $\lambda=\lambda_b$, $g=g_b$:
\[
\sum_ks_k\lambda_b^{d_k-1}+\frac{\lambda_b}2\chi''(\lambda_b)-g_b\chi'(\lambda_b)=0 .
\]
The polynomial $\sum_ks_k\lambda^{d_k-1}+\frac\lambda2\chi''(\lambda)$ has degree at most $N-1$ and takes the value $g_b\chi'(\lambda_b)=c\,g_b\prod_{b'\ne b}(\lambda_b-\lambda_{b'})$ at the $N$ points $\lambda_b$. Lagrange interpolation gives \eqref{V:eq:eta-gen}.
\end{proof}
\section{Visibility: vertices leave singularities in the density}\label{V:sec:vis}

All generating functions in Sections~\ref{V:sec:sys}--\ref{V:sec:walks} are Laplace transforms $F(t)=\langle\rho,e^{tw}\rangle$ of compactly supported distributions $\rho$ (the law of a trace, possibly twisted by a character). This section shows that every critical value of the trace that comes from a vertex of the integration domain is a point where $\rho$ fails to be smooth, and that such points bound the order of every differential equation for $F$ from below.

\subsection{Dual singular support}

\begin{lemma}[dual singular support]\label{V:lem:dual}
Let $\rho$ be a compactly supported distribution on $\R$ and $F(t)=\langle\rho_w,e^{tw}\rangle$, an entire function. Let $L=\sum_{m=0}^{r}p_m(t)\partial_t^m\ne0$ have polynomial coefficients, $d=\max_m\deg p_m$, and
\[
\chi_L(\lambda)=\sum_{m=0}^{r}[t^d]p_m\,\lambda^m ,
\]
a nonzero polynomial of degree at most $r$. If $LF=0$, then $\rho$ is real-analytic on $\R\setminus\chi_L^{-1}(0)$.
\end{lemma}

\begin{proof}
For a polynomial $p$, $p(t)\partial_t^mF(t)=\langle\rho,w^mp(\partial_w)e^{tw}\rangle=\langle p(-\partial_w)(w^m\rho),e^{tw}\rangle$. Hence $LF=\langle\hat L\rho,e^{tw}\rangle$ with $\hat L=\sum_mp_m(-\partial_w)\circ w^m$. The distribution $\hat L\rho$ has compact support and vanishing Fourier--Laplace transform, so $\hat L\rho=0$, since the Fourier transform is injective on distributions with compact support \cite[\S7.1]{Hor90}. Writing $p_m=\sum_jc_{m,j}t^j$ we have $\hat L=\sum_{j\le d}(-1)^j\partial_w^j\circ a_j(w)$ with $a_j(w)=\sum_mc_{m,j}w^m$, so the coefficient of $\partial_w^d$ in $\hat L$ is $(-1)^da_d=(-1)^d\chi_L$. If $d=0$, then $\hat L$ is multiplication by $\chi_L$ and $\rho$ vanishes off the zeros of $\chi_L$. If $d\ge1$, then on an open interval $I$ on which $\chi_L$ has no zero, $\hat L\rho=0$ is equivalent to a system $U'=A(w)U$ for $U=(\rho,\rho',\dots,\rho^{(d-1)})\in\mathcal D'(I)^d$, with $A$ real-analytic on $I$. If $\Phi$ is a fundamental matrix of $Y'=AY$, then $(\Phi^{-1}U)'=0$ in $\mathcal D'(I)$, so $\Phi^{-1}U$ is constant and $U=\Phi\cdot\mathrm{const}$ is real-analytic.
\end{proof}

\begin{corollary}\label{V:cor:visible}
Suppose $\rho$ is not $C^\infty$ at the distinct points $w_1,\dots,w_s$.
\begin{enumerate}
\item[(a)] Every nonzero operator with polynomial coefficients annihilating $F$ has order at least $s$.
\item[(b)] Every $\theta$-form in $t$ annihilating $F$ has span at least $\#\{i:w_i\ne0\}$.
\end{enumerate}
\end{corollary}

\begin{proof}
(a) $\chi_L$ is a nonzero polynomial of degree at most $\ord L$ vanishing at every $w_i$. (b) Multiplying by a power of $t$ we may assume $A=\sum_{k=k_0}^{k_1}t^kQ_k(\theta)$ with $k_0\ge0$. Let $d_*=\max_k(k+\deg Q_k)$ and $S=\{k:k+\deg Q_k=d_*\}$. The coefficient of $t^{d_*}\partial^m$ in $A$ is $\lc(Q_{d_*-m})$ if $d_*-m\in S$ and $0$ otherwise, so $\chi_A(\lambda)=\sum_{k\in S}\lc(Q_k)\lambda^{d_*-k}=\lambda^{d_*-\max S}g(\lambda)$ with $g(0)\ne0$ and $\deg g=\max S-\min S\le k_1-k_0$. The nonzero $w_i$ are roots of $g$.
\end{proof}

Part (b) is the dual form of the types lemma of Lemma~\ref{B:lem:types}: there, distinct exponential types of solutions at infinity were counted; here we count points of the singular support of the density.

\subsection{Oscillatory integrals at a vertex}

\begin{lemma}[Erd\'elyi]\label{V:lem:erdelyi}
Let $\mu\in\{1,2\}$, $\lambda>0$, $\sigma\in\{1,-1\}$ and $J\ge1$. For $g\in C^\infty([0,\infty))$ supported in $[0,R]$ and $\tau\ge1$,
\[
\int_0^\infty e^{i\sigma\tau r^\mu}r^{\lambda-1}g(r)\,dr=\sum_{j<J}\frac{g^{(j)}(0)}{j!}\,\frac1\mu\Gamma\Bigl(\frac{\lambda+j}\mu\Bigr)e^{i\sigma\pi(\lambda+j)/(2\mu)}\,\tau^{-(\lambda+j)/\mu}+E_J(\tau),
\]
where $|E_J(\tau)|\le C\,\tau^{-(\lambda+J)/\mu}$ and $C$ depends only on $\mu,\lambda,J,R$ and on $\max_{i\le M}\|g^{(i)}\|_\infty$ for some $M=M(\mu,\lambda,J)$.
\end{lemma}

This is Erd\'elyi's lemma \cite{Erd56}; see also \cite{Won89}. That the constant depends on finitely many derivatives of $g$ (how many depends on $\lambda$) can be read off from the proof, which writes $g$ as its Taylor polynomial plus a remainder vanishing to high order at $0$ and integrates by parts in the remainder term.

Let $N_1,N_2\ge0$. For $N\ge1$ we use polar coordinates $u=r\omega$ on $\R^N$, with $du=r^{N-1}dr\,d\omega$, and simplex coordinates $u=r\omega$ on $[0,\infty)^N$, with $r=\sum_ju_j$, $\omega$ in the simplex $\Sigma_N=\{\omega\in[0,1]^N:\sum\omega_j=1\}$ and $du=r^{N-1}dr\,d\sigma(\omega)$, where $d\sigma=d\omega_1\cdots d\omega_{N-1}$ is the projected Lebesgue measure.

\begin{lemma}[vertex models]\label{V:lem:vertex-model}
Let $A_i\ge0$ be measurable on $\R^{N_i}$ \textup{(}case Q\textup{)} or on $[0,\infty)^{N_i}$ \textup{(}case L\textup{)}, positively homogeneous of degree $d_i$, integrable and not a.e.\ zero on the unit sphere \textup{(}resp.\ on $\Sigma_{N_i}$\textup{)}. Let $h$ be $C^\infty$ with compact support on $\R^{N_1}\times\R^{N_2}$. Put
\[
\gamma_i=\tfrac12(N_i+d_i)\ \text{(case Q)},\qquad\gamma_i=N_i+d_i\ \text{(case L)},
\]
and assume $\gamma_1,\gamma_2>0$ when $N_1,N_2\ge1$. Then, as $\tau\to\infty$,
\[
\int e^{i\tau\phi(u,v)}A_1(u)A_2(v)h(u,v)\,du\,dv=c\,h(0)\,\tau^{-\gamma_1-\gamma_2}\bigl(1+O(\tau^{-1/2})\bigr)+O(\tau^{-\gamma_1-\gamma_2-1/2}),
\]
where in case Q the integral is over $\R^{N_1}\times\R^{N_2}$, $\phi=-|u|^2+|v|^2$ and
\[
c=\tfrac14\Gamma(\gamma_1)\Gamma(\gamma_2)e^{i\pi(\gamma_2-\gamma_1)/2}\int_{S^{N_1-1}}A_1\int_{S^{N_2-1}}A_2 ,
\]
and in case L the integral is over $[0,\infty)^{N_1}\times[0,\infty)^{N_2}$, $\phi=\sum_ju_j-\sum_jv_j$ and
\[
c=\Gamma(\gamma_1)\Gamma(\gamma_2)e^{i\pi(\gamma_1-\gamma_2)/2}\int_{\Sigma_{N_1}}A_1\int_{\Sigma_{N_2}}A_2 .
\] If $N_i=0$, the corresponding variable, factor and $\Gamma$-factor (together with one factor $\tfrac12$ in case Q) are omitted. In particular $c\ne0$.
\end{lemma}

\begin{proof}
Take case Q with $N_1,N_2\ge1$; the other cases are identical. In polar coordinates in $u$ and in $v$ separately, using homogeneity, the integral equals
\[
\int_0^\infty\!\!\int_0^\infty e^{-i\tau r^2+i\tau s^2}r^{2\gamma_1-1}s^{2\gamma_2-1}K(r,s)\,dr\,ds,\qquad K(r,s)=\int\!\!\int A_1(\omega)A_2(\omega')h(r\omega,s\omega')\,d\omega\,d\omega'.
\]
Since $h$ is smooth with compact support and the $A_i$ are integrable on the spheres, $K$ is the restriction to $[0,\infty)^2$ of a $C^\infty$ function on $\R^2$ with compact support, and $K(0,0)=h(0)\int A_1\int A_2$. Apply Lemma~\ref{V:lem:erdelyi} in $s$ with $\mu=2$, $\sigma=1$ and $J>2\gamma_1+1$; the coefficients $\partial_s^jK(r,0)$ are smooth in $r$, and the remainder is $O(\tau^{-\gamma_2-J/2})$ uniformly in $r$. Then apply Lemma~\ref{V:lem:erdelyi} in $r$ with $\sigma=-1$ to each coefficient; the remainder contributes $O(\tau^{-\gamma_2-J/2})$, which is $O(\tau^{-\gamma_1-\gamma_2-1/2})$. The leading term is $\tfrac12\Gamma(\gamma_1)e^{-i\pi\gamma_1/2}\cdot\tfrac12\Gamma(\gamma_2)e^{i\pi\gamma_2/2}K(0,0)\,\tau^{-\gamma_1-\gamma_2}$. In case L one uses simplex coordinates and $\mu=1$. Finally $\int A_i>0$.
\end{proof}

\subsection{Non-stationary points}
The weights we meet are not smooth on the diagonals, so the usual non-stationary phase argument needs a direction in which the weight is smooth. Translations of a cluster of coinciding coordinates provide one.

\begin{lemma}[cluster translation]\label{V:lem:cluster}
Let $V\ne0$ be a constant vector field on $\R^n$, $D\subseteq\R^n$ a closed set with $D+sV=D$ for all $s\in\R$, $U\subset\R^n$ open, and $\phi$ real-analytic on $U$ with $|V\phi|\ge c_0>0$ on $U$. Let $a=\Phi\,g$, where $\Phi\ge0$ is locally integrable and invariant under the translations $z\mapsto z+sV$, and $g\in C^\infty$ has compact support in $U$. Then $\int_De^{i\tau\phi}a\,dz=O(\tau^{-M})$ for every $M$.
\end{lemma}

\begin{proof}
Write $z=z'+sV$ with $z'\perp V$; then $D=D'\times\R$ in these coordinates and $\Phi$ does not depend on $s$. Integrating by parts in $s$ on almost every line gives, without boundary terms,
\[
\int_De^{i\tau\phi}\Phi g\,dz=\frac{-1}{i\tau}\int_De^{i\tau\phi}\,\Phi\,V\Bigl(\frac{g}{V\phi}\Bigr)dz,
\]
and $V(g/V\phi)$ is again smooth with compact support in $U$. Iterate $M$ times.
\end{proof}

In our applications $V=\sum_{j\in C}\partial_{z_j}$ translates a set $C$ of coordinates, and $\Phi$ is a product of factors $|\varpi(z_i-z_j)|^\beta$, with $\varpi(s)=s$ or $2\sin(s/2)$, over pairs $i,j$ both in $C$ or both outside $C$, and of factors depending only on coordinates outside $C$.

\subsection{The circular \texorpdfstring{$\beta$}{beta}-ensemble}
For $\beta>0$, $q\in\Z$ and $n\ge1$, let $\rho^{\beta}_{n,q}$ be the image of the complex measure
\[
e^{iq\sum_j\varphi_j}\prod_{i<j}|e^{i\varphi_i}-e^{i\varphi_j}|^\beta\,d\varphi_1\cdots d\varphi_n\qquad\text{on }\T^n=(\R/2\pi\Z)^n
\]
under $f(\varphi)=\sum_j2\cos\varphi_j$. It is a compactly supported distribution, and $\langle\rho^\beta_{n,q},e^{tw}\rangle$ is a constant multiple of $\E_{\mathrm{C}\beta\mathrm{E}_n}(\det U)^qe^{t\tr(U+U^{-1})}$.

\begin{proposition}\label{V:prop:vis-cbe}
For every $0\le b\le n$, the distribution $\rho^\beta_{n,q}$ is not $C^\infty$ at $\lambda_b=2(n-2b)$. More precisely, if $\psi\in C^\infty_c(\R)$ equals $1$ near $\lambda_b$ and vanishes near the other $\lambda_{b'}$, then
\[
\langle\rho^\beta_{n,q},\psi(w)e^{i\tau w}\rangle=\binom nb\,c_b\,e^{i\tau\lambda_b}\tau^{-\gamma_b}\bigl(1+O(\tau^{-1/2})\bigr),\qquad\gamma_b=\frac n2+\frac\beta2\Bigl(\binom{n-b}2+\binom b2\Bigr),
\]
with $c_b\ne0$.
\end{proposition}

\begin{proof}
The left side equals $\int_{\T^n}e^{i\tau f}\psi(f)\,a\,d\varphi$, where $a$ is the density above. The critical points of $f$ are the vertices $\varphi^S$, $S\subseteq[n]$, with $\varphi^S_j=\pi$ for $j\in S$ and $\varphi^S_j=0$ otherwise, and $f(\varphi^S)=\lambda_{|S|}$. Let $\varphi^*\in\supp\psi(f)$. If $\varphi^*$ is not a vertex, some \emph{cluster} $C=\{j:\varphi^*_j=\vartheta\}$ of coinciding coordinates has $\vartheta\notin\{0,\pi\}$. On a small neighbourhood $U$ of $\varphi^*$, the factors $|e^{i\varphi_i}-e^{i\varphi_j}|^\beta=|2\sin((\varphi_i-\varphi_j)/2)|^\beta$ with $i,j$ in different clusters are smooth, and $V=\sum_{j\in C}\partial_{\varphi_j}$ satisfies $Vf=-2\sum_{j\in C}\sin\varphi_j\ne0$. By Lemma~\ref{V:lem:cluster} (on $D=\R^n$, lifting to a chart, with $\Phi$ the product of the factors over pairs inside clusters), a smooth cutoff of $\psi(f)a$ to $U$ contributes $O(\tau^{-M})$. By compactness and a partition of unity, the integral equals the sum over $|S|=b$ of the contributions of small neighbourhoods $W_S$ of $\varphi^S$, up to $O(\tau^{-M})$; the vertices with $|S|\ne b$ lie outside $\supp\psi(f)$.

Near $\varphi^S$ put $u_j=2\sin(\varphi_j/2)$ for $j\notin S$ and $v_j=2\sin((\varphi_j-\pi)/2)$ for $j\in S$. Then $f=\lambda_b-|u|^2+|v|^2$ exactly. For $i,j\notin S$, $2\sin((\varphi_i-\varphi_j)/2)$ is a real-analytic function of $(u_i,u_j)$ near $0$ that vanishes on $u_i=u_j$ and has partial derivative $1$ in $u_i$ at $0$; hence it equals $(u_i-u_j)\,g_{ij}(u)$ with $g_{ij}$ real-analytic and $g_{ij}(0)=1$, and $|2\sin((\varphi_i-\varphi_j)/2)|^\beta=|u_i-u_j|^\beta\,|g_{ij}|^\beta$ with $|g_{ij}|^\beta$ smooth near $0$. The same holds for $i,j\in S$ in the variables $v$. For $i\notin S$, $j\in S$ the factor is smooth and equals $2^\beta$ at $\varphi^S$. The twist equals $(-1)^{qb}$ at $\varphi^S$, and $d\varphi_j=(1-u_j^2/4)^{-1/2}du_j$. Hence the contribution of $W_S$ is
\[
e^{i\tau\lambda_b}\int e^{i\tau(-|u|^2+|v|^2)}|\Delta(u)|^\beta|\Delta(v)|^\beta h_S(u,v)\,du\,dv,\qquad h_S(0)=(-1)^{qb}2^{\beta b(n-b)},
\]
with $h_S$ smooth and compactly supported. Lemma~\ref{V:lem:vertex-model} (case Q, $N_1=n-b$, $N_2=b$, $d_1=\beta\binom{n-b}2$, $d_2=\beta\binom b2$) gives $c_be^{i\tau\lambda_b}\tau^{-\gamma_b}(1+O(\tau^{-1/2}))$ with $c_b\ne0$ independent of $S$. Summing over the $\binom nb$ sets $S$ proves the formula. If $\rho^\beta_{n,q}$ were smooth on a neighbourhood of $\lambda_b$, we could choose $\psi$ supported in it; then $\psi\rho^\beta_{n,q}$ would be in $C^\infty_c$ and the left side would be $O(\tau^{-M})$ for all $M$.
\end{proof}

\subsection{The \texorpdfstring{$\beta$}{beta}-Jacobi ensemble}
For $\beta>0$ and $a,b>-1$ let $\rho^{J}$ be the image of the probability measure
\[
\frac1{Z}\prod_{j=1}^{n}s_j^{a}(1-s_j)^{b}\prod_{i<j}|s_i-s_j|^\beta\,ds\qquad\text{on }[0,1]^n
\]
under $(s_j)\mapsto\sum_js_j$.

\begin{proposition}\label{V:prop:vis-jac}
For every $0\le k\le n$, $\rho^J$ is not $C^\infty$ at $k$. If $\psi\in C^\infty_c(\R)$ equals $1$ near $k$ and vanishes near the other integers in $[0,n]$, then $\langle\rho^J,\psi e^{i\tau w}\rangle=\binom nk c'_ke^{i\tau k}\tau^{-\gamma^J_k}(1+O(\tau^{-1/2}))$ with $c'_k\ne0$ and
\[
\gamma^J_k=(n-k)(a+1)+k(b+1)+\beta\Bigl(\binom{n-k}2+\binom k2\Bigr).
\]
\end{proposition}

\begin{proof}
The vertices of the cube are the indicator vectors $\mathbf 1_S$, $S\subseteq[n]$, and $\sum_js_j=|S|$ at $\mathbf 1_S$. Let $s^*\in[0,1]^n$ with $\sum s^*_j\in\supp\psi$. If $s^*$ is not a vertex, some cluster $C$ of coinciding coordinates has its common value in $(0,1)$. Choose a small open neighbourhood $U$ of $s^*$ in $\R^n$ on which the coordinates in $C$ stay in $(0,1)$ and the factors $s_j^a(1-s_j)^b$, $j\in C$, and $|s_i-s_j|^\beta$, $i,j$ in different clusters, are smooth. Apply Lemma~\ref{V:lem:cluster} with $V=\sum_{j\in C}\partial_{s_j}$, $V\sum_js_j=|C|\ne0$, on the $V$-invariant set $D=\{s:0\le s_j\le1\text{ for }j\notin C\}$; on $U\cap D$ the weight agrees with that of the ensemble, and the factors $s_j^a(1-s_j)^b$, $j\notin C$, and $|s_i-s_j|^\beta$, $i,j$ in a common cluster, go into $\Phi$. So these regions contribute $O(\tau^{-M})$. Near $\mathbf 1_S$ with $|S|=k$, put $u_j=s_j$ for $j\notin S$ and $v_j=1-s_j$ for $j\in S$; then $\sum_js_j=k+\sum u_j-\sum v_j$, and the weight is $\prod u_j^a\,|\Delta(u)|^\beta\cdot\prod v_j^b\,|\Delta(v)|^\beta\cdot h(u,v)$ with $h$ smooth near $0$ and $h(0)=1/Z$. Lemma~\ref{V:lem:vertex-model} (case L) gives the contribution $c'_ke^{i\tau k}\tau^{-\gamma^J_k}(1+O(\tau^{-1/2}))$ with $c'_k\ne0$, where we used $\int_{\Sigma_N}\prod\omega_j^a|\Delta(\omega)|^\beta d\sigma<\infty$ for $a>-1$. Conclude as before, choosing $\psi$ supported in a neighbourhood of $k$ on which $\rho^J$ would be smooth.
\end{proof}
\section{Kac recursions: scalar operators from the first-order systems}\label{V:sec:sys}

\subsection{The Sylvester--Kac matrix}
Let $T_n$ be the $(n+1)\times(n+1)$ matrix, indexed by $0\le p,p'\le n$, with $T_n[p,p-1]=p$, $T_n[p,p+1]=n-p$ and all other entries $0$.

\begin{lemma}[Sylvester--Kac]\label{V:lem:kac}
\begin{enumerate}
\item[(a)] $T_n$ has the simple eigenvalues $n-2b$, $0\le b\le n$. Every eigenvector $v$ has $v_0\ne0$ and $v_n\ne0$.
\item[(b)] Define polynomials $r_p(\mu)$ by $r_{-1}=0$, $r_0=1$ and $r_{p+1}=(\mu r_p-p\,r_{p-1})/(n-p)$ for $0\le p<n$. Then
\[
\mu\,r_n(\mu)-n\,r_{n-1}(\mu)=\frac1{n!}\prod_{b=0}^{n}(\mu-n+2b).
\]
\item[(c)] Let $S=\mathrm{diag}\binom np^{-1/2}$. Then $S^{-1}T_nS$ is the matrix of $J_++J_-=2J_x$ in the standard basis $|j,m\rangle$, $m=j-p$, of the irreducible representation of $\mathfrak{sl}_2$ of spin $j=n/2$.
\end{enumerate}
\end{lemma}

\begin{proof}
(a) On homogeneous polynomials of degree $n$ in $X,Y$ with basis $e_p=X^{n-p}Y^p$, the operator $E=Y\partial_X+X\partial_Y$ acts by $Ee_p=(n-p)e_{p+1}+pe_{p-1}$, so its matrix is $T_n^{\mathsf T}$. In the coordinates $X\pm Y$ one has $E=(X+Y)\partial_{X+Y}-(X-Y)\partial_{X-Y}$, with eigenvectors $(X+Y)^{n-b}(X-Y)^b$ and eigenvalues $n-2b$. If $T_nv=\mu v$ and $v_0=0$, the rows $0,1,\dots,n-1$ give successively $v_1=\dots=v_n=0$, since the entries $n-p$ are nonzero; similarly for $v_n$ using the entries $p$.

(b) Rows $0,\dots,n-1$ of $(T_n-\mu)v=0$ are the recursion for $v_p=r_p(\mu)$ with $v_0=1$, and row $n$ reads $n\,r_{n-1}-\mu r_n=0$. By (a), $\mu r_n-nr_{n-1}$ vanishes at the $n+1$ eigenvalues; it has degree $n+1$ and leading coefficient $1/n!$.

(c) $S^{-1}T_nS$ has entries $\sqrt{p(n-p+1)}$ in positions $(p,p-1)$ and $(p-1,p)$, and $\langle j,m+1|J_+|j,m\rangle=\sqrt{(j-m)(j+m+1)}=\sqrt{p(n-p+1)}$ for $m=j-p$.
\end{proof}

Part (a) is Sylvester's theorem \cite{Syl54}, rediscovered by Kac \cite{Kac47}; see \cite{TT91}.

\subsection{The circular ensembles}
Let $\beta>0$, $n\ge1$ and $q>-1$, and put
\[
G(x)=G^{\beta}_{n,q}(x)={}_0F_1^{(2/\beta)}\Bigl(q+1+\frac{(n-1)\beta}2;\,(x)^n\Bigr),
\]
the hypergeometric function of $n$ equal arguments with Jack parameter $2/\beta$ \cite{Kan93,FW26}. For $q\in\Z_{\ge0}$,
\begin{equation}\label{V:eq:cbe-G}
\E_{\mathrm{C}\beta\mathrm{E}_n}\bigl[(\det U)^q\,e^{t\tr(U+U^{-1})}\bigr]=c\,t^{qn}\,G(t^2),\qquad c\ne0,
\end{equation}
by \cite[(2.7)]{FW26}; for $q=0$, $c=1$ and $G(x)=\sum_N\mu_Nx^N/(N!)^2$ with $\mu_N=\E|\tr U|^{2N}$. Define
\begin{equation}\label{V:eq:Bp}
B_p=p\Bigl(q+1+\frac\beta2(n-p)\Bigr),\qquad \varepsilon_p=B_p-p=p\Bigl(q+\frac\beta2(n-p)\Bigr)\quad(0\le p\le n),
\end{equation}
operators $R_p\in\C[x][\theta]$ by
\begin{equation}\label{V:eq:kac-rec}
R_{-1}=0,\quad R_0=1,\quad R_{p+1}=\frac{(\theta-p+B_p)R_p-p\,x\,R_{p-1}}{n-p}\quad(0\le p<n),
\end{equation}
and
\begin{equation}\label{V:eq:Ltheta}
\mathcal L=(\theta-n+B_n)R_n-n\,x\,R_{n-1}.
\end{equation}

\begin{proposition}\label{V:prop:cbe-op}
$\mathcal LG=0$.
\end{proposition}

\begin{proof}
Forrester and Wei \cite[Prop.~3.3]{FW26} consider the system
\begin{equation}\label{V:eq:fw-system}
(n-p)\,x\,I_{p+1}=B_pI_p+\theta I_p-p\,I_{p-1}\qquad(0\le p\le n)
\end{equation}
with $B_p=p\bigl(a+b+1+\frac\beta2(2n-p-1)\bigr)$, and show that its power series solution normalized by $I_0(0)=1$ has first component $I_0={}_0F_1^{(2/\beta)}(\beta(n-1)+a+b+1;(x)^n)$ \cite[(3.9)]{FW26}. The choice $a+b=q-\frac\beta2(n-1)$ gives $B_p=p(q+1+\frac\beta2(n-p))$, which is \cite[(3.13)]{FW26}, and $I_0=G$. Put $K_p=x^pI_p$. Then $\theta K_p=(p-B_p)K_p+p\,x\,K_{p-1}+(n-p)K_{p+1}$, which for $p<n$ says $K_{p+1}=R_{p+1}G$ by induction, and for $p=n$ says $\mathcal LG=0$.
\end{proof}

Forrester and Wei derive \eqref{V:eq:fw-system} from the Forrester--Kumar system of Proposition~\ref{V:prop:jac-op} below by a confluent limit ($a\to\infty$ with $a+b$ fixed, coefficientwise in the power series). Their Proposition~4.1 eliminates downwards, from $I_n$; we eliminate upwards, from $I_0$, which produces the three-term recursion \eqref{V:eq:kac-rec}.

Conversely, if $\mathcal Ly=0$ for some (formal or actual) function $y$, then $K_p:=R_py$ satisfy $\theta K_p=(p-B_p)K_p+pxK_{p-1}+(n-p)K_{p+1}$ for all $p$. So the scalar equation and the system are equivalent, with $y\mapsto(R_py)_p$ and its inverse given by the first component. We verified Proposition~\ref{V:prop:cbe-op} independently, coefficient by coefficient, on the Jack expansion of $G$ (Appendix~\ref{V:app:comp}).

\begin{proposition}[structure]\label{V:prop:cbe-structure}
Write $\mathcal L=\sum_{k\ge0}x^kP_k(\theta)$ and $K=\fl{(n+1)/2}=\cl{n/2}$. Then $P_k=0$ for $k>K$ and:
\begin{enumerate}
\item[(a)] $P_0(\theta)=\frac1{n!}\prod_{p=0}^{n}(\theta+\varepsilon_p)$;
\item[(b)] $\deg P_k\le n+1-2k$, and $\sum_{k=0}^{K}[\theta^{n+1-2k}]P_k\cdot\mu^{n+1-2k}=\frac1{n!}\prod_{b=0}^{n}(\mu-n+2b)$; in particular $\deg P_k=n+1-2k$ for $0\le k\le K$;
\item[(c)] $n!\,P_k\in\Z[B_1,\dots,B_n][\theta]$ for every $k$.
\end{enumerate}
\end{proposition}

\begin{proof}
Write $R_p=\sum_kx^kR_{p,k}(\theta)$. By \eqref{V:eq:mult}, \eqref{V:eq:kac-rec} becomes
\[
R_{p+1,k}(\theta)=\frac{(\theta+k-p+B_p)R_{p,k}(\theta)-p\,R_{p-1,k-1}(\theta)}{n-p}.
\]
By induction $\deg R_{p,k}\le p-2k$, and the coefficients $\rho_{p,k}=[\theta^{p-2k}]R_{p,k}$ satisfy $\rho_{p+1,k}=(\rho_{p,k}-p\rho_{p-1,k-1})/(n-p)$, so that $\sum_k\rho_{p,k}\mu^{p-2k}=r_p(\mu)$ in the notation of Lemma~\ref{V:lem:kac}(b). Likewise $P_k=(\theta+k-n+B_n)R_{n,k}-nR_{n-1,k-1}$, so $\deg P_k\le n+1-2k$ and $\sum_k[\theta^{n+1-2k}]P_k\mu^{n+1-2k}=\mu r_n-nr_{n-1}$, which is (b) by Lemma~\ref{V:lem:kac}(b). The polynomial $\mu^{[n\text{ even}]}\prod_{b<n/2}(\mu^2-(n-2b)^2)$ has nonzero coefficients at all $\mu^{n+1-2k}$, $0\le k\le K$, and none at higher $k$. For (a), $R_{p+1,0}=(\theta-p+B_p)R_{p,0}/(n-p)$ gives $P_0=\frac1{n!}\prod_{p=0}^{n}(\theta-p+B_p)$, and $B_p-p=\varepsilon_p$. For (c), put $\tilde R_p=\frac{n!}{(n-p)!}R_p$. Then $\tilde R_0=1$ and $\tilde R_{p+1}=(\theta-p+B_p)\tilde R_p-p(n-p+1)\,x\,\tilde R_{p-1}$, so all $\tilde R_p$ have coefficients in $\Z[B_1,\dots,B_n]$, and $n!\,\mathcal L=(\theta-n+B_n)\tilde R_n-n\,x\,\tilde R_{n-1}$.
\end{proof}

\subsection{Formal solutions at infinity}
Put $x=t^2$ and $J_p(t)=t^{-p}K_p$, where $K_p=R_py$ for a solution $y$ of $\mathcal L$. The system for $K_p$ becomes
\begin{equation}\label{V:eq:J-system}
\frac{dJ}{dt}=\Bigl(2T_n+\frac1t\,\mathcal D\Bigr)J,\qquad\mathcal D=\mathrm{diag}(p-2B_p)_{0\le p\le n},
\end{equation}
with $J_0(t)=y(t^2)$. Let $\lambda_b=2(n-2b)$ and
\begin{equation}\label{V:eq:gamma-b}
\gamma_b=\frac n2+\frac\beta4\Bigl((n-b)(n-b-1)+b(b-1)\Bigr)=\frac n2+\frac\beta2\Bigl(\binom{n-b}2+\binom b2\Bigr),\qquad0\le b\le n.
\end{equation}

\begin{proposition}\label{V:prop:cbe-formal}
For every $0\le b\le n$, the operator $\Theta_t:=\sum_kt^{2k}P_k(\theta_t/2)$, which annihilates $G(t^2)$, has a formal solution
\[
e^{\lambda_bt}\,t^{-\gamma_b-qn}\Bigl(1+\sum_{i\ge1}c_{b,i}t^{-i}\Bigr).
\]
\end{proposition}

\begin{proof}
The system \eqref{V:eq:J-system} has an irregular singular point of Poincar\'e rank $1$ at $t=\infty$ whose leading matrix $2T_n$ has the distinct eigenvalues $\lambda_b$. By the formal reduction theory \cite[Ch.~III]{Was65}, \cite[Ch.~3]{Bal00}, it has a formal fundamental matrix $\hat F(t)\,t^{\Lambda_1}e^{t\Lambda_0}$ with $\hat F=F_0+F_1t^{-1}+\cdots$, $\Lambda_0=\mathrm{diag}(\lambda_b)$, $F_0$ a matrix of eigenvectors of $2T_n$, and $\Lambda_1$ the diagonal part of $F_0^{-1}\mathcal DF_0$. The first component of the $b$-th column is a formal solution of the scalar equation, of the form $e^{\lambda_bt}t^{\rho_b}((F_0)_{0b}+O(t^{-1}))$ with $(F_0)_{0b}\ne0$ by Lemma~\ref{V:lem:kac}(a).

To compute $\rho_b=(F_0^{-1}\mathcal DF_0)_{bb}$, take $F_0=S\Psi$ with $\Psi$ an orthogonal matrix of eigenvectors $\psi_b$ of the symmetric matrix $S^{-1}T_nS$. Since $\mathcal D$ is diagonal, $\rho_b=\psi_b^{\mathsf T}\mathcal D\psi_b=\langle p\rangle_b-2\langle B_p\rangle_b$, where $\langle f(p)\rangle_b=\sum_pf(p)\psi_b(p)^2$. By Lemma~\ref{V:lem:kac}(c), $\psi_b$ is an eigenvector of $J_x$ with eigenvalue $m_x=n/2-b$, and $p=j-J_z$ with $j=n/2$. For an eigenvector of $J_x$ one has $\langle J_z\rangle=-i\langle[J_x,J_y]\rangle=0$ and, using the rotation $e^{-i\pi J_x/2}$, which commutes with $J_x$ and conjugates $J_y$ to $\pm J_z$, $\langle J_z^2\rangle=\langle J_y^2\rangle=\frac12\bigl(j(j+1)-m_x^2\bigr)$. Hence $\langle p\rangle_b=n/2$ and
\[
\langle p(n-p)\rangle_b=j^2-\langle J_z^2\rangle=\frac{(n-b)(n-b-1)+b(b-1)}4 .
\]
With $B_p=(q+1)p+\frac\beta2p(n-p)$ this gives $\rho_b=-(2q+1)\frac n2-\frac\beta4\bigl((n-b)(n-b-1)+b(b-1)\bigr)=-\gamma_b-qn$.
\end{proof}

For $\beta=2$, $\gamma_b=\frac12((n-b)^2+b^2)$, as in Sections~\ref{B:sec:unitary}--\ref{B:sec:twisted}. The values $\gamma_b$ coincide with the exponents found at the vertices in Proposition~\ref{V:prop:vis-cbe}; the formal theory and the stationary phase computation see the same vertex.

\begin{remark}[balance as a trace identity]\label{V:rem:trace}
Since $\Lambda_1$ is the diagonal part of $F_0^{-1}\mathcal DF_0$, $\sum_b\rho_b=\operatorname{tr}\mathcal D=\sum_p(p-2B_p)=-\binom{n+1}2-2\sum_p\varepsilon_p$. In the variable $t$ the exponents of $G(t^2)$ at $0$ are $-2\varepsilon_p$, so
\[
\sum_{b=0}^{n}(\gamma_b+qn)-\sum_{p=0}^{n}2\varepsilon_p=\binom{n+1}2 .
\]
This is the vanishing of the balance defect of Theorem~\ref{V:thm:balance}, written for $D^\beta_{n,q}=c\,t^{qn}G(t^2)$: with $f(p)=2\varepsilon_p-qn=\beta p(n-p)-q(n-2p)$ it reads $\sum_b\gamma_b-\sum_pf(p)=\binom{n+1}2$. In Section~\ref{V:sec:even} this identity is verified by summation and used as an input; here it is the trace of the residue of the Kac system.
\end{remark}

\subsection{The Jacobi ensembles}
Let $\beta>0$, $a,b>-1$, and let
\[
H(x)=\E_{\beta\mathrm J}\,e^{-x\sum_js_j}={}_1F_1^{(2/\beta)}\Bigl(\frac\beta2(n-1)+a+1;\ \beta(n-1)+a+b+2;\ (-x)^n\Bigr)
\]
be the Laplace transform of the trace of the $\beta$-Jacobi ensemble of Section~\ref{V:sec:vis} \cite{Kan93,FK22}. Define
\begin{equation}\label{V:eq:BD-jac}
B^J_p=p\Bigl(a+b+2+\frac\beta2(2n-p-1)\Bigr),\qquad D_p=p\Bigl(b+1+\frac\beta2(n-p)\Bigr),\qquad \varepsilon^J_p=B^J_p-p,
\end{equation}
operators $R^J_{-1}=0$, $R^J_0=1$,
\[
R^J_{p+1}=\frac{(\theta-p+B^J_p+(n-p)x)R^J_p-D_p\,x\,R^J_{p-1}}{n-p},
\]
and $\mathcal L^J=(\theta-n+B^J_n)R^J_n-D_n\,x\,R^J_{n-1}$.

\begin{proposition}\label{V:prop:jac-op}
$\mathcal L^JH=0$. Writing $\mathcal L^J=\sum_kx^kP^J_k(\theta)$, one has $P^J_k=0$ for $k>n$,
\[
P^J_0(\theta)=\frac1{n!}\prod_{p=0}^{n}(\theta+\varepsilon^J_p),\qquad\sum_{k=0}^{n}[\theta^{n+1-k}]P^J_k\cdot\mu^{n+1-k}=\frac1{n!}\,\mu\prod_{k=1}^{n}(\mu+k),\qquad P^J_n(\theta)=\theta+\gamma^J_0,
\]
with $\gamma^J_0=n(a+1)+\beta\binom n2$. In particular $\deg P^J_k=n+1-k$ for $0\le k\le n$. The operator $\mathcal L^J$ has, for every $0\le k\le n$, a formal solution $e^{-kx}x^{-\gamma^J_k}(1+O(x^{-1}))$ at $x=\infty$, with $\gamma^J_k$ as in Proposition~\ref{V:prop:vis-jac}.
\end{proposition}

\begin{proof}
For $p=0,\dots,n$ let
\[
\hat H_p(x)=\binom np^{-1}\int_{[0,1]^n}\prod_ls_l^{a}(1-s_l)^{b'-1}e^{-x\sum_ls_l}\prod_{j<k}|s_k-s_j|^\beta\,e_p(1-s_1,\dots,1-s_n)\,ds,
\]
with $e_p$ the elementary symmetric polynomial. Forrester and Kumar \cite[Cor.~2.5]{FK22} show that
\[
(n-p)\,x\,\hat H_{p+1}=\bigl((n-p)x+\tilde B_p\bigr)\hat H_p+\theta\hat H_p-\tilde D_p\hat H_{p-1},
\]
where $\tilde B_p=p\bigl(a+b'+1+\frac\beta2(2n-p-1)\bigr)$ and $\tilde D_p=p\bigl(\frac\beta2(n-p)+b'\bigr)$.
Take $b'=b+1$, so that $\tilde B_p=B^J_p$, $\tilde D_p=D_p$, and $\hat H_0$ is the Selberg integral $S_n(a,b,\beta)$ times $H$. With $H_p=\hat H_p/S_n(a,b,\beta)$ the system reads
\begin{equation}\label{V:eq:fk-system}
\theta H_p=-B^J_pH_p+D_pH_{p-1}+(n-p)\,x\,(H_{p+1}-H_p),\qquad H_0=H .
\end{equation}
With $K_p=x^pH_p$ this becomes $\theta K_p=(p-B^J_p-(n-p)x)K_p+D_pxK_{p-1}+(n-p)K_{p+1}$, and $\mathcal L^JH=0$ follows as in Proposition~\ref{V:prop:cbe-op}. The degree bounds follow from the recursion as in Proposition~\ref{V:prop:cbe-structure}, now with $\deg R^J_{p,k}\le p-k$. The top coefficients satisfy $\rho_{p+1}(\mu)=(\mu+(n-p))\rho_p(\mu)/(n-p)$ for the generating polynomials $\rho_p(\mu)=\sum_k[\theta^{p-k}]R^J_{p,k}\mu^{p-k}$ (taken at $x=1$), because the term $D_pxR^J_{p-1}$ has lower total degree; hence the top symbol of $\mathcal L^J$ is $\frac1{n!}\mu\prod_{k=1}^n(\mu+k)$, all of whose coefficients are positive. The coefficient of $x^{p}$ in $R^J_p$ is $1$ by induction, which gives $P^J_n=\theta+B^J_n-D_n$, and $B^J_n-D_n=\gamma^J_0$.

For the formal solutions, $H$-system \eqref{V:eq:fk-system} reads $H'=(X+x^{-1}Y)H$ with $X=-\mathrm{diag}(n-p)+\mathrm{diag}^+(n-p)$ and $Y=-\mathrm{diag}(B^J_p)+\mathrm{diag}^-(D_p)$. The eigenvalue $-k$ of $X$ has right eigenvector $v_p=\binom{n-k}p/\binom np$ (so $v_0=1$, $v_p=0$ for $p>n-k$) and a left eigenvector $w$ supported on $p\ge n-k$ with $w_{n-k+1}/w_{n-k}=-k$. By the formal reduction theory the exponent is $\rho_k=wYv/wv$, and only $p\in\{n-k,n-k+1\}$ contribute:
\[
\rho_k=-B^J_{n-k}+D_{n-k}\frac{v_{n-k-1}}{v_{n-k}}+D_{n-k+1}\frac{w_{n-k+1}}{w_{n-k}}=-B^J_{n-k}+(k+1)D_{n-k}-k\,D_{n-k+1}=-\gamma^J_k ,
\]
by a direct computation.
\end{proof}
\section{Circular \texorpdfstring{$\beta$}{beta}-ensembles: proofs of Theorems~\ref{V:thm:A} and~\ref{V:thm:B}}\label{V:sec:cbe}

Throughout this section $\beta>0$, $n\ge1$, $q>-\frac12$ is real, $K=\cl{n/2}$, and $\mathcal L=\sum_{k=0}^{K}x^kP_k(\theta)$ is the operator of Proposition~\ref{V:prop:cbe-structure}. Put $\hat P_k=n!\,P_k$ and $g_b=\gamma_b+qn$; since $\gamma_b\ge n/2$, we have $g_b>0$.

\subsection{Exact degrees}
Set $Q_{2k}(u)=2^{n+1}\hat P_k(u/2)$. Since $\theta_x=\theta_t/2$ for $x=t^2$, the $\theta$-form $\Theta_t=\sum_kt^{2k}Q_{2k}(\theta_t)$ is a constant multiple of the operator of Proposition~\ref{V:prop:cbe-formal}. By Proposition~\ref{V:prop:cbe-structure}(b),
\[
\chi(\lambda):=\sum_k[u^{n+1-2k}]Q_{2k}\,\lambda^{n+1-2k}=2^{n+1}\prod_{b=0}^{n}\Bigl(\frac\lambda2-n+2b\Bigr)=\prod_{b=0}^{n}(\lambda-\lambda_b).
\]
By Proposition~\ref{V:prop:cbe-formal}, $\Theta_t$ has the formal solutions $e^{\lambda_bt}t^{-g_b}(1+O(t^{-1}))$. Lemma~\ref{V:lem:eta} therefore applies with $N=n+1$ and $c=1$.

\begin{lemma}\label{V:lem:exact-degree}
Let $0\le m\le n$, let $k^*=\fl{(n+1-m)/2}$ and $d=n+1-2k^*\in\{m,m+1\}$. If $d=m+1$, then $\binom{d}{2}\lc(\hat P_{k^*})+[\theta^{d-1}]\hat P_{k^*}\ne0$.
\end{lemma}

\begin{proof}
Since $\hat P_k(\theta)=2^{-n-1}Q_{2k}(2\theta)$, the quantity in question equals $2^{d-n-2}\bigl(d(d-1)[u^d]Q_{2k^*}+[u^{d-1}]Q_{2k^*}\bigr)$. By \eqref{V:eq:eta-gen} and $[u^d]Q_{2k^*}=[\lambda^d]\chi$, this is $2^{d-n-2}$ times the coefficient of $\lambda^{d-1}$ in
\[
\frac\lambda2\chi''(\lambda)+\sum_{b=0}^{n}g_b\prod_{b'\ne b}(\lambda-\lambda_{b'}).
\]
Write $\varepsilon=[n\text{ even}]$ and $c_j=\lambda_j^2>0$ for $0\le j<n/2$, so that $\chi=\lambda^\varepsilon\prod_{j<n/2}(\lambda^2-c_j)$, and note that there are exactly $K$ indices $j<n/2$. Pairing $b$ with $n-b$ (since $\lambda_{n-b}=-\lambda_b$ and $g_{n-b}=g_b$) and keeping $b=n/2$ alone when $n$ is even, one finds that the coefficient of $\lambda^{d-1}$ equals
\[
(-1)^{k^*}\Bigl(\binom d2\sigma_{k^*}(c)+2\sum_{j<n/2}g_j\,\sigma_{k^*}(c_{\ne j})+\varepsilon\,g_{n/2}\,\sigma_{k^*}(c)\Bigr),
\]
where $\sigma_i$ are the elementary symmetric functions and $c_{\ne j}$ is the list $c$ with $c_j$ removed. Indeed, $g_j(\prod_{b'\ne j}+\prod_{b'\ne n-j})(\lambda-\lambda_{b'})=2g_j\lambda^{1+\varepsilon}\prod_{i\ne j}(\lambda^2-c_i)$ and $\frac\lambda2\chi''$ contributes $\binom d2[\lambda^d]\chi=(-1)^{k^*}\binom d2\sigma_{k^*}(c)$. All terms are $\ge0$ because $g_b>0$. If $k^*<K$, the middle sum is positive, as $c_{\ne j}$ has $K-1\ge k^*$ positive entries. If $k^*=K$, then $d=n+1-2K\le1$, so $d=1$, $n$ is even, and the last term is positive.
\end{proof}

\subsection{Proof of Theorem~\ref{V:thm:A}}
\emph{(a) and (b).} By Propositions~\ref{V:prop:cbe-op} and~\ref{V:prop:cbe-structure}, $n!\,\mathcal L=\sum_{k=0}^{K}x^k\hat P_k(\theta)$ annihilates $G$, with $\hat P_0=\prod_p(\theta+\varepsilon_p)$, $\deg\hat P_k=n+1-2k$ and $\sum_k\lc(\hat P_k)\mu^{n+1-2k}=\prod_b(\mu-n+2b)$. Since $\deg\hat P_0>\deg\hat P_k$ for $k\ge1$, the operator is Fuchsian at $0$ with indicial polynomial $\hat P_0$, so the exponents are $-\varepsilon_p$. Using $x^k\theta^i=\sum_{m\le i}S(i,m)x^{k+m}\partial^m$, write $n!\,\mathcal L=x\sum_mq_m\partial^m$. For $m\ge1$, $xq_m$ is a combination of the monomials $x^{k+m}$ coming from $x^k\hat P_k$ with $\deg\hat P_k\ge m$; for $m=0$, $xq_0=\sum_kx^k\hat P_k(0)$ and $\hat P_0(0)=0$. Hence the $q_m$ are polynomials, $x^{m-1}\mid q_m$, $q_{n+1}=x^n$, and $\deg q_m\le k^*+m-1=\fl{(n-1+m)/2}$ with $k^*$ as in Lemma~\ref{V:lem:exact-degree}. The coefficient of $x^{k^*+m}$ in $xq_m$ comes from $\hat P_{k^*}$ only; it is $\lc(\hat P_{k^*})\ne0$ if $\deg\hat P_{k^*}=m$, and $\binom{m+1}2\lc(\hat P_{k^*})+[\theta^m]\hat P_{k^*}\ne0$ if $\deg\hat P_{k^*}=m+1$, by Lemma~\ref{V:lem:exact-degree}. So $\deg q_m=\fl{(n-1+m)/2}$. The operator $L_x=\sum_mq_m\partial^m$ has no singular point in $\C\setminus\{0\}$ because $q_{n+1}=x^n$, and for $q\ge0$ it is primitive because $q_1(0)=\hat P_0(1)=\prod_p(1+\varepsilon_p)>0$. The integrality statement follows from Proposition~\ref{V:prop:cbe-structure}(c), since the Stirling numbers $S(i,m)$ are integers.

If $n$ is even, Proposition~\ref{V:prop:cbe-formal} with $b=n/2$ gives a formal solution $x^{-g_{n/2}/2}(1+O(x^{-1/2}))$ of $n!\,\mathcal L$. Applying $\sum_kx^k\hat P_k(\theta)$, the term $x^K\hat P_K(\theta)$ produces $\hat P_K(-g_{n/2}/2)\,x^{K-g_{n/2}/2}$ and all other contributions are of lower order. Hence $\hat P_K(-g_{n/2}/2)=0$, and since $\deg\hat P_K=1$, $\hat P_K(\theta)=\lc(\hat P_K)(\theta+g_{n/2}/2)$.

\emph{(c)} is Proposition~\ref{V:prop:cbe-formal} in the variable $x=t^2$.

\emph{(d)} Let $q\in\Z_{\ge0}$. By \eqref{V:eq:cbe-G}, $D(t):=c\,t^{qn}G(t^2)$ is a constant multiple of $\langle\rho^\beta_{n,q},e^{tw}\rangle$. By Proposition~\ref{V:prop:vis-cbe}, $\rho^\beta_{n,q}$ is not smooth at the $n+1$ points $\lambda_b$, so by Corollary~\ref{V:cor:visible}(a) every nonzero operator annihilating $D$ has order at least $n+1$. If $G$ satisfied an equation of order $s\le n$ in $x$, then $G(t^2)$, and hence $D$, would satisfy one of order $s$ in $t$. Therefore $L_x$ is the minimal operator of $G$. If $\tilde L$ of order $n+1$ with polynomial coefficients annihilates $G$, then $\tilde L=hL_x$ with $h=u/v\in\C(x)$ in lowest terms; $v$ divides every $q_m$, hence is constant by primitivity. In the variable $t$, conjugation by $t^{qn}$ shows that $\sum_kt^{2k}\hat P_k((\theta_t-qn)/2)$ annihilates $D$. This $\theta$-form has order $n+1$, so it is a multiple of the minimal operator of $D$ by an element of $\C(t)$; its coefficient of $\partial_t^{n+1}$ is $2^{-n-1}t^{n+1}$, so the minimal operator has no singular point in $\C\setminus\{0\}$, and its exponents at $0$ are the roots $q(n-2p)-\beta p(n-p)=qn-2\varepsilon_p$ of $\hat P_0((\theta-qn)/2)$. \qed

\begin{remark}\label{V:rem:beta0}
At $\beta=0$ the recursion \eqref{V:eq:kac-rec} with $q=0$ becomes $R_{p+1}=(\theta R_p-pxR_{p-1})/(n-p)$, which is the classical recursion for the $n$-th symmetric power of the Bessel operator $\theta^2-x$ (Section~\ref{V:sec:bessel}), in accordance with $\E_{\mathrm{C}0\mathrm{E}_n}e^{t\tr(U+U^{-1})}=I_0(2t)^n$. Thus $\mathcal L$ interpolates between the symmetric power of Bessel's equation ($\beta=0$, independent eigenvalues), the Bessel--Toeplitz operator of Sections~\ref{B:sec:unitary}--\ref{B:sec:twisted} ($\beta=2$), and the even-$\beta$ operators of Section~\ref{V:sec:even}. As $\beta\to0$ all exponents $-\varepsilon_p$ with $q=0$ tend to $0$, and the degree bound of Theorem~\ref{V:thm:B} must then come from the exponent $0$ itself; this is done in Theorem~\ref{V:thm:F}(c) with the sharper dictionary of Lemma~\ref{V:lem:dict}(b).
\end{remark}

\begin{proof}[Proof of Corollary~\ref{V:cor:FW41}]
By \cite[(3.10)]{FW26}, $I_n$ is a constant multiple of ${}_0F_1^{(2/\beta)}(\beta(n-1)+a+b+2;(x)^n)=G^\beta_{n,q'}$. Theorem~\ref{V:thm:A}(a) with $q=q'$ gives an operator of the form \cite[(4.1)]{FW26} with the asserted degrees. In terms of the parameters $B^{\mathrm{FW}}_p=p(a+b+1+\frac\beta2(2n-p-1))$ of \cite[(3.3)]{FW26} one has $B_p=p(q'+1+\frac\beta2(n-p))=B^{\mathrm{FW}}_p+p$, so the integrality statement of Theorem~\ref{V:thm:A}(a) applies. The operator obtained in \cite[\S4.1]{FW26} by eliminating \eqref{V:eq:fw-system} downwards from $I_n$ has the leading term $x^n\partial^{n+1}$, annihilates $I_n$, and, for fixed $\beta$, has coefficients that are polynomials in $q'$. The same holds for $L$. For $q'\in\Z_{\ge0}$ both are equal by the uniqueness in Theorem~\ref{V:thm:A}(d); hence they are equal for all $q'$. Minimality for $q'\in\Z_{\ge0}$ is Theorem~\ref{V:thm:A}(d).
\end{proof}

\begin{remark}\label{V:rem:FWfalse}
The range $q'>-\frac12$ cannot simply be dropped in \cite[Prop.~4.1]{FW26}. For $n=2$, $\beta=\frac12$ and $q'=-\frac12$ (for instance $a+b=-\frac74$), the elimination gives $L=x^2\partial^3+\frac74x\partial^2-4x\partial$, whose coefficient of $d^0/dx^0$ vanishes identically instead of having degree $0$; also $L$ is divisible by $x$. Here $g_1=\gamma_1+2q'=0$, which is exactly the case excluded in Lemma~\ref{V:lem:exact-degree}. 
The minimal order is still $3$, as can be proved directly.
Put $h(x)={}_0F_1(\tfrac34;4x)$ and
$R=x\partial_x^2+\tfrac34\partial_x-4$. Then $Rh=0$ and
\[
 \frac Lx=\partial_x\circ R,
 \qquad
 G^{1/2}_{2,-1/2}(x)=\frac{1+h(x)}2.
\]
For the second identity, the Jack series gives $G(0)=1$ and
$G'(0)=2/(3/4)=8/3$, which are also the initial coefficients of
$(1+h)/2$. The equation $(L/x)G=0$ determines every later analytic
coefficient from these two: if $G=\sum_{j\ge0}a_jx^j$, then
\[
 a_{j+2}=\frac{4a_{j+1}}{(j+2)(j+7/4)}\qquad(j\ge0).
\]

To prove minimality, termwise integration of the exponential series
and the beta integral give, for positive $x$,
\[
 h(x)=\frac{\Gamma(3/4)}{\sqrt\pi\,\Gamma(1/4)}
       \int_{-1}^1 e^{4\sqrt{x}\,s}(1-s^2)^{-3/4}\,ds.
\]
The integral grows faster than every polynomial as $x\to+\infty$.
The probability measures obtained by normalizing its integrand
concentrate at $s=1$: for every fixed $\delta>0$, comparison with a
fixed subinterval of $(1-\delta/2,1)$ shows that the mass of
$[-1,1-\delta]$ tends to zero exponentially in $\sqrt{x}$.
Differentiation under the integral consequently yields
\[
 \frac{h'(x)}{h(x)}\sim\frac2{\sqrt{x}}.
\]
If $A+B h+C h'=0$ with $A,B,C\in\C(x)$ and $C\ne0$, then
$h'/h=-B/C-A/(Ch)$. The second term decays faster than every power
of $x$, whereas a nonzero rational function has an integral leading
power of $x$ at infinity. This contradicts the displayed half-integral
leading power. If $C=0$, the growth of $h$ forces $A=B=0$.
Thus $1,h,h'$ are linearly independent over $\C(x)$.
Finally
\[
 h'=2G',\qquad h=\frac x2G''+\frac38G',\qquad 1=2G-h,
\]
so $G,G',G''$ are linearly independent over $\C(x)$.
The displayed third-order annihilator is therefore minimal in order.
\end{remark}

\subsection{Proof of Theorem~\ref{V:thm:B}}
Let $q=0$, so that $G(x)=\sum_N\mu_Nx^N/(N!)^2$ and $\hat P_0=\theta^2\prod_{p=1}^{n-1}(\theta+\varepsilon_p)$ with $\varepsilon_p=\frac\beta2p(n-p)>0$.

\emph{(a)} By Lemma~\ref{V:lem:dict}(b), $\sum_{k}\hat P_k(N-k)(\ff Nk)^2\mu_{N-k}=0$ for all $N\ge0$. For $N\ge1$ divide by $N^2$: the term $k=0$ becomes $R_0(N)\mu_N$, and $(\ff Nk)^2/N^2=(\ff{(N-1)}{k-1})^2$ for $k\ge1$. Each $R_k$ has degree $\deg\hat P_k+2k-2=n-1$.

\emph{(b)} We have $\lc(R_k)=\lc(\hat P_k)$, and dividing $\sum_k\lc(\hat P_k)\mu^{n+1-2k}=\mu^{\varepsilon}\prod_{j<n/2}(\mu^2-(n-2j)^2)$ by $\mu^\varepsilon$ and putting $\lambda=\mu^2$ gives the characteristic polynomial. The formula for $R_K$ follows from $\lc(\hat P_K)=[\mu^\varepsilon]\prod_b(\mu-n+2b)=(-1)^K\prod_{j<n/2}(n-2j)^2$ and Theorem~\ref{V:thm:A}(b).

\emph{(c)} By Theorem~\ref{V:thm:A}(d), $L_x$ is the minimal operator of $G$; it has leading coefficient $x^n$ and $xL_x=\sum_{k=0}^{K}x^k\hat P_k(\theta)$ has span $K$. Corollary~\ref{V:cor:degree} shows that every recurrence has order $M\ge K$ and that $\prod_{p=1}^{n-1}(\theta+\varepsilon_p)$ divides $c_M(\theta-M)$; that is, $R_0(N+M)$ divides $c_M(N)$, which is the assertion about the leading coefficient in the backward form, and $\deg c_M\ge n-1$. If $\mathcal R=\sum_{s=0}^Kc_s(N)S^s$ is the forward form of the recurrence in (a) and $\mathcal R'=\sum_{s=0}^Kc'_sS^s$ another one of order $K$, then $c_K\mathcal R'-c'_K\mathcal R$ has order $<K$ and annihilates $\mu$ for large $N$, so it vanishes. The $c_s$ have no common factor, since dividing by one would give a recurrence of degree $<n-1$. Hence $c'_K/c_K$ is a polynomial and $\mathcal R'=(c'_K/c_K)\mathcal R$.

\emph{(d)} Let $q\in\Z_{\ge1}$. Expanding $e^{t(\tr U+\overline{\tr U})}$ and using the invariance of $\mathrm{C}\beta\mathrm{E}_n$ under $U\mapsto e^{i\phi}U$, only the terms $(\tr U)^N\overline{\tr U}^{\,N+qn}$ survive in \eqref{V:eq:FW27}, so that $G(x)=c^{-1}\sum_Na_q(N)\,x^N/(N!\,(N+qn)!)$. Comparing coefficients of $x^N$ in $\sum_kx^k\hat P_k(\theta)G=0$ and multiplying by $c\,N!\,(N+qn)!$ gives $\sum_k\hat P_k(N-k)\ff Nk\ff{(N+qn)}k\,a_q(N-k)=0$. Here $\hat P_0(\theta)=\theta(\theta+qn)\prod_{i=1}^{n-1}(\theta+\varepsilon_i)$, since $\varepsilon_0=0$ and $\varepsilon_n=qn$. Dividing by $N(N+qn)$ for $N\ge1$ gives a recurrence of order $K$ with coefficients $\prod_{i=1}^{n-1}(N+\varepsilon_i)$ and $\hat P_k(N-k)\ff{(N-1)}{k-1}\ff{(N+qn-1)}{k-1}$, all of degree $n-1$.

Conversely, let $\sum_{s=0}^Mc_s(N)a_q(N+s)=0$ for $N\ge n_0$. By Lemma~\ref{V:lem:dict}(b) with $r=qn$, $G$ is annihilated by a $\theta$-form of span $M$ with lowest coefficient $\ff\theta{M+n_0}\,\ff{(\theta+qn)}M\,c_M(\theta-M)$. By Lemma~\ref{V:lem:transfer}, $M\ge K$ and $\hat P_0$ divides this coefficient. The roots of $\ff\theta{M+n_0}$ are nonnegative, and the roots of $\ff{(\theta+qn)}M$ are $-qn,-qn+1,\dots,-qn+M-1$. Hence, if $I$ is the multiset of those $\varepsilon_i$, $1\le i\le n-1$, that do not lie in $\{qn-M+1,\dots,qn\}$, then $\prod_{i\in I}(\theta+\varepsilon_i)$ divides $c_M(\theta-M)$. If no $\varepsilon_i$ lies in $\{1,\dots,qn\}$, then $I$ contains all $n-1$ of them for every $M$, which gives the degree bound and the divisibility of the leading coefficient. If no $\varepsilon_i$ lies in $\{qn-K+1,\dots,qn\}$, the same holds for $M=K$, and uniqueness of the recurrence of order $K$ follows as in (c). Finally, $\varepsilon_i\in\{1,\dots,qn\}$ means $\beta=2(k/i-q)/(n-i)$ for some $k\in\{1,\dots,qn\}$, which leaves finitely many values. \qed

\begin{remark}\label{V:rem:exceptional}
The hypothesis in Theorem~\ref{V:thm:B}(d) cannot be dropped. For $n=3$, $q=1$, $\beta=1$ one has $\varepsilon_1=2$ and $\varepsilon_2=3$, both in $\{1,2,3\}$, and the twisted moments satisfy
\[
(N+3)\,a_1(N)=(10N+6)\,a_1(N-1)-9(N-1)\,a_1(N-2),
\]
a recurrence of order $K=2$ and degree $1<n-1$. Theorem~\ref{vtd:odd-degree-drop} proves this identity and its minimality as part of an infinite family. Here $K=2$, $qn=3$, and the roots $-\varepsilon_1=-2$ and $-\varepsilon_2=-3$ of $\hat P_0$ are also roots of the factor $\ff{(\theta+3)}2=(\theta+3)(\theta+2)$ in the proof, so they need not divide the leading coefficient. For $\beta=2$, the exponent-only test is exceptional, but Theorem~\ref{vtd:even-degree} proves minimum degree $n-1$ for every integer twist.
\end{remark}

\subsection{The law of \texorpdfstring{$|\tr U|^2$}{the squared trace}}

\begin{corollary}\label{V:cor:density}
Let $\beta>0$ and $n\ge2$. The law of $|\tr U|^2$ under $\mathrm{C}\beta\mathrm{E}_n$ has a density $\rho$ on $[0,n^2]$, which is annihilated, in the sense of distributions on $\R$, by
\[
\mathcal M=\sum_{j=0}^{K}w^j\,R_{K-j}(K-j-1-\theta_w)
\]
with $R_k$ as in Theorem~\ref{V:thm:B}. The operator $\mathcal M$ has order $n-1$, and its leading coefficient is
\[
(-1)^{n-1}w^{n-1}\prod_{j<n/2}\bigl(w-(n-2j)^2\bigr).
\] Every nonzero operator with polynomial coefficients annihilating $\rho$ has order at least $n-1$.
\end{corollary}

\begin{proof}
The map $U\mapsto|\tr U|^2$ is a nonconstant real-analytic function of the eigenangles, so the law is absolutely continuous. The recurrence of Theorem~\ref{V:thm:B}(a) holds for $N\ge1$ with $\mu_{<0}=0$, so $\sum_{j=0}^{K}R_{K-j}(N+K)\mu_{N+j}=0$ for all $N\ge0$. By Lemma~\ref{B:lem:dict-mom}, a $\theta$-form $\sum_jw^jP_j(\theta)$ annihilates $\rho$ if and only if $\sum_jP_j(-N-j-1)\mu_{N+j}=0$ for all $N\ge0$; with $P_j(\theta)=R_{K-j}(K-j-1-\theta)$ this is the recurrence. The order of $\mathcal M$ is $\max_j\deg R_{K-j}=n-1$, and the coefficient of $\partial_w^{n-1}$ is $(-1)^{n-1}w^{n-1}\sum_j\lc(R_{K-j})w^j$, which is given by Theorem~\ref{V:thm:B}(b). Conversely, an annihilator of order $r$ with polynomial coefficients, multiplied by $w^r$, is a $\theta$-form $\sum_jw^jP_j(\theta)$ with $\max\deg P_j=r$; the resulting recurrence $\sum_jP_j(-N-j-1)\mu_{N+j}=0$ is nontrivial, because $\mu_N>0$ for all $N$, and has degree at most $r$. By Theorem~\ref{V:thm:B}(c), $r\ge n-1$.
\end{proof}

For $\beta=2$ this is the minimality statement of Theorem~\ref{B:thm:E}.
\section{Even \texorpdfstring{$\beta$}{beta}: the balance theorem with a multiple drift}\label{V:sec:even}

In this section $\beta=2m$ with $m\in\Z_{\ge1}$, and $q\in\Z_{\ge0}$. We give a second proof of Theorem~\ref{V:thm:A}(d) and of the $t$-statements in it for these parameters. It does not use the first-order system of Proposition~\ref{V:prop:cbe-op}, and it replaces visibility by the balance theorem. It also gives information that the first proof does not see: the exact orders at $t=0$ of the vertex integrals (Theorem~\ref{V:thm:E}). The only new ingredient compared with Sections~\ref{B:sec:unitary}--\ref{B:sec:twisted} is the multiple drift inequality, Lemma~\ref{V:lem:mdrift}.

\subsection{The integral as a multilinear form}
Let
\[
i_k(t)=I_k(2t),\qquad k_k(t)=\int_0^\infty e^{-t(s+s^{-1})}(-s)^k\,\frac{ds}s=2(-1)^kK_k(2t)
\]
be the two Bessel sequences of Section~\ref{B:sec:unitary}; both satisfy $h_k'=h_{k-1}+h_{k+1}$ and $t(h_{k-1}-h_{k+1})=kh_k$, and every $h_k$ is a $\Q(t)$-linear combination of $h_0,h_0'$, where $h_0$ solves Bessel's equation $Th:=th''+h'-4th=0$. On the torus, $|z_p-z_{p'}|^2=2-z_p/z_{p'}-z_{p'}/z_p$, and $\prod_{p<p'}(2-z_p/z_{p'}-z_{p'}/z_p)=\Delta(z)\Delta(z^{-1})$ with $\Delta(z)=\prod_{p<p'}(z_{p'}-z_p)$. Hence
\[
D^\beta_{n,q}(t):=\E_{\mathrm{C}\beta\mathrm{E}_n}(\det U)^qe^{t\tr(U+U^{-1})}=\frac1{Z}\int_{\T^n}\prod_jz_j^q\,e^{t\sum_j(z_j+z_j^{-1})}\bigl(\Delta(z)\Delta(z^{-1})\bigr)^m\prod_j\frac{d\varphi_j}{2\pi},
\]
with $z_j=e^{i\varphi_j}$ and $Z=\int_{\T^n}(\Delta(z)\Delta(z^{-1}))^m\prod_jd\varphi_j/2\pi$. Since
\[
\Delta(z)\Delta(z^{-1})=\sum_{\sigma,\tau\in S_n}\sgn(\sigma\tau)\prod_jz_j^{\sigma(j)-\tau(j)},
\]
\[
\bigl(\Delta(z)\Delta(z^{-1})\bigr)^m=\sum_{\boldsymbol\sigma,\boldsymbol\tau}\epsilon(\boldsymbol\sigma,\boldsymbol\tau)\prod_jz_j^{\delta_j(\boldsymbol\sigma,\boldsymbol\tau)},\qquad\delta_j=\sum_{r=1}^m\bigl(\sigma_r(j)-\tau_r(j)\bigr),
\]
where $\boldsymbol\sigma=(\sigma_1,\dots,\sigma_m)$ and $\boldsymbol\tau$ run over $S_n^m$ and $\epsilon=\prod_r\sgn(\sigma_r\tau_r)$. For sequences $h^{(1)},\dots,h^{(n)}$ put
\[
F^{(m)}_q\bigl(h^{(1)},\dots,h^{(n)}\bigr)=\frac1{n!}\sum_{\boldsymbol\sigma,\boldsymbol\tau}\epsilon(\boldsymbol\sigma,\boldsymbol\tau)\prod_{j=1}^nh^{(j)}_{q+\delta_j(\boldsymbol\sigma,\boldsymbol\tau)} .
\]
Replacing every $\sigma_r$ by $\sigma_r\circ\pi$ and every $\tau_r$ by $\tau_r\circ\pi$ shows that $F^{(m)}_q$ is symmetric; it is multilinear. Integrating term by term with $\frac1{2\pi}\int_{-\pi}^{\pi}e^{2t\cos\varphi}e^{ik\varphi}d\varphi=i_k(t)$ gives $D^\beta_{n,q}=\frac{n!}ZF^{(m)}_q(i,\dots,i)$. For $0\le b\le n$ define the \emph{vertex integrals}
\[
\begin{aligned}\Psi^{(m,q)}_b&=F^{(m)}_q(\underbrace{k,\dots,k}_{b},\underbrace{i,\dots,i}_{n-b})\\
&=\frac1{n!}\int_{(0,\infty)^b\times(-\pi,\pi)^{n-b}}e^{tf(z)}\prod_jz_j^q\bigl(\Delta(z)\Delta(z^{-1})\bigr)^m\prod_{j\le b}\frac{ds_j}{s_j}\prod_{j>b}\frac{d\varphi_j}{2\pi}.\end{aligned}
\]
with $z_j=-s_j$ for $j\le b$, $z_j=e^{i\varphi_j}$ for $j>b$ and $f(z)=\sum_j(z_j+z_j^{-1})$; the integral formula follows from the integral representations of $i_k$ and $k_k$ by Fubini, as in Lemma~\ref{B:lem:Psi-int}. All $\Psi^{(m,q)}_b$ are real on $(0,\infty)$.

\begin{proposition}\label{V:prop:even-basic}
Let $L$ be the minimal operator of $D^\beta_{n,q}$ over $\C(t)$, normalized to be primitive.
\begin{enumerate}
\item[(a)] $\Sol(L)=V:=\spn\{\Psi^{(m,q)}_0,\dots,\Psi^{(m,q)}_n\}$. Every element of $V$ continues analytically along all paths in $\C\setminus\{0\}$, $V\subset\mathcal N$, and $L$ is Fuchsian at $0$.
\item[(b)] For $m'\ge0$ and $t\to+\infty$, $\partial^{m'}\Psi^{(m,q)}_b=C_be^{\lambda_bt}t^{-\gamma_b}\bigl(\lambda_b^{m'}+O(t^{-1})\bigr)$, with an asymptotic expansion in integral powers of $t^{-1}$ that may be differentiated termwise; here $\gamma_b$ is given by \eqref{V:eq:gamma-b} with $\beta=2m$, and
\[
C_b=\frac{(-1)^{qb+m\binom b2}\,4^{mb(n-b)}}{n!\,(2\pi)^{n-b}}\,\mathcal M^{(m)}_{n-b}\,\mathcal M^{(m)}_b\ne0,\qquad\mathcal M^{(m)}_N=\int_{\R^N}e^{-|x|^2}\Delta(x)^{2m}\,dx>0 .
\]
In particular $\ord L=n+1$.
\item[(c)] $D^\beta_{n,q}(-t)=(-1)^{qn}D^\beta_{n,q}(t)$.
\end{enumerate}
\end{proposition}

\begin{proof}
(a) $D^\beta_{n,q}$ is a polynomial in $i_0,i_0'$ with coefficients in $\Q(t)$. The differential Galois group of Bessel's equation is $\SL_2$ (Kolchin), acting transitively on nonzero solutions, so the span of the Galois orbit of $D^\beta_{n,q}$ is spanned by the values $F^{(m)}_q(g_\bullet,\dots,g_\bullet)$, $g\in\Sol(T)$; by polarization and multilinearity this span is $V$. By Lemma~\ref{B:lem:galois-orbit} it is $\Sol(L)$. Every element of $V$ is a polynomial in $i_0,i_0',k_0,k_0'$ with coefficients in $\Q[t,t^{-1}]$, because the recurrences for $h_k$ only divide by $t$; $i_0$ is entire and $k_0=-2(\log t)i_0+(\text{entire})$, which gives the remaining assertions.

(b) This is the proof of Lemma~\ref{B:lem:vertex} with the following changes. In the global coordinates $x_j=2\sin(\varphi_j/2)$ and $y_j=s_j^{1/2}-s_j^{-1/2}$ the phase is $f=\lambda_b-|x|^2-|y|^2$. A circle--circle factor $2-z_p/z_{p'}-z_{p'}/z_p$ is $(x_p-x_{p'})^2$ times an analytic function equal to $1$ at the vertex, an axis--axis factor is $-(y_p-y_{p'})^2$ times such a function, and a mixed factor equals $4$ at the vertex; the twist $\prod_jz_j^q$ equals $(-1)^{qb}$ there. After raising to the $m$-th power, the amplitude is $\Delta(x)^{2m}\Delta(y)^{2m}H(x,y)$ with $H$ analytic and $H(0)=C_b(\mathcal M^{(m)}_{n-b}\mathcal M^{(m)}_b)^{-1}$. The weight $\Delta(x)^{2m}\Delta(y)^{2m}$ is homogeneous of degree $2\gamma_b-n$, and Laplace's method gives the expansion as in Sections~\ref{B:sec:unitary}--\ref{B:sec:twisted}. The rates are distinct and $C_b\ne0$, so $\dim V=n+1$.

(c) The substitution $z\mapsto-z$ preserves $\Delta(z)\Delta(z^{-1})$ and the Haar measure, and multiplies $\prod_jz_j^q$ by $(-1)^{qn}$.
\end{proof}

\subsection{The multiple drift inequality}

\begin{lemma}\label{V:lem:mdrift}
Let $m\ge1$, $q\ge0$, and for $0\le i\le n$ put $f_m(i)=2mi(n-i)-q(n-2i)$; let $f_{m,(0)}\le\dots\le f_{m,(n)}$ be these numbers in increasing order. Let $\sigma_1,\tau_1,\dots,\sigma_m,\tau_m\in S_n$, put $d_s=q+\sum_{r=1}^m(\sigma_r(s)-\tau_r(s))$, and let $S\subseteq[n]$ with $|S|=b$. Then
\[
E:=\sum_{s\in S}|d_s|-\sum_{s\notin S}|d_s|\le f_{m,(b)},\qquad\text{and}\qquad\sum_{b=0}^nf_{m,(b)}=\frac{m(n-1)n(n+1)}3 .
\]
\end{lemma}

\begin{proof}
Let $S_+=\{s\in S:d_s>0\}$, $j=|S_+|$, $A=\sum_{s\in S_+}d_s$ and $B=-\sum_{s\in S\setminus S_+}d_s\ge0$. Since $\sum_{s=1}^nd_s=nq$, we have $\sum_{s\notin S}|d_s|\ge|nq-A+B|$, hence $E\le A+B-|nq-A+B|=\min(2A-nq,\,2B+nq)$. For $T\subseteq[n]$ with $|T|=k$ and $\sigma,\tau\in S_n$, both $\sum_{s\in T}(\sigma(s)-\tau(s))$ and its negative are at most the sum of the $k$ largest elements of $[n]$ minus the sum of the $k$ smallest, which is $k(n-k)$. Summing over $r$ gives $A\le j\bigl(q+m(n-j)\bigr)$ and $B\le(b-j)\bigl(m(n-b+j)-q\bigr)$. Hence $2A-nq\le f_m(j)$ and $2B+nq\le f_m(j')$ with $j'=j+n-b$, so $E\le\min(f_m(j),f_m(j'))$. Since $f_m$ is concave, $f_m(i)\ge\min(f_m(j),f_m(j'))$ for the $n-b+1$ integers $i\in[j,j']$; so at most $b$ values of $f_m$ are smaller, that is, $f_{m,(b)}\ge\min(f_m(j),f_m(j'))\ge E$. Finally $\sum_if_m(i)=2m\sum_ii(n-i)$, because $\sum_i(n-2i)=0$.
\end{proof}

\begin{lemma}\label{V:lem:even-origin}
For $0<t\le1$, $|\Psi^{(m,q)}_b(t)|\le C\,t^{-f_{m,(b)}}(1+|\log t|)^b$. Hence $\ord\Psi^{(m,q)}_b\ge-f_{m,(b)}$.
\end{lemma}

\begin{proof}
By definition, $\Psi^{(m,q)}_b=\frac1{n!}\sum\epsilon(\boldsymbol\sigma,\boldsymbol\tau)\prod_{j\le b}k_{d_j}\prod_{j>b}i_{d_j}$ with $d_j=q+\delta_j=q+\sum_r(\sigma_r(j)-\tau_r(j))$. Use $|i_k(t)|\le e\,t^{|k|}$ and $|k_k(t)|\le C\,t^{-|k|}(1+|\log t|)$ on $(0,1]$ (proof of Lemma~\ref{B:lem:origin}), and Lemma~\ref{V:lem:mdrift} with $S=\{1,\dots,b\}$.
\end{proof}

\subsection{The balance}

\begin{theorem}\label{V:thm:even-t}
Let $\beta=2m$, $q\in\Z_{\ge0}$, $N=n+1$, write the minimal operator of $D^\beta_{n,q}$ as $L=\sum_{k=0}^Np_k\partial^k$ \textup{(}primitive\textup{)}, and put $e_i=q(n-2i)-\beta i(n-i)$, $0\le i\le n$.
\begin{enumerate}
\item[(i)] $L$ has order $n+1$ and no singular point in $\C\setminus\{0\}$.
\item[(ii)] The exponents of $L$ at $t=0$ are $e_0,\dots,e_n$, and $\ord\Psi^{(m,q)}_b=-f_{m,(b)}$ for every $b$.
\item[(iii)] After scaling, $p_N=t^A$ with $A=n$ if $0\in\{e_i\}$ and $A=n+1$ otherwise; $\deg p_k\le A$, $\ord_0p_k\ge k+A-N$, and $\sum_k[t^A]p_k\lambda^k=\chi(\lambda)=\prod_{b=0}^n(\lambda-\lambda_b)$.
\item[(iv)] $t^{N-A}L=\sum_{k=0}^{K}t^{2k}Q_{2k}(\theta)$, $K=\cl{n/2}$, with $Q_0(\theta)=\prod_i(\theta-e_i)$, $\deg Q_{2k}=n+1-2k$, $\lc(Q_{2k})=[\lambda^{n+1-2k}]\chi$, and \eqref{V:eq:eta-gen} holds with $g_b=\gamma_b$.
\end{enumerate}
\end{theorem}

\begin{proof}
By Proposition~\ref{V:prop:even-basic} and Lemma~\ref{V:lem:even-origin}, the basis $\Psi^{(m,q)}_b$ satisfies (B1)--(B4) of Theorem~\ref{V:thm:balance} with $\beta_b=f_{m,(b)}$. Its balance defect is
\[
\begin{aligned}
\Delta&=\binom{n+1}2-\sum_b\gamma_b+\sum_bf_{m,(b)}\\
&=\frac{n(n+1)}2-\Bigl(\frac{n(n+1)}2+\frac{m(n+1)n(n-1)}3\Bigr)+\frac{m(n-1)n(n+1)}3=0,
\end{aligned}
\]
where we used $\sum_b\bigl((n-b)(n-b-1)+b(b-1)\bigr)=2\sum_bb(b-1)=\frac23(n+1)n(n-1)$. Theorem~\ref{V:thm:balance} gives (i) and (ii), since the multiset $\{-f_{m,(b)}\}$ equals $\{-f_m(i)\}=\{e_i\}$. All $e_i$ are congruent to $qn$ modulo $2$, so $0$ and $1$ are not both exponents, and Lemma~\ref{V:lem:shape-t} gives $A$ and $\ord_0p_k\ge k+A-N$. The top degree fact (iii) of Section~\ref{V:sec:ops}, applied to the asymptotics of Proposition~\ref{V:prop:even-basic}(b), gives $\deg p_k\le A$ and the formula for $\chi$.

For (iv), $t^{N-A}L=\sum_kr_k(t)\ff\theta k$ with $r_k=t^{N-A-k}p_k\in\C[t]$ of degree at most $N-k$, so $t^{N-A}L=\sum_jt^jQ_j(\theta)$ with $\deg Q_j\le N-j$ and $Q_0=\sum_kr_k(0)\ff\theta k$, the indicial polynomial. By Proposition~\ref{V:prop:even-basic}(c), the operator obtained from $L$ by $t\mapsto-t$ is primitive, annihilates $D^\beta_{n,q}$ and has order $N$, so it is a constant multiple of $L$; comparing the terms with $\partial^N$, $t^{N-A}L$ is invariant under $t\mapsto-t$, and $Q_j=0$ for odd $j$. Comparing coefficients of $t^N\partial^k$ gives $[t^A]p_k=[\theta^k]Q_{N-k}$; as $\chi=\lambda^{[n\text{ even}]}\prod_{b<n/2}(\lambda^2-\lambda_b^2)$ has nonzero coefficients at all $\lambda^{N-2k}$, $\deg Q_{2k}=N-2k$. Applying $L$ termwise to the asymptotic expansion of $\Psi^{(m,q)}_b$ gives the expansion of $L\Psi^{(m,q)}_b=0$, so $e^{\lambda_bt}t^{-\gamma_b}(1+O(t^{-1}))$ is a formal solution for every $b$, and Lemma~\ref{V:lem:eta} gives \eqref{V:eq:eta-gen}.
\end{proof}

\begin{remark}
Theorem~\ref{V:thm:even-t}(ii) says that the crude Toeplitz bound of Lemma~\ref{V:lem:even-origin} is attained, in the power of $t$, by every single vertex integral, and consequently that the bound of Lemma~\ref{V:lem:mdrift} is attained for every $b$ by some $(\boldsymbol\sigma,\boldsymbol\tau,S)$. A direct proof would require controlling the cancellations among the $(n!)^{2m}$ terms of the expansion; the balance avoids this. For $m=1$ this is Theorem~\ref{B:thm:Dtq}.
\end{remark}
\section{Jacobi ensembles: proof of Theorem~\ref{V:thm:C}}\label{V:sec:jac}

Let $\beta>0$, $a,b>-1$, $n\ge1$, and let $\mathcal L^J=\sum_{k=0}^{n}x^kP^J_k(\theta)$ be as in Proposition~\ref{V:prop:jac-op}; put $\hat P^J_k=n!\,P^J_k$. Recall
\[
\varepsilon^J_p=p\Bigl(a+b+1+\frac\beta2(2n-p-1)\Bigr),\qquad\gamma^J_k=(n-k)(a+1)+k(b+1)+\beta\Bigl(\binom{n-k}2+\binom k2\Bigr).
\]
We write \textup{($\star$)} for the condition $a+b+1+\frac\beta2(n-1)>0$. Since $\varepsilon^J_p/p\ge a+b+1+\frac\beta2(n-1)$ for $1\le p\le n$, under \textup{($\star$)} all $\varepsilon^J_p$ with $p\ge1$ are positive. Condition \textup{($\star$)} holds whenever $a+b\ge-1$ and $n\ge2$, and for $n=1$ whenever $a+b>-1$.

\begin{proof}[Proof of Theorem~\ref{V:thm:C}]
\emph{(a)} By Proposition~\ref{V:prop:jac-op}, $n!\,\mathcal L^J=x\sum_mq_m\partial^m$ annihilates $H$, where, exactly as in Section~\ref{V:sec:cbe}, $q_m=x^{m-1}\sum_kx^k\sum_i[\theta^i]\hat P^J_kS(i,m)$ for $m\ge1$ and $q_0=\sum_{k\ge1}x^{k-1}\hat P^J_k(0)$, using $\hat P^J_0(0)=0$. Since $\deg\hat P^J_k=n+1-k$, the coefficient of $\partial^m$ ($m\ge1$) is $x^{m-1}$ times a polynomial of degree at most $n+1-m$, whose top coefficient $\lc(\hat P^J_{n+1-m})$ is a positive elementary symmetric function of $1,\dots,n$. Hence $x^{m-1}\mid q_m$, $\deg q_m=n$ for $1\le m\le n+1$, and $q_{n+1}=x^n$. The top coefficient of $q_0$ is $\hat P^J_n(0)=n!\,\gamma^J_0>0$, so $\deg q_0=n-1$. Under \textup{($\star$)}, $q_1(0)=\hat P^J_0(1)=\prod_p(1+\varepsilon^J_p)>0$, so $L^J=\sum_mq_m\partial^m$ is primitive. It is Fuchsian at $0$ with indicial polynomial $\hat P^J_0$, so its exponents are $-\varepsilon^J_p$, and its formal solutions at infinity are those of Proposition~\ref{V:prop:jac-op}.

For minimality, $F(t)=H(-t)=\langle\rho^J,e^{tw}\rangle$. By Proposition~\ref{V:prop:vis-jac}, $\rho^J$ is not smooth at the $n+1$ points $0,1,\dots,n$, and Corollary~\ref{V:cor:visible}(a) shows that no nonzero operator of order at most $n$ annihilates $F$, hence none annihilates $H$. So the minimal operator of $H$ has order $n+1$; it equals $x^{-c}L^J$ for some $c\ge0$, and $c=0$ under $(\star)$, since $L^J$ is then primitive. In that case, as in the proof of Theorem~\ref{V:thm:A}(d), every operator of order $n+1$ with polynomial coefficients annihilating $H$ is a polynomial multiple of $L^J$. Without $(\star)$, $L^J$ can be divisible by $x$: for $n=2$ one has $q_1(0)=(a+b+\beta+2)(2a+2b+\beta+3)$, $q_0(0)=(2a+\beta+2)(2a+2b+\beta+3)$ and $q_m(0)=0$ for $m\ge2$, and $2a+2b+\beta+3=0$ is possible with $a,b>-1$.

\emph{(b)} Since $H(-s)=\sum_Nm_Ns^N/N!$, the $\theta$-form $\sum_k(-1)^ks^k\hat P^J_k(\theta_s)$ annihilates $\sum_Nm_Ns^N/N!$. By Lemma~\ref{V:lem:dict}(a), $\sum_k(-1)^k\hat P^J_k(N-k)\ff Nk\,m_{N-k}=0$ for all $N\ge0$. For $N\ge1$ divide by $N$; since $\hat P^J_0(N)=N\prod_{p\ge1}(N+\varepsilon^J_p)$ and $\ff Nk/N=\ff{(N-1)}{k-1}$, this is the recurrence of (b), and $\deg c_k=\deg\hat P^J_k+k-1=n$ for all $k$. The formula for $c_n$ follows from $\hat P^J_n=n!(\theta+\gamma^J_0)$. By Proposition~\ref{V:prop:jac-op}, $\lc(c_k)=(-1)^k\lc(\hat P^J_k)=(-1)^ke_k(1,\dots,n)$, where $e_k$ is the elementary symmetric function, so the characteristic polynomial $\sum_k\lc(c_k)\lambda^{n-k}$ is $\prod_{k=1}^{n}(\lambda-k)$.

Let a recurrence of order $M$ hold for large $N$. By Lemma~\ref{V:lem:dict}(a) the exponential generating function is annihilated by a $\theta$-form of span $M$ in $s$; hence $H$ is annihilated by a $\theta$-form of span $M$ in $x$. The primitive minimal operator is $x^{-c}L^J$ for some $c\ge0$, and its leading coefficient is a nonzero monomial. Multiplication by a Laurent monomial gives the equivalent form $xL^J=\sum_{k=0}^{n}x^k\hat P^J_k(\theta)$, of span $n$, to which the transfer lemma applies. By Corollary~\ref{V:cor:degree}, $M\ge n$, and the factors $\theta+\varepsilon^J_p$ with $\varepsilon^J_p\notin\Z_{\le0}$ divide $c_M(\theta-M)$. Under \textup{($\star$)} all $n$ factors with $p\ge1$ qualify, so $\deg c_M\ge n$ and $\prod_{p\ge1}(N+\varepsilon^J_p)$ divides the leading coefficient in the backward form. Uniqueness follows as in the proof of Theorem~\ref{V:thm:B}(c).

\emph{(c)} As in Corollary~\ref{V:cor:density}, Lemma~\ref{B:lem:dict-mom} translates the recurrence of (b), shifted to hold for all $N\ge0$, into the annihilator $\mathcal M^J=\sum_{j=0}^{n}w^jc_{n-j}(n-j-1-\theta)$ of $\rho^J$. Its order is $\max_j\deg c_{n-j}=n$, and the coefficient of $\partial_w^n$ is $(-1)^nw^n\sum_k\lc(c_k)w^{n-k}=(-1)^nw^n\prod_{k=1}^n(w-k)$. An annihilator of order $r$ yields a nontrivial recurrence of degree at most $r$, since $m_N>0$; under \textup{($\star$)}, (b) gives $r\ge n$.
\end{proof}

\begin{remark}
For $\beta=1$ the trace $\sum_js_j$ is Pillai's $V$ statistic of multivariate analysis, and for this case the system \eqref{V:eq:fk-system} goes back to work of Davis; see \cite[Remark~2.6]{FK22}. Forrester and Kumar \cite[Remark~3.8]{FK22} note that eliminating their system for the density leads to a linear differential equation of order $N$ (after using an integrated form of the equation with $p=0$), and they write it out for $N=2$.
\end{remark}

\begin{remark}\label{V:rem:FK38}
Assume $(\star)$. Let $\mathcal A$ be an operator of order $r$ with polynomial coefficients that annihilates $\rho^J$ on $\R\setminus\{0,1,\dots,n\}$, for instance the equation of \cite[Remark~3.8]{FK22}. Then $\mathcal A\rho^J$ is a distribution supported in $\{0,\dots,n\}$, hence a finite combination of derivatives of Dirac masses, and $\prod_{k=0}^n(w-k)^J\mathcal A$ annihilates $\rho^J$ in the sense of distributions for large $J$. This operator has order $r$, so $r\ge n$ by Theorem~\ref{V:thm:C}(c). If $\mathcal A$ has order $n$, the same argument shows that $\mathcal A$ and $\mathcal M^J$ differ by a left factor in $\C(w)$, since otherwise a suitable combination would have order less than $n$. (They are not equal: for $N=2$ the leading coefficient in \cite[(3.35)]{FK22} is $t(t-1)(t-2)$, while that of $\mathcal M^J$ is $w^2(w-1)(w-2)$.)
\end{remark}
\section{Classical groups and Weyl chamber walks: proof of Theorem~\ref{V:thm:D}}\label{V:sec:walks}

\subsection{Shifted Jacobi traces}
Let $\beta>0$, $a,b>-1$ satisfy \textup{($\star$)}, let $n\ge1$, $\kappa\ne0$ and $c\in\R$, and consider the moments $M_N=\E(\kappa\sum_js_j+c)^N$ in the $\beta$-Jacobi ensemble. Their exponential generating function is $F(t)=e^{ct}H(-\kappa t)$, the Laplace transform of the law of $\kappa\sum s_j+c$, whose non-smooth points are $c+\kappa k$, $0\le k\le n$ (Proposition~\ref{V:prop:vis-jac}).

\begin{lemma}\label{V:lem:shifted}
Let $L_F$ be the minimal operator of $F$. Then $L_F$ has order $n+1$, leading coefficient a constant times $t^n$, and
\[
tL_F=\sum_{k=0}^{n+1}t^k\tilde P_k(\theta),\qquad\deg\tilde P_k\le n+1-k,\qquad\tilde P_0=\mathrm{const}\cdot\prod_{p=0}^n(\theta+\varepsilon^J_p),
\]
with $\sum_k[\theta^{n+1-k}]\tilde P_k\,\lambda^{n+1-k}=\mathrm{const}\cdot\prod_{k=0}^{n}(\lambda-c-\kappa k)$ and $\tilde P_{n+1}=\mathrm{const}\cdot\prod_{k=0}^n(c+\kappa k)$.
\end{lemma}

\begin{proof}
Write $L^J=\sum_mq_m(x)\partial_x^m$ as in Theorem~\ref{V:thm:C}(a): $q_{n+1}=x^n$, $\deg q_m\le n$, and $\ord_0q_m\ge m-1$ for $m\ge1$. Substituting $x=-\kappa t$ and conjugating by $e^{ct}$, which replaces $\partial_t$ by $\partial_t-c$, gives the operator $L_F=\sum_j\tilde q_j(t)\partial_t^j$ with $\tilde q_j=\sum_{m\ge j}\binom mj(-c)^{m-j}(-\kappa)^{-m}q_m(-\kappa t)$. It annihilates $F$ and is minimal because $L^J$ is. Its coefficients have degree at most $n$, $\tilde q_{n+1}$ is a constant times $t^n$, and $\ord_0\tilde q_j\ge j-1$ for $j\ge1$. Hence $t\tilde q_j\partial^j=t^{1-j}\tilde q_j\,\ff\theta j$ with $t^{1-j}\tilde q_j$ a polynomial of degree at most $n+1-j$, which gives the $\theta$-form with $\deg\tilde P_k\le n+1-k$. Since $\deg\tilde P_0=n+1$ is the largest degree, $\tilde P_0$ is the indicial polynomial at $0$, which is unchanged by the substitution and the conjugation. The top coefficients are $[\theta^{n+1-k}]\tilde P_k=[t^n]\tilde q_{n+1-k}$, so their generating polynomial is $\chi_F(\lambda)=\sum_j[t^n]\tilde q_j\lambda^j$. From the formula for $\tilde q_j$, $\chi_F(\lambda)=(-\kappa)^n\chi^J\bigl((\lambda-c)/(-\kappa)\bigr)$ with $\chi^J(\mu)=\sum_m[x^n]q_m\,\mu^m=\sum_k\lc(\hat P^J_k)\mu^{n+1-k}=\mu\prod_{k=1}^n(\mu+k)$ by Proposition~\ref{V:prop:jac-op} (note $[x^n]q_0=0$). Hence $\chi_F$ is a constant times $\prod_{k=0}^n(\lambda-c-\kappa k)$. Finally $\tilde P_{n+1}=[t^n]\tilde q_0=\chi_F(0)$.
\end{proof}

\begin{proposition}\label{V:prop:shifted}
Assume \textup{($\star$)} and $c+\kappa k\ne0$ for $0\le k\le n$. Then $(M_N)$ satisfies
\[
\begin{gathered}\sum_{k=0}^{n+1}\tilde c_k(N)M_{N-k}=0\quad(N\ge1),\\\tilde c_k(N)=\tilde P_k(N-k)\,\ff{(N-1)}{k-1}\quad(k\ge1),\qquad\tilde c_0(N)=\tilde P_0(N)/N.\end{gathered}
\]
a recurrence of order $n+1$ and degree $n$, with leading coefficient a constant times $\prod_{p=1}^n(N+\varepsilon^J_p)$, last coefficient a constant times $(N-1)\cdots(N-n)$, and characteristic polynomial $\prod_{k=0}^n(\lambda-c-\kappa k)$. Every recurrence satisfied by $(M_N)$ has order at least $n+1$ and degree at least $n$, and its leading coefficient in the backward form is divisible by $\prod_{p=1}^n(N+\varepsilon^J_p)$. Every recurrence of order $n+1$ is a polynomial multiple of this one.
\end{proposition}

\begin{proof}
By Lemma~\ref{V:lem:shifted}, $tL_F$ has span exactly $n+1$, because $\tilde P_{n+1}\ne0$. Lemma~\ref{V:lem:dict}(a) gives the recurrence, and dividing by $N$ uses $\tilde P_0(0)=0$; the degree is at most $\max_k(n+1-k+k-1)=n$, with equality for $k=0$. The last coefficient is $\tilde P_{n+1}\ff{(N-1)}{n}$. Lemma~\ref{V:lem:transfer} and Corollary~\ref{V:cor:degree}, with the $n$ negative roots $-\varepsilon^J_p$ of $\tilde P_0$, give the lower bounds, and uniqueness follows as before.
\end{proof}

\begin{proposition}\label{V:prop:shifted-sym}
Assume \textup{($\star$)}, $a=b$, $\kappa\ne0$ and $c=-\kappa n/2$, so that the law of $\kappa\sum s_j+c$ is symmetric and $M_{2N+1}=0$. Put $M'_N=M_{2N}$ and $K=\cl{n/2}$. Then $\tilde P_k=0$ for odd $k$, $tL_F$ has span $2K$, and
\[
\sum_{k=0}^{K}\hat c_k(N)M'_{N-k}=0\quad(N\ge1),\qquad\hat c_0(N)=\frac{\tilde P_0(2N)}{2N},
\]
with $\hat c_k(N)=\tilde P_{2k}(2N-2k)\,\ff{(2N-1)}{2k-1}$ for $k\ge1$,
a recurrence of order $K$ and degree $n$ with leading coefficient a constant times $\prod_{p=1}^n(2N+\varepsilon^J_p)$ and characteristic polynomial $\prod_{j<n/2}(\lambda-\kappa^2(n/2-j)^2)$. Every recurrence satisfied by $(M'_N)$ has order at least $K$ and degree at least $n$, with leading coefficient divisible by $\prod_{p=1}^n(2N+\varepsilon^J_p)$; every recurrence of order $K$ is a polynomial multiple of this one.
\end{proposition}

\begin{proof}
Since $F$ is even, the operator obtained from $L_F$ by $t\mapsto-t$ is a constant multiple of $L_F$, and comparing the terms with $\partial^{n+1}$ shows that $tL_F$ is invariant; hence $\tilde P_k=0$ for odd $k$. The rates $c+\kappa k=\kappa(k-n/2)$ are symmetric, and $\prod_k(\lambda-\kappa(k-n/2))=\lambda^{[n\text{ even}]}\prod_{j<n/2}(\lambda^2-\kappa^2(n/2-j)^2)$ has nonzero coefficients at all $\lambda^{n+1-2k}$, $0\le k\le K$; so the top coefficients of $\tilde P_{2k}$ are nonzero for $k\le K$ and the span is $2K$. Comparing coefficients of $t^{2N}/(2N)!$ in $tL_FF=0$ (Lemma~\ref{V:lem:dict}(a)) and dividing by $2N$ gives the recurrence; its degree is at most $\max_k(n+1-2k+2k-1)=n$.

Conversely, suppose that $\sum_{s=0}^{M}c_s(N)M'_{N+s}=0$ for $N\ge n_0$. Then the full sequence satisfies
\[
\sum_sc_s(j/2)\,M_{j+2s}=0\qquad(j\ge2n_0),
\] because odd moments vanish; it has order $2M$ and leading coefficient $c_M(j/2)$. By Lemma~\ref{V:lem:dict}(a), $F$ is annihilated by a $\theta$-form of span $2M$ whose lowest coefficient is $\ff\theta{2M+2n_0}c_M((\theta-2M)/2)$. Lemma~\ref{V:lem:transfer} gives $2M\ge2K$ and shows that $\prod_{p\ge1}(\theta+\varepsilon^J_p)$ divides $c_M((\theta-2M)/2)$, that is, $\prod_p(2N+2M+\varepsilon^J_p)$ divides $c_M(N)$. Uniqueness follows as before.
\end{proof}

For $a=b=\tfrac12$, $\beta=2$, $\kappa=4$, $c=-2n$ this is the symplectic case $\E_{\Sp(2n)}(\tr U)^{2N}$, and Proposition~\ref{V:prop:shifted-sym} recovers the minimal recurrence for noncrossing matchings in Theorem~\ref{B:thm:D}, with leading coefficient $\prod_p(2N+p(2n+1-p))$.

\subsection{Orthogonal groups}
Let $u_j=2\cos\varphi_j$ and $s_j=(2+u_j)/4=\cos^2(\varphi_j/2)\in[0,1]$. By Weyl's integration formula, the eigenangles $\pm\varphi_1,\dots,\pm\varphi_\ell$ of a Haar-distributed element of $\mathrm{SO}(2\ell+1)$ have density proportional to $\prod_{i<j}(\cos\varphi_i-\cos\varphi_j)^2\prod_j\sin^2(\varphi_j/2)$, and those of $\mathrm{SO}(2\ell)$ have density proportional to $\prod_{i<j}(\cos\varphi_i-\cos\varphi_j)^2$. Since $d\varphi=-ds/\sqrt{s(1-s)}$ and $\sin^2(\varphi/2)=1-s$, the variables $s_j$ form a $\beta=2$ Jacobi ensemble with
\[
(a,b)=\bigl(-\tfrac12,\tfrac12\bigr)\ \text{for }\mathrm{SO}(2\ell+1),\qquad(a,b)=\bigl(-\tfrac12,-\tfrac12\bigr)\ \text{for }\mathrm{SO}(2\ell),
\]
and $\tr U=1+\sum_ju_j=4\sum_js_j+1-2\ell$, respectively $\tr U=4\sum_js_j-2\ell$. In both cases $n=\ell$, and $\varepsilon^J_p=p(2\ell-p)$, respectively $p(2\ell-1-p)$. Condition \textup{($\star$)} holds ($a+b+1=1$, respectively $a+b+1+(\ell-1)=\ell-1>0$ for $\ell\ge2$).

\begin{proof}[Proof of Theorem~\ref{V:thm:D}]
For $\mathrm{SO}(2\ell+1)$ apply Proposition~\ref{V:prop:shifted} with $\kappa=4$, $c=1-2\ell$; the rates $c+4k=1-2\ell+4k$ are odd, hence nonzero, and equal $2\ell+1-4j$ with $j=\ell-k$. For $\mathrm{SO}(2\ell)$, $\ell\ge2$, apply Proposition~\ref{V:prop:shifted-sym} with $\kappa=4$, $c=-2\ell$, so that $\kappa^2(n/2-j)^2=4(\ell-2j)^2$. The identification of $w_B(N)=\E_{\mathrm{SO}(2\ell+1)}(\tr U)^N$ and $w_D(N)=\E_{\mathrm{SO}(2\ell)}(\tr U)^{2N}$ with numbers of sequences of dominant weights follows from $\E_G\chi_V(U)^N=\dim(V^{\otimes N})^G$ and the Pieri rules for tensoring with the defining representation $V$ \cite{GM93}: $V_\lambda\otimes V=\bigoplus V_{\lambda\pm e_i}$ over dominant $\lambda\pm e_i$, plus $V_\lambda$ itself for $\mathrm{SO}(2\ell+1)$ when $\lambda_\ell>0$. The type $B$ path model here retains the restriction on zero steps; in rank one this gives Riordan, rather than Motzkin, numbers.
\end{proof}

\begin{table}[ht]
\centering
\small
\renewcommand{\arraystretch}{1.25}
\begin{tabular}{@{}cl@{}}
\toprule
group & minimal recurrence, coefficients of $w(N),w(N-1),\dots$\\
\midrule
$\mathrm{SO}(3)$ & $N+1,\quad-2(N-1),\quad-3(N-1)$\\
$\mathrm{SO}(5)$ & $(N+3)(N+4),\quad-3(N-1)(N+2),\quad-(N-1)(13N+4),\quad15(N-2)(N-1)$\\
$\mathrm{SO}(7)$ & $(N+5)(N+8)(N+9),\quad-4(N-1)(N+4)(N+7),\quad-2(N-1)(17N^2+109N+99),$\\
& $4(N-2)(N-1)(19N+13),\quad105(N-3)(N-2)(N-1)$\\
\midrule
$\mathrm{SO}(4)$ & $(N+1)^2,\quad-4(2N-1)^2$\\
$\mathrm{SO}(6)$ & $(N+2)(N+3)^2,\quad-4N(2N-1)(5N+7),\quad36(N-1)(2N-3)(2N-1)$\\
$\mathrm{SO}(8)$ & $(N+3)(N+5)(N+6)^2,\quad-4(2N-1)(10N^3+77N^2+157N+50),$\\
& $256(N-1)(N+1)(2N-3)(2N-1)$\\
\bottomrule
\end{tabular}
\medskip
\caption{Minimal recurrences of Theorem~\ref{V:thm:D}, valid for $N\ge1$ with $w(m)=0$ for $m<0$. For $\mathrm{SO}(2\ell+1)$ the sequence is $w_B(N)$; for $\mathrm{SO}(2\ell)$ it is $w_D(N)=\E(\tr U)^{2N}$. They were derived from the $\theta$-forms and checked against the walk counts obtained by dynamic programming over dominant weights (Appendix~\ref{V:app:comp}). $\mathrm{SO}(3)$ gives the Riordan numbers, $\mathrm{SO}(4)$ the squared Catalan numbers; the $\mathrm{SO}(5)$ recurrence is also recorded in \cite[A095922]{OEIS}.}\label{V:tab:walks}
\end{table}
\section{Bessel moments: proof of Theorem~\ref{V:thm:F}}\label{V:sec:bessel}

Let $n\ge1$. The modified Bessel equation of order $0$ reads $\theta^2y=t^2y$, with solutions $I_0(t)$, $K_0(t)$. Define $S_0=1$, $S_{j+1}=(\theta S_j-j\,t^2S_{j-1})/(n-j)$ for $0\le j<n$ (with $S_{-1}=0$), and
\[
\mathcal B_n=\theta S_n-n\,t^2S_{n-1}=\sum_{k=0}^{K}t^{2k}P^{(n)}_k(\theta),\qquad K=\cl{n/2}.
\]

\begin{lemma}\label{V:lem:sympow}
$\mathcal B_n$ annihilates $y^n$ for every solution $y$ of $\theta^2y=t^2y$, and its solution space is $\spn\{I_0^iK_0^{n-i}:0\le i\le n\}$. Moreover $P^{(n)}_0=\theta^{n+1}/n!$, $\deg P^{(n)}_k=n+1-2k$, and
\[
\sum_k[\theta^{n+1-2k}]P^{(n)}_k\mu^{n+1-2k}=\frac1{n!}\prod_{b=0}^n(\mu-n+2b).
\] For even $n$, $P^{(n)}_K(\theta)=\lc(P^{(n)}_K)(\theta+n/2)$. The minimal operator of $I_0^aK_0^b$, for any $a,b\ge0$ with $a+b=n$, is $\mathcal B_n$.
\end{lemma}

\begin{proof}
For a solution $y$ put $Y_j=y^{n-j}(\theta y)^j$. Then $\theta Y_j=(n-j)Y_{j+1}+j\,t^2Y_{j-1}$, so $Y_j=S_jy^n$ by induction and $\mathcal B_ny^n=0$. The products $I_0^iK_0^{n-i}$ have the distinct exponential rates $2i-n$ at $+\infty$, so they are independent and span the $(n+1)$-dimensional solution space. The statements on $P^{(n)}_k$ follow from the recursion exactly as in Proposition~\ref{V:prop:cbe-structure}, with $x=t^2$ of weight $2$ and $\theta$ of weight $1$ (Lemma~\ref{V:lem:kac}(b)); $P^{(n)}_0=\theta^{n+1}/n!$ because $S_{j+1,0}=\theta S_{j,0}/(n-j)$. For even $n$ the solution $(I_0K_0)^{n/2}=(2t)^{-n/2}(1+O(t^{-1}))$ has rate $0$, which forces $P^{(n)}_K(-n/2)=0$ as in the proof of Theorem~\ref{V:thm:A}. Finally, by Kolchin's theorem the differential Galois group of the Bessel equation acts on its solution space as $\SL_2$ \cite{Kol68,vdPS03}; the span of the orbit of $I_0^aK_0^b$ is a nonzero $\SL_2$-stable subspace of the irreducible representation $\mathrm{Sym}^n$, hence all of it, and Lemma~\ref{B:lem:galois-orbit} gives the last assertion.
\end{proof}

\begin{lemma}[Moment uniqueness with an exponential tail]
\label{V:new:mellin-exp-tail}
Let $h\in L^1(0,\infty)$ satisfy
$\int_0^\infty e^{\eta t}|h(t)|\,dt<\infty$ for some $\eta>0$.
If all its nonnegative integer moments vanish, or if just all its
nonnegative even moments vanish, then $h=0$ almost everywhere.
\end{lemma}
\begin{proof}
For all integer moments, the function
$F(z)=\int_0^\infty e^{zt}h(t)\,dt$ is holomorphic on
$\operatorname{Re}z<\eta$ and has all derivatives zero at zero.
The identity theorem gives $F=0$, and uniqueness of the Fourier
transform on $L^1(\mathbb R)$ gives $h=0$ after extending $h$ by zero.
For even moments, extend $h$ evenly to $H(t)=h(|t|)$ on $\mathbb R$.
All moments of $H$ vanish, since its odd moments vanish by symmetry.
Its bilateral transform $\int_{\mathbb R}e^{zt}H(t)\,dt$ is
holomorphic on $|\operatorname{Re}z|<\eta$ and has all derivatives
zero at zero. The same identity and Fourier uniqueness argument applies.
\end{proof}

\begin{proof}[Proof of Theorem~\ref{V:thm:F}]
Let $f=I_0^aK_0^b$ with $b>a\ge0$, $n=a+b$, and
$\mathcal M(z)=\int_0^\infty t^{z-1}f(t)\,dt$ for
$\operatorname{Re}z>0$. Thus $c_{a,b}(k)=\mathcal M(k+1)$.
For every fixed $j$, near zero
$\theta^jf=O((1+|\log t|)^b)$, while at infinity it is bounded by
a polynomial times $e^{-(b-a)t}$. All boundary terms therefore vanish
in repeated Euler integration by parts, giving
\[
 \int_0^\infty t^{z-1}t^{2i}P(\theta)f(t)\,dt
 =P(-z-2i)\mathcal M(z+2i).
\]
Applying this to $\mathcal B_nf=0$ and setting $z=k+1$ proves (a).

For (b), suppose
$\sum_{s=0}^{M}p_s(k)c_{a,b}(k+s)=0$ for all integers $k\ge k_0$,
with $k_0\ge0$, and define
\[
 O=\sum_{s=0}^{M}t^s p_s(-\theta-1-s),\qquad g=Of.
\]
Integration by parts gives
\[
 \int_0^\infty t^kg(t)\,dt
 =\sum_{s=0}^{M}p_s(k)c_{a,b}(k+s)=0\qquad(k\ge k_0).
\]
The function $t^{k_0}g(t)$ is integrable at zero and has an
exponential tail. Lemma~\ref{V:new:mellin-exp-tail} implies $g=0$.
Thus the operator $O$, with lowest $t$-power zero and span $M$,
annihilates $f$.
By Lemma~\ref{V:lem:sympow}, the minimal operator $\mathcal B_n$
has leading differential coefficient $t^{n+1}/n!$, lowest coefficient
$\theta^{n+1}/n!$, and span $2K$, where $K=\lceil n/2\rceil$.
Lemma~\ref{V:lem:transfer} now gives $M\ge2K$ and
\[
 \theta^{n+1}\mid p_0(-\theta-1),
\]
equivalently $(k+1)^{n+1}\mid p_0(k)$.
When $M=2K$, the equality case gives $O=m(\theta)\mathcal B_n$
for a polynomial $m$. Since
$m(\theta)t^{2i}=t^{2i}m(\theta+2i)$, the corresponding recurrence
is the common polynomial $m(-k-1)$ times (a).
Its lowest coefficient has degree $n+1+\deg m$.
This proves the order, degree, and uniqueness statements in (b).

For (c), write $A_N=W_n(2N)$. The multinomial identity gives
\[
 G(x)=I_0(2\sqrt{x})^n
 =\sum_{N\ge0}\frac{A_N}{(N!)^2}x^N.
\]
To distinguish the two variable conventions, define
\[
 Q_i^{(n)}(u)=2^{\,2i-n-1}P_i^{(n)}(2u),\qquad
 \mathcal B_n^x=\sum_{i=0}^{K}x^iQ_i^{(n)}(\theta_x).
\]
Indeed, $x=t^2/4$ gives $\theta_t=2\theta_x$, so
$\mathcal B_n^x$ is $2^{-(n+1)}\mathcal B_n$ after that change of
variable. Consequently $\mathcal B_n^xG=0$, its span is $K$, and
\[
 Q_0^{(n)}(u)=\frac{u^{n+1}}{n!},\qquad
 \deg Q_i^{(n)}=n+1-2i.
\]
Its leading differential coefficient is $x^{n+1}/n!$.
It is minimal: a differential equation of order at most $n$ for $G$
would pull back to one of order at most $n$ for $I_0(t)^n$,
contradicting Lemma~\ref{V:lem:sympow}.

Suppose $\sum_{s=0}^{M}c_s(N)A_{N+s}=0$ for integers $N\ge n_0$,
with $n_0\ge0$ and $c_0c_M\ne0$.
Lemma~\ref{V:lem:dict}(b), with its factorial shift equal to zero,
gives an annihilator of $G$ with span $M$ and lowest coefficient
\[
 \theta^{\underline{M+n_0}}\theta^{\underline M}c_M(\theta-M).
\]
Here $u^{\underline j}=u(u-1)\cdots(u-j+1)$.
By Lemma~\ref{V:lem:transfer}, $M\ge K\ge1$ and the displayed
coefficient is divisible by $\theta^{n+1}$.
Each falling factorial contains $\theta$ exactly once, and therefore
\[
 \theta^{n-1}\mid c_M(\theta-M).
\]
Equivalently $(N+M)^{n-1}\mid c_M(N)$; after writing the recurrence
backward, its coefficient of $A_N$ is divisible by $N^{n-1}$.
In particular its coefficient degree is at least $n-1$.

Conversely coefficient extraction in $\mathcal B_n^xG=0$ yields
\[
 \sum_{i=0}^{K}Q_i^{(n)}(N-i)
       \bigl(N^{\underline i}\bigr)^2 A_{N-i}=0.
\]
Every coefficient is divisible by $N^2$: for $i=0$ this follows from
$n+1\ge2$, and for $i\ge1$ it follows from the square of the falling
factorial. Dividing by $N^2$ gives a recurrence of order $K$ and
degree $n-1$, with leading coefficient $N^{n-1}/n!$.
The coefficients of this recurrence have no nonconstant polynomial
common divisor, since dividing by such a divisor would contradict
the degree lower bound just established.
Ore division over $\mathbb C(N)$ shows that any two recurrences of
minimum order $K$ are rational-function multiples: a nonzero
remainder would be an eventual recurrence of lower order.
Primitivity of the displayed recurrence then makes the multiplier
polynomial whenever both recurrences have polynomial coefficients.
If both have degree $n-1$, the multiplier is constant.
This applies in particular to Verrill's recurrence.
\end{proof}

\begin{theorem}[Minimal recurrences on shifted parity lattices]
\label{V:new:mellin-parity}
Let $n\ge1$, $K=\lceil n/2\rceil$, and let
\[
 f(t)=\sum_{\substack{a+b=n\\b>a\ge0}}
       \lambda_{a,b}I_0(t)^aK_0(t)^b
\]
be a nonzero linear combination with complex coefficients.
For any real $r>-1$, put
\[
 u_r(N)=\int_0^\infty t^{2N+r}f(t)\,dt\qquad(N\ge0).
\]
Then
\begin{equation}
 \sum_{i=0}^{K}P_i^{(n)}(-2N-r-1-2i)u_r(N+i)=0.
 \label{V:new:eq:mellin-parity-recurrence}
\end{equation}
Every eventual polynomial-coefficient recurrence
$\sum_{s=0}^{M}q_s(N)u_r(N+s)=0$ with $q_0q_M\ne0$ satisfies
\[
 M\ge K,\qquad (2N+r+1)^{n+1}\mid q_0(N).
\]
Thus the minimum order is $K$ and the minimum coefficient degree,
among recurrences of any order, is $n+1$.
Every recurrence of order $K$ is a polynomial multiple of
\eqref{V:new:eq:mellin-parity-recurrence}; at order $K$ and degree
$n+1$ it is unique up to a nonzero constant.
\end{theorem}
\begin{proof}
Every nonzero vector in the solution space of $\mathcal B_n$ has
minimal operator $\mathcal B_n$: its differential-Galois orbit spans
a nonzero invariant subspace of the irreducible representation
$\operatorname{Sym}^n(\mathbb C^2)$ of $\mathrm{SL}_2$.
The argument of Lemma~\ref{V:lem:sympow} therefore applies to the
indicated linear combination. All its summands have positive
exponential decay rates at infinity, and its fixed Euler derivatives
have at most logarithmic growth at zero. Integration by parts in
$\mathcal B_nf=0$ gives
\eqref{V:new:eq:mellin-parity-recurrence}.

For a recurrence valid at integers $N\ge N_0$, with $N_0\ge0$, form
\[
 O=\sum_{s=0}^{M}t^{2s}
 q_s\!\left(\frac{-\theta-r-1-2s}{2}\right).
\]
The endpoint estimates justify integration by parts and give
\[
 \int_0^\infty t^{2N+r}(Of)(t)\,dt=0\qquad(N\ge N_0).
\]
The function $h(t)=t^{2N_0+r}(Of)(t)$ is integrable and has an
exponential tail. All its even moments vanish, so
Lemma~\ref{V:new:mellin-exp-tail} gives $Of=0$.
The operator $O$ has lowest $t$-power zero and span $2M$.
Lemma~\ref{V:lem:transfer} gives $2M\ge2K$ and
\[
 \theta^{n+1}\mid q_0\!\left(\frac{-\theta-r-1}{2}\right),
\]
which is the required divisibility. The displayed recurrence attains
both the order bound and the coefficient-degree bound.
When $M=K$, the equality case of that lemma gives
$O=m(\theta)\mathcal B_n$ for a polynomial $m$.
Commuting $m(\theta)$ across $t^{2i}$ shows that its recurrence is the
common polynomial $m(-2N-r-1)$ times
\eqref{V:new:eq:mellin-parity-recurrence}.
The lowest coefficient has degree $n+1+\deg m$, proving uniqueness.
\end{proof}

\begin{remark}
The identity \eqref{V:new:eq:mellin-parity-recurrence} is the
classical symmetric-power Mellin recurrence evaluated at a shifted
argument. The additional assertion is minimality when a relation is
known only on one shifted parity lattice, and for any nonzero vector
in the decaying solution space. Restricting the original recurrence
to that lattice does not alone prove this assertion; even-moment
uniqueness supplies the missing implication. No assertion for strides
greater than two is made here.
\end{remark}


\begin{remark}
For $n=5$ the recursion (a), multiplied by $120$, is
\[
\begin{aligned}
(k+1)^6c(k)&-(35k^4+280k^3+882k^2+1288k+731)\,c(k+2)\\
&+(259k^2+1554k+2435)\,c(k+4)-225\,c(k+6)=0,
\end{aligned}
\]
the recursion for $\int_0^\infty t^kK_0(t)^5dt$ found by Bailey, Borwein, Broadhurst and Glasser \cite{BBBG08} and proved in general by Borwein and Salvy \cite{BS08}. Theorem~\ref{V:thm:F} shows that it cannot be shortened and that no recursion of lower degree exists. The same holds for the moments $\int_0^\infty t^kI_0^aK_0^bdt$ with $b>a$, which satisfy the same recursion.
\end{remark}

\section{Even Dyson parameters: a common Bessel module and common jet fields}
\label{evb:section}

The multiple-drift argument gives more than the scalar operator at even
Dyson parameters. It supplies an everywhere invertible comparison with
the same Bessel symmetric power, independently of the Dyson parameter
and of the integer twist. Throughout this section let $n,m\ge1$ and
$q\ge0$ be integers, let $\beta=2m$, and put
\[
 y(x)=I_0(2\sqrt{x}),\qquad
 \widehat G^{2m}_{n,q}(x)=t^{-nq}D^{2m}_{n,q}(t),\quad x=t^2.
\]
The normalized function is entire: in the exponential expansion,
invariance under simultaneous rotation of all $z_j$ forces the powers
of $\sum z_j$ and $\sum z_j^{-1}$ to differ by $nq$. In particular,
if
$a^{2m}_{n,q}(N)=\mathbb E[(\det U)^q(\operatorname{tr}U)^N
\overline{\operatorname{tr}U}^{\,N+nq}]$, then
\begin{equation}\label{evb:series}
 \widehat G^{2m}_{n,q}(x)
   =\sum_{N\ge0}\frac{a^{2m}_{n,q}(N)}{N!(N+nq)!}\,x^N.
\end{equation}
Thus the displayed definition involves no choice of square root.
It differs from the Jack-normalized function of Section~\ref{V:sec:cbe}
by a nonzero rational constant:
$\widehat G^{2m}_{n,q}=c\,G^{2m}_{n,q}$.
Indeed the ensemble identity gives $c\ne0$, and the constant coefficient
of the finite Laurent-integral expansion makes $c$ rational.
All module and field assertions below therefore hold equally for
$G^{2m}_{n,q}$.

Let $R=\mathbb Q[x,x^{-1}]$ and let the linear forms
$C_k(x;X,Y)\in R[X,Y]$ be those of
\eqref{B:bv:forms}; they satisfy
$t^{-k}I_k(2t)=C_k(t^2;y(t^2),y'(t^2))$ for every integer $k$.
For the permutations and signs in Section~\ref{V:sec:even}, define
\begin{equation}\label{evb:H}
 H^{(m)}_{n,q}(x;X,Y)
   =\frac1Z\sum_{\boldsymbol\sigma,\boldsymbol\tau}
       \epsilon(\boldsymbol\sigma,\boldsymbol\tau)
       \prod_{j=1}^n C_{q+\delta_j}(x;X,Y).
\end{equation}
Here $Z$ is the positive normalization integral of that section.
It is a nonzero integer, since it is the constant coefficient of the
Laurent polynomial $(\Delta(z)\Delta(z^{-1}))^m$.

\begin{theorem}[The common even-parameter module]\label{evb:module}
For all the parameters above,
\[
 \widehat G^{2m}_{n,q}(x)=H^{(m)}_{n,q}(x;y(x),y'(x)),
\]
where $H^{(m)}_{n,q}$ is a nonzero homogeneous form of degree $n$
over $R$. The minimal differential operator is irreducible of order
$n+1$ over $\mathbb C(x)$, with differential Galois group
$\operatorname{Sym}^{n}(\mathrm{SL}_2)$ in its solution
representation. Its local monodromy at zero is a single block
$J_{n+1}(1)$. Order and irreducibility persist over every finite
algebraic extension of $\mathbb C(x)$, and ramified pullback at zero
preserves this block size.
\end{theorem}

\begin{proof}
For each tuple of permutations one has
$\sum_j\delta_j=0$. Substituting
$I_{q+\delta_j}(2t)=t^{q+\delta_j}
C_{q+\delta_j}(t^2;y,y')$ in the finite expansion of the integral
therefore produces the common factor $t^{nq}$ and proves the identity.
Every summand of \eqref{evb:H} is homogeneous of degree $n$.
The form is nonzero: Proposition~\ref{V:prop:even-basic}(b), with
$b=0$, gives a nonzero leading asymptotic coefficient for
$D^{2m}_{n,q}$.

Theorem~\ref{E:acb:module}, applied to this form with $d=n$ and
$\varepsilon=0$, now proves the assertions. Explicitly, the rank-two
equation $xy''+y'-y=0$ has group $\mathrm{SL}_2$ and local block
$J_2(1)$; polarization of the nonzero form embeds its irreducible
$n$th symmetric power into the Picard--Vessiot field. The image is
the solution space of the minimal scalar equation. The symmetric
power of $J_2(1)$ has one block of size $n+1$. Connectedness of the
group gives stability under finite algebraic base extension, and
raising a unipotent block to a positive integral power preserves its
size in characteristic zero.
\end{proof}

Put
\[
 \mathbf J^{(m)}_{n,q}
  =(\widehat G^{2m}_{n,q},(\widehat G^{2m}_{n,q})',\ldots,
                          (\widehat G^{2m}_{n,q})^{(n)})^{\mathsf T},
 \qquad
 \mathbf m_n=(y^n,y^{n-1}y',\ldots,(y')^n)^{\mathsf T}.
\]
Define $P^{(m)}_{n,q}$ by
\begin{equation}\label{evb:P}
 \mathbf J^{(m)}_{n,q}=P^{(m)}_{n,q}(x)\mathbf m_n.
\end{equation}
Equivalently, its row with index $j$, $0\le j\le n$, is the
coefficient vector of $\mathscr D^jH^{(m)}_{n,q}$, where
$\mathscr D=\partial_x+Y\partial_X+x^{-1}(X-Y)\partial_Y$.

\begin{theorem}[An everywhere invertible comparison]\label{evb:gauge}
The matrix $P^{(m)}_{n,q}$ belongs to $\mathrm{GL}_{n+1}(R)$, and
\begin{equation}\label{evb:det}
 \det P^{(m)}_{n,q}(x)
 =c^{(m)}_{n,q}\,x^{-n(n+1)(m(n-1)+3q)/6},
 \qquad c^{(m)}_{n,q}=\det P^{(m)}_{n,q}(1)\in\mathbb Q^*.
\end{equation}
Consequently, for fixed $n$ and any $m,m'\ge1$, $q,q'\ge0$,
the two complete jets are related by a matrix
$P^{(m')}_{n,q'}(P^{(m)}_{n,q})^{-1}\in\mathrm{GL}_{n+1}(R)$.
\end{theorem}

\begin{proof}
All entries belong to $R$. The determinant is nonzero, since a
dependence among its rows would give an equation of order at most
$n$ for $\widehat G^{2m}_{n,q}$, contrary to Theorem~\ref{evb:module}.
By Theorem~\ref{V:thm:even-t}, the $t$-equation has no singularity
away from zero. The unramified change $x=t^2$ and multiplication by
$t^{-nq}$ preserve ordinariness locally at every nonzero point.
The $x$-equation consequently has no finite singularity except
possibly zero.

The matrix $P^{(m)}_{n,q}$ intertwines the symmetric-power
connection and the scalar companion connection. Both connections
are holomorphic at each nonzero point. A fundamental matrix for
the former is carried to a fundamental matrix for the latter;
Abel's identity therefore excludes a zero of $\det P^{(m)}_{n,q}$
there. A nonzero Laurent polynomial with no zero in $\mathbb C^*$
is a monomial, so its determinant is $cx^s$ for an integer $s$.

The exponents $q(n-2i)-2mi(n-i)$ of the $t$-equation, given by
Theorem~\ref{V:thm:even-t}, become
\[
 \rho_i=-i\bigl(q+m(n-i)\bigr),\qquad 0\le i\le n,
\]
after multiplication by $t^{-nq}$ and passage to $x=t^2$.
To justify this passage also in resonant cases, the explicit
Laurent map \eqref{evb:H} identifies the $x$-module with the
Bessel symmetric power; its local solutions have integral-power
Laurent--log expansions. Their pullbacks, multiplied by $t^{nq}$,
span the $t$-solution space, so the indicial exponents transform
by $\nu\mapsto(\nu-nq)/2$, with multiplicity.
The trace of the connection on $\mathbf m_n$ is
$-n(n+1)/(2x)$. The residue of the trace of the companion
connection is $\sum_i\rho_i-n(n+1)/2$, by Abel's indicial formula.
Taking the trace of the gauge equation at zero thus gives
\[
 s=\sum_{i=0}^n\rho_i
   =-\frac{n(n+1)(m(n-1)+3q)}6.
\]
Evaluation at $x=1$ proves that $c$ is nonzero and rational.
The adjugate formula then shows that the inverse matrix also has
entries in $R$.
\end{proof}

\begin{corollary}[Parameter-independent jet fields and arithmetic relations]
\label{evb:arithmetic}
Let $K=\mathbb C(x)$. For every positive even Dyson parameter and
every nonnegative integer twist,
\begin{equation}\label{evb:field}
 K(\mathbf J^{(m)}_{n,q})
   =K(y^n,y^{n-1}y',\ldots,(y')^n)
   =K(y,y')^{\mu_n}.
\end{equation}
Under the change of coordinates
$\mathbf Z=(P^{(m)}_{n,q})^{-1}\mathbf J^{(m)}_{n,q}$, the ideal of
all algebraic relations is generated by
\[
 Z_iZ_{j+1}-Z_{i+1}Z_j,
 \qquad 0\le i<j\le n-1.
\]
Its transcendence degree is two, and any two distinct entries of
$\mathbf J^{(m)}_{n,q}$ are algebraically independent over $K$.

At every $\alpha\in\overline{\mathbb Q}^{\,*}$ these statements
hold over $\overline{\mathbb Q}$ for the corresponding numerical
jets, with the invertible matrix $P^{(m)}_{n,q}(\alpha)$.
In particular all the derivative values of orders $0,\ldots,n$
are nonzero and transcendental. For a fixed $n$ their complete
numerical jet field is independent of $m$ and $q$.
The function fields and numerical fields in dimensions $d,e$
obey the gcd--lcm laws of Corollary~\ref{B:bv:fieldlaw}.
\end{corollary}

\begin{proof}
The invertible Laurent gauge reduces all assertions to the
degree-$n$ monomials in $y,y'$. These two functions are
algebraically independent. Their monomials have the displayed
rational-normal-cone ideal, and generate
$K(y'/y,y^n)=K(y,y')^{\mu_n}$, as in
Theorem~\ref{B:bv:jetcone}. Any two rows of an invertible matrix
give nonproportional homogeneous forms of the same positive
degree in two independent variables; such forms are
algebraically independent.

For a nonzero algebraic $\alpha$, the two values
$y(\alpha),y'(\alpha)$ are algebraically independent by the
Siegel--Shidlovskii argument in Theorem~\ref{B:bv:values}.
This argument applies the $E$-function theorem to $I_0(2t)$ and
its $t$-derivative at a nonzero algebraic square root of $\alpha$.
Equation~\eqref{evb:det} makes the specialized coordinate change
invertible for every such $\alpha$, so the same polynomial-kernel
and homogeneous-form proofs apply to the values. The field laws
follow from the fixed fields of the scalar groups $\mu_d$ and
$\mu_e$ exactly as in Corollary~\ref{B:bv:fieldlaw}.
\end{proof}

\begin{remark}
These conclusions use the integer Laurent expansion at
$\beta=2m$. They do not assert the same Bessel module or
arithmetic independence for non-even Dyson parameters.
For $m=1$ they recover the determinant results of
Section~\ref{B:bv:section}; for $m>1$ they identify all complete
jets with those same fields, even though their scalar operators
and exponents at zero differ.
\end{remark}


\section{Exact degrees for twisted moments}
\label{vtd:section}

The shifted factorials in the twisted generating function can absorb
negative integral exponents in the recurrence dictionary.  Local
monodromy supplies a degree obstruction that is unaffected by this
cancellation.  For every positive even value of the ensemble parameter,
it settles the exceptional integer twists as well.

\begin{lemma}[The double Borel transform of a rational germ]
\label{vtd:rational-double-borel}
Let $A(z)=\sum_{N\ge0}a_Nz^N$ be nonzero, rational and holomorphic at zero.
The minimal differential equation of
\[
                 A_2(x)=\sum_{N\ge0}\frac{a_Nx^N}{(N!)^2}
\]
is regular singular or ordinary at zero, and every local monodromy
Jordan block there has size at most two.
\end{lemma}
\begin{proof}
Partial fractions give, apart from finitely many initial coefficients,
\[
                    a_N=\sum_{j=1}^s p_j(N)\lambda_j^N,
        \qquad p_j\in\mathbb C[N],\quad\lambda_j\ne0.
\]
Consequently $A_2$ is a polynomial plus a finite sum of
$p_j(\theta)I_0(2\sqrt{\lambda_jx})$, with
$\theta=x\partial_x$.  Each Bessel function belongs to a regular
singular rank-two module whose local monodromy at zero satisfies
$(T-I)^2=0$.  A rational differential image, a direct sum, and a
differential subquotient preserve this annihilating-polynomial
identity.  The polynomial term belongs to a trivial rank-one module.
These observations prove both assertions.
\end{proof}

\begin{theorem}[Minimal degree for all even-parameter integer twists]
\label{vtd:even-degree}
Let $m,n\ge1$ and $q\ge0$ be integers, set $\beta=2m$, and put
\[
 a_N=\mathbb E_{\mathrm C\beta\mathrm E_n}
       \bigl[(\det U)^q(\operatorname{tr}U)^N
                    \overline{\operatorname{tr}U}^{\,N+qn}\bigr].
\]
Every nonzero polynomial-coefficient recurrence for $(a_N)$ that holds
for all sufficiently large indices has degree at least $n-1$.
The smallest recurrence order and the smallest coefficient degree are
respectively
\[
                         K=\lceil n/2\rceil,\qquad n-1.
\]
The recurrence obtained from the Kac system attains both bounds and is
primitive.  Every polynomial-coefficient recurrence of order $K$ is a
polynomial multiple of it; in particular, the recurrence attaining
both bounds is unique up to a nonzero constant factor.
These conclusions include every exceptional integer twist at
$\beta=2$.
\end{theorem}
\begin{proof}
The case $n=1$ gives $a_N=1$ and is immediate.  Assume $n\ge2$, set
$r=qn$, and write
\[
 A(z)=\sum_{N\ge0}a_Nz^N,\qquad
 G(x)=\sum_{N\ge0}\frac{a_Nx^N}{N!(N+r)!}.
\]
The normalization of $G$ differs from the normalized Jack function
by a nonzero constant, which has no effect on its differential module.
The estimate $|a_N|\le n^{2N+r}$ proves that $A$ is convergent near zero.
The Kac recurrence in Theorem~\ref{V:thm:B} has the backward form
\[
 \sum_{k=0}^K R_k(N)a_{N-k}=0\quad(N\ge1),\qquad
 R_0(N)=\prod_{i=1}^{n-1}\bigl(N+i(q+m(n-i))\bigr),
 \qquad \deg R_k\le n-1,
\]
with negative-index terms omitted.  Thus
\[
 R A=R_0(0)a_0,\qquad
 R=\sum_{k=0}^K z^kR_k(\theta+k).
\]
The operator $R$ has order $n-1$, is Fuchsian at zero, and has lowest
$\theta$-coefficient $R_0(\theta)$.  The operator $\theta R$ is a
homogeneous annihilator of $A$.  Its lowest $\theta$-coefficient is
$\theta R_0(\theta)$, of degree $n$, and all its other
$\theta$-coefficients have degree at most $n$; hence it is Fuchsian
at zero as well.  Its minimal right factor, the minimal equation of
$A$, is therefore regular singular or ordinary there.

By Theorem~\ref{evb:module}, $G$ generates the irreducible
$n$th symmetric power of the rank-two Bessel module.  Its local
monodromy at zero is the single block $J_{n+1}(1)$.  An exact
coefficient identity gives
\begin{equation}\label{vtd:double-borel-identity}
 A_2(x):=\sum_{N\ge0}\frac{a_Nx^N}{(N!)^2}
       =\prod_{j=1}^{r}(\theta+j)G(x),
\end{equation}
where the empty product is one.  This differential image is nonzero,
since $a_0\ne0$ and the displayed coefficient multipliers are nonzero
for $N\ge0$.  A nonzero differential image of an irreducible module
is isomorphic to that module.  The minimal equation of $A_2$ thus
also has the block $J_{n+1}(1)$ at zero.

Suppose a recurrence of degree $d\ge1$ exists.
Lemma~\ref{rev:degree-monodromy} bounds every local monodromy block
of $A$ by $d$.  Theorem~\ref{B:up:borel-bound}, applied twice,
bounds the unit-eigenvalue blocks of $A_2$ by $d+2$.  It follows that
$n+1\le d+2$, and hence $d\ge n-1$.
If $d=0$, the same recurrence lemma makes $A$ rational, whereas
Lemma~\ref{vtd:rational-double-borel} bounds the blocks of $A_2$
by two.  This contradicts $n+1\ge3$.  The lower bound is proved
for recurrences of arbitrary order and with any finite initial segment
omitted.

The Kac recurrence has order $K$ and degree $n-1$, and its order is
minimal by Theorem~\ref{V:thm:B}.  Its coefficients cannot have a
nonconstant common factor, since cancelling that factor would give an
eventual recurrence of degree below $n-1$.  Two order-$K$
annihilators over $\mathbb C(N)$ are proportional: cancel their
highest-shift coefficients, and minimality excludes a nonzero
annihilator of smaller order.  If one is primitive with polynomial
coefficients, a rational proportionality factor giving polynomial
coefficients must itself be polynomial, by Bezout's identity for their
greatest common divisor.  This proves the remaining assertions.
\end{proof}

\begin{remark}
The proof gives the least coefficient degree among recurrences of
every order, as well as uniqueness at the least order.  It does not
assert that $R_0(N)$ divides the leading coefficient of every
higher-order recurrence in an exceptional case.  That stronger
divisibility statement does not follow from the Jordan-block bound.
\end{remark}

\begin{theorem}[An infinite family with a degree drop]
\label{vtd:odd-degree-drop}
Let $q\ge1$ be an integer, take $n=3$ and $\beta=2q-1$, and define
$a_N$ as in Theorem~\ref{vtd:even-degree}, using this value of $\beta$.
Then, for every $N\ge1$,
\begin{equation}\label{vtd:odd-recurrence}
 (N+4q-1)a_N
       =2(5N+6q-3)a_{N-1}-9(N-1)a_{N-2},
\end{equation}
with negative-index terms omitted.  The least recurrence order is two
and the least coefficient degree is one, even among eventual
recurrences of arbitrary order.  The displayed recurrence is primitive
and unique up to a nonzero constant among recurrences attaining both
bounds.
\end{theorem}
\begin{proof}
For $n=3$, direct elimination in the Kac system gives, for arbitrary
$\beta>0$ and integer $q\ge0$, the coefficients
\begin{align*}
 R_0(N)&=(N+\beta+q)(N+\beta+2q),\\
 R_1(N)&=-10N^2-6\beta N-30qN+10N
                -9\beta q+3\beta-18q^2+15q-3,\\
 R_2(N)&=9(N-1)(N+3q-1).
\end{align*}
At $\beta=2q-1$ they have the common factor $N+3q-1$ and become
\[
 (N+3q-1)\bigl(N+4q-1,\ -2(5N+6q-3),\ 9(N-1)\bigr).
\]
Since $N+3q-1\ne0$ for $N\ge1$, this proves
\eqref{vtd:odd-recurrence} for all such indices.

To exclude degree zero at any order, let $A(z)=\sum_{N\ge0}a_Nz^N$.
It is analytic at zero and satisfies
\begin{equation}\label{vtd:odd-ode}
 \left[z(1-z)(1-9z)\partial_z
       +(4q-1)-(12q+4)z+9z^2\right]A
                         =(4q-1)a_0.
\end{equation}
A rational solution cannot have a pole at an ordinary point of this
equation.  At its two other finite singularities $z=1$ and $z=1/9$,
the exponents of its homogeneous equation are respectively
$q-\tfrac12$ and $3q-\tfrac32$.  Neither is a negative integer,
so a rational solution of the inhomogeneous equation cannot have a
pole there either: the leading polar term would have to solve the
indicial equation of the homogeneous operator.  Analyticity excludes
a pole at zero.  A rational $A$ would therefore be a polynomial.
But if $A$ had degree $D\ge0$ and nonzero leading coefficient $c$,
the left side of \eqref{vtd:odd-ode} would have leading term
$9(D+1)c z^{D+2}$, which cannot equal the nonzero constant on the
right.  Thus $A$ is not rational, and degree zero is impossible.

The order bound in Theorem~\ref{V:thm:B} gives order at least two.
The three linear coefficients in \eqref{vtd:odd-recurrence} have
no common root, so the recurrence is primitive.  The same
minimal-order proportionality argument as in
Theorem~\ref{vtd:even-degree} proves its uniqueness.
\end{proof}

The specialization $q=1$ gives the exceptional $(n,q,\beta)=(3,1,1)$ recurrence of Remark~\ref{V:rem:exceptional}.

\begin{corollary}[The three-particle minimum-order classification]
\label{vtd:three-particle-classification}
For $n=3$, every $\beta>0$ and every integer $q\ge0$, the least
degree among recurrences of the minimum order two is
\[
 \begin{cases}
  1,&q\ge1\text{ and }\beta=2q-1,\\
  2,&\text{otherwise}.
 \end{cases}
\]
The primitive order-two recurrence is unique up to scale in every case.
\end{corollary}
\begin{proof}
Use the three quadratic coefficients $R_0,R_1,R_2$ in the preceding
proof.  Since the roots of $R_0$ are negative, $N-1$ cannot be a
common factor.  For $q=0$, this already excludes any common factor.
For $q\ge1$, the only remaining possibility is $N+3q-1$.
The condition $R_0(1-3q)=0$ forces
$\beta=2q-1$ or $\beta=q-1$.  The first is the displayed degree
drop.  For the second,
\[
              R_1(1-3q)=-3q(3q-1)\ne0,
\]
so no cancellation occurs.  The primitive quadratic recurrence
therefore has degree two except in the stated family, where it has
degree one.  Minimal order is two; hence all recurrences of that order
are rational multiples of the primitive one.  Bezout's identity makes
every polynomial-coefficient multiple a polynomial multiple, which
establishes the degree assertion and uniqueness.
\end{proof}

\section{Jacobi resonances beyond the positivity condition}
\label{V:sec:jacobi-resonances}
The condition $(\star)$ is sufficient, but not necessary, for the degree
statements. The following refinement includes zero indicial parameters and
isolates the finite collection of resonance hyperplanes on which the previous
argument can lose a degree.

\begin{lemma}[Multiplicity in the exponential-generating-function transfer]
\label{V:lem:egf-multiplicity}
In the exponential-generating-function case of Lemma~\ref{V:lem:dict}, suppose
that the lowest coefficient of a minimal $\theta$-form is $P(\theta)$ and
that the transfer lemma applies. If $\nu$ is a root of $P$ of multiplicity
$m_\nu$, every recurrence of order $M$ has leading coefficient $c_M$ satisfying
\[
 \prod_{\nu\notin\mathbb Z_{\ge0}}(\theta-\nu)^{m_\nu}
 \prod_{\nu\in\mathbb Z_{\ge0}}(\theta-\nu)^{\max(m_\nu-1,0)}
 \ \mid\ c_M(\theta-M).
\]
\end{lemma}
\begin{proof}
The lowest coefficient supplied by the dictionary is
$\ff\theta{M+n_0}c_M(\theta-M)$ and is divisible by $P(\theta)$.
The falling factorial has only simple roots, all nonnegative integers.
Comparison of the multiplicity at each root proves the assertion.
\end{proof}

For the Jacobi ensemble, define the set of distinct positive-integer
resonances and a polynomial by
\[
 J=\{j\in\mathbb Z_{>0}:\varepsilon^J_p=-j
                  \text{ for some }1\le p\le n\},\qquad
 D_n(N)=\frac{\prod_{p=1}^n(N+\varepsilon^J_p)}
                  {\prod_{j\in J}(N-j)}.
\]
Each denominator factor cancels one, and only one, copy of that factor in
the numerator, so $D_n$ is a polynomial. Moreover
$\varepsilon^J_p>-p$, because $a+b+1>-1$ and $\beta>0$; consequently
$J\subseteq\{1,\ldots,n-1\}$.

\begin{theorem}[Jacobi degrees away from, and at, the first resonance]
\label{V:thm:jacobi-resonances}
Let $n\ge1$, $\beta>0$, and $a,b>-1$, without assuming $(\star)$.
\begin{enumerate}
\item Every polynomial recurrence for the Jacobi trace moments has order
at least $n$. Its leading coefficient in backward form is divisible by
$D_n(N)$, and its degree is at least $n-|J|$.
\item If $J$ is empty, all the conclusions of Theorem~\ref{V:thm:C} hold:
$L^J$ is primitive, the moment recurrence has minimal order $n$ and
minimal degree $n$, and the density has minimal differential order $n$.
In particular zero values of $\varepsilon^J_p$ are permitted. Every
recurrence of order $n$ is a polynomial multiple of the displayed
recurrence of Theorem~\ref{V:thm:C}.
\item If $J=\{1\}$, each coefficient $c_k(N)$ of that displayed moment
recurrence is divisible by $N-1$. Dividing all coefficients by $N-1$
gives a recurrence valid for $N\ge2$, of minimal order $n$ and minimal
degree $n-1$, with leading coefficient $D_n(N)$. Every recurrence of
order $n$ is a polynomial multiple of this reduced recurrence. The
Jacobi trace density has minimal differential order $n-1$.
\end{enumerate}
\end{theorem}
\begin{proof}
Visibility gives differential order $n+1$ for $H$ throughout the stated
parameter range. The scalar operator $L^J$ has leading coefficient $x^n$;
its primitive normalization therefore still has monomial leading
coefficient. The $\theta$-form $xL^J$ has span $n$ and lowest coefficient
\[
 P_0(\theta)=\theta\prod_{p=1}^n(\theta+\varepsilon^J_p).
\]
Multiplication by a power of $x$ in passing to the primitive normalization
does not change this lowest coefficient. The transfer lemma proves the
order bound. At $\theta=0$ the first factor of $P_0$ absorbs the one zero
of the falling factorial in Lemma~\ref{V:lem:egf-multiplicity}; any additional
zero factors remain. At each distinct $j\in J$ the falling factorial can
absorb at most one factor $\theta-j$. Every other factor remains.
After passing from forward to backward form, this is precisely the
asserted divisibility by $D_n(N)$, proving (1).

If $J$ is empty, no $\varepsilon^J_p$ equals $-1$, and
$q_1(0)=\prod_{p=0}^n(1+\varepsilon^J_p)\ne0$.
Thus $L^J$ is primitive. The already constructed recurrence has degree
$n$, and (1) proves optimality. For uniqueness, the leading coefficient
of any order-$n$ recurrence is a polynomial multiple of that of the
constructed recurrence. Subtracting the corresponding polynomial
multiple gives a recurrence of smaller order, which is zero by the
order bound. A density annihilator of order $r$ gives a moment recurrence
of degree at most $r$; the constructed density equation has order $n$.
This proves (2), including the density assertion.

Suppose now that $J=\{1\}$. Then $q_1(0)=0$. For $m\ge2$,
$x^{m-1}\mid q_m$, and hence $q_m(0)=0$. Since $L^JH=0$ and $H(0)=1$,
evaluation at zero gives $q_0(0)=0$. Equivalently,
$\hat P^J_1(0)=0$. The leading recurrence coefficient satisfies
$c_0(1)=0$, and $c_1(1)=-\hat P^J_1(0)=0$.
For $k\ge2$, the formula
$c_k(N)=(-1)^k\hat P^J_k(N-k)\ff{(N-1)}{k-1}$ directly gives
$c_k(1)=0$. Division by $N-1$ therefore produces polynomial
coefficients of degree $n-1$ and does not change recurrence order.
The recurrence is asserted only for $N\ge2$: its value at $N=1$
cannot be recovered by division. Its leading coefficient is $D_n$,
so (1) proves optimality and the preceding subtraction argument gives
uniqueness at order $n$.
The vanishing of every $q_m(0)$ also shows that $x$ divides $L^J$.
Thus $L^J/x$ annihilates $H$ and all of its coefficient polynomials
have degree at most $n-1$. Its Fourier--Laplace dual annihilates
$\rho^J$ and has order at most $n-1$. Conversely every density
annihilator gives a recurrence of no greater degree, so (1) forces
order at least $n-1$.
\end{proof}

\begin{remark}
For fixed $n$, the excluded set in part (2) is the finite union
\[
 p\left(a+b+1+\frac\beta2(2n-p-1)\right)=-j,
 \qquad 2\le p\le n,\quad 1\le j\le p-1.
\]
Thus part (2) applies on both sides of $(\star)$, except on these
hyperplanes. Part (1) is a lower bound on every remaining hyperplane;
we do not assert that it is always attained. The shifted-trace and
symmetric even-moment conclusions of Propositions~\ref{V:prop:shifted}
and~\ref{V:prop:shifted-sym} also hold with $(\star)$ replaced by $J=\varnothing$:
the primitive operator and its indicial polynomial are unchanged by
analytic exponential conjugation and nonzero scaling, and the preceding
multiplicity argument replaces the use of strictly negative roots.
\end{remark}

\begin{corollary}[The one- and two-particle cases]
\label{V:cor:jacobi-small-rank}
For $n=1$ the minimal recurrence order, recurrence degree, and density
differential order are all $1$ for every $a,b>-1$. The recurrence is
\[
 (N+a+b+1)m_N-(N+a)m_{N-1}=0\qquad(N\ge1).
\]
For $n=2$ the minimal recurrence order is always $2$. The minimum
recurrence degree and density differential order both equal
\[
 \begin{cases}
 1,&2a+2b+\beta+3=0,\\
 2,&2a+2b+\beta+3\ne0.
 \end{cases}
\]
On the exceptional hyperplane, necessarily $0<\beta<1$, and the reduced
recurrence is
\[
 \left(N+\frac{\beta-1}{2}\right)m_N
 -(3N+4a+2\beta)m_{N-1}
 +2(N+2a+\beta)m_{N-2}=0\qquad(N\ge2).
\]
The density satisfies, in the sense of distributions,
\[
 w(w-1)(w-2)(\rho^J)'(w)=
 \left(\frac{\beta-3}{2}w^2-(4a+2\beta)w+4a+2\beta+2\right)\rho^J(w).
\]
On each of $(0,1)$ and $(1,2)$ it is therefore a constant multiple of
\[
 w^{2a+\beta+1}|w-1|^{-(\beta+1)/2}(2-w)^{-2a-2},
\]
where the two constants are not asserted to be equal.
\end{corollary}
\begin{proof}
For $n=1$, $\varepsilon^J_1=a+b+1>-1$, so $J$ is empty.
For $n=2$, $\varepsilon^J_1=a+b+1+\beta>-1$ and
$\varepsilon^J_2=2a+2b+2+\beta>-2$.
The only possible positive-integer resonance is therefore
$\varepsilon^J_2=-1$, equivalent to the stated hyperplane. The preceding
theorem proves the orders and degrees.
On this hyperplane the scalarization gives
\[
 \begin{split}
 \hat P^J_0(\theta)&=\theta(\theta-1)
             \left(\theta+\frac{\beta-1}{2}\right),\\
 \hat P^J_1(\theta)&=\theta(3\theta+4a+2\beta+3),\\
 \hat P^J_2(\theta)&=2(\theta+2a+\beta+2).
 \end{split}
\]
These formulas give the reduced recurrence by coefficient extraction.
They also give the primitive operator
\[
 x\partial_x^3+
 \left(3x+\frac{\beta+3}{2}\right)\partial_x^2+
 (2x+4a+2\beta+6)\partial_x+4a+2\beta+4.
\]
For $H(x)=\langle\rho^J,e^{-xw}\rangle$, replace
$\partial_x$ by multiplication by $-w$ and $x$ by $\partial_w$,
retaining operator order. This gives the displayed first-order density
equation. Partial fractions and integration on each interval yield the
power expression. Finally $a+b>-2$ and
$a+b=-(\beta+3)/2$ imply $\beta<1$.
\end{proof}

\begin{remark}[The rank-one even orthogonal case]
The preceding multiplicity refinement also includes $\mathrm{SO}(2)$,
which was outside $(\star)$. Its trace is $2\cos\varphi$ with uniform
$\varphi$, so its even moments are $w_D(N)=\binom{2N}{N}$ and obey
\[
 Nw_D(N)-2(2N-1)w_D(N-1)=0\qquad(N\ge1).
\]
This recurrence has minimal order and degree $1$. It is also the
$n=1$, $a=b=-1/2$, $\kappa=4$ case of the extended symmetric-trace
statement, for which $\varepsilon^J_1=0$.
\end{remark}

\section{Parameter boundaries and remaining analytic questions}\label{V:sec:open}
\subsection{Noninteger twists}
The scalar construction and formal solutions in Theorem~\ref{V:thm:A}
hold for real $q>-1/2$. Minimality is established here for nonnegative
integer twists. For real $q\ge0$ the following remains a question.
\begin{conjecture}
For every $\beta>0$, $n\ge1$ and real $q\ge0$, the function
$G^\beta_{n,q}$ has differential order $n+1$ over $\mathbb C(x)$.
\end{conjecture}

The integer-twist proof does not extend merely by replacing the character
$e^{iq\sum\varphi_j}$ by a branch for non-integral $q$.  Such a branch is
not a smooth function on the torus.  Moreover, when $qn\notin\Z$, the
function $t^{qn}G^\beta_{n,q}(t^2)$ is not entire at $t=0$, because
$G^\beta_{n,q}(0)=1$; consequently it cannot be the Fourier--Laplace
transform of a compactly supported distribution on the real line.
There is, however, a compact distribution representation for the
ungauged function, for every real $q\ge0$.

\begin{proposition}[A compact transform for real twists]
\label{V:prop:real-twist-compact}
Let $\beta>0$, $n\ge1$, $q\ge0$, and put
$F(t)=G^\beta_{n,q}(t^2)$.  Then
\[
  |F(t)|\le \exp(2n|\operatorname{Re}t|)\qquad(t\in\C).
\]
There is a unique distribution $\sigma^\beta_{n,q}$ supported in
$[-2n,2n]$ such that
\[
 F(t)=\langle\sigma^\beta_{n,q},e^{tw}\rangle.
\]
This distribution is real-analytic away from the points
$2(n-2b)$, $0\le b\le n$.
\end{proposition}

\begin{proof}
Let $A=S^{-1}T_nS$ be the real symmetric matrix of
Lemma~\ref{V:lem:kac}(c), and let
\[
 D=\operatorname{diag}\bigl(-p(2q+1+\beta(n-p))\bigr)_{p=0}^n.
\]
Thus $D$ is negative semidefinite and $\|A\|=n$.
Consider the power-series solution
$Y(t)=\sum_{k\ge0}y_kt^k$ of
\[
 Y'=\left(2A+\frac Dt\right)Y,\qquad y_0=e_0.
\]
Its coefficients are uniquely determined by
\[
 (kI-D)y_k=2Ay_{k-1}\qquad(k\ge1).
\]
Since $\|(kI-D)^{-1}\|\le1/k$,
$\|y_k\|\le(2n)^k/k!$, and $Y$ is entire.
The tridiagonal form of $A$ implies
$(y_k)_p=0$ unless $p\le k$ and $p\equiv k\pmod2$.
Consequently $J=SY$ satisfies the system
\eqref{V:eq:J-system}; its first component is an even entire function.
The elimination in Section~\ref{V:sec:sys} shows that
$J_0(t)=g(t^2)$, where $\mathcal Lg=0$ and $g(0)=1$.
For every integer $N\ge1$ the constant-block coefficient
\[
 P_0(N)=\frac1{n!}\prod_{p=0}^n(N+\varepsilon_p)
\]
is nonzero, because $\varepsilon_p\ge0$.
The coefficient recurrence therefore gives a unique power-series
solution normalized at zero, so $g=G^\beta_{n,q}$.

For a fixed angle $\vartheta$ put $Z(r)=Y(re^{i\vartheta})$.
For $r>0$,
\[
 \frac{d}{dr}\|Z(r)\|^2
 =4\cos\vartheta\,Z(r)^*AZ(r)
   +\frac2r Z(r)^*DZ(r)
 \le 4n|\cos\vartheta|\,\|Z(r)\|^2.
\]
Since $\|Z(0)\|=1$, integration gives
$\|Y(t)\|\le e^{2n|\operatorname{Re}t|}$.
The first diagonal entry of $S$ is $1$, hence the same bound holds
for $F=J_0$.
The Paley--Wiener--Schwartz theorem applied to $F(iz)$ gives the
unique compactly supported distribution with support in $[-2n,2n]$;
see \cite[\S7.3]{Hor90}.

Finally the operator
$\Theta_t=\sum_kt^{2k}P_k(\theta_t/2)$ annihilates $F$.
By Proposition~\ref{V:prop:cbe-structure}, its top polynomial-degree
coefficient is
\[
 \chi_{\Theta_t}(w)
  =\sum_k\operatorname{lc}(P_k)
           (w/2)^{n+1-2k}
  =\frac1{n!}\prod_{b=0}^n\left(\frac w2-n+2b\right).
\]
Lemma~\ref{V:lem:dual} now proves the asserted analyticity.
\end{proof}

Proposition~\ref{V:prop:real-twist-compact} proves existence and localizes
the possible singularities.  It does not show that all those points
are singular.  To obtain minimal order $n+1$ for arbitrary real
$q\ge0$ by visibility, one still needs a nonvanishing singular term
at every point $2(n-2b)$.  In particular, the formal existence of all
exponential types for the operator does not establish their presence
in this particular solution.  The non-integer-twist minimality
conjecture remains open here.  The Hankel-contour representation
mentioned after \cite[(2.3)]{FW26} is another possible route to these
connection coefficients, but the compact real-density vertex
argument above does not apply to those contours without further work.

\subsection{Residual exceptional degree questions}
Theorem~\ref{vtd:even-degree} determines the least coefficient degree
for all positive even $\beta$ and integer twists, at every recurrence
order. Theorem~\ref{vtd:odd-degree-drop} shows that a degree drop can
occur at other parameters. A classification of all exceptional
non-even parameters is not claimed. The Jacobi refinements above
likewise distinguish exact conclusions from the residual resonant
cases. Explicit vertex connection bases for arbitrary non-even
$\beta$, and controlled local expansions on Hankel contours, remain
outside the present proofs.


\clearpage\part{Monodromy, algebraicity, and arithmetic specialization}

\section{Galois groups, jet fields, and values of twisted determinants}
\label{B:bv:section}

The absence of nonzero finite singularities in
Theorem~\ref{B:thm:twisted}(a) has consequences beyond the order of an
operator. It makes the change from Bessel monomials to scalar jets
invertible at \emph{every} nonzero argument. We determine that change
of coordinates over the Laurent polynomial ring, its determinant, the
complete algebraic relation ideal, and the arithmetic specialization.
These statements cover every integer twist $q\ge0$.

Put $R=\mathbb Q[x,x^{-1}]$, $K=\mathbb C(x)$ and
\[
 y(x)=I_0(2\sqrt{x})=\sum_{n\ge0}\frac{x^n}{(n!)^2},
 \qquad xy''+y'-y=0.
\]
The power series, rather than a choice of square root, defines $y$.
For integers $m$ define linear forms $C_m\in R[X,Y]$ by
\begin{equation}\label{B:bv:forms}
 C_0=X,\qquad C_1=Y,\qquad
 C_{m+1}=x^{-1}(C_{m-1}-mC_m)\quad(m\ge1),
 \qquad C_{-m}=x^mC_m\quad(m\ge1).
\end{equation}
For $\ell\ge1$ and $q\ge0$ put
\begin{equation}\label{B:bv:H}
 H_{\ell,q}(x;X,Y)=\det[C_{q+i-j}(x;X,Y)]_{i,j=1}^{\ell}.
\end{equation}
The Bessel recurrence gives
$t^{-m}I_m(2t)=C_m(t^2;y(t^2),y'(t^2))$ for all $m\in\mathbb Z$.
Factoring $t^{q+i-j}$ from the determinant therefore proves
\begin{equation}\label{B:bv:representation}
 G_{\ell,q}(x)=H_{\ell,q}(x;y(x),y'(x)).
\end{equation}
Thus $H_{\ell,q}$ is a nonzero homogeneous form of degree $\ell$.
Its nonvanishing also follows from the positive constant term of
$G_{\ell,q}$ in Lemma~\ref{B:lem:twisted-expansion}.

\subsection{The full group and algebraic base changes}

\begin{theorem}[Galois group and local monodromy]
\label{B:bv:galois}
For every $\ell\ge1$ and integer $q\ge0$, the minimal operator of
$G_{\ell,q}$ is irreducible over $K$ and has differential Galois group
\[
 \operatorname{Sym}^{\ell}(\mathrm{SL}_2)
 \cong
 \begin{cases}
  \mathrm{SL}_2,&\ell\text{ odd},\\
  \mathrm{PSL}_2,&\ell\text{ even}.
 \end{cases}
\]
Here the group is specified in its $(\ell+1)$-dimensional solution
representation. Its local monodromy at $x=0$ is a single unipotent
Jordan block $J_{\ell+1}(1)$, independently of $q$.
The order $\ell+1$ and irreducibility remain unchanged over every
finite algebraic extension of $K$. A ramified algebraic pullback at
zero preserves the size of this local block.
\end{theorem}

\begin{proof}
The Bessel connection for $y,y'$ has Galois group $\mathrm{SL}_2$.
Indeed its Wronskian is a nonzero constant times $x^{-1}$, so its
Galois group is contained in $\mathrm{SL}_2$; the pullback $x=t^2$
is equation~\eqref{B:eq:T}, whose group is $\mathrm{SL}_2$ by the
Bessel result used in Section~\ref{B:sec:unitary}.
In particular $y,y'$ are algebraically independent over $K$.
To see the latter assertion directly, choose a second solution $v$.
The orbit of $y$ consists of $ay+bv$, $(a,b)\ne(0,0)$.
A polynomial relation for $y,y'$ would consequently vanish at
$(ay+bv,ay'+bv')$ for all such constants $a,b$. The nonzero Wronskian
makes this an invertible linear substitution in $a,b$, forcing the
polynomial to be zero.

Let $V$ be the two-dimensional Bessel solution space. Polarization
of $u\mapsto H_{\ell,q}(x;u,u')$ defines an equivariant linear map
from $\operatorname{Sym}^{\ell}V$ into the Picard--Vessiot field.
This map is nonzero by \eqref{B:bv:representation}. Since
$\operatorname{Sym}^{\ell}V$ is irreducible, it is injective. Its
image is finite-dimensional and the orbit of $G_{\ell,q}$ spans it.
The minimal scalar equation has precisely this solution space,
so its group is the displayed image of $\mathrm{SL}_2$. The kernel
of the symmetric-power representation is trivial for odd $\ell$
and $\{I,-I\}$ for even $\ell$.

At zero the rank-two Bessel equation has one analytic solution and
a second solution containing a nonzero logarithmic term. Its local
monodromy is a nontrivial unipotent block of size two. Its $\ell$th
symmetric power is a single block of size $\ell+1$: the logarithm
acts on homogeneous binary forms by a nonzero multiple of
$X\partial_Y$, with one chain of that length. The rational
polynomial solution map above does not change local monodromy.

The group $\operatorname{Sym}^{\ell}(\mathrm{SL}_2)$ is connected.
A finite algebraic extension of the base preserves its identity
component, hence its full group and irreducible representation.
Ramification of index $e\ge1$ raises the local monodromy to its
$e$th power; in characteristic zero a unipotent block retains its
size under this operation.
\end{proof}

\subsection{An everywhere invertible Laurent gauge}

Let
\[
 \mathbf J_{\ell,q}=(G_{\ell,q},G'_{\ell,q},\ldots,
                        G_{\ell,q}^{(\ell)})^{\mathsf T},
 \qquad
 \mathbf m_{\ell}=(y^{\ell},y^{\ell-1}y',\ldots,(y')^{\ell})^{\mathsf T}.
\]
Define $P_{\ell,q}(x)$ by the coefficient identities
\begin{equation}\label{B:bv:gauge}
 \mathbf J_{\ell,q}=P_{\ell,q}(x)\mathbf m_{\ell}.
\end{equation}
It is computed without differential elimination: the $j$th row is
the coefficient vector of $\mathscr D^jH_{\ell,q}$ in the monomials
$X^{\ell-i}Y^i$, where
\[
 \mathscr D=\partial_x+Y\partial_X+x^{-1}(X-Y)\partial_Y.
\]
Thus every entry of $P_{\ell,q}$ lies in $R$.

\begin{theorem}[Exact determinant of the jet gauge]\label{B:bv:gaugedet}
One has $P_{\ell,q}\in\mathrm{GL}_{\ell+1}(R)$ and
\begin{equation}\label{B:bv:det}
 \det P_{\ell,q}(x)=c_{\ell,q}
 x^{-\ell(\ell+1)(\ell+3q-1)/6},
 \qquad c_{\ell,q}=\det P_{\ell,q}(1)\in\mathbb Q\setminus\{0\}.
\end{equation}
In particular both $P_{\ell,q}$ and its inverse are regular and
invertible at every $x\in\mathbb C^*$.
\end{theorem}

\begin{proof}
If the rows of $P_{\ell,q}$ were dependent over $K$,
\eqref{B:bv:gauge} would give an equation of order at most $\ell$
for $G_{\ell,q}$, contradicting Theorem~\ref{B:thm:twisted}(a).
Hence its determinant is nonzero.
Write $\mathbf m_\ell'=A_\ell\mathbf m_\ell$. The diagonal entry
of $A_\ell$ belonging to $y^{\ell-i}(y')^i$ is $-i/x$, whence
\[
 \operatorname{tr}A_\ell=-\frac{\ell(\ell+1)}{2x}.
\]
Let $B_{\ell,q}$ be the companion matrix of the monic minimal
scalar equation. Theorem~\ref{B:thm:twisted}(a) says that the ratio
of its two leading polynomial coefficients is $a/x$ for a
constant $a$. Put
\[
 s_{\ell,q}=\sum_{i=0}^{\ell}-i(\ell+q-i)
       =-\frac{\ell(\ell+1)(\ell+3q-1)}6.
\]
Comparison of the coefficient of $\rho^\ell$ in its indicial
polynomial gives $a=\ell(\ell+1)/2-s_{\ell,q}$. Therefore
$\operatorname{tr}B_{\ell,q}=-a/x$.
Differentiating \eqref{B:bv:gauge} and using the independence of
its monomials yields
$P_{\ell,q}'=B_{\ell,q}P_{\ell,q}-P_{\ell,q}A_\ell$.
Taking traces after multiplying by $P_{\ell,q}^{-1}$ gives
\[
 \frac{(\det P_{\ell,q})'}{\det P_{\ell,q}}
   =\operatorname{tr}B_{\ell,q}-\operatorname{tr}A_\ell
   =\frac{s_{\ell,q}}x.
\]
Its nonzero rational solution is exactly \eqref{B:bv:det}, with
$c_{\ell,q}\in\mathbb Q^*$. Since $x^{s_{\ell,q}}$ is a unit
of $R$, the adjugate formula proves that $P_{\ell,q}^{-1}$
also has entries in $R$.
\end{proof}

\begin{proposition}[The rational gauge constant]\label{B:bv:constant}
The constant in \eqref{B:bv:det} has the closed form
\[
 c_{\ell,q}=(-1)^{\binom{\ell+1}{3}+q\binom{\ell+1}{2}}
          \prod_{j=1}^{\ell}(j!)^{2\ell-2j+1}.
\]
\end{proposition}
\begin{proof}
Put $N=\binom{\ell+1}{2}$, $t=\sqrt{x}$,
$u(x)=I_0(2t)$ and $v(x)=2K_0(2t)$.
Their $x$-Wronskian is $-x^{-1}$, so the determinant of the
symmetric-power fundamental matrix is $(-1)^Nx^{-N}$.
By Proposition~\ref{B:prop:twisted-basic}(a), its image under
$P_{\ell,q}$ has columns
\[
 \binom{\ell}{b}
 \bigl(f_b,f_b',\ldots,f_b^{(\ell)}\bigr)^{\mathsf T},
 \qquad f_b(x)=t^{-q\ell}\Psi_b^{(q)}(t),\quad 0\le b\le\ell.
\]
The vertex asymptotics in that proposition give
\[
 f_b(x)\sim(-1)^{qb}C_b e^{2(\ell-2b)t}
                   t^{-\gamma_b-q\ell}.
\]
Their differentiated expansions show that the leading Wronskian
factor is the Vandermonde determinant in $\ell-2b$, multiplied
by $t^{-N}\prod_bf_b$. This remains valid when one of these
numbers is zero: the determinant of the leading coefficient
matrix is still the nonzero Vandermonde determinant.
Its value is $(-2)^N\prod_{j=1}^{\ell}j!$.
Comparing with \eqref{B:bv:det} consequently gives
\[
 c_{\ell,q}=(-1)^{qN}2^N
  \left(\prod_{j=1}^{\ell}j!\right)
  \left(\prod_{b=0}^{\ell}\binom{\ell}{b}C_b\right).
\]
For clarity, the normalization of the Mehta integral in
Lemma~\ref{B:lem:vertex} is
\[
 \mathcal M_n=\pi^{n/2}2^{-n(n-1)/2}\prod_{j=1}^n j!.
\]
This follows by applying Andreief to monic Hermite polynomials,
whose squared norms for the weight $e^{-x^2}$ are
$\sqrt{\pi}\,2^{-j}j!$. Put $A_b=\prod_{j=1}^b j!$, $A_0=1$.
Substitution into the displayed formula for $C_b$ yields
\[
 \prod_{b=0}^{\ell}C_b
 =(-1)^{\binom{\ell+1}{3}}2^{-N}
       \frac{\prod_{b=0}^{\ell}A_b^2}{(\ell!)^{\ell+1}},
 \qquad
 \prod_{b=0}^{\ell}\binom{\ell}{b}
       =\frac{(\ell!)^{\ell+1}}{A_\ell^2}.
\]
The powers of $\pi$ cancel. Combining these identities gives
$c_{\ell,q}=(-1)^{\binom{\ell+1}{3}+qN}
 A_\ell\prod_{b=0}^{\ell-1}A_b^2$, which is the asserted product.
\end{proof}

\subsection{Algebraic jet relations and twist fields}

\begin{theorem}[The jet cone]\label{B:bv:jetcone}
Set $\mathbf Z=P_{\ell,q}^{-1}\mathbf J_{\ell,q}$.
The ideal of all polynomial relations over $K$ among the
coordinates $Z_0,\ldots,Z_\ell$ is generated by
\begin{equation}\label{B:bv:minors}
 Z_iZ_{j+1}-Z_{i+1}Z_j,
 \qquad 0\le i<j\le\ell-1.
\end{equation}
For $\ell=1$ this is the zero ideal. In particular
\[
 \operatorname{trdeg}_{K}K(\mathbf J_{\ell,q})=2,
\]
and every two distinct entries of $\mathbf J_{\ell,q}$ are
algebraically independent over $K$. These assertions persist
after finite algebraic extension of $K$.
Moreover
\begin{equation}\label{B:bv:field}
 K(\mathbf J_{\ell,q})
 =K(y^\ell,y^{\ell-1}y',\ldots,(y')^\ell)
 =K(y,y')^{\mu_\ell},
\end{equation}
where $\mu_\ell$ acts by simultaneous scalar multiplication.
The index of this field in $K(y,y')$ is exactly $\ell$.
\end{theorem}

\begin{proof}
The coordinates $Z_i=y^{\ell-i}(y')^i$ are the degree-$\ell$
monomials in two algebraically independent elements. Their
relation ideal is the ideal of the affine cone over the rational
normal curve. An elementary proof of its generators is as follows.
The minors \eqref{B:bv:minors} allow a pair of indices whose gap
is at least two to be replaced by a closer pair with the same
sum. Repeating this operation makes all indices differ by at
most one. Thus every product of $r$ coordinates is congruent
to the unique such balanced product having the same sum of
indices. These are exactly the equivalence classes of monomials
with the same image in $K[y,y']$. Grouping terms according to
their images proves that the minors generate the kernel.

Two distinct rows of $P_{\ell,q}$ are linearly independent.
The corresponding homogeneous forms $A,B$ of degree $\ell$
are therefore nonproportional, and are algebraically independent.
Indeed $A/B$ is a nonconstant rational function of $Y/X$, while
common scaling of $X,Y$ changes the degree-$\ell$ magnitude.
More formally, splitting a putative relation into homogeneous
pieces would make $A/B$ algebraic over $K$, a contradiction.
This proves pairwise independence of the scalar jets.
An algebraic extension of the coefficient field cannot decrease
their transcendence degree.

Finally, writing $r=y'/y$ shows that the middle field in
\eqref{B:bv:field} is $K(r,y^\ell)$. Since $r,y$ are algebraically
independent, $K(r,y)/K(r,y^\ell)$ is a cyclic extension of degree
$\ell$, with precisely the stated scalar action.
\end{proof}

\begin{corollary}[Twist invariance and the gcd--lcm law]
\label{B:bv:fieldlaw}
For fixed $\ell$, every integer twist $q\ge0$ gives the same
complete jet field. Denote it by $\mathcal F_\ell$.
For positive integers $d,e$,
\[
 \mathcal F_d\mathcal F_e=\mathcal F_{\gcd(d,e)},
 \qquad
 \mathcal F_d\cap\mathcal F_e=\mathcal F_{\operatorname{lcm}(d,e)}.
\]
Thus complete jets in any two coprime positive dimensions recover
$K(y,y')$. For a fixed dimension, the jets at twists $q$ and $r$
are related by a matrix
$P_{\ell,r}P_{\ell,q}^{-1}\in\mathrm{GL}_{\ell+1}(R)$.
\end{corollary}

\begin{proof}
The last assertion and twist invariance follow from
\eqref{B:bv:gauge}. Put $L=\operatorname{lcm}(d,e)$.
In the cyclic Galois extension
$K(y,y')/K(y,y')^{\mu_L}$, the two fields are fixed by
$\mu_d$ and $\mu_e$. Their compositum is fixed by the
intersection $\mu_{\gcd(d,e)}$, and their intersection is
fixed by the group they generate, $\mu_L$.
\end{proof}

\subsection{Every nonzero algebraic argument}

\begin{theorem}[Exact arithmetic jet relations for every twist]
\label{B:bv:values}
Let $\ell\ge1$, $q\ge0$ be integers and let
$\alpha\in\overline{\mathbb Q}\setminus\{0\}$. Then
\[
 \operatorname{trdeg}_{\overline{\mathbb Q}}
 \overline{\mathbb Q}\bigl(G_{\ell,q}(\alpha),
 G_{\ell,q}'(\alpha),\ldots,G_{\ell,q}^{(\ell)}(\alpha)\bigr)=2.
\]
Every two distinct numbers in this tuple are algebraically
independent. Under the invertible algebraic coordinate change
\[
 \mathbf Z=P_{\ell,q}(\alpha)^{-1}\mathbf J_{\ell,q}(\alpha),
\]
all polynomial relations are generated by \eqref{B:bv:minors}.
In particular each displayed derivative value is nonzero and
transcendental, and $G'_{\ell,q}(\alpha)/G_{\ell,q}(\alpha)$
is transcendental.
\end{theorem}

\begin{proof}
It is essential here to apply the $E$-function theorem in the
variable $t$, not to assume that $I_0(2\sqrt{x})$ is an
$E$-function in $x$. Put $b(t)=I_0(2t)$.
Its exponential-series coefficients of index $2n$ are
$\binom{2n}{n}$, and its odd coefficients are zero. They are
integers of exponential growth. The same is true of the
shifted coefficient sequence of $b'$. Thus $b,b'$ are
$E$-functions, satisfying
\[
 \frac d{dt}\begin{pmatrix}b\\b'\end{pmatrix}
 =\begin{pmatrix}0&1\\4&-1/t\end{pmatrix}
 \begin{pmatrix}b\\b'\end{pmatrix}.
\]
They are algebraically independent over $\mathbb C(t)$ by the
Bessel Galois argument in the proof of Theorem~\ref{B:bv:galois}.
The Siegel--Shidlovskii theorem
\cite[Theorem~1.1]{BVBeukers2006}, whose exceptional points are
contained in the zeros of $t$ times the common denominator of
the system, proves algebraic independence of $b(s),b'(s)$
for every nonzero algebraic $s$.

Choose $s$ with $s^2=\alpha$. Since
\[
 y(\alpha)=b(s),\qquad y'(\alpha)=\frac{b'(s)}{2s},
\]
the two base values are algebraically independent over
$\overline{\mathbb Q}$. Theorem~\ref{B:bv:gaugedet} proves that
$P_{\ell,q}(\alpha)$ is an invertible algebraic matrix.
We can therefore repeat the polynomial-kernel and homogeneous-form
arguments of Theorem~\ref{B:bv:jetcone}, over
$\overline{\mathbb Q}$ and with these two independent numbers
in place of $y,y'$. This proves the exact ideal, transcendence
degree and pairwise independence. Each individual coordinate is
a nonzero positive-degree homogeneous form in two independent
numbers and is consequently transcendental. An algebraic ratio
of two distinct jet coordinates would give a linear algebraic
relation between them, so their ratio is transcendental as well.
\end{proof}

\begin{corollary}[Zeros and arithmetic field recovery]
\label{B:bv:zeros}
Every complex zero of $G_{\ell,q}^{(j)}$, for $0\le j\le\ell$,
is transcendental. At each fixed
$\alpha\in\overline{\mathbb Q}^{\,*}$, the numerical jet fields
\[
 \mathcal F_\ell(\alpha)
 =\overline{\mathbb Q}\bigl(G_{\ell,q}^{(j)}(\alpha):0\le j\le\ell\bigr)
\]
are independent of $q$ and satisfy the same gcd--lcm laws as in
Corollary~\ref{B:bv:fieldlaw}. In particular two coprime dimensions
recover
$\overline{\mathbb Q}(y(\alpha),y'(\alpha))$ exactly.
\end{corollary}

\begin{proof}
Theorem~\ref{B:bv:values} excludes every nonzero algebraic zero.
At zero all the indicated derivatives are strictly positive
rational numbers, by the coefficient expansion in
Lemma~\ref{B:lem:twisted-expansion}; hence zero is excluded as well.
The two base values are algebraically independent, so their
field is a rational function field in two generators over
$\overline{\mathbb Q}$. The proof of
Corollary~\ref{B:bv:fieldlaw} applies verbatim to it.
\end{proof}

\begin{remark}[What the arithmetic conclusion concerns]
The result concerns the normalized Bessel determinant and its
ordinary derivatives at nonzero algebraic arguments. It does
not assert that arbitrary periods, ordinary generating functions,
or unrelated numerical constants are transcendental. The
no-exception assertion uses the Laurent-unit determinant
\eqref{B:bv:det}; abstract symmetric-power irreducibility alone
would not supply it.
\end{remark}

\section{Exact exponents at every nonzero trace threshold}
\label{B:upden:section}

The indicial exponent described heuristically in the original final
remarks can be determined from the first two coefficient layers of the
minimal recurrence.  This is a statement about the scalar equation;
it does not by itself assert a nonzero singular amplitude for a chosen
physical branch of the density.

\begin{theorem}[The radial unitary thresholds]\label{B:upden:unitary}
Let $\ell\ge2$, put $d=\ell-1$ and $K=\lceil\ell/2\rceil$, and let
$\mathcal L_\ell$ be the density operator in Theorem~\ref{B:thm:E}.
For every $0\le j<K$, the nonzero point $c_j=(\ell-2j)^2$ is a
regular singular point. Its full indicial multiset is
\[
 0,1,\ldots,\ell-3,
 \quad \frac{(\ell-j)^2+j^2-3}{2}.
\]
The list $0,\ldots,\ell-3$ is empty when $\ell=2$.
The same statement holds for the polynomial operator
$\mathcal L_\ell^\circ$ displayed in that theorem, since the removed
power of $w$ is invertible at $c_j$. Repeated exponents in this list
are counted with multiplicity.
\end{theorem}

\begin{proof}
Write $h_k=d+2-2k$ and denote the top two coefficients by
\[
 P_k(\theta)=r_k\theta^{h_k}+p_k\theta^{h_k-1}+\cdots,
 \qquad
 R_k(n)=r_kn^d+s_kn^{d-1}+\cdots.
\]
When $h_k=0$, set $p_k=0$.  The normalization of
Theorem~\ref{B:thm:B} gives, for every $0\le k\le K$,
\begin{equation}\label{B:upden:subleading}
 s_k=p_k-k(d+1-k)r_k.
\end{equation}
For $k=0$ this follows from $R_0=P_0/n^2$. For $k\ge1$, expand
$P_k(n-k)\prod_{a=1}^{k-1}(n-a)^2$ to its top two powers.

Set
\[
 H(w)=\sum_{k=0}^Kr_kw^{K-k}
     =\prod_{j=0}^{K-1}(w-c_j),\qquad
 F(w)=\sum_{k=0}^K\{p_k+(k^2-k-d)r_k\}w^{K-k}.
\]
The density operator is
$\sum_{k=0}^K w^{K-k}R_k(k-1-\theta)$, where
$\theta=w\partial_w$. Since
$\theta^d=w^d\partial_w^d+\binom d2w^{d-1}\partial_w^{d-1}+\cdots$,
its top two ordinary derivative coefficients are
\begin{align*}
 a_d(w)&=(-1)^dw^dH(w),\\
 a_{d-1}(w)&=(-1)^dw^{d-1}\{\tbinom d2H(w)-F(w)\}.
\end{align*}
At a nonzero simple root $c$ of $H$, the indicial polynomial is
proportional to
\[
 \rho^{\underline{d-1}}
 \left(\rho-d+1+\frac{a_{d-1}(c)}{a_d'(c)}\right).
\]
Thus the extra exponent is $d-1+F(c)/(cH'(c))$.

It remains to evaluate this ratio without computing all the $P_k$.
Fix $j<K$, put $\mu=\ell-2j$, $c=\mu^2$,
$\gamma=((\ell-j)^2+j^2)/2$, and
$\epsilon=\ell+1-2K\in\{0,1\}$.  Define
\[
 C(v)=\sum_{k=0}^Kr_kv^{h_k}=v^\epsilon H(v^2).
\]
The formal solution of Theorem~\ref{B:thm:A}(c), in $x=t^2$, is
$e^{2\mu t}t^{-\gamma}(1+O(t^{-1}))$. Substitution into
$\sum_k t^{2k}P_k(\tfrac12t\partial_t)$ gives at the next power
after the leading characteristic equation
\begin{equation}\label{B:upden:formalidentity}
 \sum_{k=0}^Kp_k\mu^{h_k-1}
      =\frac\gamma2C'(\mu)-\frac\mu4C''(\mu).
\end{equation}
Indeed the subleading term of
$(\tfrac12t\partial_t)^{h_k}$ on the displayed formal solution is
$\{h_k(h_k-1)/4-h_k\gamma/2\}\mu^{h_k-1}t^{h_k-1}$.
The next coefficient of $1+O(t^{-1})$ multiplies $C(\mu)=0$.

With $E=v\partial_v$ and $k=(d+2-h_k)/2$, direct expansion gives
\[
 \sum_k(k^2-k-d)r_kv^{h_k}
 =\frac14\{v^2C''(v)-(2d+1)vC'(v)+d(d-2)C(v)\}.
\]
Multiply \eqref{B:upden:formalidentity} by $\mu$ and add this
identity at $v=\mu$. The second derivative terms cancel, and
$C(\mu)=0$, so
\[
 \mu^\epsilon F(c)
   =\frac{2\gamma-2d-1}{4}\mu C'(\mu)
   =\left(\gamma-d-\frac12\right)\mu^\epsilon cH'(c).
\]
The extra exponent is therefore $\gamma-3/2$, as asserted.
\end{proof}

\begin{lemma}[A simple exponential rate gives a density exponent]
\label{B:upden:fourier}
Suppose a polynomial differential operator has the form
\[
 L=t^d\chi(\partial_t)+t^{d-1}\psi(\partial_t)
      +\text{terms of smaller degree in }t,
\]
where $d\ge1$. If $\lambda$ is a simple root of $\chi$ and
$e^{\lambda t}t^{-\gamma}(1+O(t^{-1}))$ is a formal solution,
then the Fourier operator under $t\mapsto-\partial_w$,
$\partial_t\mapsto w$ has at $w=\lambda$ the indicial multiset
\[
                     0,1,\ldots,d-2,\gamma-1.
\]
\end{lemma}
\begin{proof}
The first subleading formal coefficient gives
$\psi(\lambda)=\gamma\chi'(\lambda)$.
The two highest derivative coefficients of the Fourier operator are
$(-1)^d\chi(w)$ and $(-1)^d(d\chi'(w)-\psi(w))$.
Its leading coefficient has a simple zero at $\lambda$, hence the
point is regular singular. The simple-zero indicial formula used
above gives the remaining exponent
$d-1-(d-\gamma)=\gamma-1$.
\end{proof}

\begin{corollary}[Tableau and symplectic trace thresholds]
\label{B:upden:other}
For the tableau trace law with $k=2r+\varepsilon\ge2$,
the Fourier transform of $L_k$ has at each
$\lambda_j^{(k)}=\varepsilon+2r-4j$ the exponents
\[
 0,1,\ldots,r-2,\qquad
 (r-j)(r-j+\varepsilon-\tfrac12)+j(j+\tfrac12)-1.
\]
For the trace law on $\operatorname{Sp}(2\ell)$, the Fourier
transform of the operator of $e^{-t}Y_{2\ell+1}$ has at
$2(\ell-2j)$ the exponents
\[
 0,1,\ldots,\ell-2,\qquad
 (\ell-j)(\ell-j+\tfrac12)+j(j+\tfrac12)-1.
\]
These statements include a zero threshold whenever it occurs.
\end{corollary}
\begin{proof}
Apply Lemma~\ref{B:upden:fourier} to the degree-$r$ operator in
Theorem~\ref{B:thm:jacobi} and its degree-$\ell$ symplectic gauge
transform. Their characteristic roots are distinct and their formal
powers are the $\gamma_j^{(k)}$ in \eqref{B:eq:jacobi-data}.
\end{proof}

\begin{remark}
Integral differences among these exponents permit logarithmic
solutions, but do not force a logarithm in every solution.
In particular the indicial data do not replace a connection formula
for the actual probability density. All Fourier transforms above are
formed before any cancellation by a function vanishing at a threshold.
\end{remark}



\section{Exact recurrence widths after arithmetic sampling}
\label{B:up:subsequence-section}

For a nonzero sequence $a_n$, its \emph{eventual recurrence width} is
the least $L$ for which
\[
 \sum_{j=0}^{L}p_j(n)a_{n-j}=0
\]
holds for every sufficiently large integer $n$, with
$p_j\in\C(n)$ and $p_0p_L\ne0$.  Clearing denominators gives the
same minimum with polynomial coefficients.  The following result
concerns every sampling interval, without a genericity assumption.

\begin{theorem}[Recurrence widths under arithmetic sampling]
\label{B:up:sample-order}
Let $d\ge1$, $0\le b<d$, $\ell\ge1$, and $q\ge0$ be integers.
Then
\[
 \operatorname{width}\bigl(u_{\ell,q}(dn+b)\bigr)
 =\operatorname{width}\bigl(m_\ell(dn+b)\bigr)
 =\left\lceil\frac\ell2\right\rceil.
\]
For tableaux, put $r=\lfloor k/2\rfloor$.  For every $k\ge2$,
\[
 \operatorname{width}\bigl(y_k(dn+b)\bigr)=
 \begin{cases}
 r+1,&k\text{ odd},\\
 2\lceil r/2\rceil,&k\text{ even},\ d\text{ odd},\\
 \lceil r/2\rceil,&k\text{ even},\ d\text{ even}.
 \end{cases}
\]
All these conclusions remain valid after multiplying the sampled
sequence by $c^n h(n)$, where $c\in\C^*$ and
$0\ne h\in\C(n)$, with the finitely many poles omitted.
\end{theorem}

The untwisted assertion $q=0$ applies to permutations with bounded
increasing subsequences.  For even tableau height, sampling on an even
progression halves the original exact width.  This reduction is caused
by antipodal Bessel components; it is independent of the residue $b$.

\begin{lemma}[Sampling upper bound]\label{B:up:sample-upper}
If a sequence satisfies an eventual recurrence of width $R$, then
every arithmetic subsequence satisfies one of width at most $R$.
\end{lemma}
\begin{proof}
Outside finitely many indices, a recurrence of width $R$ gives a
first-order system $v_{n+1}=A(n)v_n$ for a vector of $R$ consecutive
terms, with $A(n)$ over $\C(n)$.  On indices $n=dm+b$ this becomes
\[
 v_{d(m+1)+b}
 =A(dm+b+d-1)\cdots A(dm+b)v_{dm+b}.
\]
Each of $R+1$ consecutive sampled outputs is a rational row vector
applied to $v_{dm+b}$.  These rows are linearly dependent over
$\C(m)$, yielding a recurrence of width at most $R$ after deleting
zero extreme coefficients.  Since the sampled sequences in the
theorem have positive terms, the relation cannot have width zero.
\end{proof}

\begin{lemma}[Distinct simple components cannot cancel]
\label{B:up:simple-filter}
Suppose $V_1,\ldots,V_s$ are pairwise nonisomorphic simple
differential modules over $\C(t)$, realized as differential
submodules of a common field of meromorphic functions, and
$0\ne f_j\in V_j$.  The differential module generated by
$\sum_j f_j$ is their direct sum.
\end{lemma}
\begin{proof}
The sum of pairwise nonisomorphic simple submodules is direct.
It is semisimple, and any submodule is the direct sum of a subset
of these simple summands: this follows inductively from simplicity
and the fact that a nonzero homomorphism between simple modules
is an isomorphism.  The submodule generated by $\sum_jf_j$ projects
nontrivially to each summand, so it contains every one of them.
\end{proof}

\begin{lemma}[A factorial budget for a sampled recurrence]
\label{B:up:sample-budget}
Let $p,h\ge1$, $c\ge0$, and $\alpha_1,\ldots,\alpha_p\ge0$ be
integers, and put
\[
 H(t)=\sum_{n\ge0}
       \frac{a_n t^{p(hn+c)}}
       {\prod_{i=1}^{p}(hn+c+\alpha_i)!}.
\]
Suppose the minimal scalar equation of $H$ has formal solutions
at infinity of the form
$e^{\beta_jt}t^{\mu_j}(1+O(t^{-1}))$, with $\nu$ distinct nonzero
numbers $\beta_j$.  Then every eventual recurrence of $(a_n)$
has width at least $\nu/(ph)$.
\end{lemma}
\begin{proof}
Clear the denominators of a recurrence
$\sum_{j=0}^{L}p_j(n)a_{n-j}=0$ and let
$D=\max_j\deg p_j$.  With $\theta=t\partial_t$, form
\[
 P=\sum_{j=0}^{L}t^{phj}
 p_j\left(\frac{\theta}{ph}+j-\frac ch\right)
 \prod_{i=1}^{p}
 \left(\frac\theta p+\alpha_i\right)^{\underline{h(L-j)}}.
\]
For sufficiently large $n$, the coefficient of
$t^{p(hn+c)}$ in $PH$ equals the proposed recurrence divided by
$\prod_i(hn+c-hL+\alpha_i)!$.  Consequently $PH$ is a polynomial.
A polynomial in $\theta$ annihilates this finite remainder.

For $P=\sum_m b_m(t)\partial_t^m$, the largest degree of the
polynomials $b_m$ is $D+phL$.  Its leading polynomial in the
derivative variable is
\[
 \sigma_P(\xi)=
 \sum_{\deg p_j=D}\operatorname{lc}(p_j)
 \left(\frac\xi{ph}\right)^D
 \left(\frac\xi p\right)^{ph(L-j)}.
\]
The powers for distinct $j$ differ, so this polynomial is nonzero.
After a power of $\xi$ is removed, its degree is at most $phL$.
The additional polynomial in $\theta$ used to remove the finite
remainder changes only this power of $\xi$.

Every homogeneous annihilator is a left multiple of the minimal
operator of $H$.  It therefore annihilates each of its formal
solutions.  Substituting
$e^{\beta_jt}t^{\mu_j}(1+O(t^{-1}))$ shows that
$\sigma_P(\beta_j)=0$.  Since the $\beta_j$ are distinct and
nonzero, $\nu\le phL$.
\end{proof}


\begin{lemma}[A Laurent Bessel form has a short parity recurrence]
\label{B:up:laurent-parity}
Let $r\ge1$, $y(t)=I_0(2t)$, and let
$\Phi\in\C[t,t^{-1}][X_0,X_1]$ be homogeneous of degree $r$.
Suppose $F(t)=\Phi(t,y(t),y'(t))$ is nonzero, holomorphic at zero and
$F(-t)=(-1)^pF(t)$, where $p\in\{0,1\}$. Write
\[
 F(t)=\sum_{n\ge0}a_n\frac{t^{2n+p}}{(2n+p)!}.
\]
Then $(a_n)$ and all its arithmetic subsequences have eventual
recurrence width at most $\lceil r/2\rceil$.
\end{lemma}
\begin{proof}
For $0\le j\le r$, differentiation of $y^r$ using
$y''=4y-t^{-1}y'$ gives
\[
 \partial_t^j(y^r)
 =r^{\underline j}y^{r-j}(y')^j
   +\sum_{i<j}a_{j,i}(t)y^{r-i}(y')^i,
 \qquad a_{j,i}\in\C[t,t^{-1}].
\]
The change-of-basis matrix is triangular with nonzero constant
 diagonal entries.  Its inverse therefore also has Laurent-polynomial
entries.  It follows that $F=R(y^r)$ for a differential operator
$R\in\C[t,t^{-1}]\langle\partial_t\rangle$ of order at most $r$.
Replacing $R$ by its parity-$p$ part under
$(t,\partial_t)\mapsto(-t,-\partial_t)$, we may take every monomial
$t^a\partial_t^j$ occurring in $R$ to satisfy $a+j\equiv p\pmod2$.

Put
\[
 W_r(n)=\sum_{n_1+\cdots+n_r=n}\binom{n}{n_1,\ldots,n_r}^2,
 \qquad y(t)^r=\sum_{n\ge0}W_r(n)\frac{t^{2n}}{(n!)^2}.
\]
Verrill's recurrence has width $K=\lceil r/2\rceil$
\cite[Theorem~1 and equation~(15)]{Ver04}; its Bessel derivation
is recalled in Remark~\ref{B:rem:walks}.  For $r=1$ this is simply
$W_1(n)=1$.
The contribution of $t^a\partial_t^j(y^r)$ to $a_n$ is
\[
 (2n+p)!\,
 \frac{(2n+2c)^{\underline j}}{((n+c)!)^2}\,W_r(n+c),
 \qquad c=\frac{p+j-a}{2}\in\Z,
\]
for every sufficiently large $n$. Dividing by
$\binom{2n}{n}$ makes its coefficient of $W_r(n+c)$ a rational
function of $n$.  Thus $a_n/\binom{2n}{n}$ is a finite
rational-function linear combination of shifts of $W_r$.

A recurrence of width $K$ expresses every such shift through
$K$ consecutive values, using forward or backward shifts outside
finitely many poles.  Hence $K+1$ consecutive values of any
rational linear combination are linearly dependent over
$\C(n)$, and its width is at most $K$.  Finally
$\binom{2n+2}{n+1}/\binom{2n}{n}=2(2n+1)/(n+1)$ is rational,
so restoring this hypergeometric factor preserves the width.
Lemma~\ref{B:up:sample-upper} handles further sampling.
\end{proof}

\begin{proof}[Proof of Theorem~\ref{B:up:sample-order}]
Let $\zeta=e^{2\pi i/d}$ and choose
$\omega_j=e^{\pi i j/d}$, $0\le j<d$, so
$\omega_j^2=\zeta^j$.  Root-of-unity filtering gives
\[
 H(t)=\frac1d\sum_{j=0}^{d-1}\zeta^{-bj}
 G_{\ell,q}(\zeta^jt^2)
 =\sum_{n\ge0}
 \frac{u_{\ell,q}(dn+b)t^{2(dn+b)}}
 {(dn+b)!(dn+b+q\ell)!}.
\]
Each term on the left generates a simple differential module:
it is a rational gauge of the $\ell$th symmetric power of the
Bessel module at scale $\omega_j$.  Simplicity follows from
the $\mathrm{SL}_2$ Galois group and irreducibility of its
symmetric powers, as in the proof of the original minimal-order
theorem.  Its formal exponential rates are
\[
 2\omega_j(\ell-2i),\qquad 0\le i\le\ell.
\]
Two such sets could coincide only if their scales differed by a
sign: compare the two extreme points of largest modulus.
The chosen $\omega_j$ never differ by a sign.  Thus the modules
are pairwise nonisomorphic, and Lemma~\ref{B:up:simple-filter}
shows that the filtered function contains all of them.
Furthermore all their nonzero rates are distinct: a ratio of two
nonzero real rates can have modulus one only when it is $1$ or
$-1$.  There are exactly
$2d\lceil\ell/2\rceil$ such rates.
Lemma~\ref{B:up:sample-budget}, with
$p=2$, $h=d$, $c=b$, and $(\alpha_1,\alpha_2)=(0,q\ell)$,
gives the required lower bound.  The recurrence of
Theorem~\ref{B:thm:twisted} and Lemma~\ref{B:up:sample-upper}
give the matching upper bound.

For matchings put
$M_\ell(t)=\sum_{N\ge0}m_\ell(N)t^{2N}/(2N)!$.
It is an even Bessel determinant generating the $\ell$th
symmetric-power module, with rates $2(\ell-2i)$.
The same filter
$d^{-1}\sum_{j<d}\zeta^{-bj}M_\ell(\omega_jt)$ selects
$N\equiv b\pmod d$.  The simple-module argument gives the
same $2d\lceil\ell/2\rceil$ distinct nonzero rates.
Apply the budget with $p=1$, $h=2d$, $c=2b$, and $\alpha_1=0$;
Theorem~\ref{B:thm:D} supplies the upper bound.

For tableaux write $k=2r+\varepsilon$, put
$R=r+1-[4\mid k]$, and use the filter
\[
 \frac1d\sum_{j=0}^{d-1}\zeta^{-bj}Y_k(\zeta^jt)
 =\sum_{n\ge0}y_k(dn+b)\frac{t^{dn+b}}{(dn+b)!}.
\]
Its components are simple, being exponential twists of
$\operatorname{Sym}^r$ of a scaled Bessel module.  Their rate
sets are
\[
 \zeta^j\{\varepsilon+2r-4i:0\le i\le r\}.
\]
If $\varepsilon=1$, the mean of the unscaled set is $1$, so
different scales give different sets.  No two unscaled rates
are negatives of one another, since all are congruent to the
same odd number modulo $4$.  All $d(r+1)$ rates in the filter
are therefore nonzero and distinct.  If $\varepsilon=0$ and
$d$ is odd, no two $d$th roots of unity are antipodal;
the same argument gives $dR$ distinct nonzero rates.
The budget with $p=1$, $h=d$, $c=b$ yields the lower bound $R$,
and Theorem~\ref{B:thm:C} gives the upper bound in these cases.

It remains to consider even $k=2r$ and even $d$.  Write
$p\equiv b\pmod2$ and put
$F_p(t)=(Y_{2r}(t)+(-1)^pY_{2r}(-t))/2$.
This is a homogeneous Laurent-polynomial Bessel form of degree $r$;
it is nonzero, since its coefficients count tableaux of sizes
$2n+p$ and are positive.  By Lemma~\ref{B:up:laurent-parity},
$y_{2r}(2n+p)$ has width at most $K=\lceil r/2\rceil$.
Our sampled sequence is a further arithmetic subsequence of it,
so it also has width at most $K$.

In the $d$-term filter above, pair indices $j$ and $j+d/2$.
Their sum is a nonzero scalar multiple of $F_p(\zeta^jt)$.
For $0\le j<d/2$ these functions generate pairwise nonisomorphic
simple modules, with rate sets
$\zeta^j\{2r-4i:0\le i\le r\}$.  The scales are distinct modulo
sign, and each module has $2K$ distinct nonzero rates.
Consequently the entire filter has $dK$ such rates.
Lemma~\ref{B:up:sample-budget}, still with $p=1$, $h=d$, $c=b$,
gives the matching lower bound $K$.  This proves the last line
of the classification.

Finally multiplication by $c^nh(n)$ is an invertible change of
unknown in recurrences over $\C(n)$; it changes their
coefficients but not their extreme shifts.
\end{proof}

\begin{corollary}[Exact width for every homogeneous parity observable]
\label{B:up:general-parity-width}
Let $r\ge1$, $p\in\{0,1\}$, and
$\Phi\in\C[t,t^{-1}][X_0,X_1]$ be homogeneous of degree $r$.
Suppose
\[
 F(t)=\Phi\bigl(t,I_0(2t),\partial_t I_0(2t)\bigr)
\]
is nonzero, holomorphic at zero, and satisfies
$F(-t)=(-1)^pF(t)$.  Define
\[
 a_n=(2n+p)!\,[t^{2n+p}]F(t).
\]
For every $d\ge1$ and $0\le b<d$, the sequence $(a_{dn+b})$
is not eventually zero and has exact eventual recurrence width
$\lceil r/2\rceil$ over $\C(n)$.
\end{corollary}
\begin{proof}
Lemma~\ref{B:up:laurent-parity} supplies the upper bound, including
all further sampling.  For the lower bound let
$\zeta=e^{2\pi i/d}$ and $\omega_j=e^{\pi i j/d}$ for $0\le j<d$.
Then
\[
 H(t)=\frac1d\sum_{j=0}^{d-1}\zeta^{-bj}\omega_j^{-p}F(\omega_jt)
 =\sum_{n\ge0}a_{dn+b}\frac{t^{2(dn+b)+p}}{(2(dn+b)+p)!}.
\]
The nonzero homogeneous form $F$ generates the entire simple
$r$th symmetric-power Bessel module: it is a nonzero element of
that module, and the $\mathrm{SL}_2$ symmetric-power representation
is irreducible.  After scaling by $\omega_j$, the formal
exponential rates are $2\omega_j(r-2i)$, $0\le i\le r$.
The scales are distinct modulo sign, so the corresponding simple
modules are pairwise nonisomorphic.  Lemma~\ref{B:up:simple-filter}
therefore makes the module generated by $H$ their direct sum.
In particular $H$ is not a polynomial: a nonzero polynomial
spans a trivial rank-one module over $\C(t)$, whereas this direct
sum consists of nontrivial simple summands of rank $r+1\ge2$.
Thus $(a_{dn+b})$ is not eventually zero.

The direct sum has $2d\lceil r/2\rceil$ distinct nonzero formal
exponential rates.  Apply Lemma~\ref{B:up:sample-budget} with one
factorial, sampling step $h=2d$, offset $c=2b+p$, and
$\alpha_1=0$.  It gives the lower bound
$\lceil r/2\rceil$, completing the proof.
\end{proof}

\begin{example}
The sequences $u_{5,7}(100n+37)$ and $m_5(100n+37)$ both have
minimum recurrence width $3$.  The sequences $y_7(dn+b)$ have
minimum width $4$ for every $d,b$, and $y_8(dn+b)$ have minimum
width $4$ for odd $d$ and width $2$ for even $d$.  These assertions concern exact
eventual recurrence width, not merely the order of a particular
elimination output.  They make no claim about the smallest
coefficient degree after sampling.
\end{example}

\section{A local Fourier proof and sharpening of the Laplace--Jordan principle}
\label{B:up:ljr:sec}

The following local Fourier argument controls the origin monodromy of
iterated factorial Borel transforms.  It uses formal stationary phase and
the nearby-cycle/vanishing-cycle description at one point.  Only the unit
eigenvalue can acquire an additional Jordan step.

\begin{lemma}[Nearby and vanishing monodromy]\label{B:up:ljr:quiver}
Let $V,W$ be finite-dimensional complex vector spaces, with maps
$u:V\to W$ and $v:W\to V$, and suppose
\[
                 T=I_V+vu,\qquad S=I_W+uv
\]
are invertible. For every polynomial $p\in\mathbb C[z]$,
\begin{equation}\label{B:up:ljr:identity}
                    (T-I_V)p(T)=v\,p(S)u.
\end{equation}
In particular, $p(S)=0$ implies $(T-I_V)p(T)=0$.
For each $\lambda\ne1$, the restrictions of $T$ and $S$ to their
$\lambda$-generalized eigenspaces are isomorphic.  Only the
$1$-eigenspace can acquire a further Jordan step in passing from $S$ to $T$.
\end{lemma}
\begin{proof}
The relation $uT=Su$ implies $up(T)=p(S)u$. Multiplication by
$v$ proves \eqref{B:up:ljr:identity}.  On a $\lambda$-generalized
eigenspace with $\lambda\ne1$, $T-I_V$ and $S-I_W$ are invertible.
Thus $u$ and $v$ give mutually inverse isomorphisms after composing
one of them with the appropriate inverse of $T-I_V$ or $S-I_W$.
The identity then proves the assertion at $\lambda=1$ as well.
\end{proof}

We use local monodromy at infinity in the positive local coordinate
$1/s$. With this convention, regular local Fourier transform preserves
its monodromy conjugacy class. A change of Fourier sign is a homothety
of the target and does not alter that class.

\begin{proposition}[The Fourier monodromy bound]\label{B:up:ljr:fourier}
Let $\mathcal M$ be a holonomic module on the complex affine $s$-line,
whose generic connection is regular singular at infinity. Let $M$ be
the monodromy of that regular germ. Write $T$ for the monodromy at
$x=0$ of the generic connection underlying its Fourier transform.
Then the transformed connection has no singularity at a finite
$x\ne0$, is regular singular or ordinary at $0$, and
\[
                 p(M)=0\quad\Longrightarrow\quad(T-I)p(T)=0.
\]
The same annihilating polynomial applies to the monodromy of every
differential subquotient of the transformed generic connection.
No middle-extension hypothesis on $\mathcal M$ is required.
\end{proposition}
\begin{proof}
Formal stationary phase gives
\[
 \Phi_x(\mathcal F\mathcal M)
  \simeq\mathcal F^{(\infty,x)}
       \bigl((\Psi_\infty\mathcal M)^{\le1,(x)}\bigr)
\]
for every finite $x$; this is \cite[Corollary~2.3(1)]{LJR_Arinkin}.
Because $\Psi_\infty\mathcal M$ is regular, its exponential-type
summand is zero unless $x=0$.  At $x=0$ its local Fourier transform
is again regular and has monodromy $M$:
\cite[Remark~5.2 and Remark~5.3(4)]{LJR_Sabbah}.
These statements concern the extended formal local Fourier functor,
which includes trivial regular summands. For example, a trivial
rank-one infinity germ contributes the trivial rank-one vanishing
cycles of a delta module; it is not discarded.

For a holonomic module on a formal disc, its irregular part is the
same in nearby and vanishing cycles. Thus the regularity of the
vanishing cycles proves regularity of the nearby germ, and their
vanishing at $x\ne0$ proves smoothness there.  On the regular part,
nearby and vanishing cycles have the description
$(V,W,u,v)$ with monodromies $I+vu$ and $I+uv$, respectively;
see \cite[\S3.3, formula~(3.3)]{LJR_Arinkin}.  The vanishing
monodromy $S$ is conjugate to $M$. Lemma~\ref{B:up:ljr:quiver}
therefore proves the asserted identity and bound. Restricting an
annihilating polynomial to an invariant subspace or its quotient
preserves the identity. The argument also works for the dual local
system of solutions, since local Fourier transform is compatible
with duality and both source and target monodromies are dualized.
\end{proof}

\begin{theorem}[Eigenvalue-resolved iterated Borel bound]\label{B:up:borel-bound}
Let $A_0(t)=\sum_{n\ge0}a_nt^n$ be a nonzero D-finite germ with
positive radius of convergence. Suppose its minimal equation is
regular singular or ordinary at $0$, and denote its monodromy by
$M$. For $r\ge0$, put
\[
                  A_r(x)=\sum_{n\ge0}\frac{a_nx^n}{(n!)^r}.
\]
Each minimal equation is regular singular or ordinary at $0$.
If $p(M)=0$, then its monodromy $T_r$ satisfies
\begin{equation}\label{B:up:ljr:iterated}
                      (T_r-I)^r p(T_r)=0.
\end{equation}
In particular, nonunit monodromy eigenvalues acquire no larger
Jordan blocks under any number of Borel transforms. The largest
Jordan-block size at eigenvalue $1$ increases by at most $r$.
If $A_0$ is algebraic, every nonunit block of $T_r$ is semisimple,
and every block at eigenvalue $1$ has size at most $r+1$.
\end{theorem}
\begin{proof}
It suffices to prove the assertion for one transform and iterate.
Put $f(s)=s^{-1}A_0(s^{-1})$ and
$\mathcal M=\mathbb C[s]\langle\partial_s\rangle\cdot f$.
Its generic connection is the rational gauge and variable transform
of the minimal equation of $A_0$.  It is regular at infinity, and
its local monodromy in the coordinate $1/s$ is $M$.

Here is the exact relation with the series transform. On
$\mathbb C[s]\oplus s^{-1}\mathbb C[[s^{-1}]]$, define
\[
 \mathcal B(s^j)=0\quad(j\ge0),\qquad
 \mathcal B(s^{-n-1})=x^n/n!\quad(n\ge0).
\]
Direct checking on monomials gives
$\mathcal B(sf)=\partial_x\mathcal B(f)$ and
$\mathcal B(\partial_sf)=-x\mathcal B(f)$.
Consequently, with Fourier convention
$s\mapsto\partial_x$, $\partial_s\mapsto-x$, there is a
surjection of holonomic modules
\[
                   \mathcal F(\mathcal M)
                       \longrightarrow
                   \mathbb C[x]\langle\partial_x\rangle\cdot A_1.
\]
The solution local system of the minimal equation of $A_1$ is
therefore a subquotient of the Fourier solution local system.
Proposition~\ref{B:up:ljr:fourier} gives $(T_1-I)p(T_1)=0$ and
regularity. Applying the same argument successively with the
polynomial $(z-1)^j p(z)$ proves \eqref{B:up:ljr:iterated}.

For algebraic $A_0$, the local monodromy of its minimal equation
has finite order: its solutions are spanned by finitely many
algebraic conjugates. Its minimal polynomial is a product of
distinct roots-of-unity factors. Multiplication by $(z-1)^r$
therefore gives exactly the stated eigenvalue-specific bounds.
\end{proof}

For an algebraic source germ, the theorem gives a unit-eigenvalue block
bound of two after one transform and three after two transforms.
Multiplication by an integer power of the variable is an invertible
rational gauge, so it preserves both the regularity hypothesis and
these monodromy bounds.

\section{Twisted Catalan identities and ordinary algebraicity}
\label{B:up:twist-algebraicity-section}

\begin{proposition}[Two-row twists as Catalan tails]
\label{B:up:twist-catalan}
For every $q,n\ge0$,
\[
 u_{1,q}(n)=1,\qquad
 u_{2,q}(n)=C_{n+q},\qquad
 C_j=\frac1{j+1}\binom{2j}{j}.
\]
Consequently, with $C(z)=(1-\sqrt{1-4z})/(2z)$,
\[
 \sum_{n\ge0}u_{2,q}(n)z^n
 =z^{-q}\left(C(z)-\sum_{j=0}^{q-1}C_jz^j\right),
\]
where the sum is empty for $q=0$ and the singularity at $z=0$
is removable.
\end{proposition}
\begin{proof}
There is exactly one standard tableau of every one-row shape,
which proves the first identity.  At $\ell=2$, the extremal
coefficient formula in Theorem~\ref{B:thm:twisted} gives
\[
 (n+q+1)u_{2,q}(n)
 =2(2n+2q-1)u_{2,q}(n-1),\qquad n\ge1.
\]
The initial value is $u_{2,q}(0)=f^{(q,q)}=C_q$ by the
two-row hook-length formula.  The sequence $C_{n+q}$ has the
same recurrence and initial value, proving the claim.
Subtracting the first $q$ coefficients of $C(z)$ gives the
ordinary generating function.
\end{proof}


\begin{theorem}[Twisted determinants under arithmetic sampling]
\label{B:up:twist-all-algebraicity}
For integers $\ell,d\ge1$, $q\ge0$, and $0\le b<d$, the
ordinary generating function
\[
 A_{\ell,q;d,b}(z)=\sum_{n\ge0}u_{\ell,q}(dn+b)z^n
\]
is algebraic over $\C(z)$ if and only if $\ell\le2$.
\end{theorem}
\begin{proof}
For $\ell=1$ the series is $(1-z)^{-1}$.  For $\ell=2$,
Proposition~\ref{B:up:twist-catalan} gives an algebraic full
series.  Every arithmetic section of an algebraic series is
algebraic: the root-of-unity filter expresses
$z^bA_{\ell,q;d,b}(z^d)$ as a finite sum of algebraic
functions, and $z$ is algebraic over $z^d$.

Suppose now $\ell\ge3$ and $A_{\ell,q;d,b}$ were algebraic.
Then
\[
 F(z)=z^bA_{\ell,q;d,b}(z^d)
     =\sum_{n\ge0}u_{\ell,q}(dn+b)z^{dn+b}
\]
is algebraic.  Write $\mathcal B(\sum a_nz^n)
=\sum a_nz^n/n!$ and $Q=q\ell$.  An exact coefficient
calculation gives
\[
 \mathcal B\bigl(z^Q\mathcal B F\bigr)
 =z^Q H(z),\qquad
 H(z)=\sum_{n\ge0}
 \frac{u_{\ell,q}(dn+b)z^{dn+b}}
 {(dn+b)!(dn+b+Q)!}.
\]
The regular-singular Borel block bound
(Theorem~\ref{B:up:borel-bound}) applies twice.  Algebraic
functions have semisimple local monodromy, so the largest
unipotent Jordan block at zero is initially of size at most
one.  Each Borel transform increases the permitted size by
at most one; multiplication by the integer power $z^Q$ is
a rational gauge and changes no local monodromy.  Thus the
minimal equation of $H$ has no Jordan block at zero larger
than three: its nonunit blocks remain semisimple, and its
unit-eigenvalue blocks have size at most three.  The hypotheses needed at the intermediate step
hold: an algebraic module is regular singular everywhere;
its Borel transform is regular singular at zero by the cited
bound; rational gauges preserve this property.

On the other hand,
\[
 H(t^2)=\frac1d\sum_{j=0}^{d-1}\zeta^{-bj}
       G_{\ell,q}(\zeta^jt^2)
\]
is exactly the direct sum of $d$ pairwise nonisomorphic simple
symmetric-Bessel modules used in the proof of
Theorem~\ref{B:up:sample-order}.  A Bessel module at zero has
one nontrivial unipotent block of size two.  Its $\ell$th
symmetric power has one unipotent block of size $\ell+1$:
if the local basis is $v,w=v\log t+\text{meromorphic}$,
then $w^\ell$ has nonzero coefficient $v^\ell$ at
$(\log t)^\ell$.  The rational gauge defining
$G_{\ell,q}$ does not change this block.
Consequently $H(t^2)$ has a block of size $\ell+1\ge4$.
The finite pullback $z=t^2$ squares monodromy.  For every
nonzero eigenvalue, squaring preserves each Jordan-block size,
since the derivative of $z^2$ at that eigenvalue is nonzero.
It therefore cannot create a block larger than three, which is
the required contradiction.  Hence $A_{\ell,q;d,b}$ is transcendental.
\end{proof}

\begin{remark}[Scope of the two assertions]
The recurrence theorem is independent of the Borel block bound:
its proof uses only the Bessel modules and the explicit
factorial budget.  The algebraicity theorem additionally uses
the local Fourier--Laplace theorem stated in
Theorem~\ref{B:up:borel-bound}.  Neither assertion is an
irrationality or transcendence assertion about evaluations
of an ordinary generating function at individual algebraic
points.
\end{remark}


\begin{corollary}[The remaining tableau algebraicity cases]
\label{B:up:three-algebraicity}
The ordinary generating function $\sum_ny_k(n)z^n$ is algebraic
over $\C(z)$ exactly when $k\le3$. The ordinary generating
function $\sum_nm_\ell(n)z^n$ is algebraic exactly when $\ell=1$.
These are the unfiltered classifications developed in the entropy--transport development; their proof is reproduced here to make
the present paper independent of that manuscript.
\end{corollary}
\begin{proof}
The scalar tableau equation has the solution representation
$e^{\varepsilon t}\operatorname{Sym}^r V$, with
$r=\lfloor k/2\rfloor$ and $V$ the rank-two Bessel solution
space. Thus its origin monodromy has the single block
$J_{r+1}(1)$. Algebraicity of the ordinary series would give
blocks of size at most two after one Borel transform, by
Theorem~\ref{B:up:borel-bound}. This excludes every $k\ge4$.
For the smaller heights the standard determinant formulas give
\[
 \sum_ny_1(n)z^n=\frac1{1-z},\qquad
 \sum_ny_2(n)z^n=
 \frac{(1+2z)/\sqrt{1-4z^2}-1}{2z},
\]
and
\[
 \sum_ny_3(n)z^n=
       \frac{1-z-\sqrt{1-2z-3z^2}}{2z^2}.
\]
The apparent singularities at zero are removable. For example,
the even and odd coefficients of $Y_2$ are
$\binom{2n}{n}$ and $\binom{2n+1}{n}$, and $Y_3$ is the
exponential generating function of the Motzkin numbers, as follows
by coefficient extraction in \eqref{B:eq:Gordon}.

For matchings, algebraicity of $\sum_nm_\ell(n)z^n$ implies
algebraicity after $z=t^2$. Its Borel transform is
$e^{-t}Y_{2\ell+1}(t)$, whose origin block has size
$\ell+1$. The bound two therefore excludes $\ell\ge2$.
For $\ell=1$, noncrossing matchings are counted by the Catalan
numbers, whose generating function is algebraic.
\end{proof}

\section{Bessel determinants and the three algebraicity classifications}
\label{E:sfb:section}

This section supplies the combinatorial definitions, Bessel reductions,
and scalar ordinary-point result used in the differential and arithmetic
theorems of Section~\ref{E:acb:section}.  The Laplace input is
Theorem~\ref{B:up:borel-bound}.  In particular, no unpublished tableau
operator theorem is needed for the arithmetic specialization.


\subsection{A coordinate and indexing dictionary}
The combinatorial sequences are those defined in
Section~\ref{B:families:main}. In this section the crossing parameter is
written $k$, so
\[
 f_k(n)=m_{k-1}(n)\quad(k\ge2),\qquad
 f_1(0)=1,\quad f_1(n)=0\quad(n\ge1).
\]
Define
\[
 \widetilde Y_k(z)=\sum_{n\ge0}y_k(n)z^n,\quad
 \widetilde U_k(z)=\sum_{n\ge0}u_k(n)z^n,\quad
 \widetilde F_k(z)=\sum_{n\ge0}f_k(n)z^n.
\]
Their normalized versions satisfy
\[
 Y_k(x)=\sum_{n\ge0}y_k(n)x^n/n!,\qquad
 U_k(s)=\sum_{n\ge0}u_k(n)s^n/(n!)^2=G_k(s),
\]
\[
 F_k(x)=\sum_{n\ge0}f_k(n)x^{2n}/(2n)!
             =e^{-x}Y_{2k-1}(x).
\]
The determinant identities are \eqref{B:eq:Gessel},
\eqref{B:eq:Gordon}, and \eqref{B:eq:Sp}, with the indicated
index shift for matchings. In particular, $U_k(s)=D_k(\sqrt{s})$.
These are the Gordon--Gessel and Weyl-chamber enumerations; see also
\cite{SFB_GrabinerMagyar,SFB_JinReidysWang}.

\subsection{Reduction to two Bessel coordinates}

Write
\[
 y(x)=I_0(2x)=\sum_{m\ge0}\frac{x^{2m}}{(m!)^2},\qquad
 y_s(s)=I_0(2\sqrt{s})=\sum_{m\ge0}\frac{s^m}{(m!)^2}.
\]
Their scalar equations and logarithmic companion solutions are
\begin{align}
 y''&=4y-x^{-1}y',&
 \kappa(x):=K_0(2x)&=-(\log x)y(x)+h(x),\label{E:sfb:besselx}\\
 y_s''&=s^{-1}(y_s-y_s'),&
 \kappa_s(s):=K_0(2\sqrt{s})
       &=-\tfrac12(\log s)y_s(s)+h_s(s),\label{E:sfb:bessels}
\end{align}
where $h,h_s$ are holomorphic at zero.  The integer-order Bessel
recurrences and these local expansions may be found in
\cite[\S\S10.29, 10.31]{DLMF}.

\begin{lemma}[Bessel group and independence]\label{E:sfb:independence}
The two Bessel equations have differential Galois group
$\mathrm{SL}_2$ over their rational function fields.  The
pairs $(y,y')$ and $(y_s,y_s')$ are algebraically independent
over $\mathbb C(x)$ and $\mathbb C(s)$, respectively.
\end{lemma}
\begin{proof}
The rank-two case of Theorem~\ref{B:bv:galois}, together with
the pullback $s=x^2$, gives both Galois groups. The orbit argument
there proves algebraic independence of the two distinguished coordinates.
\end{proof}

\begin{lemma}[Determinants are homogeneous forms]\label{E:sfb:forms}
Each determinant in \eqref{B:eq:Gessel}, \eqref{B:eq:Gordon}, and \eqref{B:eq:Sp}
is a nonzero homogeneous polynomial in the corresponding Bessel
pair, with rational Laurent-polynomial coefficients.  The degrees
are $\lfloor k/2\rfloor$ for $e^{-\varepsilon x}Y_k(x)$,
where $k=2\lfloor k/2\rfloor+\varepsilon$; $k$ for $U_k(s)$;
and $k-1$ for $F_k(x)$.
\end{lemma}
\begin{proof}
The recurrences
\[
 b_0=y,\qquad b_1=y'/2,\qquad
 b_{j+1}=b_{j-1}-(j/x)b_j
\]
express every $b_j$ as a rational Laurent-polynomial linear
form in $y,y'$.  Thus a determinant of size $d$ is homogeneous
of degree $d$.

For the $s$-variable put
$J_j(s)=s^{-j/2}I_j(2\sqrt{s})$ for $j\ge0$.
Then $J_0=y_s$, $J_1=y_s'$ and
$sJ_{j+1}=J_{j-1}-jJ_j$.
Multiplying row $i$ of \eqref{B:eq:Gessel} by $s^{-i/2}$
and column $j$ by $s^{j/2}$ does not change its determinant,
and changes its entry to
$s^{\max(0,j-i)}J_{|i-j|}(s)$.
All coefficients are therefore in $\mathbb Q[s,s^{-1}]$.
Each resulting form is nonzero because the normalized counting
series has constant coefficient $1$.
\end{proof}

\begin{lemma}[The logarithmic degree is exact]\label{E:sfb:logblock}
Let $v$ be either analytic Bessel solution above, in variable $z$,
and let $\Phi\in\mathbb C(z)[X_0,X_1]$ be nonzero homogeneous
of degree $d\ge1$.  The function $\Phi(v,v')$ has minimal
differential order $d+1$, and its minimal equation at zero has
one unipotent Jordan block of size $d+1$.
Multiplication by $e^{\varepsilon z}$ does not change these
two assertions.
\end{lemma}
\begin{proof}
This is the order and local-monodromy part of
Theorem~\ref{E:acb:module} below. Its proof uses
Lemma~\ref{E:sfb:independence}, independently of the present lemma.
\end{proof}

\subsection{Ordinary generating functions}

\begin{theorem}[Algebraicity of the three combinatorial families]
\label{E:sfb:algebraicity}
For positive integer parameters:
\begin{enumerate}
\item $\widetilde Y_k$ is algebraic if and only if $k\le3$;
\item the ordinary generating function for permutations avoiding
      an increasing subsequence of length $\ell$ is algebraic
      if and only if $\ell\le3$;
\item $\widetilde F_k$ is algebraic if and only if $k\le2$.
\end{enumerate}
\end{theorem}
\begin{proof}
For tableaux and matchings use
Corollary~\ref{B:up:three-algebraicity} and $f_k=m_{k-1}$;
the case $f_1$ has generating function $1$.
For permutations use Theorem~\ref{B:up:twist-all-algebraicity}
at twist zero and sampling interval one, with $k=\ell-1$.
Avoidance of length one permits only the empty permutation.
\end{proof}

\subsection{A published Jacobi system and a scalar tableau equation}

The arithmetic specialization below requires every nonzero
point to be ordinary for the minimal tableau scalar equation.
We retain a general-parameter Jacobi elimination and identify its specialization with the scalar equation already proved.

For $N\ge1$, $\alpha>-1$, $b>0$, $\beta>0$, and $0\le p\le N$, let
\[
 \begin{aligned}
 H_p(z)=\binom Np^{-1}\int_{[0,1]^N}
 &e^{-z\sum u_i}\prod_i u_i^\alpha(1-u_i)^{b-1}
 \prod_{i<j}|u_i-u_j|^\beta\\[-2pt]
 &\hspace{20mm}\cdot e_p(1-u_1,\ldots,1-u_N)\,du,
 \end{aligned}
\]
where $e_p$ is the elementary symmetric polynomial and $e_0=1$.
The Forrester--Kumar system
\cite[Corollary 2.5, equations (2.19)--(2.20)]{FK22}
states that
\begin{equation}\label{E:sfb:FKsystem}
 (N-p)zH_{p+1}=((N-p)z+B_p)H_p+zH_p'-D_pH_{p-1},
\end{equation}
where
\[
 B_p=p\left(\alpha+b+1+\frac\beta2(2N-p-1)\right),\qquad
 D_p=p\left(b+\frac\beta2(N-p)\right).
\]
Terms outside the range $0\le p\le N$ are omitted.
This is the only published differential-system input to the
construction; the elimination and its specialization are given
explicitly next.

\begin{lemma}[Scalar elimination]\label{E:sfb:elimination}
There is a scalar differential operator
$\mathcal A=\sum_{m=0}^{N+1}a_m(z)\partial_z^m$ annihilating $H_0$ such that
\[
 a_m\in\mathbb R[z],\qquad a_{N+1}=z^N,\qquad
 \deg a_m\le N,\qquad z^{m-1}\mid a_m\quad(m\ge1).
\]
The coefficients are rational when the parameters are rational.
\end{lemma}
\begin{proof}
Write $D=\partial_z$ and construct rational operators
\[
 P_0=1,\qquad P_1=1+D/N,
\]
and, for $1\le p<N$,
\begin{equation}\label{E:sfb:Precursion}
 (N-p)P_{p+1}=
 \left(D+N-p+\frac{B_p}{z}\right)P_p
       -\frac{D_p}{z}P_{p-1}.
\end{equation}
Equation~\eqref{E:sfb:FKsystem} gives $H_p=P_pH_0$.
Write $P_p=\sum_{m=0}^pq_{p,m}D^m$ for $p\ge1$.
Induction using $D\circ(qD^m)=qD^{m+1}+q'D^m$ gives
\[
 q_{p,p}=\frac{(N-p)!}{N!},\qquad
 q_{p,m}\in\operatorname{span}\{z^{-(p-1)},\ldots,1\},
 \qquad \operatorname{ord}_0q_{p,m}\ge m-p.
\]
For the induction, differentiation or division by $z$ lowers
the least Laurent exponent by at most one, and no positive
exponent is introduced.  The highest derivative comes only
from $D\circ P_p$.  The coefficient of $D^m$ in that product
is $q_{p,m}'+q_{p,m-1}$, whose order is at least $m-p-1$;
the other terms in \eqref{E:sfb:Precursion} have the same bound.
The case $p=1$ involving $P_0$ is immediate.

The $p=N$ equation in \eqref{E:sfb:FKsystem} now shows that
\begin{equation}\label{E:sfb:Aoperator}
 \mathcal A=N!z^{N-1}
       \bigl(zD\circ P_N+B_NP_N-D_NP_{N-1}\bigr)
\end{equation}
annihilates $H_0$.  The bracket has Laurent exponents between
$-(N-1)$ and $1$, so multiplication by $z^{N-1}$ makes its
coefficients polynomials of degree at most $N$.
Its top coefficient is $z/N!$, giving $a_{N+1}=z^N$.
The coefficient of $D^m$ in the bracket has order at least
$m-N$, proving the stated divisibility after multiplication.
The formulas also prove rationality.
\end{proof}

\begin{proposition}[No nonzero tableau scalar singularities]
\label{E:sfb:tableauordinary}
Let $k=2r+\varepsilon\ge2$, with $\varepsilon\in\{0,1\}$.
There is an order-$r+1$ operator over $\mathbb Q[x]$
annihilating $Y_k$ whose leading coefficient is $x^r$.
The minimal monic operator has order $r+1$ and is ordinary
at every $x\ne0$.
\end{proposition}
\begin{proof}
The ordinary-point and order assertions are
Theorem~\ref{B:thm:jacobi}, proved in
Appendix~\ref{B:app:jacobi}. The relation with the general
Jacobi elimination above is explicit: at
$N=r$, $\alpha=\tfrac12$, $b=\varepsilon+\tfrac12$, $\beta=2$,
\begin{equation}\label{E:sfb:JacobiY}
 Y_k(x)=\frac{4^{r(r+\varepsilon)}}{r!(2\pi)^r}
          e^{(\varepsilon-2r)x}H_0(-4x).
\end{equation}
Substituting $z=-4x$ and
$\partial_z=-\tfrac14(\partial_x-(\varepsilon-2r))$
in \eqref{E:sfb:Aoperator}, and multiplying by $-4$,
gives leading coefficient $x^r$. This retains the normalization
while avoiding a second proof of the same tableau integral.
\end{proof}

Only the ordinary-point conclusion of
Proposition~\ref{E:sfb:tableauordinary} is needed in
Theorem~\ref{E:acb:values}.  Possible cancellation of a common
power of $x$ in a primitive polynomial normalization is irrelevant
to that conclusion.  Both the general elimination and its specialized minimality proof are included here.


\clearpage

\section{Exact Bessel modules, algebraic base changes, and tableau values}
\label{E:acb:section}

The algebraicity classifications of Theorem~\ref{E:sfb:algebraicity}
admit stronger differential and arithmetic conclusions.  The general
homogeneous-form theorem below determines the normalized equations,
their complete local partition, and all algebraic relations between
their initial jets.  The ordinary-point result needed for arithmetic
specialization was proved explicitly in
Proposition~\ref{E:sfb:tableauordinary}.

\subsection{A homogeneous-form theorem}

Let $K=\mathbb C(z)$ and let $y$ denote either
\[
 y(z)=I_0(2z),\qquad y''=4y-z^{-1}y',
 \quad\hbox{or}\quad
 y(z)=I_0(2\sqrt z),\qquad y''=z^{-1}(y-y').
\]
Both rank-two connections have differential Galois group $\mathrm{SL}_2$.
Their local monodromy at zero is a nontrivial unipotent block of size two.
These Bessel inputs, including their primary references, are
established in Lemma~\ref{E:sfb:independence}.
For $d\ge1$, let $\Phi\in K[X_0,X_1]$ be a nonzero homogeneous
form of degree $d$, and put
\[
 F(z)=e^{\varepsilon z}\Phi(y(z),y'(z)),\qquad
 \varepsilon\in\{0,1\}.
\]

\begin{theorem}[Exact module and stability]\label{E:acb:module}
The minimal differential operator of $F$ over $K$ is irreducible of
order $d+1$.  Its differential Galois group, in its solution
representation, is
\[
 \begin{cases}
 \operatorname{Sym}^{d}(\mathrm{SL}_2),&\varepsilon=0,\\
 \{c\operatorname{Sym}^{d}(g):c\in\mathbb C^*,\ g\in\mathrm{SL}_2\},
       &\varepsilon=1.
 \end{cases}
\]
In the first case this is isomorphic to $\mathrm{SL}_2$ for odd $d$
and to $\mathrm{PSL}_2$ for even $d$.
At zero the complete local monodromy partition consists of one
unipotent block $J_{d+1}(1)$.

All these orders and irreducibility assertions remain valid over
every finite algebraic extension of $K$.  After a ramified algebraic
change of variable, the local block at a point above zero still has
size $d+1$.
\end{theorem}

\begin{proof}
Let $V$ be the two-dimensional Bessel solution space in a
Picard--Vessiot extension.  Polarizing $u\mapsto\Phi(u,u')$ gives
a $\mathrm{SL}_2$-equivariant map from $\operatorname{Sym}^dV$
to the differential field.  It is nonzero because $y,y'$ are
algebraically independent over $K$, by Lemma~\ref{E:sfb:independence}.
The representation $\operatorname{Sym}^dV$ is irreducible,
so this map is injective.  The Galois orbit of any nonzero image
spans the entire image.  The standard Galois-orbit characterization
of a minimal scalar operator therefore gives the claimed order,
irreducibility, and solution representation.

For $\varepsilon=1$, the compositum of the Bessel
Picard--Vessiot extension and that of $u'=u$ has group
$\mathrm{SL}_2\times\mathbb G_m$.  Indeed, a proper subdirect
product would give a nontrivial common quotient of
$\mathrm{SL}_2$ and $\mathbb G_m$, which is impossible.
The preceding argument applies to the irreducible representation
$\operatorname{Sym}^dV$ tensored with the standard character of
$\mathbb G_m$.

The symmetric power of a nontrivial $J_2(1)$ is a single
$J_{d+1}(1)$: in the usual degree-$d$ polynomial representation,
its logarithm acts by a nonzero multiple of $X_0\partial_{X_1}$,
with one chain of length $d+1$.  The exponential twist is
single-valued and nonvanishing at zero and does not alter this
partition.

Both displayed Galois groups are connected.  Finite algebraic
extension of the base leaves the identity component unchanged,
and hence leaves the full connected group and its irreducible
representation unchanged.  A ramification of index $e$ replaces
the local monodromy by its $e$th power; a unipotent Jordan block
in characteristic zero retains its size under this operation.
\end{proof}

\begin{corollary}[The three determinant families]\label{E:acb:families}
With the notation of Section~\ref{E:sfb:section}, the following orders are exact:
\[
\begin{array}{c|c|c|c}
 \text{function}&d&\text{minimal order}&\text{monodromy at zero}\\ \hline
 Y_k(x),\ k\ge2&\lfloor k/2\rfloor&\lfloor k/2\rfloor+1
                                  &J_{\lfloor k/2\rfloor+1}(1)\\
 U_k(s),\ k\ge1&k&k+1&J_{k+1}(1)\\
 F_k(x),\ k\ge2&k-1&k&J_k(1).
\end{array}
\]
The permutation and matching groups are respectively
$\operatorname{Sym}^{k}(\mathrm{SL}_2)$ and
$\operatorname{Sym}^{k-1}(\mathrm{SL}_2)$.
For tableaux, write $k=2d+\varepsilon$; the group is the
corresponding group in Theorem~\ref{E:acb:module}.
Every listed minimal order persists after finite algebraic
extension of the coefficient field.
\end{corollary}

\begin{proof}
The determinant reductions in Lemma~\ref{E:sfb:forms} give nonzero homogeneous
forms of exactly the displayed degrees.  Their nonvanishing
also follows from the constant term $1$ of each generating
function.  Apply Theorem~\ref{E:acb:module}.
\end{proof}

\subsection{Algebraic relations among initial jets}

Write $\mathbf J_F=(F,F',\ldots,F^{(d)})^{\mathsf T}$ and
\[
 \mathbf m_d=(y^d,y^{d-1}y',\ldots,(y')^d)^{\mathsf T}.
\]
Differentiating and reducing $y''$ by the Bessel equation produces
a rational matrix $P(z)$ with
\begin{equation}\label{E:acb:jetgauge}
 \mathbf J_F=e^{\varepsilon z}P(z)\mathbf m_d.
\end{equation}

\begin{theorem}[Jet cone]\label{E:acb:jets}
The matrix $P$ belongs to $\mathrm{GL}_{d+1}(K)$.
After the invertible rational change of jet coordinates
$\mathbf Z=P^{-1}\mathbf J_F$, the ideal of all polynomial
relations over $K$ is generated by
\begin{equation}\label{E:acb:minors}
 Z_iZ_{j+1}-Z_{i+1}Z_j,
             \qquad 0\le i<j\le d-1.
\end{equation}
For $d=1$ the ideal is zero.  Consequently
\[
 \operatorname{trdeg}_{K}K(F,F',\ldots,F^{(d)})=2.
\]
Every two distinct entries of this jet vector are algebraically
independent over $K$, and remain so over every finite algebraic
extension of $K$.
When $\varepsilon=0$, one has the further exact field identity
\[
 K(F,F',\ldots,F^{(d)})
   =K(y^d,y^{d-1}y',\ldots,(y')^d)
   =K(y,y')^{\mu_d},
\]
where $\mu_d$ acts by simultaneous scalar multiplication on
$y,y'$.  Thus its index in $K(y,y')$ is exactly $d$.
\end{theorem}

\begin{proof}
If the rows of $P$ were dependent, \eqref{E:acb:jetgauge} would give
a nonzero rational differential equation for $F$ of order at most
$d$, contradicting Theorem~\ref{E:acb:module}.  Therefore $P$ is
invertible.

The kernel of the substitution
$Z_j\mapsto X_0^{d-j}X_1^j$ is the homogeneous ideal
\eqref{E:acb:minors} of the affine cone over the degree-$d$ rational
normal curve.  For completeness, these minors identify all
products of $Z_j$ having the same number of factors and the same
sum of indices; their images are the same monomial in $X_0,X_1$.
Grouping monomials by those two integers proves that they generate
the entire kernel.  Algebraic independence of $y,y'$ shows that
this is exactly the relation ideal for $\varepsilon=0$.
For $\varepsilon=1$, the three functions $y,y',e^z$ are
algebraically independent: the group
$\mathrm{SL}_2\times\mathbb G_m$ acts transitively on the
nonzero Bessel solution vector and independently on the
exponential coordinate.  Multiplying every degree-$d$ monomial
by the same independent scalar preserves the same homogeneous
kernel.  This proves the ideal assertion and its dimension.

Two different rows of $P$ are linearly independent, so their
degree-$d$ forms $A,B$ are not proportional.  Two nonproportional
homogeneous forms of the same positive degree in two variables
are algebraically independent: their ratio is a nonconstant
function of $X_1/X_0$, while simultaneous scaling changes their
common magnitude.  Equivalently, their Jacobian is nonzero.
The same conclusion holds after multiplying both by the
independent scalar $e^z$.  This proves pairwise independence;
an algebraic extension of the coefficient field cannot decrease
transcendence degree.  Finally the displayed monomial field is
the fixed field of the faithful scalar action of $\mu_d$ on
the rational field $K(y,y')$, proving the index statement.
\end{proof}

\subsection{Arithmetic specialization without exceptional nonzero points}

For tableaux Proposition~\ref{E:sfb:tableauordinary} supplies an
additional essential input: an order-$d+1$ annihilator of $Y_k$
has leading coefficient $x^d$, where $d=\lfloor k/2\rfloor$,
and $d+1$ is the exact minimal order.  Thus every
$\alpha\ne0$ is an ordinary point of the minimal scalar equation.
This property is stronger than having a connection which is
merely gauge-equivalent to one regular at $\alpha$; it is what
rules out exceptional zeros of the scalar determinant form.

\begin{theorem}[Values at nonzero algebraic tableau arguments]
\label{E:acb:values}
Let $k\ge2$, $d=\lfloor k/2\rfloor$, and
$\alpha\in\overline{\mathbb Q}^{\,*}$.
Then
\[
 \operatorname{trdeg}_{\overline{\mathbb Q}}
 \overline{\mathbb Q}\bigl(Y_k(\alpha),Y_k'(\alpha),\ldots,
                                Y_k^{(d)}(\alpha)\bigr)=2.
\]
Every two distinct numbers in this displayed tuple are
algebraically independent over $\overline{\mathbb Q}$.
All their algebraic relations are obtained from
\eqref{E:acb:minors} by the invertible algebraic change of
coordinates $\mathbf Z=P(\alpha)^{-1}\mathbf J_{Y_k}(\alpha)$.
In particular $Y_k(\alpha)$ is transcendental and nonzero.
For $k=1$, $Y_1(\alpha)=e^\alpha$ is transcendental.
\end{theorem}

\begin{proof}
Put $y=I_0(2x)$.  The functions $y,y',e^x$ are $E$-functions
over $\mathbb Q$; their first-order system is
\[
 \frac{d}{dx}
 \begin{pmatrix}y\\y'\\e^x\end{pmatrix}
 =\begin{pmatrix}0&1&0\\4&-1/x&0\\0&0&1\end{pmatrix}
 \begin{pmatrix}y\\y'\\e^x\end{pmatrix}.
\]
They are algebraically independent over $\mathbb C(x)$ by
the group argument in Theorem~\ref{E:acb:jets}.
The Siegel--Shidlovskii theorem in
\cite[Theorem~1.1]{BVBeukers2006}, with common denominator $x$,
therefore says that $y(\alpha),y'(\alpha),e^\alpha$ are
algebraically independent over $\overline{\mathbb Q}$ for every
nonzero algebraic $\alpha$.

The determinant reductions use coefficients in
$\mathbb Q[x,x^{-1}]$, so $P$ in \eqref{E:acb:jetgauge} has no
pole at $\alpha$.  Moreover $P(\alpha)$ is invertible.
Indeed the symmetric-power Bessel system is ordinary there,
and so is the minimal scalar equation of $Y_k$ by the preceding
input.  Evaluation of its first $d+1$ jets is an isomorphism
from its solution space to $\mathbb C^{d+1}$.  The injective
polarized solution map of Theorem~\ref{E:acb:module} is an
isomorphism between solution spaces of that same dimension.
Their evaluation matrices are related by $P(\alpha)$, forcing
its determinant to be nonzero.

We may therefore specialize the monomial-coordinate argument
of Theorem~\ref{E:acb:jets}.  The independence of the three basic
values gives exactly the Veronese-cone relation ideal and
transcendence degree two.  Different rows of the invertible
$P(\alpha)$ are nonproportional homogeneous forms, giving
pairwise algebraic independence by the same Jacobian argument.
The case $k=1$ is the classical transcendence of $e^\alpha$.
\end{proof}

\begin{remark}[Scope of the arithmetic statement]
No corresponding absence of exceptional nonzero scalar points
is assumed here for the permutation and matching determinants.
Their functional order, Galois, and jet-cone conclusions hold
for all parameters as stated.  The arithmetic theorem above
uses Proposition~\ref{E:sfb:tableauordinary} explicitly;
functional Galois information alone would not justify its
all-points specialization.
\end{remark}


\clearpage

\section{Arithmetic subsequences and quasipolynomial coefficient weights}
\label{E:qpa:section}

The three algebraicity thresholds persist under every arithmetic
subsequence and every nonzero polynomial weighting.  In fact, one can
allow an arbitrary quasipolynomial weight.  The point requiring proof is
that the roots-of-unity filter cannot remove the long logarithmic block.
We prove this using the simple Bessel modules before applying the Borel
bound; no semisimplicity assertion about a Borel transform is needed.

A \emph{nonzero quasipolynomial} is a function $w:\mathbb Z_{\ge0}\to
\mathbb C$ for which some $q\ge1$ and polynomials
$P_0,\ldots,P_{q-1}\in\mathbb C[X]$, not all zero, satisfy
$w(n)=P_b(n)$ when $n\equiv b\pmod q$.
We use the counting sequences and normalizations of
Section~\ref{E:sfb:section}.

\begin{theorem}[Algebraicity survives no nonzero periodic polynomial filter]
\label{E:qpa:classification}
For every nonzero quasipolynomial $w$, the following equivalences hold:
\begin{align*}
 \sum_{n\ge0}w(n)y_k(n)t^n\text{ is algebraic over }\mathbb C(t)
    &\quad\Longleftrightarrow\quad k\le3,\\
 \sum_{n\ge0}w(n)u_k(n)t^n\text{ is algebraic over }\mathbb C(t)
    &\quad\Longleftrightarrow\quad k\le2,\\
 \sum_{n\ge0}w(n)f_k(n)t^n\text{ is algebraic over }\mathbb C(t)
    &\quad\Longleftrightarrow\quad k\le2.
\end{align*}
Here all parameters $k$ are positive integers.  Thus the threshold for
permutations avoiding $12\cdots\ell$ remains $\ell\le3$.
The same statements hold after changing finitely many coefficients.
\end{theorem}

\begin{corollary}[Every arithmetic progression and every polynomial weight]
\label{E:qpa:subsequence}
Let $q\ge1$, $b\ge0$, and $0\ne P\in\mathbb C[X]$.
The generating functions
\[
 \sum_{n\ge0}P(n)y_k(qn+b)t^n,\qquad
 \sum_{n\ge0}P(n)u_k(qn+b)t^n,\qquad
 \sum_{n\ge0}P(n)f_k(qn+b)t^n
\]
are algebraic exactly for $k\le3$, $k\le2$, and $k\le2$,
respectively.  No exceptional residue classes occur.
\end{corollary}

\subsection{A noncancellation lemma for simple differential modules}

Put $K=\mathbb C(z)$ and let $\mathcal E$ be a differential field
extension with field of constants $\mathbb C$.  A finite-dimensional
$K$-subspace of $\mathcal E$ stable under differentiation is viewed as a
differential module.

\begin{lemma}[Multiplicity-free sums retain their local blocks]
\label{E:qpa:noncancellation}
Let $M_1,\ldots,M_h\subset\mathcal E$ be simple differential modules
which are pairwise nonisomorphic.  Then their sum is direct.  If
$g_i\in M_i$ and $g=\sum_i g_i$, the differential module generated by
$g$ is exactly $\bigoplus_{g_i\ne0}M_i$.
Consequently, the local monodromy of the minimal equation of $g$
is the direct sum of those of the nonzero components.
\end{lemma}
\begin{proof}
The abstract direct sum of the $M_i$ is semisimple.  Every simple
submodule is isomorphic to one of the $M_i$, and projection to every
nonisomorphic summand is zero.  Because each isomorphism type occurs
once, its simple submodules are precisely the individual summands.
The kernel of the summation map to $\mathcal E$ therefore is a sum
of some summands.  The map is injective on each $M_i$, so its kernel
is zero.  The differential module generated by $g$ is a submodule of
this direct sum.  Its projection onto every index with $g_i\ne0$ is
nonzero and hence surjective.  The same multiplicity-free argument
shows that it contains exactly those summands.

The cyclic differential module of a scalar germ determines the
solution local system of its minimal equation, with the conventional
duality if necessary.  Direct sums and their Jordan block sizes are
preserved by this correspondence, giving the last assertion.
\end{proof}

\begin{lemma}[Filtered Bessel blocks]
\label{E:qpa:blocks}
Let $w$ be a nonzero quasipolynomial.  In the ranges specified below,
the minimal equations of the following entire functions are regular
singular at zero and have at least one block of the indicated size:
\[
\begin{array}{c|c|c}
 \text{function}&\text{parameter range}&\text{block size}\\ \hline
 \displaystyle\sum_nw(n)y_k(n)z^n/n!
      &k\ge2&\lfloor k/2\rfloor+1\\[2pt]
 \displaystyle\sum_nw(n)u_k(n)z^n/(n!)^2
      &k\ge1&k+1\\[2pt]
 \displaystyle\sum_nw(n)f_k(n)z^{2n}/(2n)!
      &k\ge2&k.
\end{array}
\]
More precisely, each complete local partition consists of one or more
copies of the displayed unipotent block, with no blocks of other sizes.
\end{lemma}
\begin{proof}
Write $\vartheta=z\,d/dz$.  Fourier expansion of the residue-class
indicators expresses each weighted series as a sum of terms
$Q_\zeta(\vartheta)H(\zeta z)$, where $\zeta$ ranges over roots of
unity, $Q_\zeta$ are polynomials, and $H$ is the appropriate
unweighted normalized series.  For matchings, use $2q$th roots and
$\vartheta/2$ because their normalized series is supported on even
powers.  For example, for tableaux the exact expression is
\[
 \sum_{b=0}^{q-1}P_b(\vartheta)
       \left\{\frac1q\sum_{\zeta^q=1}\zeta^{-b}Y_k(\zeta z)\right\}.
\]

For the $z$-Bessel equation, the differential module spanned by the
degree-$d$ monomials in $I_0(2az)$ and its derivative is simple,
has rank $d+1$, and has local block $J_{d+1}(1)$ at zero, by
Theorem~\ref{E:acb:module}.  Its formal exponential factors at infinity
are
\begin{equation}\label{E:qpa:untwisted-factors}
             \exp\bigl(2a(d-2j)z\bigr),\qquad 0\le j\le d.
\end{equation}
These follow directly from the two formal Bessel solutions
$e^{\pm2az}z^{-1/2}(1+O(z^{-1}))$ by taking symmetric powers.
The multiset is unchanged when $a$ is replaced by $-a$, and determines
$a$ up to that sign: its two extreme points are $\pm2da$.
Modules with distinct multisets of formal exponential factors are
nonisomorphic, since a rational gauge cannot change these factors.
Moreover, the two modules for $a$ and $-a$ are the same embedded
module, because $I_0$ is even and its derivative is odd.  We therefore
combine those two contributions before applying
Lemma~\ref{E:qpa:noncancellation}.

For even tableaux and for matchings, this is the entire argument,
with $d=\lfloor k/2\rfloor$ and $d=k-1$, respectively.  For odd
tableaux, write $k=2d+1$.  The exponential twist in
\eqref{B:eq:Gordon} changes the factors to
\begin{equation}\label{E:qpa:twisted-factors}
             \exp\bigl(a(1+2d-4j)z\bigr),\qquad0\le j\le d.
\end{equation}
These multisets distinguish every nonzero $a$: their sums of
exponents are $(d+1)a$.  The modules are again simple by
Theorem~\ref{E:acb:module}, and every homothety has local block
$J_{d+1}(1)$.

For permutations in the variable $z$, the degree-$k$ symmetric
power associated with $I_0(2\sqrt{az})$ has, after the common
ramification $z=v^2$, the formal factors
\[
               \exp\bigl(2\sqrt a(k-2j)v\bigr),\qquad0\le j\le k.
\]
Their extreme points determine $\sqrt a$ up to sign and hence
determine $a$.  Thus the modules for distinct roots of unity $a$
are pairwise nonisomorphic.  Each has rank $k+1$ and block
$J_{k+1}(1)$ before ramification as well, by
Theorem~\ref{E:acb:module}.

Application of a polynomial in $\vartheta$ keeps a term in its
simple differential module.  After grouping the equal modules
in the even cases, Lemma~\ref{E:qpa:noncancellation} applies to all
remaining terms, including any that become zero.  At least one
term remains nonzero.  Indeed $y_k(n),u_k(n)>0$ for every $n$,
and $f_k(n)>0$ for every $n$ when $k\ge2$, while a nonzero
quasipolynomial is nonzero at infinitely many integers.
Thus the weighted power series is nonzero.  Every surviving
summand has the same indicated local block, proving both the
lower bound and the complete partition assertion.
\end{proof}

\begin{proof}[Proof of Theorem~\ref{E:qpa:classification}]
Algebraic series are closed under homotheties, finite sums,
differentiation, and multiplication by rational functions.
The roots-of-unity formula and the Euler operators above show
that quasipolynomial coefficient weighting preserves algebraicity.
The algebraic directions therefore follow from
Theorem~\ref{E:sfb:algebraicity}.

For $k\ge4$ the weighted tableau exponential generating
function has a block of size $\lfloor k/2\rfloor+1\ge3$ by
Lemma~\ref{E:qpa:blocks}.  Algebraicity of its ordinary generating
function would bound that size by two after one Borel transform,
by Theorem~\ref{B:up:borel-bound}.  This is a contradiction.
For $k\ge3$, the weighted permutation function has a block of
size $k+1\ge4$, whereas two Borel transforms of an algebraic
series permit size at most three.  This is again impossible.
For matchings with $k\ge3$, algebraicity of the ordinary series
would imply algebraicity after $t\mapsto t^2$.  One Borel
transform of that substituted series is the weighted matching
function in Lemma~\ref{E:qpa:blocks}, whose block has size
$k\ge3$, contradicting the upper bound two.
Finally, changing finitely many ordinary coefficients adds a
polynomial, and therefore cannot change algebraicity.
\end{proof}

\begin{proof}[Proof of Corollary~\ref{E:qpa:subsequence}]
First suppose $0\le b<q$.  Set $w(m)=0$ outside the residue
class $b\pmod q$ and $w(m)=P((m-b)/q)$ on that class.
Then, for each of the three sequences $(a_n)$, the weighted
ordinary generating function is
\[
 \sum_{m\ge0}w(m)a_mt^m
      =t^b\sum_{n\ge0}P(n)a_{qn+b}(t^q)^n.
\]
The two series on the right and left are algebraic simultaneously,
because substitution by $t^q$ and its algebraic inverse preserve
algebraicity.  Apply Theorem~\ref{E:qpa:classification}.
For general $b=b_0+qh$, remove the first $h$ terms from the
corresponding $b_0$-subsequence and replace $P(n)$ by $P(n-h)$.
Removing a finite initial segment and multiplying by a power
of $t$ do not change algebraicity.
\end{proof}

\begin{corollary}[Exact differential orders after a single residue filter]
\label{E:qpa:orders}
Fix $q\ge1$, $0\le b<q$, and $0\ne P\in\mathbb C[X]$.
Use the inherited factorial normalization, namely
\[
 \mathcal Y(z)=\sum_{n\ge0}P(n)y_k(qn+b)
                    \frac{z^{qn+b}}{(qn+b)!},\qquad
 \mathcal U(z)=\sum_{n\ge0}P(n)u_k(qn+b)
                    \frac{z^{qn+b}}{((qn+b)!)^2},
\]
\[
 \mathcal F(z)=\sum_{n\ge0}P(n)f_k(qn+b)
                    \frac{z^{2(qn+b)}}{(2(qn+b))!}.
\]
The following minimal differential orders over $\mathbb C(z)$ are exact:
\[
\begin{array}{c|c|c}
 \text{function and parameters}&\text{order}&\text{partition at zero}\\ \hline
 \mathcal Y,\ k=2d\ge2
   &q(d+1)/\gcd(q,2)&J_{d+1}(1)^{\oplus q/\gcd(q,2)}\\
 \mathcal Y,\ k=2d+1\ge3
   &q(d+1)&J_{d+1}(1)^{\oplus q}\\
 \mathcal U,\ k\ge1
   &q(k+1)&J_{k+1}(1)^{\oplus q}\\
 \mathcal F,\ k\ge2
   &qk&J_k(1)^{\oplus q}.
\end{array}
\]
These are orders of the displayed normalized differential equations,
not assertions about the shift width of a recurrence for the sampled
coefficient sequence.
\end{corollary}
\begin{proof}
For a single residue filter, every root-of-unity coefficient is nonzero.
In the even-tableau case there are $q/\gcd(q,2)$ classes modulo sign.
If $q$ is even, each class contributes a nonzero multiple of
$Y_k(az)+(-1)^bY_k(-az)$.  This is nonzero because its surviving
parity has strictly positive counting coefficients.  If $q$ is odd,
each class has a single nonzero contribution.  The odd-tableau and
permutation cases have $q$ nonisomorphic simple summands.  For
matchings use the $2q$th-root filter and the residue $2b$; combining
opposite roots gives $q$ nonzero summands because $F_k$ is even.

It remains to justify that applying the weight polynomial kills
none of these summands.  Before weighting, each nonzero component
is holomorphic at zero and generates a simple module of rank at
least two.  It cannot be a polynomial: a nonzero polynomial
generates the rank-one differential module $\mathbb C(z)$.
For a holomorphic germ $g(z)=\sum c_nz^n$, the identity
$Q(\vartheta)g=0$ with a nonzero polynomial $Q$ forces
$c_n=0$ unless $Q(n)=0$, and would therefore make $g$ a polynomial.
Thus every component survives the relevant operator
$P((\vartheta-b)/q)$, or
$P((\vartheta/2-b)/q)$ for matchings.
Apply Lemma~\ref{E:qpa:noncancellation} and count the surviving
simple summands, each with its already known local block.
\end{proof}

\begin{remark}[Exact scope]
These statements concern algebraicity of generating functions over
$\mathbb C(t)$, not arithmetic properties of their values at particular
arguments. They extend the unfiltered classification and use the
classical enumeration identities cited above. They do not require a joint
differential-Galois group theorem for all Bessel scales.
\end{remark}


\section{An algebraicity dichotomy for homogeneous Bessel observables}
\label{fusion:observables}

Write $y(t)=I_0(2t)$.  A \emph{homogeneous Laurent Bessel observable}
of degree $r$ is a function
\[
 F(t)=\Phi(t,y(t),y'(t)),\qquad
 \Phi\in\mathbb C[t,t^{-1}][X_0,X_1],
\]
where $\Phi$ is homogeneous of degree $r$ in $X_0,X_1$.
The Laurent condition allows powers $t^j$ with integer
$j$, but no other denominators.  The following theorem uses the
holomorphic cancellation of all poles at $t=0$.

For clarity, a quasipolynomial weight is a function $w$ for which
there are an integer $h\ge1$ and polynomials
$P_0,\ldots,P_{h-1}\in\mathbb C[X]$ such that
\[
 w(n)=P_s(n)\quad\text{when }n\equiv s\pmod h.
\]
It is nonzero if at least one $P_s$ is nonzero.  No positivity of
these polynomials is assumed.

\begin{theorem}[Homogeneous-observable dichotomy]
\label{fusion:observable-dichotomy}
Let $r\ge1$ and $p\in\{0,1\}$.  Suppose that a homogeneous
Laurent Bessel observable $F$ of degree $r$ is nonzero,
holomorphic at zero, and satisfies $F(-t)=(-1)^pF(t)$.  Put
\[
 a_n=(2n+p)!\,[t^{2n+p}]F(t).
\]
For every $d\ge1$, $0\le b<d$, and every nonzero
quasipolynomial weight $w$, the series
\[
 A_{d,b,w}(z)=\sum_{n\ge0}w(n)a_{dn+b}z^n
\]
is algebraic over $\mathbb C(z)$ if and only if $r=1$.
None of these weighted sampled sequences is eventually zero.
Moreover, when $r=1$, the unweighted unsampled series satisfies
the sharper conclusion
\[
 \sum_{n\ge0}a_nz^n\in
 \mathbb C\bigl(z,\sqrt{1-4z}\bigr).
\]
\end{theorem}

The exact-width result of Corollary~\ref{B:up:general-parity-width}
and the algebraicity dichotomy measure different phenomena.
For example, degree-two observables have coefficient recurrence
width one, but their ordinary generating functions are
transcendental even after arbitrary nonzero quasipolynomial
weighting and arithmetic sampling.

\subsection{The holomorphic lattice in degree one}

For $m\ge0$, define
\[
 J_m(z)=\sum_{n\ge0}\frac{z^n}{n!(n+m)!}.
\]
Then $J_0(t^2)=y(t)$ and $y'(t)=2tJ_1(t^2)$.
Direct coefficient comparison gives
\begin{equation}\label{fusion:J-recurrence}
 zJ_{m+2}(z)=J_m(z)-(m+1)J_{m+1}(z),\qquad m\ge0.
\end{equation}

\begin{lemma}[Holomorphic division in the Bessel lattice]
\label{fusion:holomorphic-lattice}
Let
\[
 \mathcal S=\left\{\sum_{m=0}^{M}P_m(z)J_m(z):
 M<\infty,\ P_m\in\mathbb C[z]\right\}.
\]
If $g\in\mathcal S$ and $g(0)=0$, then $g(z)/z\in\mathcal S$.
Consequently every function
$a(z)J_0(z)+b(z)J_1(z)$ with
$a,b\in\mathbb C[z,z^{-1}]$ that is holomorphic at zero
belongs to $\mathcal S$.
\end{lemma}
\begin{proof}
Write $P_m(z)=P_m(0)+zQ_m(z)$.  It remains to divide a
constant linear combination $\sum_m c_mJ_m$ whose value at
zero vanishes.  Put $K_m=m!J_m$ and $d_m=c_m/m!$, so
$\sum_md_m=0$.  If $D_j=\sum_{m=0}^{j}d_m$, telescoping gives
\[
 \sum_{m=0}^{M}d_mK_m
 =\sum_{j=0}^{M-1}D_j(K_j-K_{j+1})
 =z\sum_{j=0}^{M-1}D_jj!J_{j+2},
\]
where the last identity is \eqref{fusion:J-recurrence}.
This proves the closure under division by $z$.

For the last assertion, choose $N$ so that
$z^N(aJ_0+bJ_1)\in\mathcal S$.  Its first $N$ Taylor
coefficients vanish, because the function before multiplication
is holomorphic.  Apply the division result successively $N$
times.
\end{proof}

\begin{lemma}[Algebraic factorial transforms of the lattice]
\label{fusion:lattice-algebraic}
For every $m\ge0$,
\[
 H_m(z)=\sum_{n\ge0}\frac{(2n)!}{n!(n+m)!}z^n
 \quad\text{belongs to}\quad
 \mathbb C\bigl(z,\sqrt{1-4z}\bigr).
\]
If $f\in\mathcal S$ and $p\in\{0,1\}$, the same is true of
$\sum_{n\ge0}(2n+p)![z^n]f(z)\,z^n$.
\end{lemma}
\begin{proof}
One has $H_0=(1-4z)^{-1/2}$.  For $m\ge1$, the beta
integral, evaluated term by term near $z=0$, gives
\[
 H_m(z)=\frac1{(m-1)!}
 \int_0^1\frac{(1-u)^{m-1}}{\sqrt{1-4zu}}\,du.
\]
Set $v=\sqrt{1-4zu}$, taking the branch equal to $1$ at
$z=0$.  The resulting identity is
\[
 H_m(z)=\frac1{2z(m-1)!(4z)^{m-1}}
 \int_{\sqrt{1-4z}}^1(v^2+4z-1)^{m-1}\,dv.
\]
The integrand is a polynomial in $v$, so evaluating its
polynomial antiderivative at the two endpoints proves the
claimed quadratic-field membership.  The displayed apparent
pole at zero is removable by the original power series.

For a generator $z^kJ_m(z)$ of $\mathcal S$, with $k\ge0$,
its factorial transform of parity $p$ is
\[
 z^k Q_{k,p}(\theta)H_m(z),\qquad
 Q_{k,p}(X)=\prod_{j=1}^{2k+p}(2X+j),
 \quad\theta=z\partial_z.
\]
This follows by putting $n=\nu+k$ in its coefficients.
The quadratic function field is closed under differentiation,
and finite linear combinations prove the second assertion.
\end{proof}

\subsection{Why polynomial weights cannot remove the obstruction}

\begin{lemma}[Nonzero Euler polynomials preserve the module]
\label{fusion:euler-noncancellation}
Let $V$ be a simple differential module over $\mathbb C(t)$
realized by meromorphic functions.  Suppose $V$ has a nonzero
formal exponential rate at infinity.  For every $0\ne f\in V$
and every nonzero polynomial $Q\in\mathbb C[X]$,
$Q(t\partial_t)f$ is nonzero and generates $V$.
\end{lemma}
\begin{proof}
If $Q(t\partial_t)f=0$, the minimal scalar equation of $f$
is a right divisor of this Euler operator.  Its local solutions
at infinity would therefore be finite linear combinations
of powers of $t$ times powers of $\log t$, with no nonzero
exponential rate.  This contradicts the hypothesis, since a
nonzero element of the simple module is cyclic.  Thus the
resulting element is nonzero, and simplicity makes it cyclic
as well.
\end{proof}

\begin{proof}[Proof of Theorem~\ref{fusion:observable-dichotomy}]
First take $r=1$, and write $F=A(t)y+B(t)y'$ with Laurent
polynomials $A,B$.  The functions $y,y'$ are linearly
independent over $\mathbb C(t)$, as follows from the Bessel
$\mathrm{SL}_2$ group.  The parity assumption therefore
implies
\[
 A(t)=t^pa(t^2),\qquad B(t)=t^{p-1}b(t^2),
 \qquad a,b\in\mathbb C[z,z^{-1}].
\]
Consequently
$F(t)=t^p f(t^2)$ with $f=aJ_0+2bJ_1$.
Holomorphicity of $F$ implies holomorphicity of $f$.
Lemmas~\ref{fusion:holomorphic-lattice} and
\ref{fusion:lattice-algebraic} now prove the stronger
quadratic-field assertion for the full coefficient series.

Arithmetic sections of algebraic series are algebraic by
root-of-unity filtering and algebraic substitution.
Polynomial coefficient weighting is application of a
polynomial in $z\partial_z$, which preserves algebraicity.
Finally a quasipolynomial can be expressed as a finite sum
$\sum_{\omega^h=1}\omega^n Q_\omega(n)$ with polynomial
$Q_\omega$, by Fourier inversion on its residue polynomials.
Its weighted generating function is a finite sum of
$Q_\omega(z\partial_z)A(\omega z)$, and is algebraic.
This proves all the upper assertions for $r=1$.

For the converse, first consider a polynomial weight $P(n)$
on an arithmetic section, where $P\ne0$.  Replacing $F$ by
\[
 F_P(t)=P\left(\frac{(t\partial_t-p)/2-b}{d}\right)F(t)
\]
multiplies its coefficient at $t^{2m+p}$ by
$P((m-b)/d)$.  It remains a holomorphic homogeneous Laurent
Bessel observable of degree $r$ and parity $p$.  It is
nonzero by Lemma~\ref{fusion:euler-noncancellation}, since
the underlying simple symmetric-power Bessel module has
nonzero rates $2r$ and $-2r$.
Thus it suffices to rule out algebraicity of unweighted
arithmetic sections for any observable with these properties.

Suppose $A(z)=\sum_{n\ge0}a_{dn+b}z^n$ is algebraic.
The padded series $z^{2b+p}A(z^{2d})$ is algebraic, and its
onefold factorial Borel transform is
\[
 H(t)=\sum_{n\ge0}a_{dn+b}
                 \frac{t^{2(dn+b)+p}}{(2(dn+b)+p)!}.
\]
Theorem~\ref{B:up:borel-bound} implies that all its local
Jordan blocks at zero have size at most two.  On the other
hand, the root-of-unity decomposition in the proof of
Corollary~\ref{B:up:general-parity-width} identifies its
module with a direct sum of $d$ pairwise nonisomorphic
simple $r$th symmetric-power Bessel modules.  Each has a
unipotent block of size $r+1$ at zero.  This contradicts
the Borel bound when $r\ge2$.

Now let $w$ be a nonzero quasipolynomial of period $h$ and
choose a residue $s$ with $P_s\ne0$.  If $A_{d,b,w}$ were
algebraic, so would its section at $n=hm+s$, whose coefficients
are
\[
 P_s(hm+s)\,a_{dhm+ds+b}.
\]
This is a nonzero polynomial weighting of an arithmetic
section, already excluded for $r\ge2$.  For every $r\ge1$, the same choice of
residue also proves that the weighted sequence is not
eventually zero: Corollary~\ref{B:up:general-parity-width}
gives that property for $a_{dhm+ds+b}$, and the nonzero
polynomial $P_s(hm+s)$ has only finitely many integer roots.
\end{proof}

\begin{remark}[The homogeneity assumption cannot be omitted]
The function
\[
 F(t)=\frac{I_0(2t)-1}{t}
\]
is holomorphic and odd, but is not homogeneous of degree one
in the Bessel coordinates.  Its factorial coefficient series is
\[
 \sum_{n\ge0}\frac{(2n+1)!}{((n+1)!)^2}z^n
 =\frac1z\log\left(\frac{2}{1+\sqrt{1-4z}}\right).
\]
The identity follows by differentiating after multiplication
by $z$ and matching the zero constant term.  This function
is not algebraic: at infinity, on a sector with fixed branches,
it has a nonzero $(\log z)/z$ term, whereas algebraic functions
have Puiseux expansions without logarithms.  Thus arbitrary
integration or subtraction of constants cannot replace the
holomorphic-lattice argument above.
\end{remark}

\section{Quasipolynomial robustness of the twisted classification}
\label{fusion:twisted-weights}

\begin{theorem}[Weighted sections of twisted determinants]
\label{fusion:twisted-weighted-dichotomy}
Let $\ell,d\ge1$, $q\ge0$, and $0\le b<d$ be integers, and
let $w$ be a nonzero quasipolynomial.  Then
\[
 \sum_{n\ge0}w(n)u_{\ell,q}(dn+b)z^n
\]
is algebraic over $\mathbb C(z)$ if and only if $\ell\le2$.
For $\ell=1$ it is rational.
\end{theorem}
\begin{proof}
The identities $u_{1,q}(n)=1$ and $u_{2,q}(n)=C_{n+q}$
from Proposition~\ref{B:up:twist-catalan}, together with the
algebraic closure operations just proved, handle $\ell\le2$.
The generating function of a quasipolynomial is rational,
which gives the stronger statement for $\ell=1$.

For $\ell\ge3$, extraction of a residue on which $w$ is a
nonzero polynomial reduces the claim to a polynomial weight
$P(n)$ on a section $u_{\ell,q}(dn+b)$.
Suppose its generating function $A(z)$ were algebraic.
Put $Q=q\ell$ and $F_0(z)=z^bA(z^d)$, which is algebraic.
Two Borel transforms with an intervening monomial gauge give
\[
 \mathcal B\bigl(z^Q\mathcal B F_0(z)\bigr)=z^QH(z),
 \quad
 H(z)=\sum_{n\ge0}
 \frac{P(n)u_{\ell,q}(dn+b)z^{dn+b}}
 {(dn+b)!(dn+b+Q)!}.
\]
Theorem~\ref{B:up:borel-bound} and the rational gauge $z^Q$
imply that every Jordan block of $H$ at zero has size at
most three; nonunit blocks are semisimple.

After $z=t^2$, the unweighted filtered function is the direct
sum of $d$ simple $\ell$th symmetric-power Bessel modules
at scales distinct modulo sign.  The polynomial weight is
application of
$P(((t\partial_t)/2-b)/d)$ to this filter.
Lemma~\ref{fusion:euler-noncancellation} shows that each
component remains nonzero, and
Lemma~\ref{B:up:simple-filter} preserves their direct sum.
Each component has a unipotent block of size $\ell+1\ge4$
at zero.  Squaring monodromy under $z=t^2$ cannot increase
any Jordan-block size, since every monodromy eigenvalue is
nonzero.  This contradicts the upper bound three and proves
the theorem.
\end{proof}

\begin{remark}
These classifications concern algebraicity over a rational
function field.  They do not assert irrationality or
transcendence of individual numerical values of the generating
functions.  The new scope here is uniformity in the homogeneous
observable, together with robustness under every arithmetic
sampling and every nonzero quasipolynomial weighting.
\end{remark}


\section{Multipoint arithmetic of determinant jets}
\label{mpj:section}

The all-twist Laurent gauge can be combined with the classical
transcendence theory of scaled Bessel functions. The resulting
statement describes \emph{all} relations among normalized Toeplitz
jets at any finite set of distinct nonzero algebraic arguments.
The differential-Galois independence step is proved explicitly below;
pairwise distinct exponential rates alone would not establish a
product Galois group. Multipoint independence of the scalar Bessel
values themselves is classical: Siegel's 1929 theorem for $J_0$,
recalled in \cite[\S1]{MPJDelaygue2025}, already covers distinct
nonzero squared algebraic arguments. The purpose here is the exact
all-twist jet ideal and field structure, including the derivatives.

\subsection{The product group for scaled Bessel connections}

For $a\in\mathbb C^*$ let $M_a$ be the rank-two connection
\begin{equation}\label{mpj:system}
 \frac d{dt}\begin{pmatrix}f\\f'\end{pmatrix}
 =A_a(t)\begin{pmatrix}f\\f'\end{pmatrix},
 \qquad
 A_a(t)=\begin{pmatrix}0&1\\4a^2&-1/t\end{pmatrix}.
\end{equation}
It has differential Galois group $\mathrm{SL}_2$, by the Bessel
result already used above and the nonzero linear change of variable
$t\mapsto at$. Its entire distinguished solution is
$b_a(t)=I_0(2at)$.

\begin{lemma}[Subdirect products of the rank-two groups]
\label{mpj:lie}
Let $H$ be an algebraic subgroup of $\mathrm{SL}_2^m$ whose
projection onto each factor is surjective. Then $H^\circ$ also
projects surjectively onto every factor. If
$H\ne\mathrm{SL}_2^m$, there are two distinct factors whose
three-dimensional adjoint representations restrict to isomorphic
representations of $H^\circ$.
\end{lemma}
\begin{proof}
The image of $H^\circ$ in a factor is a connected normal subgroup
with finite quotient in $\mathrm{SL}_2$. Since $\mathrm{SL}_2$
is connected, this image is the whole factor. Consequently
$\mathfrak h=\operatorname{Lie}(H^\circ)$ is a subdirect Lie
subalgebra of $\mathfrak{sl}_2^m$.
Its solvable radical projects into a solvable ideal of each
$\mathfrak{sl}_2$, hence projects to zero everywhere and is zero.
Thus $\mathfrak h$ is a direct sum of simple ideals.

The image of a simple ideal under any coordinate projection is
an ideal of $\mathfrak{sl}_2$. If nonzero, the projection is an
isomorphism onto $\mathfrak{sl}_2$: its kernel is an ideal of the
simple source. Exactly one simple ideal maps onto a given factor,
since distinct simple ideals commute and $\mathfrak{sl}_2$ does
not commute with itself. Every simple ideal maps nontrivially to
at least one factor. If none maps to two factors, there are $m$
simple ideals and $\mathfrak h=\mathfrak{sl}_2^m$; this forces
$H^\circ=H=\mathrm{SL}_2^m$.
Otherwise two factors share a simple ideal. The isomorphism
between their two copies of $\mathfrak{sl}_2$ intertwines the
adjoint actions of $\mathfrak h$, and hence those of the connected
group $H^\circ$. This proves the assertion.
\end{proof}

\begin{lemma}[Adjoint scales cannot merge algebraically]
\label{mpj:adjoint}
If $M_a$ and $M_b$ have isomorphic adjoint differential modules
after a finite algebraic extension of $\mathbb C(t)$, then
$a=b$ or $a=-b$.
\end{lemma}
\begin{proof}
The constant term of $A_a(t)$ at infinity has the distinct
eigenvalues $2a,-2a$. The induced adjoint connection on traceless
endomorphisms therefore has the three distinct leading eigenvalues
\[
                         0,\quad 4a,\quad -4a.
\]
A formal change of basis over $\mathbb C((t^{-1}))$ diagonalizes
this connection with diagonal entries
$\lambda+O(t^{-1})$, one for each displayed $\lambda$.
Indeed, after diagonalizing the constant term, successive
coefficients of the off-diagonal gauge are solved uniquely by
dividing by the nonzero eigenvalue differences. The same applies
to the adjoint connection at scale $b$.

Choose a place above infinity in the finite algebraic extension.
Its completion embeds in a Puiseux field
$\mathbb C((t^{-1/e}))$ for some $e\ge1$. After the two formal
diagonalizations, any entry $h$ of an intertwining matrix obeys
\[
 h'=(\mu-\lambda+O(t^{-1}))h,
 \qquad \lambda\in\{0,4a,-4a\},\quad
        \mu\in\{0,4b,-4b\}.
\]
For a nonzero Puiseux Laurent series, $h'/h=O(t^{-1})$.
Thus such an entry can be nonzero only when $\lambda=\mu$.
An invertible intertwiner requires equality of the two multisets
of leading eigenvalues. Since $a,b\ne0$, this gives $a=\pm b$.
\end{proof}

\begin{theorem}[Full product group and functional independence]
\label{mpj:product}
Let $a_1,\ldots,a_m\in\mathbb C^*$ satisfy
$a_i\ne\pm a_j$ for $i\ne j$. The differential Galois group of
$M_{a_1}\oplus\cdots\oplus M_{a_m}$ over $\mathbb C(t)$ is
\[
                              \mathrm{SL}_2^m.
\]
The $2m$ functions
$b_{a_1},b_{a_1}',\ldots,b_{a_m},b_{a_m}'$
are algebraically independent over $\mathbb C(t)$.
\end{theorem}
\begin{proof}
Let $E$ be a Picard--Vessiot field for the direct sum and $H$ its
Galois group. Restriction to each summand makes $H$ a subdirect
algebraic subgroup of $\mathrm{SL}_2^m$. If it were proper,
Lemma~\ref{mpj:lie} would supply two adjoint representations
isomorphic on $H^\circ$. Let $F=E^{H^\circ}$; the differential
Galois correspondence makes $F/\mathbb C(t)$ finite algebraic.
The two adjoint differential modules would then be isomorphic
over $F$. Explicitly, if $U_i,U_j$ are their fundamental
matrices and $C$ is a constant intertwiner of the
$H^\circ$-representations, the matrix $U_jCU_i^{-1}$ is fixed
by $H^\circ$ and is the required invertible differential gauge
over $F$. Lemma~\ref{mpj:adjoint} contradicts the condition on
the scales. Hence $H=\mathrm{SL}_2^m$.

For each $i$, complete $b_{a_i}$ to a solution basis
$b_{a_i},v_i$ of $M_{a_i}$. The product group moves these
first solution columns independently to
\[
 c_i\begin{pmatrix}b_{a_i}\\b_{a_i}'\end{pmatrix}
 +d_i\begin{pmatrix}v_i\\v_i'\end{pmatrix},
 \qquad (c_i,d_i)\ne(0,0).
\]
A polynomial relation among the $2m$ functions would therefore
vanish after these substitutions for all constants $c_i,d_i$.
Their product parameter set is Zariski dense in $\mathbb C^{2m}$,
and each two-by-two substitution matrix has nonzero Wronskian.
The relation must be the zero polynomial. This proves functional
independence.
\end{proof}

\begin{remark}[Why central or projective identifications do not escape]
The proof passes to $H^\circ$ and compares adjoint modules over
its finite algebraic fixed field. A finite central ambiguity or
a rank-one projective twist disappears in the adjoint comparison;
it cannot change the leading-rate multiset used in
Lemma~\ref{mpj:adjoint}. Once the Lie algebra is the full product,
$H^\circ$ is already the full connected ambient group. No separate
assumption that the original subdirect group is connected is used.
\end{remark}

\subsection{Simultaneous arithmetic specialization}

Recall $y(x)=I_0(2\sqrt{x})$ and the Laurent matrices
$P_{\ell,q}$ of \eqref{B:bv:gauge}. For the remainder of this section
put $k=\overline{\mathbb Q}$.

\begin{theorem}[Independent Bessel coordinates at distinct arguments]
\label{mpj:basevalues}
If $\alpha_1,\ldots,\alpha_m\in k^*$ are pairwise distinct,
then the $2m$ numbers
\[
                y(\alpha_i),\ y'(\alpha_i),\qquad 1\le i\le m,
\]
are algebraically independent over $k$.
\end{theorem}
\begin{proof}
Choose $a_i\in k^*$ with $a_i^2=\alpha_i$ and set
$b_i(t)=I_0(2a_it)=y(\alpha_i t^2)$. Distinctness of the
$\alpha_i$ is exactly the condition $a_i\ne\pm a_j$.
Each $b_i,b_i'$ is an $E$-function: its coefficients are
algebraic scalar multiples of the central binomial coefficients,
with heights and common-denominator growth bounded exponentially.
The block-diagonal first-order system \eqref{mpj:system} has
common denominator $t$ and no other finite singularity.
By Theorem~\ref{mpj:product} its $2m$ distinguished coordinate
functions are algebraically independent. The
Siegel--Shidlovskii theorem, in the formulation
\cite[Theorem~1.1]{BVBeukers2006}, applied at $t=1$, proves
independence of $b_i(1),b_i'(1)$ for all $i$ simultaneously.
Finally
\[
 b_i(1)=y(\alpha_i),\qquad b_i'(1)=2\alpha_i y'(\alpha_i).
\]
The nonzero algebraic factors preserve independence.
\end{proof}

\begin{theorem}[The multipoint relation ideal]
\label{mpj:alljets}
Let the $\alpha_i\in k^*$ be pairwise distinct. Independently
choose integers $\ell_i\ge1$ and $q_i\ge0$, and put
\[
 \mathbf J_i=
 \bigl(G_{\ell_i,q_i}(\alpha_i),G_{\ell_i,q_i}'(\alpha_i),
                 \ldots,G_{\ell_i,q_i}^{(\ell_i)}(\alpha_i)\bigr)^{\mathsf T}.
\]
Then
\[
 \operatorname{trdeg}_k k(\mathbf J_1,\ldots,\mathbf J_m)=2m.
\]
For each $i$, use the invertible algebraic coordinate change
$\mathbf Z_i=P_{\ell_i,q_i}(\alpha_i)^{-1}\mathbf J_i$.
The ideal of all polynomial relations among the combined
coordinates is the sum of the $m$ ideals generated by
\begin{equation}\label{mpj:ideal}
 Z_{i,r}Z_{i,s+1}-Z_{i,r+1}Z_{i,s},
 \quad 0\le r<s\le\ell_i-1,\quad 1\le i\le m.
\end{equation}
Thus every relation is generated by the individual-argument
Veronese relations; no additional generator connecting distinct
arguments is needed. Choosing any two distinct initial jet
entries at each argument gives $2m$ algebraically independent
numbers over $k$.
\end{theorem}
\begin{proof}
By Theorem~\ref{B:bv:gaugedet}, every specialized gauge is
invertible. Set $u_i=y(\alpha_i)$ and $v_i=y'(\alpha_i)$.
Then
\[
                  Z_{i,r}=u_i^{\ell_i-r}v_i^r.
\]
The $2m$ base numbers are independent by
Theorem~\ref{mpj:basevalues}. The polynomial substitution above
uses disjoint pairs of independent variables. Within each pair,
the kernel is the Veronese ideal proved in
Theorem~\ref{B:bv:jetcone}. Tensoring these monomial subalgebras
introduces no new relation: distinct products of their monomial
bases have distinct exponent vectors in the independent variables
$u_1,v_1,\ldots,u_m,v_m$. Hence the combined kernel is exactly
\eqref{mpj:ideal}, and its dimension is $2m$.

Two distinct jet entries at argument $\alpha_i$ are
nonproportional homogeneous forms of the same positive degree
in $u_i,v_i$, because the specialized gauge has independent rows.
Their two-variable Jacobian is nonzero, as in the proof of
Theorem~\ref{B:bv:jetcone}. The Jacobian for all the chosen pairs
is block diagonal with these nonzero determinants. In
characteristic zero this proves their joint algebraic independence.
\end{proof}

\begin{corollary}[Arbitrarily many independent determinant values]
\label{mpj:individual}
With the hypotheses above, the $m$ values
$G_{\ell_i,q_i}(\alpha_i)$ are algebraically independent over $k$.
The same holds if one chooses, independently at each argument,
any one of the derivatives of orders $0,\ldots,\ell_i$.
In particular any family of such choices indexed by distinct
nonzero algebraic arguments is algebraically independent in the
sense that every finite subfamily is independent.
\end{corollary}
\begin{proof}
Complete each choice to a pair of distinct initial jet entries
and apply Theorem~\ref{mpj:alljets}. Algebraic independence of
an arbitrary family is, by definition, a finite-subfamily property.
\end{proof}

\subsection{The exact field laws across arguments and dimensions}

Fix distinct $\alpha_1,\ldots,\alpha_m\in k^*$ and write
\[
 \mathcal E=k(u_1,v_1,\ldots,u_m,v_m),\qquad
 (u_i,v_i)=(y(\alpha_i),y'(\alpha_i)).
\]
For a vector $\mathbf d=(d_1,\ldots,d_m)$ of positive integers,
let $\mathcal F_{\mathbf d}$ be the field generated by complete
initial jets of dimension $d_i$ at $\alpha_i$, choosing arbitrary
nonnegative integer twists at the different arguments.

\begin{theorem}[Product cyclic fields and componentwise gcd--lcm]
\label{mpj:fields}
The field $\mathcal F_{\mathbf d}$ is independent of all chosen
twists, and
\[
 \mathcal F_{\mathbf d}
 =k\bigl(u_i^{d_i},v_i/u_i:1\le i\le m\bigr)
 =\mathcal E^{\mu_{d_1}\times\cdots\times\mu_{d_m}}.
\]
The $i$th cyclic group scales $u_i,v_i$ simultaneously and fixes
the other pairs. The extension
$\mathcal E/\mathcal F_{\mathbf d}$ is Galois of degree
$\prod_i d_i$, with precisely that product group. For any two
positive dimension vectors $\mathbf d,\mathbf e$,
\[
 \mathcal F_{\mathbf d}\mathcal F_{\mathbf e}
       =\mathcal F_{\gcd(\mathbf d,\mathbf e)},\qquad
 \mathcal F_{\mathbf d}\cap\mathcal F_{\mathbf e}
       =\mathcal F_{\operatorname{lcm}(\mathbf d,\mathbf e)},
\]
where gcd and lcm are taken componentwise.
In particular dimension choices coprime at every argument recover
all $2m$ basic Bessel values exactly.
\end{theorem}
\begin{proof}
The exact single-argument field identity \eqref{B:bv:field},
specialized by the invertible gauge, gives the displayed
monomial generators. Theorem~\ref{mpj:basevalues} identifies
$\mathcal E$ with a rational function field in $2m$ variables
over $k$. Therefore the independent scalar actions have the
stated fixed field and degree. Put
$L_i=\operatorname{lcm}(d_i,e_i)$ and work in the finite Galois
extension $\mathcal E/\mathcal F_{\mathbf L}$.
The intersection of the two fixing subgroups is
$\prod_i\mu_{\gcd(d_i,e_i)}$, and the subgroup they generate
is $\prod_i\mu_{L_i}$. The Galois correspondence gives the two
field identities.
\end{proof}

\begin{remark}[Scope and attribution]
The use of scaled Bessel equations and Siegel--Shidlovskii
specialization belongs to classical transcendence theory. The
contribution of this section is the explicit combination with the
all-twist, all-nonzero-argument Laurent gauge: it supplies complete
multipoint determinant-jet ideals and exact field laws. These
statements do not assert a new transcendence criterion for unrelated
constants, nor do they certify external priority for the classical
Bessel independence input.
\end{remark}


\clearpage\part{Entropy transport and certified finite-dimensional recovery}

\section{The constant-one Chebyshev--Bernstein theorem}
\label{E:cba:sec}

We prove the constant-one coefficient bound by an exact rational
certificate, printed below, and then derive an energy refinement.
Section~\ref{E:cbx:sec} determines the optimal exponential constant when a
fixed universal prefactor is allowed. These are different normalizations
of the coefficient inequality.

Write $B_j^t(x)=\binom tjx^j(1-x)^{t-j}$ and define $C(t,m,j)$ by
\[
 T_m(2x-1)=\sum_{j=0}^t C(t,m,j)B_j^t(x),\qquad 0\le m\le t.
\]

\begin{theorem}[The coefficient conjecture]
\label{E:cba:constantone}
For all integers $t\ge1$, $0\le m\le t$, and $0\le j\le t$,
\begin{equation}\label{E:cba:conjecture}
 |C(t,m,j)|\le \exp(m^2/t).
\end{equation}
The inequality is strict for $m>0$. The proof below is a finite-certificate proof: its only computational part
uses exact rational arithmetic for $t<512$. The complete verifier is
included in the text.
\end{theorem}

This is the precise constant-one assertion posed in
\cite[Supplement, Remark~4.1, equation~(23)]{BT_VKVK2019}, rather than a
bound with an unspecified multiplicative constant. The distinction is discussed in
\cite[Appendix~C.2]{BT_MMRS2025}. The universal exponent with an unspecified
fixed prefactor is determined in
Theorem~\ref{E:cbx:main}; the prefactor-one optimum remains separate.

\begin{lemma}[Three complementary coefficient bounds]
\label{E:cba:threebounds}
For $1\le m\le t$,
\begin{align}
 |C(t,m,j)|&\le 2^m,\label{E:cba:degreebound}\\
 |C(t,m,j)|&\le (t+1)\exp\{tH(m/t)\}
       \le(t+1)\exp\{(\log2)m^2/t\},\label{E:cba:contourbound}
\end{align}
where $H(a)=\frac12[(1+a)\log(1+a)+(1-a)\log(1-a)]$ and $0\log0=0$.
For $m<t$ there is also the bound
\begin{equation}\label{E:cba:dettebound}
 |C(t,m,j)|\le D(t,m):=
          \prod_{l=1}^{m-1}\frac{t+l}{t-l}.
\end{equation}
In particular, $D(t,m)\le e^{m^2/t}$ whenever
\begin{equation}\label{E:cba:integercriterion}
 (m+6)(m-1)^2\le6t^2.
\end{equation}
\end{lemma}
\begin{proof}
At degree $m$ the coefficients are
$(-1)^{m-i}\binom{2m}{2i}/\binom mi$. The successive absolute-value ratio
is $(2m-2i-1)/(2i+1)$, so their maximum occurs at a middle index.
Vandermonde's identity bounds the numerator at that index by
$2^m\max_i\binom mi$, proving \eqref{E:cba:degreebound} at degree $m$.
Degree elevation replaces coefficients by convex combinations and proves
the bound at every $t\ge m$.

For \eqref{E:cba:contourbound}, put
\[
 F(z)=(1+z)^tT_m\left(\frac{z-1}{z+1}\right)
     =\sum_{j=0}^t\binom tj C(t,m,j)z^j.
\]
Endpoint coefficients have absolute value one. For $0<j<t$, set
$r=j/(t-j)$ and, on $z=re^{i\theta}$, set
$u=2\sqrt r\,|\sin(\theta/2)|/(1+r)$. The larger inverse Joukowski
modulus of $(z-1)/(z+1)$ is
$\sqrt{(1+u)/(1-u)}$. Consequently
\[
 \frac{|F(re^{i\theta})|}{(1+r)^t}
 \le(1-u^2)^{t/2}
          \left(\frac{1+u}{1-u}\right)^{m/2}
 \le \exp\{tH(m/t)\}.
\]
The last maximum is attained at $u=m/t$; limiting values handle $u=1$.
Cauchy's formula and
$\binom tj r^j/(1+r)^t\ge1/(t+1)$ give the first inequality in
\eqref{E:cba:contourbound}. This probability is a mode of a
$\operatorname{Binomial}(t,j/t)$ law, which proves the stated lower bound.
Finally
\[
 H(a)=\sum_{k\ge1}\frac{a^{2k}}{2k(2k-1)}
       \le a^2\sum_{k\ge1}\frac1{2k(2k-1)}=a^2\log2.
\]

To apply the older bound \eqref{E:cba:dettebound}, combine the
hypergeometric formula with the identity
\[
 x^r=\sum_j(j)_{\underline r}B_j^t(x)/(t)_{\underline r}.
\]
This gives
\[
 C(t,m,j)=(-1)^m\,{}_3F_2
       \left(\begin{matrix}-m,m,-j\\1/2,-t\end{matrix};1\right)
       =(-1)^mQ_m(j;-1/2,-1/2,t).
\]
Dette's Hahn bound \cite[Corollary~3.6(b)]{CB_Dette} is exactly
\eqref{E:cba:dettebound}. Its stipulated range $m<t$ is respected here.
Thus this inequality itself is prior art. For the final assertion,
using the positive series for $\operatorname{arctanh}$ gives
\begin{align*}
 \log D(t,m)
 &=2\sum_{l=1}^{m-1}\operatorname{arctanh}(l/t)\\
 &\le \frac{m(m-1)}t+
 \frac{m^2(m-1)^2}
 {6t^3\{1-(m-1)^2/t^2\}}.
\end{align*}
The second term is at most $m/t$ precisely under
\eqref{E:cba:integercriterion}.
\end{proof}

\begin{proof}[Proof of Theorem~\ref{E:cba:constantone}]
The case $m=0$ is immediate. Suppose first $t\ge512$.
If $m^2/t\ge\frac{10}{3}\log(t+1)$, the contour bound proves the
assertion because $\log2<7/10$. Otherwise,
\[
 m^2<\frac{10}{3}t\log(t+1)\le\frac{25}{9}t^{4/3},
 \qquad m<\frac53t^{2/3}.
\]
Here $\log(t+1)\le(5/6)t^{1/3}$ for $t\ge512$: at $512$ it follows
from $513<729$, $\log3<11/10$, and $33/5<20/3$; the difference
$(5/6)t^{1/3}-\log(t+1)$ has positive derivative thereafter.
In particular $m<t$, and
\[
 (m+6)(m-1)^2\le m^3+6m^2
 \le\left(\frac{125}{27}+\frac{25}{96}\right)t^2<6t^2.
\]
The last estimate uses $t^{2/3}\ge64$. The Hahn bound therefore applies.
The elementary logarithmic inequalities just used can, for example, be
verified by positive Taylor lower bounds for $e^{7/10}$ and $e^{11/10}$.

For $1\le t<512$, the following verifier covers every pair $(t,m)$.
Its first branch uses \eqref{E:cba:integercriterion}; its second uses
$2^m\le e^{7m/10}\le e^{m^2/t}$; its third uses the contour bound with
$\log2<7/10$. Every other pair is checked with the exact identity
\begin{equation}\label{E:cba:exactcoeff}
 C(t,m,j)=\frac1{\binom tj}
 \sum_{i=\max(0,j-t+m)}^{\min(m,j)}
 (-1)^{m-i}\binom{2m}{2i}\binom{t-m}{j-i}.
\end{equation}
The function \texttt{exp\_lower} returns a rational number no larger
than $e^x$ for $x\ge0$, so each successful assertion is a rigorous
inequality in the required direction. There are no floating-point
comparisons. The complete output is
\[
 (34784,39518,56329,185),\qquad
 12734\ \hbox{direct coefficient comparisons},\quad
 \max t_{\rm direct}=145.
\]
The four counts sum to $511\cdot512/2$, as required.
\end{proof}

\noindent\begin{minipage}{\textwidth}\small
\begin{verbatim}
from fractions import Fraction
from math import comb

def exp_lower(x, degree=32):
    term = total = Fraction(1)
    for k in range(1, degree + 1):
        term *= x / k
        total += term
    return total

counts = [0, 0, 0, 0]
comparisons = largest_direct = 0
for t in range(1, 512):
    for m in range(1, t + 1):
        if m < t and (m+6)*(m-1)**2 <= 6*t*t:
            counts[0] += 1
            continue
        if 10*m >= 7*t:
            counts[1] += 1
            continue
        if exp_lower(Fraction(3*m*m, 10*t)) >= t+1:
            counts[2] += 1
            continue
        lower = exp_lower(Fraction(m*m, t))
        for j in range(t + 1):
            numerator = sum(
                (-1)**(m-i)*comb(2*m, 2*i)*comb(t-m, j-i)
                for i in range(max(0,j-t+m), min(m,j)+1)
            )
            assert Fraction(abs(numerator), comb(t,j)) <= lower
            comparisons += 1
        counts[3] += 1
        largest_direct = max(largest_direct, t)
print(counts, comparisons, largest_direct)
assert sum(counts) == 511*512//2
\end{verbatim}
\end{minipage}

\subsection{Krawtchouk identities and the energy refinement}

The following identities also justify the weighted norm used in the next
section. Define
\[
 K_k(x;N)=\sum_h(-1)^h\binom xh\binom{N-x}{k-h},\qquad
 \sum_{k=0}^NK_k(x;N)u^k=(1-u)^x(1+u)^{N-x}.
\]
\begin{lemma}\label{E:cba:kraw}
For $0\le m,j\le t$,
\[
 \binom tjC(t,m,j)=(-1)^{m+j}K_{2j}(t+m;2t).
\]
Moreover
\[
 \sum_{x=0}^N\binom NxK_k(x;N)K_\ell(x;N)
                       =2^N\binom Nk\mathbf1_{k=\ell}.
\]
For fixed $t,j$, write $y=2j/t-1$ and $C_m=C(t,m,j)$. Then
\[
 C_0=1,\quad C_1=y,\qquad
 (t-m)C_{m+1}=2tyC_m-(t+m)C_{m-1}
                           \quad(1\le m\le t-1).
\]
\end{lemma}
\begin{proof}
In the coefficient generating polynomial $F(z)$ above, put $z=w^2$.
The even part of $(1+iw)^{t+m}(1-iw)^{t-m}$ is
\[
 \frac12\{(1+iw)^{t+m}(1-iw)^{t-m}
          +(1-iw)^{t+m}(1+iw)^{t-m}\}.
\]
Multiplication by $(-1)^m$ gives $F(z)$. Comparing its even powers
with the Krawtchouk generating function, and using
$K_k(N-x;N)=(-1)^kK_k(x;N)$, gives the first identity.
Multiplying two generating functions and summing with weight
$\binom Nx$ gives $(2+2uv)^N$; comparison of coefficients gives
orthogonality. Finally, the same generating function verifies
\[
 (N-x)K_k(x+1;N)+xK_k(x-1;N)=(N-2k)K_k(x;N).
\]
Take $N=2t$, $k=2j$, $x=t+m$ and insert the signs in the first
identity to obtain the recurrence.
\end{proof}

A one-term consequence of this orthogonality identity is the
squared bound
\[
 C(t,m,j)^2\le
       \frac{4^t\binom{2t}{2j}}
            {\binom{2t}{t+m}\binom tj^2}.
\]
It is not needed by the rational-certificate proof above.

\begin{proposition}[An energy bound]
\label{E:cba:energy}
For fixed $t\ge2,j$, set $y=2j/t-1$, $C_m=C(t,m,j)$, and
\[
 \alpha_0=1,\qquad
 \alpha_{m+1}=\alpha_m\frac{(t-m)^2}{(t+m)(t-m+1)}.
\]
For $1\le m\le t$ define
\[
 Q_m=\alpha_m C_m^2+\alpha_{m-1}C_{m-1}^2
       -\frac{2t\alpha_m}{t-m+1}yC_mC_{m-1},\qquad
 \kappa_m=\frac{t^2}{(t+m-1)(t-m+2)}.
\]
Then, for $1\le m<t$,
\[
 Q_{m+1}-Q_m=
 -\frac{2(m-1)\alpha_{m-1}}{(t+m-1)(t-m+2)}C_{m-1}^2\le0,
 \qquad Q_1=1-\frac{t}{t+1}y^2.
\]
If $\kappa_my^2<1$, this yields
\[
 C(t,m,j)^2\le D(t,m)\frac{t+1}{t-m+1}
                \frac{1-ty^2/(t+1)}{1-\kappa_my^2},
\]
where $D(t,m)=\prod_{i=1}^{m-1}(t+i)/(t-i)$ also at $m=t$.
In particular, for each fixed $0<\delta<1$, uniformly for
$|y|\le1-\delta$ and $m\le\delta t/2$,
\[
 |C(t,m,j)|\le K_\delta
       \exp\{m(m-1)/(2t)+O(m^4/t^3)\}.
\]
\end{proposition}
\begin{proof}
Substitute the three-term recurrence into $Q_{m+1}$ and cancel its
$y^2C_m^2$ terms. The result is
\[
 Q_{m+1}=\alpha_m C_m^2
       -\frac{2t\alpha_m}{t-m+1}yC_mC_{m-1}
       +\alpha_m\frac{t+m}{t-m+1}C_{m-1}^2.
\]
Subtracting $Q_m$ and using the ratio of consecutive $\alpha$'s
gives the displayed difference. Completing the square gives
\[
 Q_m=\alpha_{m-1}
 \left(C_{m-1}-\frac{t\alpha_m y}{(t-m+1)\alpha_{m-1}}C_m\right)^2
       +\alpha_m C_m^2(1-\kappa_my^2).
\]
Telescoping the defining product gives
$\alpha_m^{-1}=D(t,m)(t+1)/(t-m+1)$.
These identities and $Q_m\le Q_1$ prove the bound.
In the stated bulk range,
$(t+m-1)(t-m+2)\ge t^2-m^2$, so
$\kappa_my^2\le(1-\delta)^2/(1-\delta^2/4)<1$ uniformly.
The remaining rational factors are bounded, and the positive
series for $\operatorname{artanh}$ gives
$\log D(t,m)=m(m-1)/t+O(m^4/t^3)$ for $m\le t/2$.
\end{proof}



\clearpage

\section{The optimal exponent with a universal prefactor}
\label{E:cbx:sec}

Let $M_{t,m}=\max_{0\le j\le t}|C(t,m,j)|$ and put
\[
 H(a)=\tfrac12\{(1+a)\log(1+a)+(1-a)\log(1-a)\},
 \qquad 0\le a\le1,
\]
where $0\log0=0$. The following improves the exponential constant in
Theorem~\ref{E:cba:constantone} when a universal multiplicative prefactor
is allowed. The distinction between that prefactor and prefactor one
is essential.

\begin{theorem}[A sharp exponential constant]\label{E:cbx:main}
There is an absolute constant $K$ such that, for all
$t\ge1$, $0\le m\le t$ and $0\le j\le t$,
\begin{align}
 |C(t,m,j)|&\le K(1+m^2/t)^{1/4}\exp\{tH(m/t)\},
                                      \label{E:cbx:profileupper}\\
 |C(t,m,j)|&\le K\,2^{m^2/t}.           \label{E:cbx:optimalupper}
\end{align}
The exponent $\log2$ in the second inequality is optimal among
inequalities $|C(t,m,j)|\le K_c\exp(c m^2/t)$ with a finite constant
$K_c$ independent of $t,m,j$. This does not prove the same inequality
with prefactor exactly one.
\end{theorem}

The improvement comes from integrating the coefficient contour,
rather than replacing its integrand by its supremum. The resulting
prefactor depends on $m^2/t$, instead of on $t$ alone.

\begin{lemma}[A Gaussian integral against the arcsine law]
\label{E:cbx:arcsine}
For every $A,B\ge0$,
\[
 (1+B)\int_0^{\pi/2}
       \exp\{-\tfrac12(B\sin v-A)^2\}\,dv
       \le K\sqrt{1+A}.
\]
\end{lemma}
\begin{proof}
The case $B\le1$ is immediate. Suppose $B>1$. The substitution
$y=B\sin v$ gives the measure $dy/\sqrt{B^2-y^2}$ on $[0,B]$.
If $B\le2(A+1)$, use
$\sqrt{B^2-y^2}\ge\sqrt{B(B-y)}$ to bound the integral by
\[
 B^{-1/2}\int_0^B e^{-(y-A)^2/2}(B-y)^{-1/2}\,dy
 \le B^{-1/2}(2+\sqrt{2\pi}).
\]
For the last inequality put $u=B-y$, split at $u=1$, and use
$u^{-1/2}\le1$ on the remaining interval. Multiplication by
$1+B\le2B$ proves the claim in this case.

If $B>2(A+1)$, the part $0\le y\le3B/4$ is at most $K/B$,
because the denominator is at least $\sqrt7 B/4$ and the Gaussian
integral is bounded. On $3B/4\le y\le B$ we have $y-A>B/4$;
this part is at most $(\pi/2)e^{-B^2/32}$. Multiplying by
$1+B$ gives a uniform constant, which is at most
$K\sqrt{1+A}$.
\end{proof}

\begin{proof}[Proof of \eqref{E:cbx:profileupper}]
Endpoint coefficients have absolute value one. For $0<j<t$, put
$p=j/t$, $r=p/(1-p)$, $v=p(1-p)$, $a=m/t$, and
\[
 F(z)=(1+z)^tT_m\left(\frac{z-1}{z+1}\right)
      =\sum_{j=0}^t\binom tjC(t,m,j)z^j.
\]
For $z=re^{i\theta}$ set $u=2\sqrt v\,|\sin(\theta/2)|$.
The inverse Joukowski formula and $|T_m(w)|\le|\zeta|^m$ when
$w=(\zeta+\zeta^{-1})/2$, $|\zeta|\ge1$, give
\[
 \frac{|F(re^{i\theta})|}{(1+r)^t}
 \le \exp\{t f_a(u)\},\qquad
 f_a(u)=\tfrac12\log(1-u^2)+a\operatorname{artanh}u.
\]
Explicitly, one may take
$\zeta=(\sqrt z+i)/(\sqrt z-i)$ or its reciprocal; choosing
the larger modulus gives $|\zeta|=\sqrt{(1+u)/(1-u)}$.
Limits cover the removable endpoint configurations. Since
$f_a'(u)=(a-u)/(1-u^2)$ and $f_a(a)=H(a)$,
integration of this derivative on the interval between $u$ and $a$
shows
\[
                 f_a(u)\le H(a)-\tfrac12(u-a)^2.             \tag{*}
\]
This uses $1/(1-w^2)\ge1$ and also holds at $a=1$ by a limit.

The binomial mass $\pi_j=\binom tjp^j(1-p)^{t-j}$ is a mode of
$\operatorname{Bin}(t,p)$. Chebyshev's probability inequality on
the interval of radius $2\sqrt{tv}$ around its mean $j$ implies
\[
 \pi_j\ge \frac{3/4}{4\sqrt{tv}+1}
                     \ge\frac{c}{1+\sqrt{tv}}.
\]
Cauchy's coefficient formula, (*) and Lemma~\ref{E:cbx:arcsine}, with
$A=m/\sqrt t$ and $B=2\sqrt{tv}$, therefore yield
\[
 |C(t,m,j)|\le K e^{tH(a)}\sqrt{1+m/\sqrt t}
      \le K(1+m^2/t)^{1/4}e^{tH(a)}.
\]
The final inequality follows from $(1+u)^2\le2(1+u^2)$.
\end{proof}

\begin{proof}[Proof of \eqref{E:cbx:optimalupper} and optimality]
Write $x=m^2/t$ and
$\Delta(a)=(\log2)a^2-H(a)$. The positive power series
\[
 H(a)=\sum_{h\ge1}\frac{a^{2h}}{2h(2h-1)}
\]
shows that $H(a)/a^2$ increases strictly on $(0,1)$ and tends to
$\log2$ at one. Hence for $0\le a\le99/100$ there is an absolute
$c_0>0$ such that $t\Delta(a)\ge c_0x$. The factor
$(1+x)^{1/4}e^{-c_0x}$ is bounded, proving the assertion in this range.

Suppose $99/100\le a<1$ and write $\delta=1-a$, $k=t-m$.
Then $k$ is an integer with $1\le k\le t/100$. Since
$H'(a)=\operatorname{artanh}a$,
\begin{align*}
 \Delta(1-\delta)
 &\ge -2(\log2)\delta
            +\tfrac12\int_0^\delta\log(1/u)\,du\\
 &=\delta\{\tfrac12\log(1/\delta)+\tfrac12-2\log2\}
 \ge\tfrac14\delta\log(1/\delta).
\end{align*}
The last inequality follows from $\delta\le1/100$; for example
$\log100>4$ and $\log2<7/10$ suffice. The function
$k\mapsto k\log(t/k)$ increases for $1\le k\le t/100$, so
\[
 t\Delta(a)\ge\tfrac14 k\log(t/k)\ge\tfrac14\log t.
\]
Since $x\le t$, this absorbs the prefactor in
\eqref{E:cbx:profileupper}. At $m=t$, the elementary degree bound
$|C(t,t,j)|\le2^t$ proves the result directly.

Finally, for even $t$ the exact central coefficient is
\[
 |C(t,t,t/2)|=\frac{\binom{2t}{t}}{\binom t{t/2}}
                 =2^{t-1/2}(1+o(1)).
\]
Thus no fixed prefactor permits an exponent $c<\log2$.
\end{proof}

\subsection{Exact weighted energy and matching exponential profiles}

The next identity is the classical Hahn orthogonality in this
normalization. We prove it directly from the Krawtchouk identities;
its sharp-profile consequences are the additional conclusions needed here.
Put
\[
 w_{t,j}=4^{-t}\binom{2j}{j}\binom{2t-2j}{t-j},\qquad
 A_{t,m}=\frac{(t+m)!(t-m)!}{(t!)^2}.
\]

\begin{proposition}[Exact weighted inverse norm]\label{E:cbx:parseval}
The weights $w_{t,j}$ sum to one. If
$P(x)=\sum_{m=0}^t a_mT_m(2x-1)=\sum_{j=0}^t b_jB_j^t(x)$, with
real coefficients, then
\[
 \sum_{j=0}^t w_{t,j}b_j^2
       =a_0^2+\frac12\sum_{m=1}^t A_{t,m}a_m^2.
\]
In particular, for $m\ge1$,
\begin{equation}\label{E:cbx:lower}
                      M_{t,m}\ge\sqrt{A_{t,m}/2}.
\end{equation}
\end{proposition}
\begin{proof}
The generating series of $4^{-j}\binom{2j}{j}$ is
$(1-z)^{-1/2}$, so convolution proves the assertion about the weights.
Write $K_k(x;N)$ for the Krawtchouk polynomials defined in the preceding
section. Their combinatorial formula gives the duality
\[
 \frac{K_k(x;N)}{\binom Nk}=
 \frac{K_x(k;N)}{\binom Nx}.
\]
Indeed, multiplying either side by $\binom Nx\binom Nk$ counts
ordered subsets of sizes $x,k$, with the sign of their intersection.
Together with the orthogonality identity from that section this gives
\[
 \sum_{k=0}^N\frac{K_k(x;N)K_k(y;N)}{\binom Nk}
                   =\frac{2^N}{\binom Nx}\mathbf1_{x=y}.
\]
Using $K_k(N-y;N)=(-1)^kK_k(y;N)$, its even-$k$ part is half the
sum of the right sides for $y$ and $N-y$.

Take $N=2t$, $x=t+m$, $y=t+\ell$, and use
$\binom tjC(t,m,j)=(-1)^{m+j}K_{2j}(t+m;2t)$ and
\[
 \frac{w_{t,j}}{\binom tj^2}
      =\frac{(2t)!}{4^t(t!)^2}\frac1{\binom{2t}{2j}}.
\]
For $0\le m,\ell\le t$, the resulting weighted inner product is
zero unless $m=\ell$. Its value is one for $m=\ell=0$ and
$A_{t,m}/2$ otherwise. Expanding $b_j=\sum_m a_mC(t,m,j)$ proves
the identity. Averaging $C(t,m,j)^2\le M_{t,m}^2$ proves the lower bound.
\end{proof}

\begin{corollary}[Sharp proportional-degree and moderate-degree scales]
\label{E:cbx:asymptotic}
If $t\to\infty$ and $m/t\to a\in[0,1]$, then
\[
                    \frac1t\log M_{t,m}\longrightarrow H(a).
\]
If instead $m=o(t)$ and $m^2/t\to\infty$, then
\[
                    \log M_{t,m}=(\tfrac12+o(1))m^2/t.
\]
The latter maximum is over every coefficient, without a restriction
keeping $j/t$ away from zero and one.
\end{corollary}
\begin{proof}
The elementary factorial estimate
$\log(n!)=n\log n-n+O(\log(n+1))$, valid also at $n=0$ with
$0\log0=0$, gives
$\log A_{t,m}=2tH(m/t)+O(\log(t+1))$ uniformly.
Combine \eqref{E:cbx:lower} with \eqref{E:cbx:profileupper}; the case
$m=0$ is immediate. For $m=o(t)$, the product
\[
 A_{t,m}=\prod_{i=1}^m\frac{t+i}{t-i+1}
\]
and $\log(1+u)=u+O(u^2)$ give
$\log A_{t,m}=m^2/t+O(m^3/t^2)$.
Also $tH(m/t)=m^2/(2t)+O(m^4/t^3)$.
The logarithm of the prefactor in \eqref{E:cbx:profileupper} is
$O(1+\log(1+m^2/t))=o(m^2/t)$, proving the second assertion.
\end{proof}

\begin{remark}[What remains open in this argument]
The conjecturally stronger statement $M_{t,m}\le2^{m^2/t}$ with
prefactor one is not proved here. Conversely, a global exponent
$1/2$ with either a fixed or polynomial prefactor is impossible:
the degree-$t$ central coefficient grows as $2^{t-1/2}$.
The exponent $1/2$ is the sharp moderate-degree limit, whereas
$\log2$ is the sharp universal exponent when a fixed prefactor is allowed.
\end{remark}

\section{An entropy--transport inequality for mixture likelihood}
\label{E:et:section}

Let $\mathcal X$ be a countable metric space, let $K_\theta$ be component
probability kernels, and write $p_H=\int K_\theta\,dH(\theta)$.
For an empirical law $q$, assume the exact mixing NPMLE $H$ exists
and has positive probability on the observed support. Zero empirical
cells are omitted in score sums.

\begin{theorem}[Component concentration transfers through likelihood contact]
\label{E:et:transfer}
Suppose that $\psi:[0,\infty)\to[0,\infty)$ is nondecreasing and
concave and that
\[
                 W_1(R,K_\theta)\le\psi(D(R\Vert K_\theta))
\]
for every component and every probability law with finite entropy.
Then
\[
             W_1(q,p_H)\le\psi(D(q\Vert p_H))
                       \le\psi(D(q\Vert p_G))
\]
for every admissible comparator mixing law $G$ for which the final
entropy is finite.
\end{theorem}
\begin{proof}
The directional derivative toward a point mass gives
$S(\theta)=\sum_xq(x)K_\theta(x)/p_H(x)\le1$.
Its $H$-integral equals one, so $S=1$ almost surely under $H$.
Thus $R_\theta(x)=q(x)K_\theta(x)/p_H(x)$ is a probability kernel
$H$-almost surely, and $\int R_\theta\,dH=q$. The entropy chain rule gives
\[
 \int D(R_\theta\Vert K_\theta)\,dH(\theta)=D(q\Vert p_H).
\]
Mixing couplings, applying the component bound, and then Jensen's
inequality prove the first conclusion. Likelihood optimality gives
$D(q\Vert p_H)\le D(q\Vert p_G)$, proving the second.
All applications below have finite right sides; truncation of transport
costs gives the same conclusion before passing to an unbounded cost.
\end{proof}

For normalized binomial counts, the component bound has $\psi(u)=\sqrt{u/(2t)}$ for normalized count cost.
For Poisson mixtures, a parameter-free occupancy channel will import
that same inequality while preserving the original likelihood fit.

The argument does not normalize every stratum of an arbitrary
heterogeneous joint likelihood: only the aggregate score is normalized.
Consequently an arbitrary-profile theorem does not follow verbatim.
The results proved below concern common known trial counts or common
known exposures.



\clearpage

\section{Closure of the Poisson critical window}
\label{E:pc:sec}

Consider the experiment
\[
 p_G(j)=\int_0^B e^{-\lambda}\frac{\lambda^j}{j!}\,dG(\lambda),
 \qquad X_1,\ldots,X_N\overset{\mathrm{iid}}\sim p_G,
\]
where the support radius $B>0$ is known.  Let $\widehat G$ be any
measurable maximizer of the original Poisson mixture likelihood over
all probability laws on $[0,B]$.  There is no restriction on the
number of mixing atoms.  Set $L=\log(eN)$ and, for $N\ge1$, define
\begin{equation}\label{E:pc:rate}
 r(N,B)=
 \begin{cases}
 B\min\{1,L^{-1}\log(e+L/B)\},&0<B\le L,\\[2mm]
 \max\{\sqrt{B/L},B/\sqrt N\},&B>L.
 \end{cases}
\end{equation}
The two expressions at the breakpoint agree up to absolute constants.

\begin{theorem}[Global Poisson population recovery]\label{E:pc:main}
There exist absolute constants $0<c<C<\infty$ such that, for every
$N\ge1$ and every $B>0$,
\[
 c\,r(N,B)\le
 \inf_{\widetilde G}\sup_{G\in\mathcal P([0,B])}
       \mathbb E_G W_1(\widetilde G,G)
 \le \sup_{G\in\mathcal P([0,B])}
       \mathbb E_G W_1(\widehat G,G)
 \le C\,r(N,B).
\]
In particular, the ordinary NPMLE is minimax up to absolute factors
throughout the entire support/sample-size plane.  For observations
$X_i\mid\Theta_i\sim\operatorname{Poi}(t\Theta_i)$ with
$G\in\mathcal P([0,M])$, the corresponding rate is $r(N,tM)/t$.
\end{theorem}

The growing-support crossover in Theorem~\ref{E:pc:main} is
$B\asymp N/\log(eN)$: below that scale the deconvolution term
$\sqrt{B/L}$ dominates, and above it the latent-sampling term
$B/\sqrt N$ dominates. The statement uses the ordinary likelihood
estimator throughout, with no logarithmic loss in this transition.

Miao et al.\ \cite[Theorem 6.1(a)]{PC_Miao2021} established the
small-support rate for the NPMLE, and their Theorem 6.1(b) gives
$\sqrt{B/\log N}$ under $B\le C_0N^{1/3-\varepsilon_0}$.
The rate includes diameter clipping for arbitrarily small $B$:
an unclipped formula cannot be extrapolated beyond the maximum
possible loss $B$.
All ingredients used below are proved here, including the two-scale
binomial inverse inequality in Proposition~\ref{E:pc:bc:twoscale}.
The estimator is never replaced by a transformed-data MLE.

For $N\in\{1,2\}$ the rate in \eqref{E:pc:rate} is comparable to $B$.
The diameter bound and the two-point latent-sampling lower bound given
below therefore suffice. In the proof of the nontrivial regimes we
may assume $N\ge3$.

\subsection{A parameter-free occupancy channel}

For an integer $m\ge1$, define a Markov kernel $A_m$ as follows.  Given
a count $x$, independently assign $x$ balls uniformly to $m$ boxes,
and output the number of occupied boxes.  Poisson splitting proves
\begin{equation}\label{E:pc:occupancy}
 A_m\operatorname{Poi}(\lambda)
 =\operatorname{Bin}(m,\phi_m(\lambda)),\qquad
 \phi_m(\lambda)=1-e^{-\lambda/m}.
\end{equation}
Indeed the box counts are independent Poisson random variables of
mean $\lambda/m$, so their nonempty indicators are independent
Bernoulli variables.  In particular, if $F=\phi_m\#G$, then
$A_mp_G$ is exactly the binomial mixture with mixing law $F$.
On $[0,\phi_m(B)]$ the inverse map has derivative at most
$m e^{B/m}$, and consequently
\begin{equation}\label{E:pc:backmap}
 W_1(G,H)\le m e^{B/m}
                 W_1(\phi_m\#G,\phi_m\#H).
\end{equation}
This channel is a proof device.  We do not transform the data before
fitting, or assert that a transformed MLE remains an MLE.

We use $W_{1,m}$ for transportation on $\{0,\ldots,m\}$ with cost
$|i-j|/m$, and $h(p,q)^2=\sum_j(\sqrt{p_j}-\sqrt{q_j})^2$.

\begin{lemma}[Empirical entropy on the infinite alphabet]\label{E:pc:kl}
Let $q$ be the empirical law of $N$ observations from $p_G$.
Uniformly over $G\in\mathcal P([0,B])$,
\begin{equation}\label{E:pc:klbound}
                    \mathbb E D(q\Vert p_G)\le C(B+L)/N.
\end{equation}
\end{lemma}
\begin{proof}
For every coordinate with probability $p_j>0$, convexity and
$\log u\le u-1$ give
\[
 0\le\mathbb E\{q_j\log(q_j/p_j)-q_j+p_j\}
 \le\min\{(1-p_j)/N,\ p_j\log(1/p_j)\}.
\]
The second bound also uses $q_j\log q_j\le0$.  The expression inside
braces is nonnegative, so its infinite sum can be averaged by
Tonelli's theorem.  Coordinates with $p_j=0$ contribute zero almost
surely.  For $j\ge B$, the Poisson mass is increasing in its mean
throughout $[0,B]$, and hence
\[
 p_G(j)\le e^{-B}B^j/j!\le(eB/j)^j\le e^{-j}
                         \quad(j\ge e^2B).
\]
Take $J=\lceil e^2B+3L\rceil$.  The first $J+1$ coordinates
contribute at most $(J+1)/N$.  Since $u\log(1/u)$ is increasing
on $(0,e^{-1}]$, the remaining contribution is bounded by
$\sum_{j>J}j e^{-j}\le C(J+1)e^{-J}\le C/N$.
This proves \eqref{E:pc:klbound}.  Notice that the expected empirical
chi-square divergence on the full infinite alphabet would not give
this conclusion.
\end{proof}

\begin{lemma}[Root-$N$ prediction after occupancy]\label{E:pc:transport}
For $m=\lceil\max(B,L)\rceil$ and every Poisson NPMLE $\widehat G$,
\begin{equation}\label{E:pc:transportbound}
       \mathbb E W_{1,m}(A_mp_{\widehat G},A_mp_G)\le C/\sqrt N.
\end{equation}
\end{lemma}
\begin{proof}
Existence follows from compactness of the mixing parameter space
and upper semicontinuity of the likelihood; the admissible fit
$\delta_B$ is strictly positive at every observed count.  Write
$H=\widehat G$.  The likelihood directional derivative gives
\[
 S(\lambda)=\sum_{j:q_j>0}\frac{q_jK_\lambda(j)}{p_H(j)}\le1
                    \quad(0\le\lambda\le B).
\]
Since $\int S\,dH=1$, equality holds $H$-almost everywhere.
Thus $R_\lambda(j)=q_jK_\lambda(j)/p_H(j)$ is a probability kernel
$H$-almost everywhere and $\int R_\lambda\,dH=q$.  Its entropy
chain rule and the likelihood comparison yield
\begin{equation}\label{E:pc:posteriorchain}
 \int D(R_\lambda\Vert K_\lambda)\,dH(\lambda)
       =D(q\Vert p_H)\le D(q\Vert p_G).
\end{equation}
The binomial transport-entropy inequality is
\[
 W_{1,m}(R,\operatorname{Bin}(m,p))^2
          \le D(R\Vert\operatorname{Bin}(m,p))/(2m).
\]
For completeness, a normalized-count 1-Lipschitz function changes
by at most $1/m$ on changing one of the $m$ Bernoulli bits.
The conditional Hoeffding bound gives centered log-mgf at most
$u^2/(8m)$; entropy duality and Wasserstein duality give the
displayed inequality, including degenerate endpoints whenever the
entropy is finite.  Apply it to $A_mR_\lambda$, use entropy data
processing, and integrate \eqref{E:pc:posteriorchain}.  Convexity of
transportation under mixtures and Cauchy--Schwarz give
\[
 W_{1,m}(A_mq,A_mp_H)
                  \le\sqrt{D(q\Vert p_G)/(2m)}.
\]
For each count threshold, the CDF of $A_mq$ is an average of
independent $[0,1]$-valued functions of the $X_i$.  Its expected
absolute error is at most $1/(2\sqrt N)$.  The CDF formula for
$W_{1,m}$ therefore gives
$\mathbb E W_{1,m}(A_mq,A_mp_G)\le1/(2\sqrt N)$.
Now use Lemma~\ref{E:pc:kl} and $(B+L)/m\le2$.
\end{proof}

\subsection{A growing-support Hellinger bound}

\begin{lemma}[Poisson mixture entropy and likelihood control]\label{E:pc:hellinger}
For every fixed absolute $b_0>0$, uniformly over
$b_0L\le B\le N$,
\begin{equation}\label{E:pc:hellingerbound}
       \mathbb E h(p_{\widehat G},p_G)
                       \le C_{b_0}B^{1/4}L^{3/4}/\sqrt N.
\end{equation}
\end{lemma}
\begin{proof}
For complex $z$ the vector $K_z=(e^{-z}z^j/j!)_{j\ge0}$ belongs
to $\ell_1$ and satisfies
$\|K_z\|_1=e^{|z|-\operatorname{Re}z}$.  On the ellipse
$z=B(1+\cos(\theta+iu))/2$, direct calculation gives
\[
 |z|-\operatorname{Re}z
   =\tfrac B2(\cosh u-1)(1-\cos\theta)
   \le B(\cosh u-1).
\]
The $k$th vector Chebyshev coefficient $v_k$ of
$x\mapsto K_{Bx}$ therefore satisfies, by contour shifting,
\[
 \|v_k\|_1\le2\exp\{B(\cosh u-1)-ku\}
       \le2\exp\{-c\min(k^2/B,k)\}.
\]
For the second inequality take $u=k/(2B)$ if $k\le2B$, using
$\cosh u-1\le u^2$ for $0\le u\le1$, and take $u=1$ otherwise.
It follows by summing the Gaussian and exponential tails that
truncating the expansion at
\[
 d=\left\lceil C\left\{
 \sqrt{B\log(C(B+1)/\beta)}+
                   \log(C(B+1)/\beta)\right\}\right\rceil
\]
has uniform $\ell_1$ error at most $\beta/8$.
Every Chebyshev moment of a probability law on $[0,B]$ belongs to
$[-1,1]$, and $\|v_k\|_1\le2$ also follows directly from its real
integral formula.  Grid the first $d$ moments with mesh
$\beta/(16(d+1))$.  Choose one actual mixture in each occupied grid
box.  This gives an $\ell_1$ cover of radius $\beta$, with centers
in the mixture class and cardinality $M$ satisfying
\begin{equation}\label{E:pc:cover}
                       \log M\le Cd\log(C(d+1)/\beta).
\end{equation}

Take $J=\lceil e^2B+6L\rceil$ and $\beta=N^{-3}/(J+1)$.
In the indicated range, $d\le C_{b_0}\sqrt{BL}$ and
$\log M\le C_{b_0}\sqrt B L^{3/2}$.  For every cover center $p_a$,
define the pointwise envelope
\[
 u_a(j)=\begin{cases}p_a(j)+\beta,&j\le J,\\ e^{-j},&j>J.
                     \end{cases}
\]
Every mixture is bounded coordinatewise by its associated envelope,
whose total mass $U_a$ lies in $[1,1+\eta]$ for
$\eta=CN^{-3}$, whenever that envelope is used for at least one
mixture.  Unused envelopes can be discarded.  Normalize
$r_a=u_a/U_a$.  An associated mixture $p$ satisfies
$h(p,r_a)\le\sqrt{2\eta}$.

Here is the elementary likelihood argument, to specify the treatment
of the infinite alphabet.  If $\widehat p$ is a likelihood maximizer
and $u_a\ge\widehat p$, then
\[
 \prod_{i=1}^N\frac{r_a(X_i)}{p_G(X_i)}\ge e^{-N\eta}.
\]
Under $p_G$, the expectation of the square root of this product is
the $N$th power of the affinity, at most
$\exp\{-Nh(p_G,r_a)^2/2\}$.  Markov's inequality followed by a
union bound shows, with probability at least $1-e^{-x}$,
\[
 h(p_G,\widehat p)\le\sqrt{2\eta}
       +\sqrt{\eta+2(\log M+x)/N}.
\]
Integration gives
$\mathbb E h(p_G,\widehat p)^2\le C\eta+C(\log M+1)/N$,
which proves \eqref{E:pc:hellingerbound}.
\end{proof}

\subsection{The two-scale argument closes every growing-support regime}

The binomial inverse inequality proved below states that if $F,H$ are
mixing laws on $[0,1]$, $b_F,b_H$ their $m$-trial binomial mixtures,
and $1\le S\le K\le m/2$, then
\begin{equation}\label{E:pc:bininverse}
 W_1(F,H)\le C/K+Ce^{4S^2/m}W_{1,m}(b_F,b_H)
                      +Ce^{2K^2/m}h(b_F,b_H)/S.
\end{equation}
Its small high-frequency amplitude $1/S$ is indispensable.
The proof approximates a Lipschitz test function at two polynomial
degrees, preserves a bounded derivative at the lower degree, and
uses the degree-elevated Bernstein coefficient bounds.  This is
Proposition~\ref{E:pc:bc:twoscale}.

Suppose $b_0L\le B\le N$, with $b_0>0$ an absolute constant to
be fixed below.  Set $m=\lceil\max(B,L)\rceil$, so $m\asymp_{b_0}B$.
For sufficiently large absolute $N$, take
\[
 S=\lfloor\sqrt m/16\rfloor,\qquad
 K=\left\lfloor\min\{m/2,\sqrt{a mL}\}\right\rfloor,
 \qquad a=1/64.
\]
These are admissible, $S\asymp\sqrt m$, $e^{4S^2/m}\le C$, and
$m/K\le C\sqrt{m/L}$ because $m\ge L$.
Apply \eqref{E:pc:bininverse} to the pushforwards under $\phi_m$,
then use \eqref{E:pc:backmap}, Lemmas~\ref{E:pc:transport} and
\ref{E:pc:hellinger}, and Hellinger data processing.  This gives
\[
 \mathbb E W_1(\widehat G,G)
 \le C_{b_0}\left\{
 \sqrt{B/L}+\frac B{\sqrt N}
       +\frac{e^{2aL}B^{3/4}L^{3/4}}{\sqrt N}\right\}.
\]
The ratio of the last term to $\sqrt{B/L}$ is at most
$C e^{2aL}N^{-1/4}L^{5/4}$, which is uniformly bounded since
$B\le N$ and $2a<1/4$.  We have proved
\begin{equation}\label{E:pc:medium}
 \mathbb E W_1(\widehat G,G)
                  \le C_{b_0}\{\sqrt{B/L}+B/\sqrt N\}.
\end{equation}
Bounded $N$ is absorbed in the absolute constant by the diameter
bound; this is uniform in $B$ because $B\ge b_0L$ in this step.

If $B>N$, no Hellinger estimate or approximation is necessary.
For any binomial mixing law, coupling its parameter $p$ with its
normalized count gives expected cost at most $1/(2\sqrt m)$.
Thus
\[
 W_1(F,H)\le m^{-1/2}+W_{1,m}(b_F,b_H).
\]
Together with \eqref{E:pc:backmap} and Lemma~\ref{E:pc:transport}, this
gives $\mathbb E W_1(\widehat G,G)\le C(\sqrt B+B/\sqrt N)
\le CB/\sqrt N$.  Every $B\ge L$ is now covered.

\subsection{The small-support regime, including arbitrarily small radii}

We include a direct proof to ensure uniformity as $B\downarrow0$.
Let $f$ be 1-Lipschitz on $[0,B]$ and normalize $f(0)=0$.
Jackson approximation gives a degree-$k$ polynomial
$P(\lambda)=\sum_{l=0}^k a_l\lambda^l$ such that
\[
 \|f-P\|_{[0,B]}\le CB/k,\qquad \|P\|_{[0,B]}\le CB.
\]
Its monomial coefficients satisfy
\begin{equation}\label{E:pc:monomial}
                         |a_l|\le BC_1^kB^{-l}.
\end{equation}
Indeed expand $P(Bx)$ in shifted Chebyshev polynomials: each
Chebyshev coefficient has magnitude at most $2\|P\|$, and the sum
of absolute monomial coefficients of $T_j(2x-1)$ is at most $7^j$
by its three-term recurrence.  Summing the finite geometric series
proves \eqref{E:pc:monomial} after enlarging $C_1$.

Take $k\ge2eB$ and put $Q_k(\lambda)=\sum_{r=0}^k\lambda^r/r!$.
The exponential remainder implies
\[
 \|e^{-\lambda}P(\lambda)Q_k(\lambda)-f(\lambda)\|_{[0,B]}
                                                    \le CB/k.
\]
For example, $0\le1-e^{-\lambda}Q_k(\lambda)
\le C(eB/(k+1))^{k+1}\le C2^{-k}$, which suffices.
Write this weighted polynomial as $\sum_{j=0}^{2k}b_jK_\lambda(j)$.
Since
\[
 b_j=\sum_{\substack{0\le l\le k\\0\le j-l\le k}}
                      a_l\frac{j!}{(j-l)!},
\]
equation \eqref{E:pc:monomial} gives
\begin{equation}\label{E:pc:smallcoeff}
             \max_{0\le j\le2k}|b_j|\le B(C_2 k/B)^k.
\end{equation}
The factors $k+1$ and $C_1^k$ have been absorbed into $C_2^k$;
$2k/B\ge1$ by hypothesis.  Wasserstein duality consequently gives
the deterministic inequality
\begin{equation}\label{E:pc:smallinverse}
 W_1(F,H)\le CB/k+
                B(C_2 k/B)^k\|p_F-p_H\|_1.
\end{equation}

For the MLE, likelihood comparison, $h^2\le D$, the Hellinger
triangle inequality and $\|p-q\|_1\le2h(p,q)$ imply
\[
 \mathbb E\|p_{\widehat G}-p_G\|_1
              \le4\sqrt{\mathbb E D(q\Vert p_G)}
              \le C\sqrt{L/N}\qquad(B\le L).
\]
Put $d_0=\log(e+L/B)$ and $k=\lfloor\gamma L/d_0\rfloor$,
where $\gamma>0$ is a sufficiently small absolute constant.
If $\gamma L/d_0<2$, the diameter bound $W_1\le B$ is already at
most a constant times $B\min(1,d_0/L)$.  Otherwise
$k\ge\gamma L/(2d_0)$.  Choose a fixed $b_0>0$ sufficiently small
that $B\le b_0L$ implies
$\gamma(L/B)/(2\log(e+L/B))\ge2e$; then $k\ge2eB$.
Furthermore
\[
 k\log(C_2k/B)\le C_3\gamma L\le L/8
\]
after decreasing $\gamma$ once, since
$\log(C_2\gamma(L/B)/d_0)\le C_3d_0$.
The stochastic term in \eqref{E:pc:smallinverse} is therefore at most
$CB e^{L/8}\sqrt{L/N}\le CB/k$: here $k\le\gamma L$ and
$e^{-3L/8}L^{3/2}$ is bounded.  This proves the required rate for
$B\le b_0L$.  On $b_0L\le B\le L$, equation \eqref{E:pc:medium}
is bounded by an absolute constant, whereas
$B\min(1,d_0/L)$ is bounded below by a positive absolute constant.
This proves the upper bound in Theorem~\ref{E:pc:main} for every $B$.

\subsection{Matching lower bounds without support-range restrictions}

We first record a convenient moment-matching pair.  On an interval
$[a-M,a+M]$, push forward the probability measures
\[
 (1\pm\tfrac12T_k(x))\frac{dx}{\pi\sqrt{1-x^2}},\qquad -1<x<1,
\]
under $x\mapsto a+Mx$.  Denote them by $F_k^+,F_k^-$.
They match all moments of degree less than $k$.  Their CDF difference
is, up to orientation, $\sin(k\theta)/(\pi k)$ for $x=\cos\theta$.
Consequently
\begin{equation}\label{E:pc:witness}
 W_1(F_k^+,F_k^-)
   =\frac M{\pi k}\int_0^\pi|\sin(k\theta)|\sin\theta\,d\theta
   \ge\frac M{\pi k}.
\end{equation}
The last inequality follows from $|\sin u|\ge\sin^2u$ and
$\int_0^\pi\sin^2(k\theta)\sin\theta\,d\theta
=1+(4k^2-1)^{-1}$.

For $B\le L$, use this pair on $[0,B]$.  Expanding $e^{-\lambda}$
and using the matched moments gives the elementary bound
\begin{equation}\label{E:pc:smalltv}
 \|p_{F_k^+}-p_{F_k^-}\|_1
 \le2\sum_{r\ge k}\frac{(2B)^r}{r!}.
\end{equation}
Indeed the term with observed count $j$ and exponential-expansion
index $l$ vanishes if $j+l<k$; summing the absolute values of all
remaining terms gives \eqref{E:pc:smalltv}.
Set $d_0=\log(e+L/B)$ and $k=\lceil100L/d_0\rceil$.
Writing $t=L/B\ge1$, we have
$t=e^{d_0}-e\ge e^{d_0}/(e+1)$ and
$e^{d_0/2}/d_0\ge e/2$.  These inequalities show
\[
 \log\frac{k}{2eB}\ge\frac{d_0}2,
 \qquad k\log\frac{k}{2eB}\ge50L.
\]
They also ensure $k\ge4B$.  Thus the tail in
\eqref{E:pc:smalltv} is at most $4(2eB/k)^k\le4e^{-50L}$.
The two $N$-sample laws have TV at most $2Ne^{-50L}$, whereas
\eqref{E:pc:witness} is at least
$cB\min\{1,d_0/L\}$.  Le Cam's two-point inequality proves the
small-support lower bound, including arbitrarily tiny $B$.

For $B\ge L$, set $a=B/2$, $M=\sqrt{BL}/8$, and
$k=\lceil20L\rceil$.  The moment-matching interval lies inside
$[0,B]$.  The exact Poisson identity
\[
 \sum_{x=0}^\infty\frac{K_\lambda(x)K_\mu(x)}{K_a(x)}
                   =\exp\{(\lambda-a)(\mu-a)/a\}
\]
shows, upon integrating against the signed difference $F_k^+-F_k^-$,
that
\[
 \sum_x\frac{(p_{F_k^+}(x)-p_{F_k^-}(x))^2}{K_a(x)}
 \le4\sum_{r\ge k}\frac{(M^2/a)^r}{r!}
 =4\sum_{r\ge k}\frac{(L/32)^r}{r!}.
\]
All exchanges are justified by boundedness of $\lambda,\mu$ on
the integration interval and the absolutely convergent exponential
series.  The last tail is at most $Ce^{-40L}$, for instance by
$r!\ge(r/e)^r$ and a geometric-tail bound.  Cauchy--Schwarz bounds
the predictive TV by half the square root of this quantity.
Le Cam's inequality and \eqref{E:pc:witness} therefore give the
lower rate $c\sqrt{B/L}$.

Finally consider mixing laws on $\{0,B\}$ with respective weights
$1/2\pm c_0/\sqrt N$ at $B$, for a sufficiently small absolute
$c_0>0$.  Their latent $N$-sample laws have TV bounded away from
one, by the Bernoulli KL formula and Pinsker's inequality.  The
Poisson sample is a Markov image of that latent sample, so its TV
is no larger.  The mixing-law distance is $2c_0B/\sqrt N$.
Le Cam's inequality yields the remaining lower rate $cB/\sqrt N$.
Taking the maximum completes the proof of Theorem~\ref{E:pc:main}.

The proof uses a common known exposure and a known compact support.
It makes no assertion for heterogeneous exposures, unknown support
radii, or an MLE constrained to a prescribed number of components.
The gain over a direct one-scale inverse estimate is structural:
the occupancy channel imports a binomial concentration inequality,
likelihood contact supplies posterior entropy, and the two-scale
inverse places predictive transportation and high-frequency
Hellinger error in different terms.

\subsection{Self-contained binomial inversion tools}
\label{E:pc:binomialtools}

We give the approximation tools used in the proof. Write
$B_j^t(x)=\binom tjx^j(1-x)^{t-j}$ and
$p_F(j)=\int B_j^t\,dF$. The integer $t$ is a binomial trial count,
$W_{1,t}$ uses cost $|i-j|/t$, and
$h(p,q)^2=\sum_j(\sqrt{p_j}-\sqrt{q_j})^2$.

\begin{lemma}[Stable approximation and bounded-polynomial inversion]
\label{E:pc:bc:approx}
There is an absolute $C$ with the following properties.
For every $1$-Lipschitz $f:[0,1]\to\mathbb R$ and integer $k\ge1$, there
is a polynomial $P_k$ of degree at most $k$ such that
\[
 \|P_k-f\|_\infty\le C/k,\qquad \|P_k'\|_\infty\le2.
\]
If $P=\sum_{j=0}^t b_jB_j^t$ has degree $k\le t/2$, then
\begin{equation}\label{E:pc:bc:coeff}
 \max_j|b_j|\le Ce^{2k^2/t}\|P\|_\infty.
\end{equation}
For the derivative-stable $P_k$, its coefficient vector also satisfies
\begin{equation}\label{E:pc:bc:coefflip}
 \max_{j<t}t|b_{j+1}-b_j|\le Ce^{4k^2/t}.
\end{equation}
\end{lemma}
\begin{proof}
Subtract $f(0)$, reflect $f$ on $[0,1]$ to a period-two $1$-Lipschitz
function, and convolve with the normalized positive Jackson kernel
$c_s[\sin(s\pi u/2)/\sin(\pi u/2)]^4$. Its first absolute moment is
$O(1/s)$, so the resulting trigonometric polynomial $Q_s$ has error
$O(1/s)$ and derivative norm at most one. Its frequencies have magnitude
at most $2s$, and its Fourier coefficients have modulus at most one.
Replacing each exponential by its Taylor polynomial at $1/2$ of degree
$100s$ changes the function and derivative by at most $e^{-100s}$.
Indeed $|\pi l(x-1/2)|\le\pi s$ and the factorial inequality
$M!\ge(M/e)^M$ bounds both remainders, including the $4s+1$ terms and the
extra derivative factor. Take $s=\lfloor k/100\rfloor$ for $k\ge100$;
use a constant polynomial and enlarge $C$ for smaller $k$.

For \eqref{E:pc:bc:coeff}, put $M=\|P\|_\infty$ and
$F(z)=(1+z)^tP(z/(1+z))$. The maximum principle in the inverse Joukowski
variable gives
$|P((2+w+w^{-1})/4)|\le M|w|^k$ for $|w|\ge1$.
For $0<j<t$, set $r=j/(t-j)$, $v=r/(1+r)^2$, and
$u=2\sqrt v|\sin(\phi/2)|$. The larger inverse Joukowski modulus for
$re^{i\phi}/(1+re^{i\phi})$ is $(1+u)/\sqrt{1-u^2}$. Hence
\[
 \frac{|F(re^{i\phi})|}{M(1+r)^t}
 \le(1-u^2)^{(t-k)/2}(1+u)^k
 \le e^{2k^2/t}e^{-tu^2/8}.
\]
The last inequality follows from $k\le t/2$ and completing the square.
The average of the final Gaussian factor is at most
$C/\sqrt{1+tv}$. Cauchy's coefficient formula, together with
$\binom tj(j/t)^j(1-j/t)^{t-j}\ge c/\sqrt{1+tv}$, proves
\eqref{E:pc:bc:coeff}. The last lower bound follows directly from Stirling's
factorial bounds; endpoint coefficients are $P(0)$ and $P(1)$.
Finally differentiate the Bernstein expansion:
$P_k'=t\sum_{j<t}(b_{j+1}-b_j)B_j^{t-1}$.
Apply \eqref{E:pc:bc:coeff} at degree $k-1\le(t-1)/2$ and use
$2(k-1)^2/(t-1)\le4k^2/t$.
\end{proof}

\begin{proposition}[Two-scale deterministic inversion]
\label{E:pc:bc:twoscale}
For integers $1\le s\le k\le t/2$ and any two mixing laws $G,H$,
\begin{equation}\label{E:pc:bc:twoscaleineq}
 W_1(G,H)\le \frac Ck
       +Ce^{4s^2/t}W_{1,t}(p_G,p_H)
       +\frac{Ce^{2k^2/t}}s h(p_G,p_H).
\end{equation}
\end{proposition}
\begin{proof}
For any $1$-Lipschitz test $f$, take the polynomials $P_k,P_s$ from
Lemma~\ref{E:pc:bc:approx}. The approximation error in integrating $P_k$
instead of $f$ against $G-H$ is at most $C/k$. The Bernstein coefficients
of $P_s$ have normalized-grid Lipschitz constant at most $Ce^{4s^2/t}$.
On the other hand,
\[
 \|P_k-P_s\|_\infty\le C/k+C/s\le C/s.
\]
Thus the Bernstein coefficients of $P_k-P_s$ have supremum norm at most
$Ce^{2k^2/t}/s$. Test $P_s$ by count transportation and $P_k-P_s$ by
$\|p_G-p_H\|_1\le2h(p_G,p_H)$. All constants are uniform over $f$,
so taking the Wasserstein dual supremum proves the result.
\end{proof}



\clearpage

\section{Binomial mixture estimation and testing across the replication range}
\label{E:p06:bm:sec}

The critical replication scale is $t\asymp N/\log N$.
The estimation proof uses two statistical norms on different
parts of a test polynomial. The low-degree part is controlled by
likelihood transportation. The remaining polynomial has small amplitude,
and is controlled by a predictive Hellinger bound. This yields the full
replication curve, including its transition and its parametric plateau.

For a probability measure $G$ on $[0,1]$, put
\[
 K_x(j)=B_j^t(x)=\binom tj x^j(1-x)^{t-j},\qquad
 p_G(j)=\int B_j^t(x)\,dG(x),\quad 0\le j\le t.
\]
We observe $N$ independent counts from $p_G$. Their empirical law is $q$.
The unrestricted mixing-measure MLE is any maximizer $\widehat G$ of
$\ell(H)=\sum_jq_j\log p_H(j)$. Write $\widehat p=p_{\widehat G}$,
$L_N=\log(eN)$, and
\[
 \mathcal R_{N,t}=\inf_{\widetilde G}\sup_G
       \mathbb E_G W_1(\widetilde G,G),\qquad
 r_{N,t}=\min\left\{1,\frac1t+
           \frac1{\sqrt{tL_N}}+\frac1{\sqrt N}\right\}.
\]
The infimum allows randomized estimators. Wasserstein distance on the
count space uses the normalized cost $|j-l|/t$ and is denoted $W_{1,t}$.
Our Hellinger convention is
$h(p,v)^2=\sum_j(\sqrt{p_j}-\sqrt{v_j})^2$, also for nonnormalized
nonnegative vectors $v$.

\begin{theorem}[All-replication optimality of the unrestricted MLE]
\label{E:p06:bm:full}
There are absolute constants $0<c<C<\infty$ such that, for every pair of
integers $N,t\ge1$,
\begin{equation}\label{E:p06:bm:fullrate}
 c r_{N,t}\le \mathcal R_{N,t}
 \le \sup_G\mathbb E_G W_1(\widehat G,G)
 \le C r_{N,t}.
\end{equation}
The upper bound holds for every measurable choice of an unrestricted MLE.
In particular, at $t\asymp N/\log(eN)$ its risk is $O(N^{-1/2})$,
and the same order holds for every larger $t$.
\end{theorem}

The statistical question extends the replication-range conjectures of
Vinayak, Kong, Valiant and Kakade \cite{BT_VKVK2019} and the discussion
in Section~6.2 of Musco, Musco, Rosenblatt and Singh
\cite{BT_MMRS2025}. The preceding near-linear result already answered the
fixed-polynomial-range question in that discussion. The estimation theorem is reproduced from the preceding compilation;
its consequence for the full testing curve is derived below. It is not
counted as a new estimation result here.

\subsection{An elementary finite-envelope bound for predictive likelihood}

The following entropy argument is stated with its details because it
must apply even when the true distribution has zero cells. Neither an
interior mixing density nor a regular parametric likelihood is assumed.

\begin{lemma}[Small covers of the binomial moment body]
\label{E:p06:bm:cover}
For $0<\beta<1$ the set $\{p_G:G\in\mathcal P([0,1])\}$ has a finite
$\ell_1$ cover of radius $\beta$, with centers in that same set, whose
cardinality $M$ satisfies
\begin{equation}\label{E:p06:bm:entropy}
 \log M\le m\log\frac{C(m+1)}\beta,
 \qquad
 m=\min\left\{t,
       \left\lceil\sqrt{2t\log(16t/\beta)}\right\rceil\right\}.
\end{equation}
\end{lemma}
\begin{proof}
View $K_x=(B_0^t(x),\ldots,B_t^t(x))$ as a vector-valued polynomial.
Its shifted Chebyshev expansion is
\[
 K_x=\sum_{k=0}^t v_kT_k(2x-1).
\]
We claim $\|v_k\|_1\le2\exp\{-k^2/(2t)\}$ for $k\ge1$.
Indeed, the even periodic function
$F(z)=K_{(1+\cos z)/2}$ has Fourier coefficient $v_k/2$.
For $z=\theta+iu$, the ellipse identity gives
\[
 \left|\frac{1+\cos z}{2}\right|
 +\left|\frac{1-\cos z}{2}\right|=\cosh u.
\]
The binomial formula therefore yields
$\|F(\theta+iu)\|_1=(\cosh u)^t\le e^{tu^2/2}$.
Shift the Fourier contour in the appropriate direction and take
$u=k/t$. This proves the claim. In particular, truncating after degree
$m<t$ has vector $\ell_1$ error at most
$2t e^{-(m+1)^2/(2t)}\le\beta/8$, uniformly in $x$.
For $m=t$ the error is zero.

Each mixing law contributes its first $m$ shifted Chebyshev moments in
$[-1,1]^m$. Partition that cube into cells of coordinate diameter at most
$\beta/[8(m+1)]$, and choose one mixing-law representative from every
nonempty cell. If $G,H$ lie in one cell, their truncated transforms differ
in $\ell_1$ norm by at most
$2m\beta/[8(m+1)]\le\beta/4$. The two tails contribute at most
$\beta/4$. Thus their predictive laws have distance at most $\beta$.
The number of cells is at most $[C(m+1)/\beta]^m$, as asserted.
\end{proof}


\begin{lemma}[Likelihood envelopes, including zero cells]
\label{E:p06:bm:envelopes}
Let $\varepsilon\ge0$ and $\eta>0$. Let $p$ be the true law on a finite alphabet with $d$ cells, and suppose a
model admits an $\ell_1$ cover of radius $\eta/d$, with $M$ probability
centers in the model. Let $\widetilde p$ be any fitted model law satisfying
\[
 \sum_{i=1}^N\log\widetilde p(X_i)
 \ge\sum_{i=1}^N\log p(X_i)-N\varepsilon.
\]
For every $x\ge0$, simultaneously for all such fitted laws, except on an
event of probability at most $e^{-x}$,
\begin{equation}\label{E:p06:bm:htail}
 h(p,\widetilde p)
 \le\sqrt\eta+
       \sqrt{\varepsilon+\eta+2(\log M+x)/N}.
\end{equation}
Consequently
\begin{equation}\label{E:p06:bm:hexpect}
 \mathbb E h(p,\widetilde p)^2
 \le C\left\{\varepsilon+\eta+\frac{\log M+1}{N}\right\}.
\end{equation}
\end{lemma}
\begin{proof}
For each center $p_a$, form the nonnegative vector
$u_a(j)=p_a(j)+\eta/d$. Every model law is bounded coordinatewise by at
least one $u_a$, and $\sum_j u_a(j)=1+\eta$. If $\widetilde p\le u_a$,
then
\[
 h(u_a,\widetilde p)^2\le\sum_j(u_a(j)-\widetilde p(j))=\eta.
\]
Thus $h(p,\widetilde p)>\sqrt\eta+r$ implies $h(p,u_a)>r$.
Moreover the likelihood comparison implies
$\prod_i u_a(X_i)/p(X_i)\ge e^{-N\varepsilon}$.
The ratios need only be defined on the support of $p$, where all sampled
observations lie almost surely. Markov's inequality with exponent $1/2$
gives
\begin{align*}
 \Pr\left\{\prod_i\frac{u_a(X_i)}{p(X_i)}
                      \ge e^{-N\varepsilon}\right\}
 &\le e^{N\varepsilon/2}
        \left(\sum_j\sqrt{p(j)u_a(j)}\right)^N\\
 &=e^{N\varepsilon/2}
        \left(1+\frac\eta2-\frac{h(p,u_a)^2}{2}\right)^N\\
 &\le\exp\{N(\varepsilon+\eta-r^2)/2\}
\end{align*}
when $h(p,u_a)>r$. Taking the union over the $M$ envelopes proves
\eqref{E:p06:bm:htail}. Integrating that tail, and using
$(a+b)^2\le2a^2+2b^2$, proves \eqref{E:p06:bm:hexpect}.
\end{proof}

\begin{corollary}[A sufficient predictive Hellinger rate]
\label{E:p06:bm:hellinger}
For $1\le t\le N$, every exact MLE satisfies
\begin{equation}\label{E:p06:bm:hbound}
 \mathbb E_G h(p_G,\widehat p)
       \le \frac{C t^{1/4}L_N^{3/4}}{\sqrt N}.
\end{equation}
More generally, if $\ell(\widetilde G)\ge\ell^*-\varepsilon$, the right
side can be increased by $C\sqrt\varepsilon$.
\end{corollary}
\begin{proof}
Set $\eta=(eN)^{-2}$ and $\beta=\eta/(t+1)$ in the preceding lemmas.
For $t\le N$, \eqref{E:p06:bm:entropy} gives
$m\le C\sqrt{tL_N}$ and
$\log M\le C\sqrt t\,L_N^{3/2}$. The MLE likelihood is at least that of
the true law. The approximate likelihood condition loses at most
$N\varepsilon$. Apply \eqref{E:p06:bm:hexpect} and Jensen's inequality.
\end{proof}

The logarithmic factors in \eqref{E:p06:bm:hbound} need not be sharp. The
argument below only needs the strict improvement from effective dimension
$t$ to a quantity of order $\sqrt t$ up to logarithms.

\begin{lemma}[Likelihood transport, recalled]
\label{E:p06:bm:transport}
For every binomial-mixture MLE,
\begin{equation}\label{E:p06:bm:trans}
 W_{1,t}(q,\widehat p)^2\le\frac{D(q\Vert\widehat p)}{2t}
                    \le\frac{D(q\Vert p_G)}{2t}.
\end{equation}
Consequently
\begin{equation}\label{E:p06:bm:predictive}
 \mathbb E_G W_{1,t}(\widehat p,p_G)\le C/\sqrt N.
\end{equation}
\end{lemma}
\begin{proof}
In score sums, terms with $q_j=0$ are set to zero; every observed
cell has positive fitted probability. The likelihood score condition gives
$S(x)=\sum_jq_jK_x(j)/\widehat p(j)\le1$, and its integral against
$\widehat G$ is one. Thus $S=1$ almost surely under $\widehat G$.
The posterior vectors $R_x(j)=q_jK_x(j)/\widehat p(j)$ are therefore
probability laws, with mixture $q$, and
\[
 \int D(R_x\Vert K_x)\,d\widehat G(x)=D(q\Vert\widehat p).
\]
Every normalized-count binomial kernel satisfies
$W_{1,t}(R,K_x)^2\le D(R\Vert K_x)/(2t)$: a $1$-Lipschitz test of the
normalized sum of $t$ Bernoulli variables has centered logarithmic moment
generating function at most $\lambda^2/(8t)$ by bounded differences;
entropy duality and optimization over $\lambda$ give the inequality.
Mix the couplings, use Cauchy--Schwarz, and then use likelihood optimality
to prove \eqref{E:p06:bm:trans}. Also
$\mathbb E D(q\Vert p_G)\le t/N$, by
$D(q\Vert p_G)\le\sum_{j:p_G(j)>0}(q_j-p_G(j))^2/p_G(j)$.
Finally $\mathbb E W_{1,t}(q,p_G)\le1/(2\sqrt N)$ by integrating the
standard deviation of the empirical distribution function. The triangle
inequality proves \eqref{E:p06:bm:predictive}.
\end{proof}

\begin{proof}[Upper bound in Theorem~\ref{E:p06:bm:full}]
It suffices first to treat $t\le N$. Bounded $N$ or bounded $t$ are handled
by the loss bound one. Fix a sufficiently small absolute $a>0$, for
example $a=1/64$, and take
\[
 k=\left\lfloor\min\{t/2,\sqrt{atL_N}\}\right\rfloor,
 \qquad s=\lfloor\sqrt t/4\rfloor.
\]
For sufficiently large $N$ and $t$, these satisfy
$1\le s\le k\le t/2$, $s\asymp\sqrt t$, and
\[
 1/k\le C\{1/t+(tL_N)^{-1/2}\},\qquad e^{4s^2/t}\le C.
\]
Use Proposition~\ref{E:pc:bc:twoscale}, Lemma~\ref{E:p06:bm:transport}, and
Corollary~\ref{E:p06:bm:hellinger}. This gives
\begin{equation}\label{E:p06:bm:master}
 \mathbb E_G W_1(\widehat G,G)
 \le \frac Ck+\frac C{\sqrt N}
       +\frac{Ce^{2k^2/t}L_N^{3/4}}{\sqrt N\,t^{1/4}}.
\end{equation}
The final term, multiplied by $\sqrt{tL_N}$, is at most
\[
 C(eN)^{2a}N^{-1/2}t^{1/4}L_N^{5/4}
 \le C(eN)^{2a}N^{-1/4}L_N^{5/4}\le C.
\]
The last bound is uniform because $2a<1/4$. This proves the desired
upper bound when $t\le N$.

For $t\ge N$, couple every mixing parameter $x$ to its normalized
binomial count. The expected absolute difference is at most
$\sqrt{x(1-x)/t}\le1/(2\sqrt t)$. Hence, after identifying count laws
with measures on $[0,1]$,
\[
 W_1(\widehat G,G)
 \le 1/\sqrt t+W_{1,t}(\widehat p,p_G).
\]
Take expectations and apply \eqref{E:p06:bm:predictive}. This is $O(N^{-1/2})$
for $t\ge N$, and completes the proof of the upper bound.
\end{proof}

\subsection{Matching lower bounds}

\begin{proof}[Lower bound in Theorem~\ref{E:p06:bm:full}]
The $N^{-1/2}$ term follows from mixing laws supported on $\{0,1\}$ with
weights $1/2\pm c/\sqrt N$. Each count then reveals exactly one endpoint
label, regardless of $t$. Their $N$-sample total variation stays bounded
away from one for sufficiently small absolute $c$.

For the other terms let $\pi(dx)=dx/[\pi\sqrt{x(1-x)}]$ and
$h_k(x)=T_k(2x-1)$. Define
$G_k^\pm=(1\pm h_k/2)\pi$. Integration of their signed cumulative
function gives
\[
 W_1(G_k^+,G_k^-)
  =\frac1{2\pi k}\int_0^\pi|\sin(k\theta)|\sin\theta\,d\theta
  \asymp1/k.
\]
If $k=t+1$, the count laws are identical, proving the $1/t$ lower bound.
For $k\le t$, draw $X\sim\pi$, $J\mid X\sim\operatorname{Binomial}(t,X)$
and put $Qh(j)=\mathbb E[h(X)\mid J=j]$. The posterior is
$\operatorname{Beta}(j+1/2,t-j+1/2)$. The self-adjoint operator $Q^*Q$
preserves polynomial degrees and has leading degree-$k$ coefficient
\[
 \lambda_{t,k}=\prod_{i=0}^{k-1}\frac{t-i}{t+1+i}.
\]
Orthogonality of the Chebyshev polynomials therefore gives
$Q^*Qh_k=\lambda_{t,k}h_k$ and
$\mathbb E[Qh_k(J)^2]=\lambda_{t,k}/2$.
This is the classical beta--binomial diagonalization
\cite[Section~4.1]{BT_DKS2008}.
The count laws are $p_\pi(1\pm Qh_k/2)$, so
\[
 \chi^2(p_{G_k^+}\Vert p_{G_k^-})
 \le2\mathbb E[Qh_k(J)^2]=\lambda_{t,k}
 \le e^{-k^2/(2t)}.
\]
Take $k=\lceil\sqrt{8t\log(8N)}\rceil$ if this is at most $t$.
The chi-squared divergence of the $N$-fold products is
$(1+\chi^2)^N-1\le e^{N(8N)^{-4}}-1$, so their total variation is
bounded by an absolute constant smaller than one. The two-point testing
bound yields risk at least $c/k\ge c'/\sqrt{tL_N}$.
If this choice has $k>t$, then $t\le C L_N$ apart from bounded $t$,
and $1/\sqrt{tL_N}\le C'/t$; the exact moment-matching lower bound
already suffices. The maximum of the three lower bounds is comparable
to their sum clipped at one, proving \eqref{E:p06:bm:fullrate}.
\end{proof}


\subsection{The full worst-null testing curve}

Fix $\alpha,\beta\in(0,1)$ with $\alpha+\beta<1$. A null-dependent test
$\psi_{G_0}$ has level $\alpha$ if
$\mathbb E_{G_0}\psi_{G_0}\le\alpha$. Define
$\varepsilon^*_{N,t}(\alpha,\beta)$ as the infimum of the separations
$\varepsilon$ for which such tests exist for every $G_0$, with
$\mathbb E_G(1-\psi_{G_0})\le\beta$ whenever
$W_1(G,G_0)\ge\varepsilon$. Thus the rate is uniform in the null law;
it need not describe the difficulty at each particular fixed null.

\begin{theorem}[Testing across the replication range]
\label{E:bt:testing}
For all integers $N,t\ge1$,
\[
 \varepsilon^*_{N,t}(\alpha,\beta)
 \asymp_{\alpha,\beta}
 \min\left\{1,\frac1t+\frac1{\sqrt{t\log(eN)}}+\frac1{\sqrt N}\right\}.
\]
An upper bound is attained by the plug-in test based on any unrestricted
NPMLE. In particular the parametric plateau begins at
$t\asymp N/\log(eN)$, with no logarithmic loss in the transition.
\end{theorem}
\begin{proof}
Let $r=r_{N,t}$ and let $C$ be the upper-bound constant in
Theorem~\ref{E:p06:bm:full}. Reject when
$W_1(\widehat G,G_0)>Cr/\alpha$. Markov's inequality gives level
$\alpha$. If $W_1(G,G_0)>Cr(\alpha^{-1}+\beta^{-1})$, the triangle
inequality bounds acceptance probability by
$\Pr_G\{W_1(\widehat G,G)>Cr/\beta\}\le\beta$.
Using a slightly larger constant handles non-strict separation.

For the lower bound use the same three pairs constructed in the lower
proof of Theorem~\ref{E:p06:bm:full}, taking one member as the null and
the other as the alternative. The degree-$t+1$ pair gives identical
sample laws and separation at least $c/t$. For the middle term choose
$k=\lceil\sqrt{C_{\alpha,\beta}t\log(eN)}\rceil$ when $k\le t$;
the displayed chi-square estimate makes the product-law total variation
smaller than $1-\alpha-\beta$ by taking $C_{\alpha,\beta}$ large.
If $k>t$, the $1/t$ pair already gives the required order. Finally
choose the endpoint-weight perturbation of size
$c_{\alpha,\beta}/\sqrt N$ with $c_{\alpha,\beta}$ small, giving
the $N^{-1/2}$ obstruction. For any two sample laws $P_0,P_1$,
the sum of type I and type II errors is at least
$1-\|P_0-P_1\|_{\rm TV}$. These pairs therefore prove the maximum
of the three lower bounds. Their maximum is comparable to their sum
clipped at one, as required.
\end{proof}

Kania, Wasserman and Balakrishnan~\cite[Section 3, Theorem 3]{BT_KWB2025}
give the corresponding worst-null lower curve and discuss the gap in
the intermediate replication range. Theorem~\ref{E:bt:testing}
combines the estimation bound of Theorem~\ref{E:p06:bm:full}
with an elementary Markov argument to obtain the stated testing curve.



\clearpage

\section{Certified approximate likelihood fits}
\label{E:upst:sec}

This section extends the rates in Theorems~\ref{E:pc:main} and
\ref{E:p06:bm:full} to approximate likelihood maximizers. The entropy
chain rule, posterior reweighting, and likelihood monotonicity are
classical EM ingredients \cite{UP_DLR1977}; the contribution here is their quantitative
combination with the two-scale inverse inequalities and the resulting
uniform computational tolerance. The result concerns the known compact
parameter interval. It does not assert analogous stability on an
unbounded parameter space.

\subsection{An exact posterior-defect identity}

Let $q$ be a finitely supported probability law, let $H$ be a mixing
law with $p_H>0$ on the support of $q$, and put
\[
 \ell(H)=\sum_xq(x)\log p_H(x),\qquad
 S_H(\theta)=\sum_x\frac{q(x)K_\theta(x)}{p_H(x)},
 \qquad dH^+(\theta)=S_H(\theta)\,dH(\theta).
\]
The supremum $\ell^*$ below is over the same admissible mixing class,
which is assumed to contain $H^+$ and all point masses. Assume all
component first moments are finite and
$D_K:=\sup_{\theta,\vartheta}W_1(K_\theta,K_\vartheta)<\infty$.

\begin{theorem}[Transport with an explicit likelihood defect]
\label{E:upst:repair}
Suppose the component transport bound of Theorem~\ref{E:et:transfer}
holds with concave nondecreasing function $\psi$. Put
$D_0=D(q\Vert p_H)$ and $d_0=D(H^+\Vert H)$. Then
\begin{align}
 d_0&\le\ell(H^+)-\ell(H)\le\ell^*-\ell(H),
       \label{E:upst:gain}\\
 W_1(q,p_H)&\le\psi(D_0-d_0)+D_K\sqrt{d_0/2}.
       \label{E:upst:repairbound}
\end{align}
In particular, if $\ell^*-\ell(H)\le\delta$, then for every comparator
$G$ in the model,
\begin{equation}\label{E:upst:gapbound}
 W_1(q,p_H)\le
 \psi\bigl(D(q\Vert p_G)+\delta\bigr)+D_K\sqrt{\delta/2}.
\end{equation}
The conclusions also hold after a fixed Markov channel $A$, with
$q,K_\theta,p_H$ replaced in the transport distances by
$Aq,AK_\theta,Ap_H$, provided the transport hypothesis and component
diameter refer to $AK_\theta$. The entropy on the right remains that
of the original data.
\end{theorem}

\begin{proof}
Since $\int S_H\,dH=1$, $H^+$ is a probability law. On $S_H>0$, set
\[
 R_\theta(x)=\frac{q(x)K_\theta(x)}{p_H(x)S_H(\theta)}.
\]
These are probability laws, their $H^+$ mixture is $q$, and the entropy
chain rule gives the exact identity
\begin{equation}\label{E:upst:chain}
 D(q\Vert p_H)=D(H^+\Vert H)
       +\int D(R_\theta\Vert K_\theta)\,dH^+(\theta).
\end{equation}
For every observed $x$, the posterior
$\Pi_x(d\theta)=K_\theta(x)\,dH(\theta)/p_H(x)$ satisfies
$p_{H^+}(x)/p_H(x)=\int S_H\,d\Pi_x$. Jensen's inequality gives
\[
 \ell(H^+)-\ell(H)
 \ge\sum_xq(x)\int\log S_H\,d\Pi_x
 =\int S_H\log S_H\,dH=d_0.
\]
Zero-score points carry no mass in the posterior joint law; the
formulas follow by the convention $0\log0=0$ or by truncation.
This is the usual EM ascent argument, specialized to mixing measures.
Mix the component couplings and use concavity of $\psi$ in
\eqref{E:upst:chain} to obtain
$W_1(q,p_{H^+})\le\psi(D_0-d_0)$.
Couple the common mass of $H,H^+$ identically. The remaining component
mass costs at most $D_K\operatorname{TV}(H,H^+)$, which is at most
$D_K\sqrt{d_0/2}$ by Pinsker's inequality. This proves
\eqref{E:upst:repairbound}. Finally
$D_0\le D(q\Vert p_G)+\delta$ by likelihood comparison.
For a channel, apply entropy data processing to each $R_\theta$ before
mixing; its mixture remains $Aq$. The posterior reweighting and its
likelihood bound are still computed from the original observations.
\end{proof}

\begin{proposition}[A global stopping certificate]
\label{E:upst:certificate}
For every such fit,
\begin{equation}\label{E:upst:scoregap}
 \ell^*-\ell(H)\le\log\sup_\theta S_H(\theta).
\end{equation}
Thus $\sup_\theta S_H(\theta)\le\exp(A/N)$ certifies a total
log-likelihood deficit of at most $A$ for $N$ observations.
\end{proposition}
\begin{proof}
For every $G$, Jensen's inequality yields
\[
 \ell(G)-\ell(H)
 \le\log\sum_xq(x)\frac{p_G(x)}{p_H(x)}
 =\log\int S_H\,dG\le\log\sup_\theta S_H(\theta).
\]
Take the supremum over $G$. The integral of $S_H$ under $H$ is one,
so its supremum is at least one.
\end{proof}

For binomial data $S_H$ is a polynomial of degree at most $t$ on
$[0,1]$; its maximum is among the endpoints and zeros of its derivative.
For Poisson data it is $e^{-\lambda}P(\lambda)$, where $P$ is a
polynomial of degree at most the largest observed count. Its maximum
on $[0,B]$ is among the endpoints and real zeros of $P'-P$.
These are finite candidate sets for the mathematical supremum certificate.
A numerical implementation must bound the candidate values with controlled
error. No polynomial-time complexity or certification from a small change
between consecutive iterations is asserted.

\subsection{The same estimation and testing curves at finite precision}

\begin{theorem}[A bounded total likelihood error costs only a constant]
\label{E:upst:approxrates}
Let $A\ge0$. In the binomial model, let $\widetilde G$ be any
measurable fit on $[0,1]$ satisfying, for every data set,
\[
 \sum_{i=1}^N\log p_{\widetilde G}(X_i)
 \ge\sup_H\sum_{i=1}^N\log p_H(X_i)-A.
\]
Then, for every $N,t\ge1$,
\begin{equation}\label{E:upst:binrate}
 \sup_G\mathbb E_GW_1(\widetilde G,G)
       \le C(1+\sqrt A)\,r_{N,t}.
\end{equation}
In the Poisson model with known $[0,B]$, the same total likelihood
condition over that model implies, for every $N\ge1$ and $B>0$,
\begin{equation}\label{E:upst:poirate}
 \sup_G\mathbb E_GW_1(\widetilde G,G)
       \le C(1+\sqrt A)\,r(N,B).
\end{equation}
Here $r_{N,t}$ and $r(N,B)$ are exactly the complete curves defined
before Theorems~\ref{E:p06:bm:full} and~\ref{E:pc:main}; $C$ is absolute.
In particular, a fixed total log-likelihood tolerance preserves
minimax optimality throughout both parameter planes.
\end{theorem}

\begin{proof}
Write $\delta=A/N$.
For normalized binomial counts, $\psi(u)=\sqrt{u/(2t)}$ and
$D_K=1$. Equation~\eqref{E:upst:gapbound}, the bound
$\mathbb E D(q\Vert p_G)\le t/N$, and the empirical CDF bound give
\begin{equation}\label{E:upst:binpredict}
 \mathbb E W_{1,t}(p_{\widetilde G},p_G)
 \le \sqrt{\frac{t/N+\delta}{2t}}+
       \sqrt{\delta/2}+\frac1{2\sqrt N}
 \le\frac{C(1+\sqrt A)}{\sqrt N}.
\end{equation}
For $t\le N$, Corollary~\ref{E:p06:bm:hellinger}, which explicitly
allows approximate fits, gives
\[
 \mathbb E h(p_{\widetilde G},p_G)
 \le C\{t^{1/4}L_N^{3/4}/\sqrt N+\sqrt\delta\}
 \le C(1+\sqrt A)t^{1/4}L_N^{3/4}/\sqrt N.
\]
Use the same $s,k$ as in the proof of
Theorem~\ref{E:p06:bm:full}. Every term in
\eqref{E:p06:bm:master} increases by at most $C(1+\sqrt A)$, while the
deterministic approximation term is unchanged. This proves
\eqref{E:upst:binrate} for $t\le N$. For $t\ge N$, use
$W_1(\widetilde G,G)\le t^{-1/2}
 +W_{1,t}(p_{\widetilde G},p_G)$ and \eqref{E:upst:binpredict}.

For Poisson data, take $m=\lceil\max(B,L_N)\rceil$ and use the
occupancy channel of \eqref{E:pc:occupancy}. The transformed components
are binomial, with $\psi(u)=\sqrt{u/(2m)}$ and diameter at most one.
The channel version of Theorem~\ref{E:upst:repair} gives
\[
 \mathbb E W_{1,m}(A_mp_{\widetilde G},A_mp_G)
 \le\sqrt{\frac{C(B+L_N)/N+\delta}{2m}}
       +\sqrt{\delta/2}+\frac1{2\sqrt N}
 \le C(1+\sqrt A)/\sqrt N.
\]
The finite-envelope proof of Lemma~\ref{E:pc:hellinger} also applies:
the likelihood ratio threshold becomes $e^{-N\eta-A}$, so the
Hellinger upper bound acquires the additional term $C\sqrt\delta$.
For $b_0L_N\le B\le N$ this is absorbed into
$C(1+\sqrt A)B^{1/4}L_N^{3/4}/\sqrt N$. Substitution into the same
two-scale bound proves \eqref{E:pc:medium} with that factor.
For $B>N$, the coupling and occupancy argument following
\eqref{E:pc:medium} uses only the just-proved predictive bound.

Finally, when $B\le b_0L_N$, likelihood comparison gives
$D(q\Vert p_{\widetilde G})\le D(q\Vert p_G)+\delta$.
The Hellinger triangle inequality therefore implies
\[
 \mathbb E\|p_{\widetilde G}-p_G\|_1
       \le C\sqrt{L_N/N}+2\sqrt\delta
       \le C(1+\sqrt A)\sqrt{L_N/N}.
\]
The small-support inverse inequality \eqref{E:pc:smallinverse} with
the same $k$ as before proves \eqref{E:upst:poirate}, including the
diameter clipping for arbitrarily small $B$. All intermediate parameter
regions and bounded sample sizes were already covered by those same
choices and the diameter bound.
\end{proof}

For fixed testing levels $\alpha,\beta$, the plug-in binomial test
based on $\widetilde G$ therefore has the full worst-null separation
curve $r_{N,t}$ for every fixed $A$. Indeed use the threshold
$C(1+\sqrt A)r_{N,t}/\alpha$ and the same Markov/triangle proof as
in the exact-fit testing theorem. The existing minimax testing lower
bound is unchanged. This closes the finite-computation version as well
as the exact-optimizer version, without claiming EM convergence at a
particular iteration count.

\section{Arbitrary known replication profiles: exact reduction and lower certificates}
\label{E:upst:heterosec}

Let independent units have $P_i\sim G$ and
$X_i\mid P_i\sim\operatorname{Binomial}(t_i,P_i)$, where the positive
integer depths $t_i$ are known. Write $N_m=\#\{i:t_i\ge m\}$.
The following result concerns an explicitly constructed estimator,
not the unmodified heterogeneous joint likelihood.

\begin{theorem}[Thinning upper bound and spectral lower bound]
\label{E:upst:profile}
Let $\mathcal R(t_1,\ldots,t_N)$ be the minimax expected $W_1$ risk
over all mixing laws on $[0,1]$. Then
\begin{equation}\label{E:upst:profileupper}
 \mathcal R(t_1,\ldots,t_N)
 \le C\min_{m:\,N_m>0}r_{N_m,m}.
\end{equation}
Set
\[
 \lambda_{t,k}=\begin{cases}
 \displaystyle\prod_{j=0}^{k-1}\frac{t-j}{t+1+j},&1\le k\le t,\\
 0,&k>t,
 \end{cases}
 \qquad \Lambda_k=\sum_{i=1}^N\lambda_{t_i,k}.
\]
For an absolute $c>0$,
\begin{equation}\label{E:upst:profilelower}
 \mathcal R(t_1,\ldots,t_N)
 \ge c\max\left\{N^{-1/2},
             \sup_{k\ge1}\frac1{k\sqrt{1+\Lambda_k}}\right\}.
\end{equation}
These are design-specific upper and lower certificates. Equality up
to an absolute factor for every arbitrary profile is not asserted.
\end{theorem}

\begin{proof}
Fix $m$ and discard units with $t_i<m$. For each retained unit,
independently draw
\[
 Y_i\mid X_i\sim\operatorname{Hypergeometric}(t_i,X_i,m).
\]
Viewing $X_i$ as the total of $t_i$ independent
Bernoulli trials shows that $Y_i\mid P_i\sim\operatorname{Binomial}(m,P_i)$.
Thus the retained transformed observations are exactly $N_m$ independent
observations from the common-depth mixture. Theorem~\ref{E:p06:bm:full}
proves \eqref{E:upst:profileupper}. The minimization is deterministic;
it can be restricted to the distinct observed depths, since $N_m$
is constant between successive depths and $r_{N_m,m}$ decreases in $m$.

For the lower bound, use the arcsine law $\pi$ and
$h_k(x)=T_k(2x-1)$. As in the proof of Theorem~\ref{E:p06:bm:full},
the beta--binomial posterior operator has squared norm
$\mathbb E[\mathbb E(h_k(P_i)\mid X_i)^2]=\lambda_{t_i,k}/2$;
for $k>t_i$ it is zero by polynomial orthogonality. Take
$G_\pm=(1\pm a h_k)\pi$ with
$a=1/(4\sqrt{1+\Lambda_k})$. These are probability laws.
The individual predictive chi-squared divergences satisfy
\[
 \chi^2(p_{G_+,t_i}\Vert p_{G_-,t_i})
 \le4a^2\lambda_{t_i,k}.
\]
Consequently the product chi-squared divergence is at most
$\exp(4a^2\Lambda_k)-1\le e^{1/4}-1$, and product total variation
is bounded away from one. Direct integration gives
\[
 W_1(G_+,G_-)=\frac a{\pi k}
        \int_0^\pi|\sin(k\theta)|\sin\theta\,d\theta
 \ge\frac a{\pi k}.
\]
For the last inequality, use $|\sin(k\theta)|\ge\sin^2(k\theta)$
and
$\int_0^\pi\sin^2(k\theta)\sin\theta\,d\theta
 =1+(4k^2-1)^{-1}\ge1$.
Le Cam's two-point inequality proves the second lower certificate.
Endpoint mixing laws with weights $1/2\pm c/\sqrt N$ give the
$N^{-1/2}$ certificate, since every positive depth reveals the
endpoint label exactly.
\end{proof}

\begin{corollary}[Any number of bounded-depth observations]
\label{E:upst:shallow}
Fix $d\ge1$. Suppose $n\ge1$ units have a common depth $T>d$ and
an arbitrary additional number of units have arbitrary depths in
$\{1,\ldots,d\}$. There exist $c_d>0$ and absolute $C$ such that
\[
 c_d r_{n,T}\le\mathcal R(t_1,\ldots,t_N)\le C r_{n,T}.
\]
Thus an arbitrarily large bounded-depth sample does not change the
worst-case Wasserstein rate supplied by the deep sample, up to constants
depending on the fixed depth bound $d$. An estimator ignoring the shallow
sample attains this rate. No claim about the heterogeneous joint NPMLE
is needed.
\end{corollary}
\begin{proof}
The upper bound discards the shallow sample. In the lower bound,
every $k>d$ makes all shallow contributions to $\Lambda_k$ vanish.
Choose $k=T+1$ to obtain $c/T$, and $k=d+1$ to obtain
$c/((d+1)\sqrt n)$ since $\Lambda_k\le n$.
For the remaining term, let
$k=\max\{d+1,\lceil\sqrt{8T\log(8n)}\rceil\}$.
If $k\le T$, the estimate
$\lambda_{T,k}\le e^{-k^2/(2T)}$ gives $\Lambda_k\le1$,
so \eqref{E:upst:profilelower} is at least
$c_d/\sqrt{T\log(en)}$. If $k>T$, then
$T\le C_d\log(en)$, and $1/\sqrt{T\log(en)}\le C_d/T$.
The maximum of these three lower bounds is comparable to $r_{n,T}$.
\end{proof}

\begin{proposition}[Known heterogeneous Poisson exposures]
\label{E:upst:poithin}
Suppose $\Theta_i\sim G$ on the known interval $[0,M]$ and,
conditionally, $X_i\sim\operatorname{Poi}(u_i\Theta_i)$, independently,
where $u_i>0$ are known. For $u>0$ let $N_u=\#\{i:u_i\ge u\}$.
There is a randomized estimator with risk at most
\[
 C\inf_{u:\,N_u>0}\frac{r(N_u,uM)}u.
\]
Any fixed choice of $u$ attains its corresponding upper bound.
\end{proposition}
\begin{proof}
Retain the units with $u_i\ge u$ and draw
$Y_i\mid X_i\sim\operatorname{Binomial}(X_i,u/u_i)$ independently.
Poisson thinning gives $Y_i\mid\Theta_i\sim\operatorname{Poi}(u\Theta_i)$.
Fit the common-exposure model on $[0,uM]$ and divide the estimated
mixing parameter by $u$. Theorem~\ref{E:pc:main} gives the bound.
The design is fixed and known, so choose a threshold whose displayed
bound is within a factor of two of the infimum, absorbing that factor
into $C$.
\end{proof}

\subsection{Why separate likelihood strata cannot be normalized silently}

For stratum weights $w_j>0$ summing to one, empirical laws $q_j$,
and an exact joint fit $H$, write
$S_j(\theta)=\sum_xq_j(x)K_{j,\theta}(x)/p_{H,j}(x)$ and
$dH_j^+=S_j\,dH$. The joint score gives only
\[
             \sum_jw_jH_j^+=H.
\]
Applying \eqref{E:upst:chain} separately yields the exact decomposition
\[
 \sum_jw_jD(q_j\Vert p_{H,j})
 =\underbrace{\sum_jw_jD(H_j^+\Vert H)}_{I(J;\Theta)}
 +\sum_jw_j\int D(R_{j,\theta}\Vert K_{j,\theta})\,dH_j^+(\theta).
\]
The first term is the mutual information of the stratum label and
the reweighted parameter under the joint law $w_jH_j^+$. It is the
precise normalization defect. For example, take two equally weighted
Bernoulli strata, $q_1=\delta_0$, $q_2=\delta_1$, and
$H=(\delta_0+\delta_1)/2$. This is a joint likelihood maximizer,
but $H_1^+=\delta_0$, $H_2^+=\delta_1$ and the defect is $\log2$.
Each stratum's within-component entropy is zero, whereas
$W_1(q_j,p_{H,j})=1/2$. Therefore the common-depth transfer theorem
cannot be applied separately to the unaltered strata. The thinning
results above avoid this defect through an exact experiment reduction.

\section{Data-independent finite grids with certified full likelihood accuracy}
\label{fg:section}

The preceding approximate-likelihood theorem turns a bounded total
log-likelihood error into the complete binomial and known-support Poisson
estimation rates. We now give explicit, data-independent finite parameter
grids for which ordinary finite-dimensional likelihood maximization has
exactly this accuracy. The mechanism is a quadratic loss at an interior
score maximum, expressed by the relative entropy between two component
laws. The endpoint nodes are essential to the statement.

Finite-grid nonparametric likelihood estimation and quantitative
discretization guarantees have an established literature. Dicker and
Zhao~\cite{FG_DickerZhao2016} analyze a convex approximation to the NPMLE
in empirical Bayes classification; Soloff, Guntuboyina and
Sen~\cite[Section~2.3]{FG_SoloffGuntuboyinaSen2025} give quantitative
fixed-data grid optimality-gap bounds in Gaussian mixture estimation.
The purpose here is a self-contained stationary-score derivation of the
three explicit grids below and their combination with this manuscript's
mixing-distribution risk bounds. The use of a finite grid itself is not
claimed as a new method.

\subsection{A stationary-score inequality for exponential families}

Let $K_\theta$ have mass function or density $k_\theta$ relative to a
fixed dominating measure. The parameter interval is $I=[a,b]$ and
the functions $\theta\mapsto k_\theta(x_i)$ are continuous for every
observation $x_i$. In the interior, suppose this is a regular
one-parameter exponential family in a smooth coordinate $\eta(\theta)$:
\[
 k_\theta(x)=h(x)
     \exp\{\eta(\theta)T(x)-B(\eta(\theta))\},\qquad
 \eta'(\theta)\ne0.
\]
The endpoint distributions may be continuous limiting distributions of
this family. For a mixing law $H$ on $I$, define
\[
 p_H(x)=\int_I k_\theta(x)\,dH(\theta),\qquad
 \mathcal L_N(H)=\sum_{i=1}^N\log p_H(x_i),\qquad
 S_H(\theta)=\frac1N\sum_{i=1}^N
                         \frac{k_\theta(x_i)}{p_H(x_i)}.
\]
Throughout, $p_H(x_i)>0$ and $k_\theta(x_i)$ is finite. Repeated
observations are allowed. All three models considered below satisfy
these assumptions, and $\int S_H\,dH=1$.

\begin{lemma}[Stationary scores lose at most component relative entropy]
\label{fg:stationary}
If $\theta_*\in(a,b)$ is an interior stationary point of $S_H$ with
$S_H(\theta_*)>0$, then for every $v\in I$ whose component belongs
to the same regular exponential family (including a regular endpoint),
\begin{equation}\label{fg:jensen}
 \frac{S_H(v)}{S_H(\theta_*)}
       \ge \exp\{-D(K_{\theta_*}\Vert K_v)\}.
\end{equation}
Moreover, for any admissible $H$,
\begin{equation}\label{fg:scoregap}
 \sup_G\mathcal L_N(G)-\mathcal L_N(H)
                  \le N\log\sup_{\theta\in I}S_H(\theta).
\end{equation}
\end{lemma}

\begin{proof}
Set $\eta_* =\eta(\theta_*)$ and
\[
 w_i=\frac{k_{\theta_*}(x_i)}{N p_H(x_i)S_H(\theta_*)}.
\]
These weights are nonnegative and sum to one. Since
$\eta'(\theta_*)\ne0$, differentiating the score at its stationary
point gives
$\sum_iw_iT(x_i)=B'(\eta_*)$. Thus Jensen's inequality yields
\begin{align*}
 \frac{S_H(v)}{S_H(\theta_*)}
 &=\sum_iw_i\exp\{(\eta(v)-\eta_*)T(x_i)
                           -B(\eta(v))+B(\eta_*)\}\\
 &\ge\exp\{(\eta(v)-\eta_*)B'(\eta_*)
                           -B(\eta(v))+B(\eta_*)\}\\
 &=\exp\{-D(K_{\theta_*}\Vert K_v)\}.
\end{align*}
Here the last equality is obtained by integrating the exponential-family
log-density ratio under $K_{\theta_*}$.
For any mixing law $G$, a second application of Jensen gives
\[
 \frac{\mathcal L_N(G)-\mathcal L_N(H)}N
 \le\log\left\{\frac1N\sum_i\frac{p_G(x_i)}{p_H(x_i)}\right\}
 =\log\int S_H\,dG\le\log\sup_I S_H.
\]
If $p_G$ vanishes at an observation, its likelihood is minus infinity
and the same inequality is immediate. Taking the supremum proves
\eqref{fg:scoregap}.
\end{proof}

\begin{proposition}[Finite grids certify the continuous likelihood]
\label{fg:gridprinciple}
Let $V\subset I$ be finite and contain both endpoints. Suppose every
interior $\theta$ has a node $v\in V$ belonging to the same regular
exponential family and satisfying
$D(K_\theta\Vert K_v)\le\rho$. Then every admissible fit $H$ satisfies
\begin{equation}\label{fg:globalcertificate}
 \sup_G\mathcal L_N(G)-\mathcal L_N(H)
       \le N\log\max_{v\in V}S_H(v)+N\rho.
\end{equation}
It is enough that the stated relative-entropy comparison hold for every
interior global score maximizer, instead of every interior parameter.

Suppose some mixture supported on $V$ has positive likelihood. Every
exact likelihood maximizer $H_V$ over mixtures on $V$ exists, has
positive fitted probability at each observation, and satisfies
\begin{equation}\label{fg:exactgrid}
                 \sup_G\mathcal L_N(G)-\mathcal L_N(H_V)\le N\rho.
\end{equation}
More generally, for $A_{\rm grid}\ge0$, either of the following conditions suffices for a bound
$A_{\rm grid}+N\rho$ on the continuous-model likelihood deficit:
\begin{enumerate}
\item the fit satisfies the finite score certificate
$\max_{v\in V}S_H(v)\le\exp(A_{\rm grid}/N)$;
\item the fit is supported on $V$ and its total log likelihood is
within $A_{\rm grid}$ of the finite-grid optimum.
\end{enumerate}
\end{proposition}

\begin{proof}
Continuity and compactness give a global score maximizer $\theta_*$.
Its score is positive since $\int S_H\,dH=1$. If it is an endpoint,
it belongs to $V$. Otherwise its score derivative is zero, so
Lemma~\ref{fg:stationary} gives
$\sup_I S_H\le e^\rho\max_VS_H$. This proves
\eqref{fg:globalcertificate}.

The finite-grid parameter space is a compact probability simplex.
The likelihood product is continuous and attains a positive maximum by
the assumed positive candidate. Hence its maximizers are positive at all
observations. Differentiating the normalized log likelihood towards
the point mass at any grid node gives $S_{H_V}(v)-1\le0$.
Since the score has $H_V$-average one, $\max_VS_{H_V}=1$.
Substitute into \eqref{fg:globalcertificate} to obtain
\eqref{fg:exactgrid}. Condition (1) follows directly from that same
inequality. Under (2), add the finite-grid optimization deficit to
\eqref{fg:exactgrid}. The two certificates are sufficient conditions;
we do not infer a score bound merely from an objective-value gap.
\end{proof}

\subsection{Binomial angular grids}

\begin{theorem}[A deterministic grid for every binomial trial count]
\label{fg:binomial}
For integers $N,t,m\ge1$, use the grid
\[
 V_m^{\rm bin}=
 \left\{\sin^2\left(\frac{j\pi}{4m}\right):j=0,1,\ldots,2m\right\}.
\]
Every exact maximum-likelihood fit $H_m$ on this grid satisfies
\begin{equation}\label{fg:bindeficit}
 \sup_{G\in\mathcal P([0,1])}\mathcal L_N(G)-\mathcal L_N(H_m)
                         \le\frac{\pi^2Nt}{8m^2}.
\end{equation}
A finite-grid objective or score certificate as in
Proposition~\ref{fg:gridprinciple} increases the right side by at most
$A_{\rm grid}$. The grid is fixed before observing the data and
contains $0,1/2,1$.
\end{theorem}

\begin{proof}
We first prove the needed one-sided Bernoulli divergence bound.
For $0\le a\le b\le\pi/4$, put $p=\sin^2a$ and $q=\sin^2b$.
Differentiation at $a<v\le b$ gives
\[
 \frac d{dv}D(\operatorname{Ber}(p)\Vert
                       \operatorname{Ber}(\sin^2v))
 =\frac{4\sin(v+a)\sin(v-a)}{\sin(2v)}
 \le4(v-a).
\]
Indeed $0\le v+a\le2v\le\pi/2$ implies
$\sin(v+a)\le\sin(2v)$, while $\sin(v-a)\le v-a$.
Integration from $a$ to $b$ proves
\begin{equation}\label{fg:berkl}
 D(\operatorname{Ber}(p)\Vert\operatorname{Ber}(q))
                \le2(b-a)^2\quad(0\le p\le q\le1/2).
\end{equation}
The value at $a=0$ follows by the same integral or continuity.
Reflection $p\mapsto1-p$, $q\mapsto1-q$ gives the corresponding
bound when $1/2\le q\le p\le1$.

Let $p_*\in(0,1)$ maximize the score. If $p_*\le1/2$, round
$\arcsin\sqrt{p_*}$ up to the next angular grid node; if $p_*\ge1/2$,
round it down. Thus rounding is always towards $1/2$, and the chosen
node lies in $(0,1)$ with angular displacement at most $\pi/(4m)$.
The binomial log-density ratio, integrated using $\mathbb E X=tp_*$,
gives
\[
 D(\operatorname{Bin}(t,p_*)\Vert\operatorname{Bin}(t,q))
 =tD(\operatorname{Ber}(p_*)\Vert\operatorname{Ber}(q))
 \le\frac{\pi^2t}{8m^2}.
\]
Interior stationarity is legitimate because the natural parameter is
$\log(p/(1-p))$. Endpoint maxima are already grid nodes. The grid
node $1/2$ assigns positive probability to every count $0,\ldots,t$,
so finite-grid likelihood maximizers exist and are positive on the
observed counts. Proposition~\ref{fg:gridprinciple} proves the claim.
In particular, all-zero or all-$t$ data cause no exceptional case:
the relevant endpoint score maximum is captured exactly.
\end{proof}

\subsection{Poisson square-root grids and Gaussian uniform grids}

\begin{theorem}[Poisson square-root discretization]
\label{fg:poisson}
For $M>0$ and integers $N,m\ge1$, take
\[
             V_m^{\rm poi}=\{M(j/m)^2:j=0,1,\ldots,m\}.
\]
Every exact finite-grid Poisson mixture MLE on the known interval
$[0,M]$ satisfies
\begin{equation}\label{fg:poideficit}
 \sup_{G\in\mathcal P([0,M])}\mathcal L_N(G)-\mathcal L_N(H_m)
                                      \le\frac{2NM}{m^2}.
\end{equation}
Either finite-grid certificate in Proposition~\ref{fg:gridprinciple}
adds at most $A_{\rm grid}$ to this bound.
\end{theorem}

\begin{proof}
For $0\le a\le b$,
\[
 D(\operatorname{Poi}(a^2)\Vert\operatorname{Poi}(b^2))
 =b^2-a^2+a^2\log(a^2/b^2)\le2(b-a)^2.
\]
For $a>0$, the derivative of the left side with respect to $b$ is
$2(b-a)(b+a)/b\le4(b-a)$, so integration from $a$ proves the
inequality. At $a=0$ it follows directly, with $0\log0=0$.
At an interior score maximum $\lambda_*>0$, choose the next grid
node to its right in square-root coordinates. Its displacement is at
most $\sqrt M/m$, and it is positive. The score is stationary in
the natural parameter $\log\lambda$, equivalently its score-weighted
observations have mean $\lambda_*$. The relative-entropy loss is at
most $2M/m^2$. Both endpoints are included. The grid node $M>0$
assigns positive probability to every nonnegative count, which proves
existence and positivity of finite-grid maximizers. This also covers
all-zero data: a maximum at $\lambda=0$ is an included endpoint.
Apply Proposition~\ref{fg:gridprinciple}.
\end{proof}

\begin{theorem}[Gaussian discretization on a known mean interval]
\label{fg:gaussian}
For known $M,\sigma>0$, integers $N,m\ge1$, Gaussian components
$K_\mu=\mathcal N(\mu,\sigma^2)$ on $\mu\in[-M,M]$, and the grid
\[
              V_m^{\rm gau}=\{-M+2Mj/m:j=0,1,\ldots,m\},
\]
every exact finite-grid MLE satisfies
\begin{equation}\label{fg:gaudeficit}
 \sup_{G\in\mathcal P([-M,M])}\mathcal L_N(G)-\mathcal L_N(H_m)
                                 \le\frac{NM^2}{2\sigma^2m^2}.
\end{equation}
Either finite-grid certificate adds at most $A_{\rm grid}$.
\end{theorem}

\begin{proof}
Every Gaussian kernel is positive at every real observation, so the
finite-grid likelihood maximizer exists and has positive fitted values.
At an interior maximum the score is stationary in the natural parameter
$\mu/\sigma^2$. The nearest grid node $v$ has
$|v-\mu_*|\le M/m$, and
\[
 D(\mathcal N(\mu_*,\sigma^2)\Vert\mathcal N(v,\sigma^2))
                         =\frac{(v-\mu_*)^2}{2\sigma^2}
                         \le\frac{M^2}{2\sigma^2m^2}.
\]
Here a nearest node may be an endpoint, but both endpoint Gaussian
components are regular members of the same exponential family; the
proof of Lemma~\ref{fg:stationary} applies unchanged to them.
Endpoint maxima are captured exactly. Proposition~\ref{fg:gridprinciple}
therefore gives the result. No restriction on the values of the
observations is required.
\end{proof}

\subsection{Finite-dimensional fits and population recovery}

\begin{corollary}[Minimax recovery on explicit grids]
\label{fg:rates}
For binomial mixtures, take $m=\lceil\sqrt{Nt}\rceil$ in
Theorem~\ref{fg:binomial}. Any measurable exact grid MLE, or a grid
fit with fixed total objective deficit $A_{\rm grid}$, satisfies
\[
 \sup_G\mathbb E_GW_1(H_m,G)
 \le C(1+\sqrt{A_{\rm grid}})
       \min\left\{1,\frac1t+\frac1{\sqrt{t\log(eN)}}
                                      +\frac1{\sqrt N}\right\}.
\]
This uses $2\lceil\sqrt{Nt}\rceil+1$ deterministic grid nodes and
attains the full minimax curve, including the transition and plateau.

For Poisson mixtures on known $[0,M]$, take
$m=\lceil\sqrt{NM}\rceil$. Any such finite-grid fit satisfies
\[
             \sup_G\mathbb E_GW_1(H_m,G)
                  \le C(1+\sqrt{A_{\rm grid}})\,r(N,M),
\]
where, with $L=\log(eN)$,
\[
 r(N,M)=\begin{cases}
 M\min\{1,L^{-1}\log(e+L/M)\},&0<M\le L,\\
 \max\{\sqrt{M/L},M/\sqrt N\},&M>L.
 \end{cases}
\]
This uses $\lceil\sqrt{NM}\rceil+1$ nodes. Both conclusions also
hold when the corresponding finite score certificate replaces the
finite-grid objective-gap certificate.
\end{corollary}

\begin{proof}
The binomial total continuous-model likelihood deficit is at most
$A_{\rm grid}+\pi^2/8$; the Poisson deficit is at most
$A_{\rm grid}+2$. Apply Theorem~\ref{E:upst:approxrates}, absorbing the
fixed constants into $C(1+\sqrt{A_{\rm grid}})$. The matching minimax
lower bounds already hold for all estimators, so they apply to these
grid estimators as well. This argument does not assume that the true
mixing distribution is supported on the grid.
\end{proof}

The Gaussian grid similarly has a bounded total likelihood deficit
when $m\ge\lceil M\sqrt N/\sigma\rceil$. Only likelihood accuracy
is asserted for that Gaussian conclusion here. The grid sizes above
are sufficient, not claimed minimal, and no iteration-complexity bound
is claimed for a particular numerical optimizer. These are consequences
proved within this manuscript; no claim of external historical priority
is made for the discretization principle.


\clearpage\part{Further arithmetic classifications and scope}

\section{Global Franel recurrences and arithmetic Ap\'ery limits}
\label{E:ac:section}

Definite critical values give one criterion for recurrence complexity.
The Franel classification below instead uses multiplicities of formal
exponential factors, which also cover complex weights and colliding
critical values. We classify arithmetic subsequences and determine
the space of Ap\'ery limits over the algebraic numbers.

\subsection{The critical-value criterion and its scope}

The width of an eventual nonzero recurrence
$\sum_{j\in J}p_j(n)u_{n+j}=0$, with $p_j\in\mathbb C(n)$, is
$\max J-\min J$ after zero coefficients are removed. Here $J$ is a
finite set of integers; gaps are allowed. The zero sequence has minimum
width zero. A sequence with positive minimum width is understood modulo
finitely many initial terms. Width is not the number of its nonzero
terms and is not the order of an ordinary differential equation.

\begin{proposition}[Visibility of same-side definite critical values]
\label{E:ac:visibility}
Let $\Phi$ be real analytic on a compact real analytic manifold of
positive dimension, equipped with a positive analytic density $d\mu$.
Set $u_n=\int\Phi^n\,d\mu$. Suppose that a nonzero critical value $v$
has finitely many critical points above it, all nondegenerate local
maxima, or all nondegenerate local minima. Regular points in the same
fiber are allowed. Then some analytic continuation of
$\sum_{n\ge0}u_nz^n$ is singular at $z=1/v$. Consequently, $s$
distinct such nonzero critical values force minimum recurrence width
at least $s$.
\end{proposition}
\begin{proof}
The analytic Morse lemma and radial integration give, near $v$, a
pushforward density with singular part
$\mathbf1_{\{\epsilon u>0\}}(\epsilon u)^{d/2-1}h(\epsilon u)$,
where $u=x-v$, $d$ is the manifold dimension, $\epsilon$ is the
common side and $h$ is analytic with $h(0)>0$. The regular contribution
is smooth. Positivity and the common side prevent cancellation of the
first nonsmooth term. This remains true when $d/2-1$ is an integer:
the one-sided polynomial then has a derivative jump. The Cauchy
transform of this measure cannot be holomorphic through $v$, because
its boundary imaginary part is a nonzero constant multiple of this
nonsmooth density. Inverting the Cauchy variable gives the asserted
singularity of the moment generating function.

For the recurrence assertion, translate shifts to $0,\ldots,L$ and
clear denominators. With $\theta=z\partial_z$, the coefficient
recurrence gives an inhomogeneous differential equation whose operator
has the form $\sum_{j=0}^Lz^{L-j}p_j(\theta-j)$. Its highest derivative
coefficient is $z^D H(z)$ with $\deg H\le L$. An operator polynomial
in $\theta$ removes the polynomial inhomogeneous term without creating
new nonzero finite zeros in that leading coefficient. Every branch is
regular away from those zeros and zero. Thus at most $L$ distinct
nonzero finite singularities are possible.
\end{proof}

This proposition requires the stated real positivity and Hessian
conditions. It does not assert that every complex critical value is
visible. None of those conditions is needed in the next theorem.

\subsection{Every weight and every arithmetic subsequence}

For $r\ge2$ and $\lambda\in\mathbb C^*$ write
\[
 a_{r,\lambda}(N)=\sum_{k=0}^N\binom Nk^r\lambda^k,
 \qquad u_{r,\lambda;q,b}(n)=a_{r,\lambda}(qn+b),
 \quad q\ge1,\quad 0\le b<q.
\]

\begin{theorem}[Weighted Franel subsequences]
\label{E:ac:subsequence}
For every integer $r\ge2$, every $\lambda\in\mathbb C^*$, and all
$q,b$ as above, the minimum eventual recurrence width of
$u_{r,\lambda;q,b}$ over $\mathbb C(n)$ is
\[
\begin{array}{c|c}
\text{parameter stratum}&\text{minimum width}\\ \hline
\lambda\notin\{1,-1\}&r\\
\lambda=1&\lceil r/2\rceil\\
\lambda=-1,\ q\text{ odd}&2\lfloor r/2\rfloor\\
\lambda=-1,\ q\text{ even},\ b\text{ even}&\lfloor r/2\rfloor\\
\lambda=-1,\ q\text{ even},\ b\text{ odd}&0.
\end{array}
\]
In the last case the sequence is identically zero. The same minimum
holds over every field $K(n)$ with
$\mathbb Q(\lambda)\subseteq K\subseteq\mathbb C$.
\end{theorem}

The $q=2,b=0$ case contains the all-weight centered classification.
For arbitrary $q,b$, the conclusion also rules out a loss of width
under sampling, except for the stated alternating parity symmetry.
Stoll's and McIntosh's constructions provide the symmetry-weight
upper bounds \cite{AC_Stoll,AC_McIntosh}. Weighted low-degree
recurrences and the higher-Bessel product interpretation have
established precedents \cite{AC_MatalaAho,AC_GRvS}. The theorem
above is proved here by a uniform lower bound, including multiplicity
at exceptional weights; it is not presented as a previously named
conjecture about subsequences.

\subsection{The differential Galois input}


Define the entire function
\[
 \mathcal F_r(x)=\sum_{n\ge0}\frac{x^n}{(n!)^r}.
\]
It solves $((x\partial_x)^r-x)y=0$. Work over $K=\mathbb C(t)$,
and write $M_a$ for the rank-$r$ differential module whose scalar
solutions include
\[
 f_a(t)=\mathcal F_r(a t^r),\qquad a\in\mathbb C^*.
\]
Equivalently its scalar equation is
\[
 \left\{\left(\frac{t\partial_t}{r}\right)^r-a t^r\right\}y=0.
\]
We use the following precise classical input. The differential Galois
group of $M_a$ in its standard representation is
\begin{equation}\label{E:ac:katz}
 G_r=\begin{cases}\mathrm{SL}_r(\mathbb C),&r\text{ odd},\\
                    \mathrm{Sp}_r(\mathbb C),&r\text{ even}.
       \end{cases}
\end{equation}
This is Katz's Kloosterman-group calculation
\cite[Corollary 4.4.8]{AC_Katz}; the connection and its $r$-fold
pullback are identified explicitly in
\cite[Proposition 2.4 and the proof of Proposition 2.7]{AC_FSY}.
Scaling $t$ by a nonzero constant gives the assertion for every $a$.
In particular these groups are connected and their standard
representations are irreducible.

The formal decomposition at $t=\infty$ consists of $r$ rank-one
factors with exponential parts
\begin{equation}\label{E:ac:exp}
 \exp(r\alpha\zeta t),\qquad \zeta^r=1,\quad\alpha^r=a.
\end{equation}
Their remaining parts are regular singular. Direct substitution gives
formal solutions
\[e^{r\alpha\zeta t}t^{-(r-1)/2}(1+O(t^{-1}));\]
the distinct leading exponential roots recursively determine the
formal series. Only the exponential parts, with their multiplicities,
will be used below. The duality relation is
\begin{equation}\label{E:ac:dual}
 M_a^*\simeq M_{(-1)^r a}.
\end{equation}
For example, this follows by reversing the cyclic basis for the
companion system and alternating its signs, or by formal adjunction
with respect to $dx/x$ before the pullback.

\begin{lemma}[Independence of distinct scale classes]
\label{E:ac:independence}
Choose finitely many nonzero scales $a_1,\ldots,a_s$.
If $r$ is even, assume they are pairwise distinct. If $r$ is odd,
assume $a_i\ne\pm a_j$ for $i\ne j$. The connected differential
Galois group of $\bigoplus_i M_{a_i}$ is the full product
$G_r^s$.
\end{lemma}
\begin{proof}
The connected joint group projects onto each individual group:
its image has finite index in the connected group $G_r$.
Its Lie algebra is therefore a subdirect product of simple Lie
algebras. By the Lie-algebra form of Goursat's lemma, a proper
subdirect product identifies two of these simple factors by an
automorphism. Automorphisms of $\mathfrak{sp}_r$ are inner;
automorphisms of $\mathfrak{sl}_r$ are inner or the standard duality
automorphism. Consequently the two standard representations on the
connected joint group would be isomorphic, or in the odd case dual.

Passing to the finite algebraic extension corresponding to the
identity component makes this an isomorphism of the associated
differential modules. But an isomorphism preserves the formal
exponential parts. After any common finite ramification at infinity,
\eqref{E:ac:exp} shows that $M_a$ and $M_b$ can be isomorphic only if
$a=b$, and that $M_a$ and $M_b^*$ can be isomorphic only if
$a=(-1)^r b$. Indeed equality of the multisets of coefficients of
the exponential terms, followed by taking their $r$th powers, gives
these equalities. Both possibilities contradict the scale hypotheses.
Thus no pair of simple factors is identified and the Lie algebra is
the full product. A connected closed subgroup with this Lie algebra
is the full product group.
\end{proof}


\subsection{A factorial-normalization budget with an arbitrary sampling interval}

For a nonzero differential operator $P=\sum_i b_i(t)\partial_t^i$
with formal Laurent coefficients at infinity, let $v$ be the largest
leading power of $t$ among its coefficients and put
\[
 \sigma_\infty(P)(\xi)
 =\sum_i\left(\lim_{t\to\infty}t^{-v}b_i(t)\right)\xi^i.
\]
The limit means the constant coefficient after Laurent normalization.
Differentiating a Laurent coefficient lowers its degree, so
\begin{equation}\label{E:ac:symbol-product}
 \sigma_\infty(PQ)=\sigma_\infty(P)\sigma_\infty(Q)
\end{equation}
(up to a harmless nonzero constant if another normalization is used).

\begin{lemma}[Nonzero exponential multiplicities bound width]
\label{E:ac:budget}
Let
\[
 H(t)=\sum_{n\ge0}\frac{u_n}{((qn+b)!)^r}t^{r(qn+b)}.
\]
Suppose the differential module generated by this scalar function has
at infinity a decomposition into rank-one factors with exponential
parts $e^{\beta t}$ and regular-singular remaining parts. If $\nu$
of these factors have $\beta\ne0$, counted with multiplicity, then
every eventual recurrence of $u_n$ has width at least $\nu/(qr)$.
\end{lemma}
\begin{proof}
Write a proposed width-$L$ recurrence as
$\sum_{j=0}^L p_j(n)u_{n-j}=0$, with $p_j\in\mathbb C[n]$.
Let $D=\max\deg p_j$ and let $c_j$ be the leading coefficient when
$\deg p_j=D$. Multiplication by
$((qn+b)_{\underline{qL}})^r/((qn+b)!)^r$ gives
$PH=Q(t)$, with $Q$ polynomial and
\[
 P=\sum_{j=0}^L t^{qrj}
 p_j\left(\frac{t\partial_t}{qr}+j-\frac bq\right)
 \prod_{\ell=0}^{q(L-j)-1}
       \left(\frac{t\partial_t}{r}-\ell\right)^r.
\]
Indeed the coefficient contributed by the $j$th term at exponent
$r(qn+b)$ is $p_j(n)u_{n-j}/((qn+b-qL)!)^r$ for sufficiently
large $n$. The finitely many other coefficients account for $Q$.
The largest coefficient degree in $t$ is $D+qrL$, and
\begin{equation}\label{E:ac:symbol}
 \sigma_\infty(P)(\xi)
 =\sum_{\deg p_j=D}c_j
       \left(\frac\xi{qr}\right)^D
       \left(\frac\xi r\right)^{qr(L-j)}.
\end{equation}
Distinct $j$ give distinct powers, so this is nonzero and has at most
$qrL$ nonzero roots with multiplicity. Left multiplication by a
polynomial in $t\partial_t$ killing $Q$ makes the equation homogeneous
and adds only zero roots to the principal symbol.

Let $R$ be the monic minimal scalar annihilator of $H$. Its formal
rank-one decomposition gives a first-order factorization with factors
$\partial_t-\beta_i+O(t^{-1})$: take a formal rank-one solution,
right-divide by its logarithmic-derivative factor, and repeat on the
quotient. Laurent gauge changes add only $O(t^{-1})$.
Thus $\sigma_\infty(R)=\prod_i(\xi-\beta_i)$.
The homogeneous annihilator is a left multiple of $R$.
Equation~\eqref{E:ac:symbol-product} then forces its principal symbol
to contain all $\nu$ nonzero factors of $R$, with multiplicity.
Consequently $\nu\le qrL$.
\end{proof}

\subsection{Cyclicity after root-of-unity filtering}

The Cauchy-product identity gives
\[
 \mathcal F_r(x)\mathcal F_r(\lambda x)
 =\sum_{N\ge0}\frac{a_{r,\lambda}(N)}{(N!)^r}x^N.
\]
Writing $\mu_q=\{\epsilon:\epsilon^q=1\}$, its root-of-unity
filter is
\begin{equation}\label{E:ac:filter}
 H_{q,b}(t)=\frac1q\sum_{\epsilon\in\mu_q}
                 \epsilon^{-b} f_\epsilon(t)f_{\lambda\epsilon}(t)
 =\sum_{n\ge0}\frac{a_{r,\lambda}(qn+b)}{((qn+b)!)^r}
                   t^{r(qn+b)}.
\end{equation}
All products lie in one Picard--Vessiot field. Multiplication of scalar
solutions is an equivariant linear map from their tensor product.
On an irreducible representation a nonzero such map is injective.
Distinct nonisomorphic irreducible images have direct sum, and a vector
with a nonzero component in each generates the whole direct sum.
These facts establish cyclicity for the specific filtered function,
not just for an ambient system that annihilates it.

\begin{proof}[Lower bounds in Theorem~\ref{E:ac:subsequence}]
First suppose $\lambda\ne\pm1$. Each product in
\eqref{E:ac:filter} belongs to an irreducible rank-$r^2$ tensor
representation. For even $r$ the two scales are distinct independent
symplectic factors. For odd $r$ they are distinct modulo sign, and
hence independent special-linear factors. Use
Lemma~\ref{E:ac:independence} for the finite set of all scale classes
occurring in the filter.

The $q$ product representations are pairwise nonisomorphic even
when their scales overlap. For even $r$, their two-factor supports
agree only if the unordered pairs
$\{\epsilon,\lambda\epsilon\}$ agree. For odd $r$, a representation
also records whether it is standard or dual at each signed scale
class: these are inequivalent for odd $r\ge3$. Hence isomorphism
again requires equality of the unordered pairs of actual scales.
For two distinct indices $\epsilon,\eta$, such equality would require
$\eta=\lambda\epsilon$ and $\lambda\eta=\epsilon$, forcing
$\lambda^2=1$, which is excluded. Since every coefficient
$\epsilon^{-b}/q$ is nonzero, the filter is cyclic in all $q$
irreducible summands. A formal exponential coefficient in
$M_a\otimes M_c$ can vanish only if $c=(-1)^r a$.
This too is excluded. Therefore $\nu=qr^2$ and the budget gives
$L\ge r$.

For $\lambda=1$, the summands are
$\operatorname{Sym}^2 M_\epsilon$. Their representations are
irreducible and pairwise nonisomorphic. For even $r$ this follows
from distinct symplectic scale factors. For odd $r$, opposite scales
if present give $\operatorname{Sym}^2 V$ and
$\operatorname{Sym}^2 V^*$, which are nonisomorphic for $r\ge3$.
In each summand the formal exponents are the unordered sums of two
$r$th-root exponents. There are $r/2$ zero sums when $r$ is even,
and none when $r$ is odd. The nonzero rank per summand is thus
$r^2/2$ for even $r$, and $r(r+1)/2$ for odd $r$.
The $q$ nonzero filter components give respectively
$L\ge r/2$ and $L\ge(r+1)/2$.

Now let $\lambda=-1$. The products indexed by $\epsilon$ and
$-\epsilon$ coincide. If $q$ is even, their filter coefficients sum
to zero for odd $b$, and to $2\epsilon^{-b}/q$ for even $b$.
This proves the asserted vanishing in the first case. In the second
case there are $q/2$ surviving products. If $q$ is odd, none of these
pairs repeats and there are $q$ surviving products.

For even $r$, each surviving product is an irreducible tensor of two
independent symplectic modules. Distinct surviving products have
disjoint scale pairs, so they are nonisomorphic. Each has $r^2$
nonzero formal factors, since zero would require the scales to be
equal rather than opposite.

For odd $r$, each product lies in
$V\otimes V^*=\mathbf1\oplus\operatorname{Ad}$ for the special-linear
factor indexed by the class $\{\epsilon,-\epsilon\}$.
The multiplication map kills the trivial summand. To check this
without assuming that invariant scalar functions are constant, put
$\delta=(t/r)\partial_t$. For solutions $y,z$ of
$\delta^r y=\epsilon t^r y$, $\delta^r z=-\epsilon t^r z$, the
nondegenerate concomitant
\[
 \mathcal B(y,z)=\sum_{j=0}^{r-1}(-1)^j
                  (\delta^j y)(\delta^{r-1-j}z)
\]
is constant. Indeed its derivative telescopes to zero.
Choose dual solution bases $y_i,z_i$ for this pairing and jet matrices
$Y,Z$ whose rows are $\delta^j$, $0\le j<r$. Then
$Y^{\mathsf T}JZ=I$, with $J$ invertible and anti-diagonal, whence
\[
 \sum_i y_i z_i=(YZ^{\mathsf T})_{00}=(J^{-\mathsf T})_{00}=0
 \qquad(r\ge2).
\]
The selected product is nonzero, so the multiplication map is
injective on the irreducible adjoint summand. Different surviving
products use different independent scale classes. Each therefore
contributes the $r(r-1)$ nonzero ordered differences of distinct
$r$th-root exponents. There are $q$ such products for odd $q$, or
$q/2$ for even $q$ and even $b$. Dividing their nonzero ranks by $qr$
gives respectively $2\lfloor r/2\rfloor$ and
$\lfloor r/2\rfloor$, in both parities of $r$.
\end{proof}

\begin{proof}[Upper bounds and coefficient fields]
For $\lambda\ne\pm1$ set
$\mathcal L_N[p]=\sum_{k=0}^Np(k)\binom Nk^r\lambda^k$.
An index shift gives
\[
 \mathcal L_N[k^{r+h}-\lambda(N-k)^r(k+1)^h]=0
 \qquad(h\ge0).
\]
The polynomial displayed here has leading coefficient
$1-(-1)^r\lambda\ne0$. Successive elimination reduces every moment
to $\mathcal L_N[1],\ldots,\mathcal L_N[k^{r-1}]$.
The identity
$\binom{N-1}{k}^r=N^{-r}(N-k)^r\binom Nk^r$
therefore gives a backward $r$-dimensional state relation over
$\mathbb Q(\lambda)(N)$. Apply it $q$ times between observations:
$r+1$ consecutive observations on the progression $qn+b$ are rows
on the same $r$-dimensional state, so they satisfy a recurrence of
width at most $r$. This argument needs no matrix inverse.

For $\lambda=1$, Stoll's ordinary Franel recurrence has width at
most $m=\lceil r/2\rceil$ over $\mathbb Q(n)$
\cite{AC_Stoll}. A rational companion system has dimension $m$;
its $q$-step system gives the same upper bound on every progression.
For $\lambda=-1$, symmetry gives $a_{r,-1}(N)=0$ for odd $N$.
McIntosh's recurrence \cite[Theorem 1]{AC_McIntosh} gives width
at most $h=\lfloor r/2\rfloor$ for
$A_r(n)=(-1)^n a_{r,-1}(2n)$. Geometric rescaling preserves width,
so even progressions have the claimed upper bound $h$. Interlacing
this recurrence with the zero odd sequence gives a width-$2h$
recurrence for $a_{r,-1}(N)$: a recurrence on its even entries with
shifts $j$ becomes one with shifts $2j$, valid also on the zero
odd entries. Taking a $q$-step companion system gives the upper
bound $2h$ for odd $q$. All upper constructions have coefficients
in $\mathbb Q(\lambda)(n)$. Since the lower bounds hold over
$\mathbb C(n)$, extending the coefficient field anywhere between
these fields cannot change the minimum.
\end{proof}

\begin{corollary}[Franel and centered alternating sums]
\label{E:ac:oldEF}
For every $r\ge2$,
\[
 \operatorname{minwidth}_{\mathbb C(n)}
       \sum_{k=0}^n\binom nk^r=\lceil r/2\rceil,
 \qquad
 \operatorname{minwidth}_{\mathbb C(n)}
       \sum_{k=-n}^n(-1)^k\binom{2n}{n+k}^r=\lfloor r/2\rfloor.
\]
More generally, the centered weighted sequence
$\lambda^{-n}a_{r,\lambda}(2n)$ has width
$\lceil r/2\rceil$ at $\lambda=1$, $\lfloor r/2\rfloor$ at
$\lambda=-1$, and $r$ elsewhere.
\end{corollary}
\begin{proof}
Use Theorem~\ref{E:ac:subsequence} with $(q,b)=(1,0)$ or $(2,0)$.
The centered alternating sum equals $(-1)^na_{r,-1}(2n)$, and
geometric rescaling preserves width.
\end{proof}

The alternating conclusion does not require a critical-value visibility
argument. The unweighted conclusion removes the restriction to
telescoping recurrences in
\cite[Theorem 1.1]{AC_SZ}. Critical-value collisions do not create
exceptions in Theorem~\ref{E:ac:subsequence}: the budget counts formal
exponential factors with multiplicity, not distinct critical values.

\subsection{Arithmetic spaces of Ap\'ery limits}

Put $A_r(N)=a_{r,1}(N)$ and $m=\lceil r/2\rceil$. In this
subsection the notation $A_r$ denotes the ordinary, not the alternating,
Franel sequence. Define rational sequences $A_{r,j}(N)$ by
\begin{equation}\label{E:ac:deformation}
 \sum_{j\ge0}A_{r,j}(N)z^{2j}
 =\sum_{k=0}^N\binom Nk^r
 \left\{\prod_{\ell=1}^k\left(1-\frac z\ell\right)
        \prod_{\ell=1}^{N-k}\left(1+\frac z\ell\right)\right\}^{-r}.
\end{equation}
The right side is even by replacing $k$ with $N-k$, and
$A_{r,0}=A_r$. Let
\[
 \phi_{r,j}=[z^{2j}]\left(\frac z{\sin z}\right)^r
 \in\mathbb Q_{>0}.
\]
The input from Straub--Zudilin \cite[Theorem 1.3]{AC_SZ} is that
$A_{r,j}$, $0\le j<m$, satisfy every telescoping recurrence for
$\binom Nk^r$ on a sufficiently late tail, and have limits
\begin{equation}\label{E:ac:known-limits}
 \lim_{N\to\infty}\frac{A_{r,j}(N)}{A_r(N)}
             =\phi_{r,j}\pi^{2j}.
\end{equation}
Van Hoeij's invariant-submodule calculation
\cite[Remark 6]{AC_vH} gives a telescoper of order $m$ for this
summand. These are external inputs; the transfer to the unrestricted
minimal recurrence is supplied by the preceding theorem.

For $q\ge1$ and $0\le b<q$, let $L_{r;q,b}$ denote the monic
minimal eventual recurrence of $A_r(qn+b)$, with shifts normalized
to $0,\ldots,m$. If $K$ is a field between $\mathbb Q$ and
$\overline{\mathbb Q}$, write $\mathcal S_{r;q,b}(K)$ for the
$K$-valued solutions on a sufficiently late regular tail, identifying
sequences that agree eventually.

\begin{theorem}[The Ap\'ery-limit map on algebraic-valued solutions]
\label{E:ac:apery}
For every $r\ge2$, $q\ge1$, $0\le b<q$ and
$\mathbb Q\subseteq K\subseteq\overline{\mathbb Q}$, the following
statements hold.
\begin{enumerate}
\item $L_{r;q,b}$ has order $m$ and coefficients in $\mathbb Q(n)$.
For $(q,b)=(1,0)$ it is also the monic minimal telescoper for the
summand $\binom nk^r$.
\item The sequences $A_{r,j}(qn+b)$, $0\le j<m$, are a basis of
$\mathcal S_{r;q,b}(K)$. Every complex tail solution also has a
unique expansion in this same basis, with complex constants.
\item Every $X\in\mathcal S_{r;q,b}(K)$ has an Ap\'ery limit, and
\begin{equation}\label{E:ac:limit-isomorphism}
 X\longmapsto \lim_{n\to\infty}\frac{X(n)}{A_r(qn+b)}
 \quad\text{is a $K$-linear isomorphism onto}\quad
 \bigoplus_{j=0}^{m-1}K\pi^{2j}.
\end{equation}
In particular two algebraic-number-valued tail solutions with the
same limit are equal eventually.
\end{enumerate}
\end{theorem}
\begin{proof}
For $q=1,b=0$, the minimal recurrence has order $m$ by
Corollary~\ref{E:ac:oldEF}. Stoll provides an annihilator of this order
over $\mathbb Q(n)$, so its monic normalization is the unique minimal
operator even over $\mathbb C(n)$. The minimal telescoper supplied
by van Hoeij has the same order and annihilates the sequence. Right
division by the monic minimal recurrence leaves a lower-order
annihilator as remainder, necessarily zero. Equality of orders thus
makes the two monic operators equal. Equation~\eqref{E:ac:known-limits}
therefore applies to solutions of this minimal recurrence.

Pass to a tail beyond all poles and integer zeros of its leading
and trailing coefficients, and beyond the finitely many exceptional
indices for the $A_{r,j}$. A companion system has dimension $m$.
Its $q$-step system yields one common rational-coefficient recurrence
of order at most $m$ for every sampled solution. By
Theorem~\ref{E:ac:subsequence}, the sampled sequence $A_r(qn+b)$
has minimum order exactly $m$. Hence this common operator has order
$m$ and, after normalization, is $L_{r;q,b}$.

The $m$ rational sequences $A_{r,j}(qn+b)$ are linearly independent
over $\mathbb C$. Indeed a complex dependence among rational
columns of an infinite matrix implies a rational dependence: the
row space is finite-dimensional and can be spanned by finitely many
rational rows, whose nullspace is defined over $\mathbb Q$.
A rational dependence on any tail would, after division by
$A_r(qn+b)$ and passage to the limit, contradict the transcendence
of $\pi$ in \eqref{E:ac:known-limits}. The solution space of a regular
order-$m$ tail recurrence has dimension $m$, so these columns are a
basis over $\mathbb C$.

Evaluation at $m$ consecutive sufficiently late indices is an
isomorphism from this solution space to $\mathbb C^m$. The resulting
basis matrix has rational entries and nonzero determinant, hence
rational inverse. A solution taking values in $K$ consequently has
all its basis coefficients in $K$, proving the basis assertion over
$K$. Its limit is
$\sum_{j=0}^{m-1}c_j\phi_{r,j}\pi^{2j}$.
The numbers $1,\pi^2,\ldots,\pi^{2m-2}$ are linearly independent
over $\overline{\mathbb Q}$, and every $\phi_{r,j}$ is nonzero
rational. This proves injectivity and the stated image.
\end{proof}

An algebraic-number-valued solution need not be assumed in advance
to lie in one number field: its $m$ initial tail values generate a
finite extension of $\mathbb Q$, and the rational recurrence then
keeps its entire tail in that field. Thus the theorem covers all such
solutions, including those initially described only as
$\overline{\mathbb Q}$-valued.

\begin{corollary}[Normalized Ap\'ery limits and uniqueness]
\label{E:ac:normalized}
For $1\le j<m$ put
\[
 D_{r,j}(N)=\frac{A_{r,j}(N)}{2\binom{r+2j-1}{2j}}.
\]
Then $D_{r,j}(0)=0$, $D_{r,j}(1)=1$, and it is a solution of the
minimal Franel recurrence on a tail, with
\[
 \lim_{N\to\infty}\frac{D_{r,j}(N)}{A_r(N)}
 =\frac{\phi_{r,j}\pi^{2j}}{2\binom{r+2j-1}{2j}}.
\]
It is the unique algebraic-number-valued tail solution with that
limit. The same uniqueness holds for its sampled sequence
$D_{r,j}(qn+b)$ relative to $L_{r;q,b}$ and $A_r(qn+b)$.
In particular, for $r\ge3$ and $r\ge5$, respectively,
\[
 \lim\frac{D_{r,1}(N)}{A_r(N)}=\frac{\zeta(2)}{r+1},
 \qquad
 \lim\frac{D_{r,2}(N)}{A_r(N)}
 =\frac{3(5r+2)\zeta(4)}{(r+1)(r+2)(r+3)}.
\]
\end{corollary}
\begin{proof}
At $N=0$ the deformation is $1$, while at $N=1$ it is
$(1-z)^{-r}+(1+z)^{-r}$. The stated initial values follow by
coefficient extraction. The limits and uniqueness follow from
Theorem~\ref{E:ac:apery}. Finally
$\phi_{r,1}=r/6$ and $\phi_{r,2}=r(5r+2)/360$ give the two
special formulas, already computed in \cite[Examples 1.5 and 1.7]{AC_SZ}.
\end{proof}

\subsection{Exact scope of the conjectural statements}

Chamberland--Straub \cite[Conjectures 9 and 11]{AC_CS} state the
normalized $\zeta(2)$ and $\zeta(4)$ limits, respectively for
$r\ge3$ and $r\ge5$, and include uniqueness. Their wording does not
specify a coefficient field for the secondary solution. The numerical
values of the limits and the deformed rational sequences are already
in Straub--Zudilin. What the argument here supplies is their membership
in the unrestricted minimal recurrence, complete uniqueness over
$\overline{\mathbb Q}$, and the same conclusions for every arithmetic
subsequence. It does not establish unrestricted complex uniqueness.

In fact, on the complex tail solution space the limit map has a
kernel of dimension $m-1$. For $r\ge7$ this obstruction persists
even with the two displayed initial normalizations in the eventual
sense. To make it explicit, set $c_j=A_{r,j}(1)$ and
$\ell_j=\phi_{r,j}\pi^{2j}$, $j=1,2,3$, and put
\[
 Z=(c_2\ell_3-c_3\ell_2)A_{r,1}
   +(c_3\ell_1-c_1\ell_3)A_{r,2}
   +(c_1\ell_2-c_2\ell_1)A_{r,3}.
\]
Then $Z(0)=Z(1)=0$, $Z/A_r\to0$, and $Z$ is not the zero tail
solution: its coefficient of $A_{r,3}$ is nonzero by the transcendence
of $\pi$, and the basis is independent. Thus $D_{r,1}+cZ$ has the
same two initial values and limit for every $c\in\mathbb C$, and
satisfies the recurrence eventually. An assertion requiring the
recurrence to hold at every small nonnegative index needs a separate
regularity argument; no such assertion is inferred here from eventual
annihilation. The coefficient-degree prediction in
\cite[Remark 4.1]{AC_SZ} is likewise not settled by a width bound.


\clearpage

\section{An arithmetic classification for split beta orthogonality}
\label{E:aa:section}

The critical-value argument uses positivity to prevent cancellation of
local contributions. A different use of positivity gives an arithmetic
classification for multiple orthogonal polynomials: once a logarithmic
moment identity forces orthogonality on both sides of a cut, the required
sign changes can exceed the degree. The classification below is uniform
in both orthogonality indices, in the algebraic cut, and in every positive
algebraic polynomial modification of the weight.

\subsection{Parameter and index classification}

Let $a,b>0$ be rational with $a+b\in\mathbb Z$, let
$\xi\in\overline{\mathbb Q}\cap(0,1)$, and let
$w\in(\overline{\mathbb Q}\cap\mathbb R)[X]$ satisfy
$w(x)>0$ for $0<x<1$. Define the positive integrable weight
\[
 \rho(x)=w(x)x^{a-1}(1-x)^{b-1}.
\]
For $m,n\ge0$, the monic type-II Angelesco polynomial
$P_{m,n}$ has degree $m+n$ and satisfies
\begin{align}
 \int_0^\xi P_{m,n}(x)x^j\rho(x)\,dx&=0
       &&(0\le j<m),\label{E:aa:left}\\
 \int_\xi^1 P_{m,n}(x)x^j\rho(x)\,dx&=0
       &&(0\le j<n).\label{E:aa:right}
\end{align}
Empty lists of conditions are allowed, and $P_{0,0}=1$.
Existence and uniqueness will be proved below.

\begin{theorem}[Exact algebraic-coefficient classification]
\label{E:aa:main}
For the preceding family, the following are equivalent:
\begin{enumerate}
\item Every coefficient of $P_{m,n}$ is algebraic.
\item Every zero of $P_{m,n}$ is algebraic.
\item Either $a,b\in\mathbb Z$, or $m=n$.
\end{enumerate}
In the algebraic cases the coefficients belong to
\[
 K=\mathbb Q(\xi,w_0,\ldots,w_d),
 \qquad w(X)=\sum_{j=0}^d w_jX^j.
\]
Consequently, when $a,b$ are nonintegers and $m\ne n$,
$P_{m,n}$ has at least one transcendental coefficient and at least
one transcendental zero. All its zeros are nevertheless real and simple,
with precisely $m$ in $(0,\xi)$ and $n$ in $(\xi,1)$.
\end{theorem}

The rationality of $a,b$ and the integer-sum condition are essential to
the stated arithmetic reduction. The theorem includes the arcsine pair
$a=b=1/2$; no exclusion of that pair is needed. We assert that at least
one zero is transcendental off the diagonal, not that every zero is.

\subsection{Positivity and normality}

\begin{lemma}[Sign changes and existence]\label{E:aa:normality}
For a positive integrable weight on $(0,1)$, each system
\eqref{E:aa:left}--\eqref{E:aa:right} has a unique monic solution of
degree $m+n$. That solution has exactly the zero distribution stated
in Theorem~\ref{E:aa:main}.
\end{lemma}

\begin{proof}
A nonzero real polynomial $Q$ orthogonal to
$1,x,\ldots,x^{k-1}$ on an interval with a positive weight has at
least $k$ sign changes in its interior. Indeed, if the sign-change
points were $t_1,\ldots,t_s$ with $s<k$, the polynomial
$q(x)=\prod_{i=1}^s(x-t_i)$ would make $Qq$ have one constant
nonzero sign except at finitely many points. Its weighted integral
could not vanish, contradicting orthogonality to $q$.

A nonzero polynomial satisfying both lists therefore has at least
$m+n$ distinct real zeros in the two disjoint interiors. In particular,
there is no such polynomial of degree at most $m+n-1$. The linear map
from that $(m+n)$-dimensional coefficient space to its $m+n$ moments
is injective, hence bijective. It follows that the lower coefficients
of a monic degree-$m+n$ solution exist uniquely. The same sign-change
count exhausts its degree, proving simplicity and the precise locations
of all zeros. The case $m+n=0$ is immediate.
\end{proof}

\subsection{A logarithmic moment forces orthogonality on both sides}

The transcendence input is the classical consequence of Baker's theorem:
a nonzero linear form, with algebraic coefficients, in fixed logarithms
of nonzero algebraic numbers is transcendental
\cite{aa:BakerIII}. We reproduce the needed integral reduction, so the
argument does not require any unproved period-independence assertion.

For $r\in\mathbb Q\cap(0,1)$ write
\[
 b_r(x)=x^{r-1}(1-x)^{-r},\qquad
 U_r(y)=\int_0^y b_r(x)\,dx.
\]

\begin{lemma}[Two transcendental base integrals]\label{E:aa:base}
If $\xi\in\overline{\mathbb Q}\cap(0,1)$, then both
$U_r(\xi)$ and $\int_\xi^1 b_r(x)\,dx$ are transcendental.
This includes $r=1/2$.
\end{lemma}

\begin{proof}
Write $r=p/q$ with $1\le p<q$ and put
$s=(\xi/(1-\xi))^{1/q}>0$. The substitution
$x=z^q/(1+z^q)$ and partial fractions give
\[
 U_r(\xi)=q\int_0^s\frac{z^{p-1}}{1+z^q}\,dz
 =-\sum_{\zeta^q=-1}\zeta^p\log(1-s/\zeta).
\]
For each root $\zeta$, take the continuous logarithm along positive
real $s$, starting at zero when $s=0$. Its argument never vanishes
because no root of $\zeta^q=-1$ is positive real. The displayed
identity follows by differentiation and the common value at $s=0$.
All logarithm arguments at the specified $s$ are nonzero algebraic
numbers, and all coefficients are algebraic. The integral is strictly
positive, so Baker's theorem proves its transcendence. Finally,
\[
 \int_\xi^1 b_r(x)\,dx=U_{1-r}(1-\xi),
\]
and the same argument applies.
\end{proof}

\begin{lemma}[Arithmetic synchronization of moments]
\label{E:aa:synchronization}
Let $r\in\mathbb Q\cap(0,1)$ and
$S\in(\overline{\mathbb Q}\cap\mathbb R)[X]$. Then
\[
 \int_0^\xi S(x)b_r(x)\,dx=0
 \quad\Longleftrightarrow\quad
 \int_\xi^1 S(x)b_r(x)\,dx=0.
\]
More precisely, define the unique constant $c(S)$ and polynomial
$R_S$ by
\begin{equation}\label{E:aa:reduction}
 S(X)=c(S)+X(1-X)R_S'(X)+(r-X)R_S(X).
\end{equation}
Either vanishing condition is equivalent to the two algebraic
linear conditions
\begin{equation}\label{E:aa:algebraic-conditions}
 c(S)=0,\qquad R_S(\xi)=0.
\end{equation}
If $S$ has coefficients in a field $K\supset\mathbb Q$
containing $\xi$, then $c(S)$ and all coefficients of $R_S$
belong to $K$.
\end{lemma}

\begin{proof}
The operator in \eqref{E:aa:reduction} sends $X^j$ to
$(j+r)X^j-(j+1)X^{j+1}$. Descending elimination therefore gives
a unique $R_S$ and constant $c(S)$, over the coefficient field.
Put $H_S(x)=x^r(1-x)^{1-r}R_S(x)$. Direct differentiation and
$H_S(0)=H_S(1)=0$ give
\begin{align*}
 \int_0^\xi S b_r&=c(S)U_r(\xi)+H_S(\xi),\\
 \int_\xi^1 S b_r&=c(S)\{U_r(1)-U_r(\xi)\}-H_S(\xi).
\end{align*}
Here $H_S(\xi)$ is algebraic. Either integral can vanish only if
$c(S)=0$, by Lemma~\ref{E:aa:base}. It then forces
$H_S(\xi)=0$, equivalently $R_S(\xi)=0$ since
$0<\xi<1$. Conversely, these two conditions make both integrals
zero. All operations in descending elimination take place over $K$.
\end{proof}

\subsection{Proof of the classification}

\begin{proof}[Proof of Theorem~\ref{E:aa:main}]
Suppose first that $a,b$ are integers. Since they are positive,
$\rho$ is a polynomial with coefficients in $K$. Each moment on
either interval belongs to $K$, by termwise integration and
$\xi\in K$. The invertible linear system in
Lemma~\ref{E:aa:normality} consequently has its monic solution in $K[X]$.

Otherwise neither shape is an integer, since their sum is an integer.
Write
\[
 a=A+r,\qquad b=B+1-r,\qquad A,B\in\mathbb Z_{\ge0},
 \quad r\in\mathbb Q\cap(0,1).
\]
Then $\rho=g b_r$, where
$g(X)=w(X)X^A(1-X)^B\in K[X]$ is positive on $(0,1)$.

Assume that $P_{m,n}$ has algebraic coefficients and set
$h=\max(m,n)$. For every $0\le j<h$, at least one of
\eqref{E:aa:left} and \eqref{E:aa:right} holds. Apply
Lemma~\ref{E:aa:synchronization} to the algebraic polynomial
$S=X^jP_{m,n}g$. Both integrals then vanish. Thus $P_{m,n}$
satisfies $h$ orthogonality conditions on each interval. Positivity
forces at least $2h$ distinct zeros, so
\[
 m+n=\deg P_{m,n}\ge 2\max(m,n).
\]
This is possible only when $m=n$.

Conversely let $m=n=k$. The $2k$ equations
\begin{equation}\label{E:aa:balanced-system}
 c(X^jPg)=0,\qquad R_{X^jPg}(\xi)=0,
 \qquad 0\le j<k,
\end{equation}
are homogeneous linear equations over $K$ in the $2k+1$
coefficients of $P$ of degree at most $2k$. They have a nonzero
solution in $K[X]$. Lemma~\ref{E:aa:synchronization} shows that
this solution satisfies both orthogonality lists. By the sign-change
argument it has degree exactly $2k$. Divide by its nonzero leading
coefficient. Uniqueness identifies the resulting monic polynomial
with $P_{k,k}$, proving $P_{k,k}\in K[X]$.

Finally, a monic polynomial has all roots algebraic if and only if
all its coefficients are algebraic: one direction uses elementary
symmetric functions, and the other uses the algebraic closedness of
$\overline{\mathbb Q}$. The zero distribution follows from
Lemma~\ref{E:aa:normality}.
\end{proof}

\subsection{Quadrature, an explicit example, and the literature boundary}

\begin{corollary}[Algebraic cuts destroy all-algebraic Gaussian nodes]
\label{E:aa:gauss}
If $a,b>0$ are rational nonintegers with integer sum, then for every
algebraic $\xi\in(0,1)$ and every $N\ge1$, the $N$-node Gaussian
quadrature rule for $\rho$ restricted to either $[0,\xi]$ or
$[\xi,1]$ has at least one transcendental node. In contrast, all
zeros of every balanced $P_{N,N}$ are algebraic. These assertions
remain true under every positive algebraic polynomial factor $w$.
\end{corollary}

\begin{proof}
The Gaussian nodes are the zeros of the degree-$N$ monic orthogonal
polynomial, namely $P_{N,0}$ or $P_{0,N}$. Apply the theorem.
The balanced assertion is the case $(m,n)=(N,N)$.
\end{proof}

For $w=1$ and $(a,b)=(r,1-r)$, the classical diagonal Rodrigues
formula specializes to
\begin{equation}\label{E:aa:Rodrigues}
 P_{k,k}(x)=\frac{(-1)^k}{(2k)_k\,b_r(x)}
 \frac{d^k}{dx^k}
 \left[x^{k+r-1}(1-x)^{k-r}(x-\xi)^k\right],\qquad k\ge1.
\end{equation}
Leibniz' rule makes the right side a polynomial of degree at most
$2k$. Integration by parts $k$ times on each interval proves the
orthogonality: derivatives of order at most $k-1$ of the expression
in brackets vanish at all relevant endpoints, and
$\deg(x^j)<k$. Its leading coefficient before normalization is
$(-1)^k(2k)_k$, as can be read from the highest-degree terms.
This establishes \eqref{E:aa:Rodrigues} and its monic normalization.
For example,
\[
 P_{1,1}(x)=x^2-\frac{1+r+\xi}{2}x+\frac{r\xi}{2}.
\]
At $r=1/3$ and $\xi=1/2$ its two zeros are
\[
 \frac{11-\sqrt{73}}{24},\qquad
 \frac{11+\sqrt{73}}{24},
\]
one on each side of the cut. Every off-diagonal polynomial for the
same measures has a transcendental coefficient and a transcendental
zero, including the degree-one polynomial for either single interval.

The diagonal formula is classical. After an affine change of variable,
this is the Jacobi--Angelesco system with zero exponent at the common
endpoint; see Van Assche and Coussement \cite[\S3.5,
Eq.~(3.45)]{aa:VanAsscheCoussement}. Their equations (3.48)--(3.50)
also express near-diagonal polynomials through incomplete moment
ratios. We do not count the diagonal Rodrigues formula, Angelesco
normality, polynomial beta reduction, or Baker's theorem as new.
The addition proved here is the exact arithmetic exclusion of every
unbalanced index throughout the specified parameter family, including
positive polynomial modifications. We make no historical-priority
claim. Nor does this result assert algebraic independence of different
zeros or settle an unrestricted problem for arbitrary rational beta
shapes with noninteger sum.


\clearpage

\section{Recurrence scope and remaining questions}\label{B:sec:remarks}

\begin{remark}[Random walks]\label{B:rem:walks}
For the $\ell$-step planar walk with unit steps, the even moments $W_\ell(2n)=\sum_{a_1+\dots+a_\ell=n}\binom{n}{a_1,\dots,a_\ell}^2$ have the generating function $\sum_nW_\ell(2n)t^{2n}/(n!)^2=I_0(2t)^\ell$. Its minimal operator is the $\ell$-th symmetric power of \eqref{B:eq:T}, by the Galois argument of Proposition~\ref{B:prop:solspace}. The solution space is spanned by $i_0^{\ell-b}k_0^b$, with rates $2(\ell-2b)$, exponents $\ell/2$ at infinity and exponent $0$ (with logarithms) at the origin.

The identity $\binom{\ell+1}2-0-(\ell+1)\frac{\ell}2=0$ again excludes apparent singularities. In the variable $x$ one finds $P_0=\theta^{\ell+1}$ and $\deg P_k\le\ell+1-2k$. This recovers the recurrence of order $\cl{\ell/2}$ and degree $\ell-1$, with leading coefficient $n^{\ell-1}$, found by Verrill \cite{Ver04}; see \cite[Thms.~1 and~2]{BNSW11} and also \cite{BSWZ12}. Lemma~\ref{B:lem:types} excludes recurrences of smaller order.
The exponent-divisibility argument alone does not determine the least
degree, because all exponents at the origin vanish. The logarithmic
monodromy argument in Theorem~\ref{rev:verrill-min-degree} below proves
that the least degree is $\ell-1$, even among recurrences of arbitrary
order valid only for sufficiently large indices.

The unitary and the random-walk operators thus have the same generic structure: both have solution spaces isomorphic to $\mathrm{Sym}^\ell$ of the Bessel module. The difference is entirely in the lattice at the origin. For the walk the local monodromy at $0$ is a single unipotent Jordan block and all exponents are $0$. For $U(\ell)$ the same Jordan type is realized with exponents
$-b(\ell-b)$, which provide the direct divisibility argument for
degree $\ell-1$. For planar walks, the proof below instead uses
Jordan-block length.
\end{remark}

\subsection{The polynomial degree in the planar-walk recurrence}
\label{rev:verrill-degree}

The ordinary generating function retains information that is lost by
using only the numerical exponents of the double exponential generating
function. The relevant invariant is the size of its unipotent monodromy
block.

\begin{lemma}[Recurrence degree and local monodromy]
\label{rev:degree-monodromy}
Let $A(z)=\sum_{n\ge0}a_nz^n$ be a convergent germ whose minimal
differential equation is regular singular or ordinary at zero.
Suppose $(a_n)$ satisfies an eventual nonzero polynomial-coefficient
recurrence of degree at most $d$. Then every Jordan block of the local
monodromy of the minimal equation of $A$ has size at most
$\max\{d,1\}$. If $d=0$, $A$ is rational.
\end{lemma}
\begin{proof}
Write the recurrence as
$\sum_{s=0}^M c_s(n)a_{n+s}=0$ for all sufficiently large $n$,
where $\deg c_s\le d$, and put $\theta=z\partial_z$. The nonzero
operator
\[
 R=\sum_{s=0}^M z^{M-s}c_s(\theta-s)
\]
has order $e\le d$, and $RA=P$ for a polynomial $P$.
If $P=0$, the minimal equation of $A$ is a right factor of $R$,
so its solution space has dimension at most $e$.
If $P\ne0$, set $Q=\partial_z-P'/P$. On a sufficiently small
punctured disc about zero, the solution space of $QR$ is
\[
                 \operatorname{Sol}(QR)
                  =\operatorname{Sol}(R)\oplus\mathbb C A.
\]
Indeed, $QRf=0$ gives $Rf=cP$, and hence $f-cA\in\operatorname{Sol}(R)$;
the sum is direct since $RA=P\ne0$.
The second summand is fixed by local monodromy because $A$ is
holomorphic at zero. The first has dimension $e$. Thus every Jordan
block on this direct sum has size at most $\max\{e,1\}$.
The minimal equation of $A$ right divides $QR$, so its local solution
space is an invariant subspace and obeys the same bound.
When $d=0$, $R$ is multiplication by a nonzero polynomial, and
$A=P/R$ is rational.
\end{proof}

\begin{theorem}[Minimal degree of the planar-walk recurrence]
\label{rev:verrill-min-degree}
For $\ell\ge1$, let
\[
 a_n=W_\ell(2n)=
 \sum_{a_1+\cdots+a_\ell=n}\binom{n}{a_1,\ldots,a_\ell}^{\!2}.
\]
Every nonzero polynomial-coefficient recurrence for $(a_n)$ valid for
all sufficiently large $n$ has degree at least $\ell-1$.
Together with the Verrill recurrence and the order bound in
Remark~\ref{B:rem:walks}, the smallest possible recurrence order and
the smallest possible coefficient degree are respectively
$\lceil\ell/2\rceil$ and $\ell-1$; the same recurrence attains both.
\end{theorem}
\begin{proof}
Put $A(z)=\sum_{n\ge0}a_nz^n$. The estimate $a_n\le\ell^{2n}$
gives convergence near zero. The recurrence in
Remark~\ref{B:rem:walks}, originally due to Verrill
\cite{Ver04,BNSW11}, has the form
\[
 \sum_{k=0}^K R_k(n)a_{n-k}=0\quad(n\ge1),\qquad
 R_0(n)=n^{\ell-1},\qquad \deg R_k\le\ell-1,
 \quad K=\lceil\ell/2\rceil,
\]
with negative-index terms omitted. For $\ell\ge2$ it gives the
homogeneous equation
\[
 \left(\theta^{\ell-1}
       +\sum_{k=1}^K z^kR_k(\theta+k)\right)A=0.
\]
The $n=0$ coefficient vanishes as well. This operator has order
$\ell-1$ and is Fuchsian at zero, since its lowest
$\theta$-coefficient is $\theta^{\ell-1}$ and all the other
$\theta$-coefficients have degree at most $\ell-1$.
Its minimal right factor, the minimal equation of $A$, is consequently
regular singular or ordinary there.

The second factorial Borel transform is
\[
 A_2(x)=\sum_{n\ge0}\frac{a_nx^n}{(n!)^2}
                 = I_0(2\sqrt{x})^\ell.
\]
Its minimal differential module is the $\ell$th symmetric power of
the rank-two Bessel module, as in Remark~\ref{B:rem:walks}.
The base module has a nontrivial unipotent monodromy block of size two
at zero. Its $\ell$th symmetric power therefore has a single block
$J_{\ell+1}(1)$.

If a recurrence of degree $d\ge1$ exists, Lemma~\ref{rev:degree-monodromy}
bounds the unit-eigenvalue Jordan blocks of $A$ by $d$.
Theorem~\ref{B:up:borel-bound}, applied twice, bounds those of $A_2$
by $d+2$. Hence $\ell+1\le d+2$, or $d\ge\ell-1$.
For $d=0$, $A$ would be rational and the same Borel theorem would
bound every block of $A_2$ by three; this excludes $\ell\ge3$.
For $\ell=2$, the identity
$a_n=\binom{2n}{n}$ gives $A(z)=(1-4z)^{-1/2}$, which is not
rational and thus also excludes degree zero.
Finally, for $\ell=1$ the sequence is constantly one and has a
degree-zero recurrence. These cases prove the lower bound for every
$\ell$, and the stated Verrill recurrence supplies the upper bound.
\end{proof}

\begin{remark}
This degree bound uses the logarithmic monodromy at the origin, whereas
the exponent-divisibility argument alone sees only the repeated
exponent zero. The sharper dictionary in Theorem~\ref{V:thm:F}(c) additionally proves that $N^{\ell-1}$ divides the leading coefficient of every backward recurrence, at arbitrary order.
\end{remark}

\begin{remark}[Scope of the recurrence classifications]
\begin{enumerate}
\item Theorem~\ref{V:thm:A} proves minimality for all positive $\beta$
and nonnegative integer twists. The Bessel-module and jet-field
classification extends to every positive even $\beta$ by
Theorems~\ref{evb:module}--\ref{evb:arithmetic}; no such classification
is asserted for all non-even $\beta$.
\item The minimum degree of every twisted Toeplitz sequence is
$\ell-1$, including the exceptional twists, by
Theorem~\ref{vtd:even-degree}. For non-even $\beta$, the generic
degree theorem and the explicit degree-drop family leave the residual
exceptional cases specified in Section~\ref{V:sec:open}.
\item Remark~A.9 of \cite{FW26} concerns auxiliary functions of the matrix system and is not addressed here.
\item The exponent at every nonzero radial-unitary threshold is determined in Theorem~\ref{B:upden:unitary}. Its value is $\frac12((\ell-j)^2+j^2-3)$. This indicial calculation does not determine which logarithmic terms occur in a particular physical density branch.
\item The ordinary tableau algebraicity classification is established by
Corollary~\ref{B:up:three-algebraicity}. The proof uses local monodromy and
the fully stated Borel theorem. The twisted permutation classification
is stronger and includes every arithmetic section.
\item Theorem~\ref{B:up:sample-order} includes every tableau height,
sampling interval and residue. The sampled coefficient degree is not
classified. Likewise, a polynomial factor of a global distributional
annihilator cannot be cancelled at its zeros without controlling possible
point-supported residues.
\end{enumerate}
\end{remark}



\section{Statistical and arithmetic scope}
\label{E:end:scope}

The statistical results specify estimation rates across the stated
parameter ranges and a worst-null testing consequence. The local Fourier
argument controls Jordan blocks under explicit regularity hypotheses.
The Franel classification includes arithmetic subsequences. The
Angelesco and tableau-value theorems have their own parameter
restrictions, recorded below.

The following distinctions remain part of the mathematical conclusions.
\begin{enumerate}[leftmargin=*]
\item The inequality $|C(t,m,j)|\le2^{m^2/t}$ with prefactor exactly
one is not proved. The exponent $\log2$ is optimal when an absolute
prefactor is allowed. An exponent $1/2$ is valid in the moderate-degree
asymptotic regime but impossible globally with a polynomial prefactor.
\item The mixture theorems use a known support interval and either an ordinary
unrestricted maximum-likelihood maximizer or a certified approximate fit. Approximate fits require the certificates in Section~\ref{E:upst:sec}.
They do not establish an unknown-support adaptation theorem. The testing rate is worst-null,
not a claim that every fixed null is equally difficult. Arbitrary
heterogeneous observation kernels still require a separate normalization
or comparison argument; the homogeneous posterior identity cannot simply
be applied within each stratum.
\item Recurrence width and differential order are different invariants.
The Franel results classify width over rational-function coefficient
fields. They do not determine the least possible coefficient degrees
of every minimum-width recurrence.
\item The Ap\'ery results concern regular eventual tails. Initial
singular indices of an arbitrary chosen polynomial recurrence are not
silently filled in. Complex linear combinations can share a limit;
injectivity is asserted over subfields of the algebraic numbers.
\item The iterated Fourier statement assumes the specified regularity
at the source point. It is not a bound for arbitrary irregular source
connections. For Bessel determinant values, absence of exceptional nonzero
points is proved for tableaux and every integer-twisted Toeplitz family.
The multipoint theorem retains distinct nonzero algebraic arguments.
These statements do not remove exceptional points for arbitrary rational
homogeneous observables.
\item For Angelesco polynomials, a noninteger parameter sum, an
arbitrary nonpolynomial tilt, or more than two intervals lies outside
the theorem. In the nonalgebraic cases the proof gives at least one
transcendental coefficient and at least one transcendental zero, not
transcendence of every zero.
\end{enumerate}

The finite verification in Section~\ref{E:cba:sec} is a proof certificate,
not numerical evidence: every comparison uses rational arithmetic and a
one-sided Taylor bound. It checks all $130816$ pairs with $1\le m\le
t<512$ through four exhaustive branches, including $12734$ direct
coefficient comparisons. Symbolic checks of representative Angelesco
formulas supplement, but do not replace, the all-parameter proofs.
The analytic, arithmetic and statistical arguments underwent separate internal checks. These checks improve the manuscript's auditability; they do not
constitute external peer review or establish historical priority.


\clearpage
\appendix

\section{Computational verification}\label{B:app:comp}

The proofs do not depend on the numerical checks listed here. The following
paragraphs preserve the computational record reported in the supplied
original manuscript. The original supplementary scripts were not supplied
with the two attachments, so these historical runs are not presented as
newly executed checks. The new exact verification for this revision is
reported separately below. The original record uses the prime
$p=2^{31}-1$, Bessel determinant series, and linear algebra over
$\mathbb F_p$ unless otherwise specified.

\begin{enumerate}
\item \emph{The operator $L^t_\ell$, $1\le\ell\le10$.} We found no operator of order at most $\ell$ with coefficients of degree at most $3\ell+2$. For order $\ell+1$ there is a one-dimensional space of operators with coefficients of degree $\le\ell$, and none with degree $\le\ell-1$. The operator has leading coefficient $t^\ell$, its $\theta$-form contains only even powers of $t$, and the indicial roots at $0$ are exactly $\{-2\fl{j/2}(\ell-\fl{j/2})\}_{j=0}^{\ell}$.
\item \emph{The operator $L_\ell$ in $x$, $1\le\ell\le13$.} The coefficient degrees are exactly $\fl{(\ell-1+m)/2}$. The leading coefficients of the $\theta$-coefficients $P_k$ agree with Theorem~\ref{B:thm:A}(b), and the subleading coefficients with \eqref{B:eq:eta}. For even $\ell$, the root of $P_K$ is $-\ell^2/8$.
\item \emph{Orders of the vertex integrals at $0$, $2\le\ell\le6$.} The coefficients of $\det(\alpha T_I+\beta T_K)$ were evaluated with $400$ significant digits at $t=10^{-5},10^{-7},10^{-9}$. The logarithmic slopes agree with $-2\fl{b/2}(\ell-\fl{b/2})$ up to the expected logarithmic corrections.
\item \emph{Recurrences for $u_\ell$, $2\le\ell\le12$.} The recurrence of Theorem~\ref{B:thm:B}(a), built from the computed $P_k$, holds for all $n<260$. The formulas for $R_0$, $R_K$ and the characteristic polynomial hold. The space of recurrences of order $K$ and degree $\ell-1$ is one-dimensional and is spanned by it.
\item \emph{Exact recurrences, $2\le\ell\le7$.} The polynomials $P_k$ were reconstructed over $\Q$ from two primes. The resulting recurrences (Table~\ref{B:tab:perm}) were checked in exact arithmetic against $u_\ell(n)=\sum_\lambda(f^\lambda)^2$, computed by the hook length formula independently of \eqref{B:eq:Gessel}, for $n<26$ ($\ell\le5$) and $n<20$ ($\ell=6,7$).
\item \emph{Tableaux, $2\le k\le9$, exactly.} Minimal recurrences were guessed modulo two primes and reconstructed over $\Q$. They were checked exactly against hook length sums for $n<56$ ($k\le6$) and $n<46$ ($k=7,8,9$). Theorem~\ref{B:thm:C} predicts the order, the degree, $c_0$ and $c_R$, and all predictions are confirmed. On the same data there is no recurrence of order $R-1$ and degree $\le r+1$, and none of order $R$ and degree $r-1$. Modulo $p$ the predicted orders, degrees and extreme coefficients were also confirmed for $k\le14$.
\item \emph{Matchings, $1\le\ell\le6$, exactly.} The values $m_\ell(n)$ were computed by counting oscillating tableaux of height at most $\ell$, independently of \eqref{B:eq:Sp}. The minimal recurrences have the order, degree and extreme coefficients predicted by Theorem~\ref{B:thm:D}. Modulo $p$, from \eqref{B:eq:Sp}, the same was confirmed for $\ell\le10$.
\item \emph{Twisted determinants, variable $x$.} For $1\le\ell\le8$ and $0\le q\le5$, for $q=\ell\in\{6,7,8\}$ and for $(\ell,q)\in\{(3,9),(4,7),(5,10)\}$, the series $G_{\ell,q}$ was computed from the determinant. Its first nine coefficients agree, in exact arithmetic, with the expansion of Lemma~\ref{B:lem:twisted-expansion} evaluated by the hook length formula. In every case there is no operator of order $\ell$ with coefficients of degree $\le3\ell+2$, and there is a unique operator of order $\ell+1$ and degree $\ell$, none of degree $\ell-1$. Its coefficients satisfy $q_{\ell+1}=x^\ell$, $x^{m-1}\mid q_m$ and $\deg q_m=\fl{(\ell-1+m)/2}$, and $P_0$, the $\lc(P_k)$ and $P_K$ are as in Theorem~\ref{B:thm:twisted}(a). The Taylor coefficients satisfy no recurrence of order $K-1$ and degree $\le3\ell+3$.
\item \emph{Twisted determinants, variable $t$.} For $21$ pairs $(\ell,q)$ with $\ell\le8$, among them the exceptional pairs $(3,4)$, $(4,3)$, $(5,12)$, $(6,8)$ and $(8,6)$, the minimal operator of $D_{\ell,q}$ has order $\ell+1$, leading coefficient $t^\ell$ or $t^{\ell+1}$ as predicted by Theorem~\ref{B:thm:twisted}(b), and the predicted exponents at $0$.
\item \emph{Twisted recurrences.} For $2\le\ell\le7$ with $1\le q\le4$, and for $q=\ell\in\{5,6,7\}$, the sequence $u_{\ell,q}$ satisfies a unique recurrence of order $K$ and degree $\ell-1$, with leading coefficient $\prod_{i=1}^{\ell-1}(n+i(\ell+q-i))$; there is none of order $K-1$ and degree $\le3\ell+3$, and none of order $K$ and degree $\ell-2$.
\item \emph{Displacement with drift.} For $\ell\le6$ and $q\le8$, the maximum of $E_q(\pi,S)$ over $\pi\in S_\ell$ and $|S|=b$ was computed by enumeration. It equals $f_{(b)}$ for every $b$.
\end{enumerate}

A second set of scripts, written independently of the first, computes all sequences by dynamic programming over Young diagrams and oscillating tableaux instead of the Bessel determinants. It confirmed the statements of Sections~\ref{B:sec:unitary}--\ref{B:sec:densities} for small parameters, including the vertex constants $C_b$, the orders of the $\Psi^{(q)}_b$ at $0$, and Lemma~\ref{B:lem:disp-q} for $\ell\le9$, $q\le12$. This sentence records the original manuscript's reported supplementary checks; those scripts are not part of the present source-and-PDF deliverable.


\subsection*{Exact checks executed for this revision}
The finite continuant and path-matching formulas were compared symbolically
for $1\le\ell\le8$. Independently computed hook-length sums verified
300 twisted recurrence instances with $1\le\ell\le5$, $0\le q\le3$
and $1\le n\le15$. The exponent-zero condition and the Pythagorean
square criterion agreed in all $12,800$ cases with
$1\le\ell\le80$, $1\le q\le160$. A separate exact determinant
calculation checked the Catalan-tail identity in 99 pairs $(q,n)$.
These finite checks detect normalization errors; the all-parameter claims
are established by the preceding proofs.

Additional independent checks use rational arithmetic and symbolic
polynomial identities in SymPy. They make no numerical claim of
algebraic independence. The checks were executed for this revision,
with the following results.
\begin{itemize}
\item The Laurent jet determinant, including the closed constant in
Proposition~\ref{B:bv:constant}, was checked for
$1\le\ell\le3$, $0\le q\le2$, and $(\ell,q)=(4,0)$: ten cases.
For the nine cases with $\ell\le3$, the inverse matrices were checked
exactly, and all their denominators were confirmed to be monomials
in $x$. The coordinate conversions between twists $q=0,1,2$
were also verified exactly.
\item The full indicial polynomials in
Theorem~\ref{B:upden:unitary} were checked at all nonzero thresholds
for $2\le\ell\le8$: nineteen thresholds. The operators were built
from the integral Ore continuant and the exact recurrence-to-density
dictionary.
\item The tableau indicial polynomials in
Corollary~\ref{B:upden:other} were checked for $2\le k\le7$:
eighteen thresholds. The symplectic ones were checked for
$1\le\ell\le3$: nine thresholds. In these checks the scalar
operators were constructed independently from the homogeneous
Bessel determinant forms, and their Fourier transforms were
computed by the Weyl-algebra product rule.
\end{itemize}
All these finite checks passed. They confirm normalization and
indexing in the formulas, while the preceding proofs establish
the stated results for all parameters.

\section{Proof of Theorem~\ref{B:thm:jacobi}}\label{B:app:jacobi}

For completeness we give the proof of Theorem~\ref{B:thm:jacobi}, following \cite[Part~I]{Kei26}. Let $k=2r+\varepsilon$ with $r\ge1$, $a=\varepsilon-\frac12$, $w(u)=(2-u)^a(2+u)^{1/2}$ on $(-2,2)$ and $\tilde w(u)=(2-u)^a|2+u|^{1/2}$ on $(-\infty,-2)$.

\subsection*{B.1. One-variable functions}
For $t>0$ put
\[
\varphi_I(t)=\int_{-2}^2e^{tu}w(u)\,du,\qquad\varphi_K(t)=\int_{-\infty}^{-2}e^{tu}\tilde w(u)\,du.
\]

Let $g(u)=(4-u^2)w(u)$. Then $g'=w\cdot\bigl((2-2\varepsilon)-(2+\varepsilon)u\bigr)$ and $g(\pm2)=0$. Integrating $(e^{tu}g)'$ over $[-2,2]$, and similarly over $(-\infty,-2]$, shows that both functions satisfy
\[
B_\varepsilon f:=tf''+(2+\varepsilon)f'-(4t+2-2\varepsilon)f=0 .
\]
They are independent: $\varphi_I(0)$ is finite, while $\varphi_K(t)\to\infty$ as $t\to0^+$. A direct computation gives the operator identities
\[
B_0\circ(1+\tfrac12\partial)=(\tfrac12\partial+1)\circ T,\qquad B_1\circ\bigl((2t)^{-1}\partial\bigr)=\bigl(\partial\circ(2t)^{-1}\bigr)\circ T,
\]
with $T$ as in \eqref{B:eq:T}. Hence $B_\varepsilon$ is gauge equivalent to $T$ over $\C(t)$, and its differential Galois group acts on $\Sol(B_\varepsilon)$ as $\SL_2(\C)$ by Kolchin's theorem.

\subsection*{B.2. The Hankel form}
For $g_1,\dots,g_r\in\Sol(B_\varepsilon)$ put
\[
F(g_1,\dots,g_r)=\sum_{\sigma,\tau\in S_r}\sgn(\sigma)\sgn(\tau)\prod_{p=1}^rg_p^{(\sigma(p)+\tau(p)-2)} .
\]
It is symmetric and multilinear. Suppose $g_p=\int_{\gamma_p}e^{tu}\omega_p(u)\,du$, where each pair $(\gamma_p,\omega_p)$ is either $([-2,2],w)$ or $((-\infty,-2],\tilde w)$. Expanding $\Delta(u)^2=\det[u_p^{i-1}]^2$ gives
\[
F(g_1,\dots,g_r)=\int_{\gamma_1\times\cdots\times\gamma_r}e^{t\sum u_p}\Delta(u)^2\prod_p\omega_p(u_p)\,du .
\]
In particular $F(\varphi_I,\dots,\varphi_I)=r!\,(2\pi)^re^{-\varepsilon t}Y_k$ by Lemma~\ref{B:lem:jacobi-int}. Put $\Psi_b=F(\varphi_K^{(b\text{ times})},\varphi_I^{(r-b\text{ times})})$ for $0\le b\le r$. As in Proposition~\ref{B:prop:solspace}, Lemma~\ref{B:lem:galois-orbit} and polarization show that the solution space of the minimal operator of $e^{-\varepsilon t}Y_k$ is $V=\spn\{\Psi_0,\dots,\Psi_r\}$. The minimal operator of $Y_k$ then has solution space $e^{\varepsilon t}V$.

\subsection*{B.3. Vertices}
On $(-\infty,-2]^b\times[-2,2]^{r-b}$ the function $\sum_pu_p$ attains its maximum $\mu_b=2r-4b$ only at the vertex with first $b$ coordinates $-2$ and the others $2$. In the local coordinates $y_p=-2-u_p\ge0$ ($p\le b$) and $x_p=2-u_p\ge0$ ($p>b$), the integrand is
\[
e^{\mu_bt}e^{-t(\sum x+\sum y)}\,\Delta(x)^2\Delta(y)^2\prod x_p^a\prod y_p^{1/2}\,H(x,y),
\]
with $H$ real-analytic near $0$ and $H(0)=4^{2b(r-b)+(r-b)/2+ab}>0$. The weight in front of $H$ is positively homogeneous of degree $\gamma^{(k)}_b-r$, and Watson's lemma on the orthant gives
\[
\Psi_b^{(m)}(t)=C_be^{\mu_bt}t^{-\gamma^{(k)}_b}\bigl(\mu_b^m+O(t^{-1})\bigr),\qquad C_b=H(0)\,S_{r-b}(a)\,S_b(\tfrac12)>0,
\]
with $S_n(c)=\int_{(0,\infty)^n}e^{-\sum x}\Delta(x)^2\prod x^c\,dx$. The expansion is in integral powers of $t^{-1}$, and the expansions of the derivatives are the termwise derivatives. Since the $\mu_b$ are distinct, $\dim V=r+1$, and $y_b=e^{\varepsilon t}\Psi_b$ has rate $\lambda^{(k)}_b=\varepsilon+\mu_b$.

\subsection*{B.4. The origin}
Elements of $V$ are polynomials in $\varphi_I,\varphi_K$ and their derivatives with coefficients in $\C(t)$. The exponents of $B_\varepsilon$ at $0$ are $0$ and $-1-\varepsilon$, so $V\subset\mathcal N$.

On the rays substitute $u_p=-s_p/t$. Collecting powers of $t$, using $b+\varepsilon b+b(b-1)+2b(r-b)=-e_b$, gives
\[
t^{-e_b}\Psi_b(t)=\int\!\!\int e^{t\sum u'}e^{-\sum s}\Delta(s)^2\prod_{p,i}(s_p+tu'_i)^2\prod_p(s_p+2t)^a(s_p-2t)^{1/2}\,\Delta(u')^2\prod w(u'_i)\,ds\,du',
\]
over $s_p\ge2t$ and $u'\in[-2,2]^{r-b}$. For $0<t\le1$ the integrand is dominated by an integrable function. By dominated convergence,
\[
\lim_{t\to0^+}t^{-e_b}\Psi_b(t)=S_b(2r-2b+\varepsilon)\int_{[-2,2]^{r-b}}\Delta(u)^2\prod w(u_i)\,du>0 .
\]
Hence $y_b$ has leading part $C'_bt^{e_b}$ with $C'_b\ne0$. The $e_b$ are distinct, $e_0=0>e_1>\dots>e_r$, so by Lemma~\ref{B:lem:filtration} they are exactly the exponents at $0$.

\subsection*{B.5. Conclusion}
As in the proof of Theorem~\ref{B:thm:Dt}, Abel's identity and B.3 give
\[
\sum_{\tau\ne0}\kappa_\tau=\binom{r+1}2-\sum_j\gamma^{(k)}_j-\sum_be_b .
\]
Using $a+\frac12=\varepsilon$,
\[
\sum_j\gamma^{(k)}_j=\frac{r(r+1)(2r+1)}3+\varepsilon\frac{r(r+1)}2,\qquad\sum_be_b=\frac{r(r+1)(2r+1)}6-r^2(r+1)-\varepsilon\frac{r(r+1)}2,
\]
whose sum is $\frac{r(r+1)}2=\binom{r+1}2$. Hence there are no apparent singularities. The leading coefficient is $ct^A$, and Lemma~\ref{B:lem:top} gives $\deg p_m\le A$ and the formula for $\chi_k$. Since all exponents at $0$ are $\le0$ and $0$ is an exponent, the argument in the proof of Theorem~\ref{B:thm:Dt} gives $A=r$ and $t^{m-1}\mid p_m$. \qed
\clearpage

\section{Verification of the ensemble and moment formulas}\label{V:app:comp}
The proofs of Sections~\ref{V:overview}--\ref{V:sec:open} are analytic
and algebraic. The supplementary calculations associated with the
ensemble formulas provide finite checks, not general minimality proofs.
The following checks were reproduced for this integrated version.
\begin{enumerate}[leftmargin=*]
\item The symbolic circular operators for $n\le3$ and Jacobi operators
for $n\le2$ agree exactly with the supplied symbolic output.
\item Nine circular and nine Jacobi parameter cases with $n\le3$
satisfy the scalar coefficient recursion through degree $17$ modulo
$2^{31}-1$. These include fractional positive $\beta$, integer and
fractional twists, and Jacobi endpoint exponents $-1/2$.
\item For $n=1,2$, the Bessel moment recurrence agrees exactly with
the gamma-integral ratios
$c_1(k+2)/c_1(k)=(k+1)^2$ and
$c_2(k+2)/c_2(k)=(k+1)^3/[4(k+2)]$.
\item At $(n,\beta,q)=(2,1/2,-1/2)$ the scalar operator is
$x^2\partial^3+(7/4)x\partial^2-4x\partial$, verifying the stated
failure of the degree formula at that boundary.
\item Polynomiality of the sharpened dictionary coefficients was
checked for $1\le M\le4$, $0\le n_0\le3$, $0\le r\le3$.
\end{enumerate}
The supplied archive additionally reports larger finite sweeps,
quadrature tests and recurrence searches. These were not all rerun.
An operator search with bounded coefficient degree and a finite
Taylor prefix establishes only a statement about that tested ansatz;
it does not exclude operators of arbitrary degree or recurrences
beginning at an arbitrary later index. In particular the archived
circular search tests coefficient degree at most $2n+2$, and the
planar-walk search uses a $161$-term prefix. The general lower bounds
in this volume come from visibility, Laurent transfer and monodromy.

\clearpage\phantomsection
\addcontentsline{toc}{section}{References}
\begin{thebibliography}{99}
\raggedright
\bibitem{BR01}
J.~Baik and E.~M.~Rains, \emph{Algebraic aspects of increasing subsequences}, Duke Math.\ J.\ \textbf{109} (2001), 1--65.

\bibitem{Ber86}
A.~Berele, \emph{A Schensted-type correspondence for the symplectic group}, J.\ Combin.\ Theory Ser.\ A \textbf{43} (1986), 320--328.

\bibitem{BFK95}
F.~Bergeron, L.~Favreau and D.~Krob, \emph{Conjectures on the enumeration of tableaux of bounded height}, Discrete Math.\ \textbf{139} (1995), 463--468.

\bibitem{BG00}
F.~Bergeron and F.~Gascon, \emph{Counting Young tableaux of bounded height}, J.\ Integer Seq.\ \textbf{3} (2000), Article 00.1.7.

\bibitem{BNSW11}
J.~M.~Borwein, D.~Nuyens, A.~Straub and J.~Wan, \emph{Some arithmetic properties of short random walk integrals}, Ramanujan J.\ \textbf{26} (2011), 109--132.

\bibitem{BSWZ12}
J.~M.~Borwein, A.~Straub, J.~Wan and W.~Zudilin, \emph{Densities of short uniform random walks}, Canad.\ J.\ Math.\ \textbf{64} (2012), 961--990.

\bibitem{CDDSY07}
W.~Y.~C.~Chen, E.~Y.~P.~Deng, R.~R.~X.~Du, R.~P.~Stanley and C.~H.~Yan, \emph{Crossings and nestings of matchings and partitions}, Trans.\ Amer.\ Math.\ Soc.\ \textbf{359} (2007), 1555--1575.

\bibitem{CQRZ08}
W.~Y.~C.~Chen, J.~Qin, C.~M.~Reidys and D.~Zeilberger, \emph{Efficient counting and asymptotics of $k$-noncrossing tangled-diagrams}, preprint (2008), arXiv:0802.3491.

\bibitem{DLMF}
\emph{NIST Digital Library of Mathematical Functions}, \url{https://dlmf.nist.gov}, Chapter~10.

\bibitem{FK22}
P.~J.~Forrester and S.~Kumar, \emph{Differential recurrences for the distribution of the trace of the $\beta$-Jacobi ensemble}, Phys.\ D \textbf{434} (2022), 133220; arXiv:2011.00787.

\bibitem{FW26}
P.~J.~Forrester and F.~Wei, \emph{Higher-order linear differential equations for unitary matrix integrals: applications and generalisations}, SIGMA \textbf{22} (2026), 015; arXiv:2508.20797.

\bibitem{Ges90}
I.~M.~Gessel, \emph{Symmetric functions and P-recursiveness}, J.\ Combin.\ Theory Ser.\ A \textbf{53} (1990), 257--285.

\bibitem{Inc56}
E.~L.~Ince, \emph{Ordinary Differential Equations}, Dover, New York, 1956.

\bibitem{Joh97}
K.~Johansson, \emph{On random matrices from the compact classical groups}, Ann.\ of Math.\ (2) \textbf{145} (1997), 519--545.

\bibitem{Kei26}
Keio, \emph{Minimal differential and recurrence orders: Bessel periods, bounded-height tableaux, and banana integrals}, manuscript, October 2026.

\bibitem{KW23}
J.~P.~Keating and F.~Wei, \emph{Joint moments of higher order derivatives of CUE characteristic polynomials II: Structures, recursive relations, and applications}, Nonlinearity \textbf{37} (2024), 085009; arXiv:2307.02831.

\bibitem{Kol68}
E.~R.~Kolchin, \emph{Algebraic groups and algebraic dependence}, Amer.\ J.\ Math.\ \textbf{90} (1968), 1151--1164.

\bibitem{Mis19}
M.~Mishna, \emph{On standard Young tableaux of bounded height}, in: Recent Trends in Algebraic Combinatorics, Springer, 2019; arXiv:1805.08130.

\bibitem{Rai98}
E.~M.~Rains, \emph{Increasing subsequences and the classical groups}, Electron.\ J.\ Combin.\ \textbf{5} (1998), R12.

\bibitem{Reg81}
A.~Regev, \emph{Asymptotic values for degrees associated with strips of Young diagrams}, Adv.\ Math.\ \textbf{41} (1981), 115--136.

\bibitem{Sun90}
S.~Sundaram, \emph{Tableaux in the representation theory of the classical Lie groups}, in: Invariant Theory and Tableaux, IMA Vol.\ Math.\ Appl.\ \textbf{19}, Springer, 1990, 191--225.

\bibitem{Ver04}
H.~A.~Verrill, \emph{Sums of squares of binomial coefficients, with applications to Picard--Fuchs equations}, preprint (2004), arXiv:math/0407327.

\bibitem{vdPS03}
M.~van der Put and M.~F.~Singer, \emph{Galois Theory of Linear Differential Equations}, Grundlehren \textbf{328}, Springer, 2003.

\bibitem{BVBeukers2006}
F.~Beukers, \emph{A refined version of the Siegel--Shidlovskii theorem},
Ann. of Math. (2) \textbf{163} (2006), 369--379,
\url{https://doi.org/10.4007/annals.2006.163.369}.

\bibitem{LJR_Arinkin}
D.~Arinkin, \emph{Fourier transform and middle convolution for irregular
D-modules}, 2008, especially Corollary~2.3 and \S3.3.
\url{https://arxiv.org/abs/0808.0699}.

\bibitem{LJR_Sabbah}
C.~Sabbah, \emph{An explicit stationary phase formula for the local
formal Fourier--Laplace transform}, Contemp. Math. \textbf{474}
(2008), 309--330, especially Remarks~5.2 and~5.3(4).
\url{https://arxiv.org/abs/0706.3570}.

\bibitem{aa:BakerIII}
A. Baker, Linear forms in the logarithms of algebraic numbers. III,
\emph{Mathematika} \textbf{14} (1967), 220--228.
\href{https://doi.org/10.1112/S0025579300003843}{doi:10.1112/S0025579300003843}.

\bibitem{aa:VanAsscheCoussement}
W. Van Assche and E. Coussement, Some classical multiple orthogonal
polynomials, \emph{Journal of Computational and Applied Mathematics}
\textbf{127} (2001), 317--347.
\href{https://doi.org/10.1016/S0377-0427(00)00503-3}{doi:10.1016/S0377-0427(00)00503-3}.
Author manuscript: \url{https://wis.kuleuven.be/analyse/members/walter/clasmult.pdf}.

\bibitem{BT_KWB2025}
L. Kania, L. Wasserman and S. Balakrishnan,
\emph{Testing Random Effects for Binomial Data}, arXiv:2504.13977v3,
2025, Section 3 and Theorem 3.
\url{https://arxiv.org/pdf/2504.13977}.

\bibitem{BT_VKVK2019}
R. Korlakai Vinayak, W. Kong, G. Valiant and S. Kakade,
\emph{Maximum likelihood estimation for learning populations of parameters},
Proceedings of ICML 2019, PMLR \textbf{97}, 6448--6457.
\url{https://proceedings.mlr.press/v97/vinayak19a.html}.

\bibitem{BT_MMRS2025}
C. Musco, C. Musco, L. Rosenblatt and A. V. Singh,
\emph{Sharper Bounds for Chebyshev Moment Matching, with Applications},
Proceedings of COLT 2025, PMLR \textbf{291}, 4309--4358.
\url{https://proceedings.mlr.press/v291/musco25a.html}.

\bibitem{BT_DKS2008}
P. Diaconis, K. Khare and L. Saloff-Coste,
\emph{Gibbs sampling, exponential families and orthogonal polynomials},
Statistical Science \textbf{23} (2008), 151--178.

\bibitem{CB_Dette}
H.~Dette, \emph{New bounds for Hahn and Krawtchouk polynomials},
SIAM Journal on Mathematical Analysis \textbf{26} (1995), 1647--1659.
\url{https://arxiv.org/abs/math/9406223}; Corollary~3.6(b).

\bibitem{AC_Katz}
N.~M.~Katz, On the calculation of some differential Galois groups,
\emph{Invent. Math.} \textbf{87} (1987), 13--61, Corollary 4.4.8.
\url{https://web.math.princeton.edu/~nmk/old/calcdiffgal.pdf}.

\bibitem{AC_FSY}
J.~Fres\'an, C.~Sabbah and J.-D.~Yu, Hodge theory of Kloosterman
connections, \emph{Duke Math. J.} \textbf{171} (2022), 1649--1747.
\url{https://perso.pages.math.cnrs.fr/users/claude.sabbah/articles/kloosterman-revised10.pdf}.

\bibitem{AC_Stoll}
M.~Stoll, Bounds for the length of recurrence relations for
convolutions of P-recursive sequences,
\emph{European J. Combin.} \textbf{18} (1997), 707--712.
\url{https://doi.org/10.1006/eujc.1996.0123}.

\bibitem{AC_McIntosh}
R.~J.~McIntosh, Recurrences for alternating sums of powers of
binomial coefficients, \emph{J. Combin. Theory Ser. A}
\textbf{63} (1993), 223--233.
\url{https://doi.org/10.1016/0097-3165(93)90058-G}.

\bibitem{AC_MatalaAho}
T.~Matala-Aho, On recurrences for some hypergeometric type polynomials,
\emph{Ast\'erisque} \textbf{198--199--200} (1991), 237--244.
\url{https://www.numdam.org/item/AST_1991__198-199-200__237_0/}.

\bibitem{AC_GRvS}
I.~Gaiur, V.~Rubtsov and D.~van Straten, Product formulas for the
higher Bessel functions, \emph{Selecta Math. (N.S.)} \textbf{32}
(2026), article 12. \url{https://doi.org/10.1007/s00029-025-01114-w}.
Preprint: \url{https://arxiv.org/abs/2405.03015}.

\bibitem{AC_SZ}
A.~Straub and W.~Zudilin, Sums of powers of binomials, their Ap\'ery
limits, and Franel's suspicions, \emph{Int. Math. Res. Not.}
\textbf{2023} (2023), no.~11, 9861--9879.
\url{https://arminstraub.com/downloads/pub/franel-aperylimits.pdf}.

\bibitem{AC_CS}
M.~Chamberland and A.~Straub, Ap\'ery limits: experiments and proofs,
\emph{Amer. Math. Monthly} \textbf{128} (2021), no.~9, 811--824.
\url{https://arminstraub.com/downloads/pub/aperylimits.pdf}.

\bibitem{AC_vH}
M.~van Hoeij, Submodule approach to creative telescoping,
preprint, 2024, Remark~6.
\url{https://arxiv.org/abs/2401.08455}.

\bibitem{PC_Miao2021}
Z. Miao, W. Kong, R. Korlakai Vinayak, W. Sun and F. Han,
\emph{Fisher--Pitman permutation tests based on nonparametric Poisson
mixtures with application to single cell genomics},
author version dated 8 June 2021, arXiv:2106.03022v1,
Theorems 6.1--6.2 and Remark 6.1.
\url{https://ramyakv.github.io/testing-poisson-npmle.pdf}.

\bibitem{SFB_GrabinerMagyar}
D.~J.~Grabiner and P.~Magyar,
\emph{Random walks in Weyl chambers and the decomposition of tensor powers},
J. Algebraic Combin. \textbf{2} (1993), 239--260.

\bibitem{SFB_JinReidysWang}
E.~Y.~Jin, C.~M.~Reidys, and R.~R.~Wang,
\emph{Asymptotic analysis of $k$-noncrossing matchings},
arXiv:0803.0848 (2008), equation (1.4).
\url{https://arxiv.org/abs/0803.0848}.

\bibitem{UP_DLR1977}
A.~P. Dempster, N.~M. Laird and D.~B. Rubin,
Maximum likelihood from incomplete data via the EM algorithm,
\emph{Journal of the Royal Statistical Society, Series B}
\textbf{39} (1977), 1--22.
\href{https://doi.org/10.1111/j.2517-6161.1977.tb01600.x}
{doi:10.1111/j.2517-6161.1977.tb01600.x}.

\bibitem{FG_DickerZhao2016}
L.~H. Dicker and S.~D. Zhao,
High-dimensional classification via nonparametric empirical Bayes and
maximum likelihood inference,
\emph{Biometrika} \textbf{103} (2016), 21--34.
\href{https://doi.org/10.1093/biomet/asv067}
{doi:10.1093/biomet/asv067}.

\bibitem{FG_SoloffGuntuboyinaSen2025}
J.~A. Soloff, A.~Guntuboyina and B.~Sen,
Multivariate, heteroscedastic empirical Bayes via nonparametric maximum
likelihood,
\emph{Journal of the Royal Statistical Society, Series B}
\textbf{87} (2025), 1--32.
\href{https://doi.org/10.1093/jrsssb/qkae040}
{doi:10.1093/jrsssb/qkae040}.

\bibitem{MPJDelaygue2025}
E.~Delaygue, \emph{A Lindemann--Weierstrass theorem for $E$-functions},
arXiv:2210.12046v2, 7 March 2025,
\url{https://arxiv.org/abs/2210.12046}.
\bibitem{BBC06}
D.~H.~Bailey, J.~M.~Borwein and R.~E.~Crandall, \emph{Integrals of the Ising class}, J.\ Phys.\ A \textbf{39} (2006), 12271--12302.

\bibitem{BBBC07}
D.~H.~Bailey, D.~Borwein, J.~M.~Borwein and R.~E.~Crandall, \emph{Hypergeometric forms for Ising-class integrals}, Experiment.\ Math.\ \textbf{16} (2007), 257--276.

\bibitem{BBBG08}
D.~H.~Bailey, J.~M.~Borwein, D.~Broadhurst and M.~L.~Glasser, \emph{Elliptic integral evaluations of Bessel moments and applications}, J.\ Phys.\ A \textbf{41} (2008), 205203.

\bibitem{Bal00}
W.~Balser, \emph{Formal Power Series and Linear Systems of Meromorphic Ordinary Differential Equations}, Universitext, Springer, New York, 2000.

\bibitem{BS08}
J.~M.~Borwein and B.~Salvy, \emph{A proof of a recursion for Bessel moments}, Experiment.\ Math.\ \textbf{17} (2008), 223--230.

\bibitem{Erd56}
A.~Erd\'elyi, \emph{Asymptotic Expansions}, Dover, New York, 1956.

\bibitem{GM93}
D.~J.~Grabiner and P.~Magyar, \emph{Random walks in Weyl chambers and the decomposition of tensor powers}, J.\ Algebraic Combin.\ \textbf{2} (1993), 239--260.

\bibitem{Hor90}
L.~H\"ormander, \emph{The Analysis of Linear Partial Differential Operators I}, 2nd ed., Springer, Berlin, 1990.

\bibitem{Kac47}
M.~Kac, \emph{Random walk and the theory of Brownian motion}, Amer.\ Math.\ Monthly \textbf{54} (1947), 369--391.

\bibitem{Kan93}
J.~Kaneko, \emph{Selberg integrals and hypergeometric functions associated with Jack polynomials}, SIAM J.\ Math.\ Anal.\ \textbf{24} (1993), 1086--1110.

\bibitem{OEIS}
OEIS Foundation, \emph{The On-Line Encyclopedia of Integer Sequences}, \url{https://oeis.org}.

\bibitem{Syl54}
J.~J.~Sylvester, \emph{Th\'eor\`eme sur les d\'eterminants de M.~Sylvester}, Nouv.\ Ann.\ Math.\ \textbf{13} (1854), 305.

\bibitem{TT91}
O.~Taussky and J.~Todd, \emph{Another look at a matrix of Mark Kac}, Linear Algebra Appl.\ \textbf{150} (1991), 341--360.

\bibitem{Tit39}
E.~C.~Titchmarsh, \emph{The Theory of Functions}, 2nd ed., Oxford University Press, 1939.

\bibitem{Was65}
W.~Wasow, \emph{Asymptotic Expansions for Ordinary Differential Equations}, Interscience, New York, 1965.

\bibitem{Won89}
R.~Wong, \emph{Asymptotic Approximations of Integrals}, Academic Press, Boston, 1989.

\end{thebibliography}
\end{document}
